\documentclass[12pt,a4paper,openany]{book} %openany elimina las hojas en blano que queda al crear los capitulos (Chapter)
\usepackage[utf8]{inputenc}
\usepackage[T1]{fontenc}
\usepackage{amsmath}
\usepackage{amsfonts}
\usepackage{amssymb}
\usepackage{graphicx}
\usepackage{amsthm}%para usar proof
\usepackage{mathrsfs}
\usepackage{color}
\usepackage{appendix}
\usepackage{stmaryrd }
\usepackage{booktabs}
\renewcommand{\arraystretch}{1.5} % Ajusta el factor de estiramiento de las filas tablas
\usepackage{multirow} % Para \multirow
\usepackage{enumitem}
%\author{Henrri Quino Huerta}
%%%%%%%%%%%%%%%%%%%%%%%%%%%%%%%%%%%%%%%%%%%%%%%%%%%%%%%%%%%%%%%%%%%%%%%%%%%%%%%%%%%%%%
\usepackage[spanish]{babel} %paquete de idioma español
\usepackage[left=2.54cm,right=2.54cm,top=2.54cm,bottom=2.54cm]{geometry} %bordes con normas APA

\usepackage{setspace}%para dar espacios como en las barras de la cartula
\usepackage{float} %para fijar graficos
\usepackage[colorlinks=true,linkcolor=black,urlcolor=blue]{hyperref}
%%%%%%%%%%%%%%%%%%%%%%%%%%%%%%%%%ENCABEZADO Y PIE DE PAGINA
%%%%%%%%%%%%%%% pero creo que en book no funciona

%%%%%%%%%%%%%%%%%%%%%%%%%%%%%%% DEFINICIONES
\newtheorem{teorema}{Teorema}[section]
\newtheorem{lema}{Lema}[section]
\newtheorem{corolario}{Corolario}[section]
\newtheorem{afirmacion}{Afirmación}[section]
\newtheorem{definicion}{Definición}[section]
\newtheorem{proposicion}{Proposición}[section]
\newtheorem{ejemplo}{Ejemplo}[section]
\newtheorem{observacion}{Observación}[section]
\newtheorem{propiedades}{Propiedades}[section]
\newtheorem{ejercicio}{Ejercicio}[section]
\providecommand{\abs}[1]{\lvert#1\rvert}
\providecommand{\norm}[1]{\lVert#1\rVert}

\usepackage{tikz-cd}
\usepackage{tikz}
\usetikzlibrary{positioning}
\usepackage{mathrsfs}
% Redefinir la numeración de las notas al pie para que use estrellas
\usepackage{footmisc}

\renewcommand{\thefootnote}{\fnsymbol{footnote}}
\makeatletter
\renewcommand*{\@fnsymbol}[1]{
	\ifcase#1
	\or $\star$
	\or $\star\star$
	\or $\star\star\star$
	\or $\star\star\star\star$
	\or $\star\star\star\star\star$
	\or $\star\star\star\star\star\star$
	\or $\star\star\star\star\star\star\star$
	\else\@ctrerr\fi
}
\makeatother
\usepackage{mdframed}
\usepackage{xcolor}
\setlength{\footnotesep}{15pt}

\usepackage{pdfpages}
\newcommand{\inner}[2]{\left\langle #1, #2 \right\rangle}
\newcommand{\norma}[1]{\left\lVert #1 \right\rVert}

\usepackage{array}
\usepackage{float}
\usepackage{booktabs}
\usepackage{caption}


\usepackage[all,cmtip]{xy} %para diagramas de algebra
%\usepackage{extpfeil}% para flechas largas
\usepackage{mathtools}

\usepackage{tikz-cd}
\usetikzlibrary{arrows.meta, decorations.markings}




\begin{document}
\frontmatter
	\begin{titlepage}
	\begin{center}
		
		

		

		\begin{spacing}{2}
			
		 {\huge \textbf{ALGEBRA}}\\[10pt]
		 
		 
		 {\Large Henrri Quino Huerta}
			
		\end{spacing}

		
		\vspace{3cm}
	
		{\Large IMCA, PERÚ}
		
		\vspace{0.3cm}
		{\Large 2025}
	\end{center}

\end{titlepage}


	\tableofcontents

	\mainmatter
	
	\chapter{Estructuras algebraicas fundamentales, y Categorías.}
	
\begin{mdframed}[backgroundcolor=gray!30, hidealllines=true, innertopmargin=2mm, innerbottommargin=2mm, innerleftmargin=2mm, innerrightmargin=2mm]
	\begin{center}
		Tenga en cuenta, en adelante, que
		$		\mathbb{N} := \{0, 1, 2, 3, \ldots\}.$
		
	\end{center}
\end{mdframed}



\section{Estructuras algebraicas fundamentales}

Sean $A$, $B$ y $C$ conjuntos no vacíos.
\begin{definicion}
	\upshape{
		Una \textit{ley de composición} es una función $A \times B \longrightarrow C$.
	}
\end{definicion}

\begin{ejemplo}[Ley de composición de funciones]
	\upshape{
	Sean $X$ y $Y$ conjuntos no vacíos. Definimos el conjunto de todas las funciones de $X$ a $Y$ como
	\[
	Y^X := \left\{ f : X \to Y \mid f \text{ es una función} \right\}.
	\]
	Sea $Z$ un tercer conjunto no vacío. Consideramos la siguiente ley de composición:
	\begin{align*}
		Y^X \times Z^Y & \longrightarrow Z^X \\
		(f, g) & \longmapsto g \circ f.
	\end{align*}
	}
\end{ejemplo}

\begin{definicion}
	\upshape{
		Se llama:
		\begin{enumerate}
			\item \textit{Operación externa} (acción) a una función $A \times A \longrightarrow B$,
			\item \textit{Operación interna} a una función $A \times A \longrightarrow A$.
		\end{enumerate}
	}
\end{definicion}

\begin{definicion}
	\upshape{
		Sean $M$ un conjunto y $\ast : M \times M \longrightarrow M$ una ley de composición interna (es decir, una operación interna). Entonces:
		\begin{enumerate}
			\item La pareja $(M,\ast)$ se llama un \textit{magma}.
			
			\item La ley $\ast$ se dice \textit{asociativa} si
			\[
			a \ast (b \ast c) = (a \ast b) \ast c \quad \text{para todo } a,b,c \in M.
			\]
			
			\item Un \textit{semigrupo} es un magma cuya ley es asociativa.
			
			\item Un elemento $e \in M$ es una \textit{identidad} si
			\[
			e \ast a = a = a \ast e \quad \text{para todo } a \in M.
			\]
			
			\item Un \textit{monoide} es un semigrupo que posee una identidad.
			
			\item Un \textit{grupo} es un monoide en el que todo elemento $a \in M$ admite un \textit{inverso} $a^{-1} \in M$ tal que
			\[
			a \ast a^{-1} = e = a^{-1} \ast a.
			\]
			
			\item Decimos que la ley $\ast$ es \textit{conmutativa} si
			\[
			a \ast b = b \ast a \quad \text{para todo } a,b \in M.
			\]
			
			\item Una estructura $(M,\ast)$ (sea magma, semigrupo, monoide o grupo) se dice \textit{conmutativa} si la ley $\ast$ es conmutativa. En el caso de grupos, también se dice que es un \textit{grupo abeliano}.
		\end{enumerate}
	}
\end{definicion}


\begin{ejemplo}
	{\color{white} Texto en azul.}
	\upshape{		
		\begin{itemize}
			\item $(\mathbb{Z}^+, +)$ es un semigrupo que no es un monoide, ya que no contiene el $0$, que sería el neutro para la suma.
			
			\item $(\mathbb{N}, \cdot)$ es un monoide. El neutro es el número $1$, ya que $1 \cdot a = a \cdot 1 = a$ para todo $a \in \mathbb{N}$.
			
			
			\item $M_{n \times n}(\mathbb{N})$, con la multiplicación de matrices, es un monoide. La matriz identidad actúa como elemento neutro.
			
			\item $M_{n \times n}(\mathbb{R}^+)$, también con multiplicación de matrices, es un monoide.
			
			\item El conjunto $\mathcal{P}(X)$ de subconjuntos de un conjunto $X$, con la operación unión, es un monoide. La identidad es el conjunto vacío $\emptyset$.
			
			\item El conjunto $\mathcal{P}(X)$ con la operación intersección también forma un monoide, siendo $X$ el elemento neutro.
			
			\item El conjunto de funciones de $X$ en sí mismo, con la composición de funciones, forma un monoide. La función identidad de $X$ actúa como neutro.
			
			\item $(\mathbb{Z}, +)$ es un grupo. Todo entero tiene inverso aditivo.
			
			\item $(\mathbb{Q}^\ast, \cdot)$ es un grupo, donde $\mathbb{Q}^\ast := \mathbb{Q} \setminus \{0\}$, con la multiplicación de racionales. Todo elemento tiene inverso multiplicativo.
			
			\item El conjunto de matrices invertibles $GL_n(\mathbb{R})$, con la multiplicación de matrices, es un grupo. El neutro es la matriz identidad, y toda matriz invertible tiene inversa.
		\end{itemize}
	}
\end{ejemplo}



\begin{ejemplo}
	\upshape{
		Sea \(n \in \mathbb{N}\), \(n \geq 1\). El conjunto \(\mathbb{Z}_n := \{[0], [1], \dots, [n{-}1]\}\), formado por clases de enteros módulo \(n\), con \([a]\) denotando la clase de \(a \in \mathbb{Z}\), es un grupo abeliano bajo la suma: la identidad es \([0]\) y el inverso de \([a]\) es \([-a]\).
	}
\end{ejemplo}



\begin{proposicion} \label{jkhgdvnjf}
	\upshape{
		Sea \(n \geq 2\) un número entero. Definimos
		\[ 
		(\mathbb{Z}_n)^\ast := \{ [a] \in \mathbb{Z}_n \mid [a] \neq [0] \},
		\]
		donde la operación en \((\mathbb{Z}_n)^\ast\) es la multiplicación de clases módulo \(n\), es decir, 
		\[ 
		[a] \cdot [b] := [a \cdot b] \quad \text{para} \; a, b \in \mathbb{Z}.
		\]Entonces, \((\mathbb{Z}_n)^\ast\), con esta operación de multiplicación, es un grupo si y solo si \(n\) es primo. Además, es un grupo abeliano.
	}
\end{proposicion}

\begin{proof}
	{\color{white} Texto en azul.}\\
	
	\noindent [$\Rightarrow$] Supongamos que \(n\) no es primo, entonces existen \(a, b \in \mathbb{Z}\) con \(1 < a, b < n\) tales que \(n = ab\). Tomando clases, vemos que \([ab] = [0]\). Por otro lado, como \(1 < a, b < n\), se tiene que \([a], [b] \in (\mathbb{Z}_n)^\ast\). Como \((\mathbb{Z}_n)^\ast\) tiene estructura de grupo, debe cumplirse que \([ab] \in (\mathbb{Z}_n)^\ast\), lo cual contradice el hecho de que \([ab] = [0]\). Así, \(n\) es primo.\\
	
	\noindent [$\Leftarrow$] Verifiquemos que \((\mathbb{Z}_n)^\ast\) es un grupo:
	\begin{itemize}
		\item La multiplicación es una operación interna: Sean \( [a], [b] \in (\mathbb{Z}_n)^\ast\). Supongamos que \( [a][b] \notin (\mathbb{Z}_n)^\ast \). Entonces \( [ab] = [0] \), lo que implica que \( n \mid ab \). Como \( n \) es primo, tenemos que \( n \mid a \) o \( n \mid b \), por lo tanto, \( [a] = [0] \) o \( [b] = [0] \). Esto contradice nuestra suposición de que ambos \( [a] \) y \( [b] \) son diferentes de \( [0] \). Así, la multiplicación es una operación interna.
		
		\item La multiplicación es asociativa: La multiplicación en \( (\mathbb{Z}_n)^\ast \) es asociativa, ya que la multiplicación en \( \mathbb{Z} \) es asociativa.
		
		\item El elemento identidad es \( [1] \): El elemento identidad en \( (\mathbb{Z}_n)^\ast \) es \( [1] \), ya que para cualquier \( [a] \in (\mathbb{Z}_n)^\ast \), se cumple que \( [1][a] = [a] = [a][1] \).
		
		\item Todo elemento es inversible: Sea \( [a] \in (\mathbb{Z}_n)^\ast \). Como \( n \nmid a \) y \( n \) es primo, se cumple que \( \gcd(a, n) = 1 \). Por el algoritmo de Euclides, existen \( \alpha, \beta \in \mathbb{Z} \) tales que \( \alpha a + \beta n = 1 \). Tomando clases módulo \(n\), obtenemos \( [a][\beta] = [1] \). Por lo tanto, \( [\beta] \) es el inverso de \( [a] \) en \( (\mathbb{Z}_n)^\ast \), y claramente \( [\beta] \neq [0] \).
	\end{itemize}
	Así, \( (\mathbb{Z}_n)^\ast \) es un grupo.
\end{proof}



\section{Aplicación a la teoría de categorías}

\begin{definicion}
	\upshape{
		Una \textit{categoría} $\mathcal{C}$ consta de:
		\begin{enumerate}
			\item Una clase (o colección) de objetos, denotada por $\mathrm{Ob}(\mathcal{C})$. Si $A \in \mathrm{Ob}(\mathcal{C})$, se escribe simplemente $A \in \mathcal{C}$.
			
			\item Para cualesquiera objetos $A, B \in \mathrm{Ob}(\mathcal{C})$, un conjunto de morfismos de $A$ a $B$, denotado por
			$$
			\mathrm{Hom}_{\mathcal{C}}(A, B) := \{ f : A \to B \}.
			$$
			Los elementos de este conjunto también se denominan flechas.
			
			\item Una ley de composición de morfismos: para todos $A, B, C \in \mathrm{Ob}(\mathcal{C})$, se define una aplicación
		\begin{align*}
			\mathrm{Hom}_{\mathcal{C}}(A, B) \times \mathrm{Hom}_{\mathcal{C}}(B, C) &\longrightarrow \mathrm{Hom}_{\mathcal{C}}(A, C), \\
			(f, g) &\longmapsto g \circ f.
		\end{align*}
			Usualmente se escribe $gf = g \circ f$. Cuando no hay riesgo de ambigüedad, también se puede escribir $f \in \mathcal{C}$. Además, como la composición es una función, el morfismo compuesto $g \circ f$ es único.
			
			\item Para cada objeto $A \in \mathrm{Ob}(\mathcal{C})$, existe un morfismo identidad $id_A \in \mathrm{Hom}_{\mathcal{C}}(A, A)$ tal que, para todo objeto $B \in \mathrm{Ob}(\mathcal{C})$, se cumple:
			$$
			f \circ id_A = f,\quad \forall f \in \mathrm{Hom}_{\mathcal{C}}(A, B),
			$$
			$$
			id_A \circ g = g,\quad \forall g \in \mathrm{Hom}_{\mathcal{C}}(B, A).
			$$  \label{4}
			
			\item La composición de morfismos es asociativa: dados
			$$
			f \in \mathrm{Hom}_{\mathcal{C}}(A, B), \quad g \in \mathrm{Hom}_{\mathcal{C}}(B, C), \quad h \in \mathrm{Hom}_{\mathcal{C}}(C, D),
			$$
			se cumple
			$$
			h \circ (g \circ f) = (h \circ g) \circ f.
			$$
		\end{enumerate}
	}
\end{definicion}
\begin{observacion}
	\upshape{
	En ítem  \ref{4}, el morfísmo identidad es único.
	}
\end{observacion}

\begin{proof}[Justificación]
	Supongamos que $u, v \in \mathrm{Hom}_{\mathcal{C}}(A, A)$ son dos morfismos identidad. Entonces, para todo $f : A \to B$, se cumple $f \circ u = f$ y $f \circ v = f$; y para todo $g : B \to A$, se cumple $u \circ g = g$ y $v \circ g = g$. Tomando $f = v$, se tiene $v = v \circ u$; y tomando $g = u$, se tiene $u = v \circ u$. Por lo tanto, $u = v$.
\end{proof}

\begin{proposicion}
	\upshape{
	Sea $f : A \to B$ un morfismo en una categoría. Supongamos:
	\begin{itemize}
		\item Existe un morfismo $g : B \to A$ tal que $g \circ f = \mathrm{id}_A$ \ (inverso por la izquierda),
		\item Existe un morfismo $h : B \to A$ tal que $f \circ h = \mathrm{id}_B$ \ (inverso por la derecha).
	\end{itemize}
	Entonces $g = h$, y  este morfismo se denomina el\textit{ morfismo inverso} de $f$.}
\end{proposicion}

\begin{proof}
    Sabemos que $g = g \circ \mathrm{id}_B$. Como $f \circ h = \mathrm{id}_B$, sustituimos:
	$$
	g = g \circ (f \circ h) = (g \circ f) \circ h.
	$$
	Como $g \circ f = \mathrm{id}_A$, se obtiene:
	$$
	g = \mathrm{id}_A \circ h = h.
	$$
	Por lo tanto, $g = h$, y $f$ tiene un inverso único.
\end{proof}

\begin{definicion}
	\upshape{
		Un morfismo $f: A \to B$
		se llama \textit{isomorfismo} si existe un morfismo $ g: B \to A $
		tal que se cumple $$ g \circ f = id_A, \quad f \circ g = id_B.$$ En este caso, los objetos $A$ y $B$ se denominan \textit{isomorfos}.
	}
\end{definicion}


\begin{ejemplo} 
	\upshape{
	Sea $\mathbb{K}$ un cuerpo. A continuación se muestran algunos ejemplos de categorías, con sus respectivos objetos, morfismos e isomorfismos:\\
	
	\begin{table}[H]
		\centering
		{\large
			\renewcommand{\arraystretch}{1.6}
			\resizebox{160mm}{!}{
				\begin{tabular}{
						>{\centering\arraybackslash}p{4cm}
						>{\raggedright\arraybackslash}p{5.5cm}
						>{\raggedright\arraybackslash}p{5.5cm}
						>{\raggedright\arraybackslash}p{5.5cm}
					}
					\toprule
					\multicolumn{1}{c}{$\mathcal{C}$} & 
					\multicolumn{1}{c}{$\mathrm{Ob}(\mathcal{C})$} & 
					\multicolumn{1}{c}{Morfismos} & 
					\multicolumn{1}{c}{Isomorfismos} \\
					\midrule
					$\mathbf{Vect}_\mathbb{K}$ & Espacios vectoriales sobre $\mathbb{K}$ & Aplicaciones lineales & Isomorfismos lineales (biyecciones lineales) \\
					$\mathbf{Top}$ & Espacios topológicos & Funciones continuas & Homeomorfismos \\
					$\mathbf{Mag}$ & Magmas (conjunto con operación binaria) & Homomorfismos de magmas & Isomorfismos de magmas (homomorfismos biyectivos) \\
					$\mathbf{Set}$ & Conjuntos & Funciones & Biyeciones \\
					\bottomrule
				\end{tabular}
		}}
		\caption{}
	\end{table}
	}
\end{ejemplo}

\begin{ejemplo}
	\upshape{
		Categoría $\mathbf{Set}$ (conjuntos):
		
		\begin{enumerate}
			\item \textbf{Objetos:} Son conjuntos, por ejemplo, $A = \{1, 2, 3\}$, $B = \{a, b\}$ y $C = \{\alpha, \beta\}$, por lo que $A, B, C \in \mathbf{Set}$.
			
			\item \textbf{Morfismos:} Para $A, B \in \mathbf{Set}$, los morfismos son funciones de $A$ a $B$, es decir,
			\[
			\mathrm{Hom}_{\mathbf{Set}}(A, B) = \{ f : A \to B \mid f \text{ es función} \}.
			\]
			Por ejemplo, si $A = \{1,2\}$ y $B = \{a,b,c\}$, una función puede ser $f(1) = a$, $f(2) = c$.
			
			\item \textbf{Composición:} Dados $f : A \to B$ y $g : B \to C$, con $C = \{\alpha, \beta\}$, la composición $g \circ f : A \to C$ está definida por $(g \circ f)(x) = g(f(x))$ para todo $x \in A$. Por ejemplo, si $f(1) = a$, $f(2) = b$ y $g(a) = \alpha$, $g(b) = \beta$, entonces $(g \circ f)(1) = \alpha$, $(g \circ f)(2) = \beta$.
			
			\item \textbf{Identidad:} Para cada conjunto $A$ existe la función identidad $\mathrm{id}_A : A \to A$ dada por $\mathrm{id}_A(x) = x$ para todo $x \in A$, que cumple
			\[
			f \circ \mathrm{id}_A = f, \quad \forall f : A \to B, \quad \text{y} \quad \mathrm{id}_A \circ g = g, \quad \forall g : B \to A.
			\]
			
			\item \textbf{Asociatividad:} Para $f : A \to B$, $g : B \to C$ y $h : C \to D$ se cumple
			\[
			h \circ (g \circ f) = (h \circ g) \circ f,
			\]
			porque ambas funciones asignan a cada $x \in A$ el valor $h(g(f(x)))$.
		\end{enumerate}
	}
\end{ejemplo}


\begin{definicion}
	\upshape{
	Sea $(M, *)$ y $(N, *')$ dos magmas. Un \emph{homomorfismo de magmas} es una función $ f : M \to N $
	tal que, para todos $x, y \in M$, se cumple
	$$
	f(x * y) = f(x) *' f(y).
	$$
	}
\end{definicion}



\begin{ejemplo}
	\upshape{
		Categoría $\mathbf{Mag}$ (magmas):
		
		\begin{enumerate}
			\item \textbf{Objetos:} Son magmas, es decir, conjuntos no vacíos con una operación binaria. Por ejemplo,
			$$ A = (\mathbb{Z}, +), \quad B = (\mathbb{N}, \cdot), \quad C = (\{0,1\}, \max).$$		
			\item \textbf{Morfismos:} Para magmas $A, B \in \mathbf{Mag}$, los morfismos son homomorfismos de magmas, es decir, funciones
			$$
			f : A \to B
			$$
			tales que para todos $x,y \in A$ se cumple
			$$
			f(x * y) = f(x) *' f(y),
			$$
			donde $*$ es la operación en $A$ y $*'$ la operación en $B$.
			
			\item \textbf{Composición:} Dados homomorfismos
			$$
			f : A \to B, \quad g : B \to C,
			$$
			la composición usual de funciones
			$$
			(g \circ f)(x) = g(f(x))
			$$
			es también un homomorfismo, pues
			\begin{align*}
				(g \circ f)(x * y) &= g(f(x * y)) = g(f(x) *' f(y)) \\
				&= g(f(x)) *'' g(f(y)) = (g \circ f)(x) *'' (g \circ f)(y),
			\end{align*}
			donde $*''$ es la operación en $C$.
			
			\item \textbf{Identidad:} Para cada magma $A$, la función identidad
			$$
			\mathrm{id}_A : A \to A, \quad \mathrm{id}_A(x) = x,
			$$
			es un homomorfismo, ya que
			$$
			\mathrm{id}_A(x * y) = x * y = \mathrm{id}_A(x) * \mathrm{id}_A(y).
			$$
			
			\item \textbf{Asociatividad:} La composición de funciones es asociativa, esto es,
			$$
			h \circ (g \circ f) = (h \circ g) \circ f,
			$$
			para homomorfismos $f, g, h$.
		\end{enumerate}
	}
\end{ejemplo}




\begin{ejemplo}
	\upshape{
		Categoría $\mathbf{Vect}_\mathbb{K}$ (espacios vectoriales sobre $\mathbb{K}$):
		
		\begin{enumerate}
			\item \textbf{Objetos:} Son espacios vectoriales sobre $\mathbb{K}$. Por ejemplo,
			$$
			A = \mathbb{K}^2, \quad B = \mathbb{K}^3, \quad C = \{ \mathbf{0} \} \text{ (espacio trivial)}.
			$$
			
			\item \textbf{Morfismos:} Para $A, B \in \mathbf{Vect}_\mathbb{K}$, los morfismos son transformaciones lineales,
			$$
			f : A \to B,
			$$
			tales que para todos $x,y \in A$ y $\lambda \in \mathbb{K}$ se cumple
			$$
			f(x + y) = f(x) + f(y), \quad f(\lambda x) = \lambda f(x).
			$$
			
			\item \textbf{Composición:} Dados
			$$
			f : A \to B, \quad g : B \to C,
			$$
			la composición
			$$
			(g \circ f)(x) = g(f(x))
			$$
			es lineal porque, para $x,y \in A$ y $\lambda \in \mathbb{K}$,
			\begin{itemize}
				\item Para $x, y \in A$,
				\begin{align*}
					(g \circ f)(x + y) &= g(f(x + y)) = g(f(x) + f(y)) \\[8pt]
					&= g(f(x)) + g(f(y)) \\[8pt]
					&= (g \circ f)(x) + (g \circ f)(y).
				\end{align*}
				\item Para $\lambda \in \mathbb{K}$ y $x \in A$,
				\begin{align*}
					(g \circ f)(\lambda x) &= g(f(\lambda x)) = g(\lambda f(x)) = \lambda g(f(x)) = \lambda (g \circ f)(x).
				\end{align*}
			\end{itemize}
			
			
			\item \textbf{Identidad:} Para cada espacio vectorial $A$, la transformación identidad
			$$
			\mathrm{id}_A : A \to A, \quad \mathrm{id}_A(x) = x,
			$$
			es lineal y cumple
			\begin{align*}
				\mathrm{id}_A(x + y) &= x + y = \mathrm{id}_A(x) + \mathrm{id}_A(y), \\[8pt]
				\mathrm{id}_A(\lambda x) &= \lambda x = \lambda \,\mathrm{id}_A(x).
			\end{align*}
			
			
			\item \textbf{Asociatividad:} La composición de funciones es asociativa, es decir,
			$$
			h \circ (g \circ f) = (h \circ g) \circ f,
			$$
			para transformaciones lineales $f, g, h$.
		\end{enumerate}
	}
\end{ejemplo}



\begin{ejemplo}
	\upshape{
		Ejemplo de la categoría $\mathbf{Top}$ (espacios topológicos):
		
		\begin{itemize}
			\item \textbf{Objetos:} Son espacios topológicos, es decir, pares $(X, \tau)$ donde $X$ es un conjunto y $\tau$ es una topología sobre $X$.
			
			\item \textbf{Morfismos:} Para dos objetos $(X, \tau)$ y $(Y, \sigma)$, los morfismos son las funciones continuas $f : X \to Y$.
			
			\item \textbf{Composición:} La composición de funciones continuas es continua. Si $f : X \to Y$ y $g : Y \to Z$ son continuas, entonces $g \circ f : X \to Z$ también es continua.
			
			\item \textbf{Identidad:} Para cada espacio topológico $(X, \tau)$, la función identidad $\mathrm{id}_X : X \to X$ es continua, ya que la preimagen de cualquier abierto es el mismo abierto.
			
			\item \textbf{Asociatividad:} La composición de funciones es asociativa:
			$$
			h \circ (g \circ f) = (h \circ g) \circ f.
			$$
		\end{itemize}
	}
\end{ejemplo}


\begin{ejemplo}
	\upshape{
		Ejemplo de la categoría $\mathbf{Vect}_{\mathbb{K}}^{\|\cdot\|}$ (espacios vectoriales normados):
		
		\begin{itemize}
			\item \textbf{Objetos:} Son espacios vectoriales sobre $\mathbb{K}$ equipados con una norma. Es decir, pares 
			$$
			(V, \|\cdot\|),
			$$
			donde $V$ es un espacio vectorial sobre $\mathbb{K}$ y $\|\cdot\|\colon V \to \mathbb{R}_{\ge 0}$ es una norma.
			
			\item \textbf{Morfismos:} Para dos objetos $(V, \|\cdot\|_V)$ y $(W, \|\cdot\|_W)$, los morfismos son las transformaciones lineales continuas
			$$
			f \colon V \to W
			$$
			(es decir, $f$ lineal y existe $C>0$ tal que $\|f(v)\|_W \le C\,\|v\|_V$ para todo $v\in V$).
			
			\item \textbf{Composición:} La composición de dos transformaciones lineales continuas es nuevamente lineal y continua. Si
			$$
			f \colon V \to W, \quad g \colon W \to X
			$$
			son continuas, entonces
			$$
			g \circ f \colon V \to X
			$$
			es lineal y continua.
			
			\item \textbf{Identidad:} Para cada objeto $(V, \|\cdot\|)$, la función identidad
			$$
			\mathrm{id}_V \colon V \to V
			$$
			es lineal y continua, ya que $\|\mathrm{id}_V(v)\| = \|v\|$ para todo $v \in V$.
			
			\item \textbf{Asociatividad:} La composición de funciones es asociativa, es decir, si
			$$
			f \colon U \to V,\quad g \colon V \to W,\quad h \colon W \to X,
			$$
			entonces
			$$
			h \circ (g \circ f) \;=\; (h \circ g) \circ f.
			$$
		\end{itemize}
	}
\end{ejemplo}




	\chapter{Functores y Transformaciones naturales}
	
%\xymatrix{
%	A \ar[r]^f \ar[dr]_g & B \ar[d]^h \\
%	& C 
%} $\longrightarrow$  \xymatrix{
%A \ar[r]^f \ar[dr]_g & B \ar[d]^h \\
%& C 
%}



%$$
%\xymatrix{
%	A \ar[rr]^f \ar[ddrr]_{f\circ g}& & B \ar[dd]^g &          &  & F(A) \ar[rr]^{F(f)} \ar[ddrr]_{F(g\circ f)}& & F(B) \ar[dd]^{F(g)} \\
%	                              & &             & \ar[rr]^F &  &                                            & &\\
%                                  & & C           &          &  &                                            & &  C'
%} $$       \ar@(dr,dl)[]



\section{Functores}

\begin{definicion}\label{ppoiujhgf}
	\upshape{
		Dadas dos categorías $\mathcal{C}$ y $\mathcal{D}$, un \textit{functor} de $\mathcal{C}$ a $\mathcal{D}$ es una asignación $F : \mathcal{C} \to \mathcal{D}$ que consiste en:
		
		\begin{enumerate}
			\item Una función entre las clases de objetos:
			\begin{align*}
				\mathrm{Ob}(\mathcal{C}) &\longrightarrow \mathrm{Ob}(\mathcal{D}) \\
				A &\longmapsto F(A)
			\end{align*}
			
			\item Para cada par de objetos $A, B \in \mathrm{Ob}(\mathcal{C})$, una función entre las clases de morfismos:
			\begin{align*}
				\mathrm{Hom}_{\mathcal{C}}(A, B) &\longrightarrow \mathrm{Hom}_{\mathcal{D}}\left(F(A), F(B)\right) \\
				f &\longmapsto F(f)
			\end{align*}
			que respeta:
			\begin{itemize}
				\item La identidad: $F(\mathrm{id}_A) = \mathrm{id}_{F(A)}$ para todo $A \in \mathrm{Ob}(\mathcal{C})$.
				
				\item La composición: Si $f \in \mathrm{Hom}_{\mathcal{C}}(A, B)$ y $g \in \mathrm{Hom}_{\mathcal{C}}(B, C)$, entonces
				$$
				F(g \circ f) = F(g) \circ F(f).
				$$
				Esto significa que a un diagrama conmutativo, el functor $F$ le asigna otro diagrama también conmutativo.
				
				$$
					{\Large
						\xymatrix{
								A \ar[rr]^f \ar[ddrr]_{f\circ g}& & B \ar[dd]^g &          &  & F(A) \ar[rr]^{F(f)} \ar[ddrr]_{F(g\circ f)}& & F(B) \ar[dd]^{F(g)} \\
								                              & &             & \ar[rr]^F &  &                                            & &\\
							                                & & C           &          &  &                                            & &  F(C)
							}
					} 
				$$
				
			\end{itemize}
		\end{enumerate} 
	}
\end{definicion}

\begin{observacion}
	\upshape{
		Las categorías se relacionan entre sí mediante functores.
		
		\begin{center}
			\begin{tikzpicture}[scale=0.7]
				% Categoría izquierda
				\begin{scope}[shift={(-6,0)}]
					\draw[thick] (0,0) ellipse (3cm and 1.5cm);
					\draw[thick] (-3,0) -- (3,0);
					\node at (0,2) {{\small CATEGORÍA}};
					\node at (0,0.75) {Objetos};
					\node at (0,-0.75) {Morfismos};
				\end{scope}
				
				% Categoría derecha
				\begin{scope}[shift={(6,0)}]
					\draw[thick] (0,0) ellipse (3cm and 1.5cm);
					\draw[thick] (-3,0) -- (3,0);
					\node at (0,2) {{\small CATEGORÍA}};
					\node at (0,0.75) {Objetos};
					\node at (0,-0.75) {Morfismos};
				\end{scope}
				
				% Flecha
				\draw[->, thick] ( -1.5, 0 ) -- ( 1.5, 0 ) node[midway, above] {\(F\)};
			\end{tikzpicture}
		\end{center}
		
		\noindent A todo objeto $A$ de la primera categoría se le hace corresponder un único objeto $F(A)$ de la segunda categoría, y a cada morfismo $f$ se le hace corresponder un único morfismo $F(f)$, preservando identidad y composición.
	}
\end{observacion}



\begin{definicion}
	\upshape{
		La \textit{categoría opuesta} de una categoría $\mathcal{C}$ se denota por $\mathcal{C}^{\mathrm{op}}$ y se define de la siguiente manera:
		
		\begin{enumerate}
			\item Los objetos de $\mathcal{C}^{\mathrm{op}}$ son los mismos que los de $\mathcal{C}$:
			$$
			\mathrm{Ob}(\mathcal{C}^{\mathrm{op}}) := \mathrm{Ob}(\mathcal{C}).
			$$
			
			\item Dados dos objetos $A$ y $B$, el conjunto de morfismos de $A$ a $B$ en $\mathcal{C}^{\mathrm{op}}$ se define como
			$$
			\mathrm{Hom}_{\mathcal{C}^{\mathrm{op}}}(A, B) := \mathrm{Hom}_{\mathcal{C}}(B, A).
			$$
			Es decir, un morfismo $f \in \mathrm{Hom}_{\mathcal{C}}(B, A)$ se considera en $\mathcal{C}^{\mathrm{op}}$ como un morfismo de $A$ a $B$. Aunque se trata del mismo símbolo $f$, se interpreta con dirección opuesta.
			
			\item La composición de morfismos en $\mathcal{C}^{\mathrm{op}}$ se define por
			$$
			f^{\mathrm{op}} \circ g^{\mathrm{op}} := (g \circ f)^{\mathrm{op}},
			$$
			donde $f \in \mathrm{Hom}_{\mathcal{C}}(B, C)$ y $g \in \mathrm{Hom}_{\mathcal{C}}(A, B)$. Esto garantiza que la composición en $\mathcal{C}^{\mathrm{op}}$ corresponde a la composición en $\mathcal{C}$ pero en orden inverso.
			
			\item Para cada objeto $A$, el morfismo identidad en $\mathcal{C}^{\mathrm{op}}$ es el mismo que en $\mathcal{C}$:
			$$
			\mathrm{id}_A^{\mathrm{op}} := \mathrm{id}_A.
			$$
			
			\item La composición en $\mathcal{C}^{\mathrm{op}}$ es asociativa, ya que la composición en $\mathcal{C}$ lo es.
		\end{enumerate}
	}
\end{definicion}






\begin{observacion}
	\upshape{
		Si $F: \mathcal{C}^{\mathrm{op}} \to \mathcal{D}$ es un functor, entonces para todo par de  morfismos $f: A \to B$ y $g: B \to C$ en $\mathcal{C}$, se tiene
		$$
		F(g \circ f) = F(f) \circ F(g).
		$$
		Es decir, $F$ invierte el orden de la composición.
	}
\end{observacion}

\begin{proof}[Justificación]
	Dado que en $\mathcal{C}^{\mathrm{op}}$ el morfismo \( f: A \to B \) en $\mathcal{C}$ se interpreta como \( f: B \to A \), y análogamente \( g: B \to C \) como \( g: C \to B \), en $\mathcal{C}^{\mathrm{op}}$ se tiene
	$$
	f: B \to A, \quad g: C \to B, \quad \text{de modo que } f \circ g: C \to A,
	$$
	lo cual, por definición de composición en $\mathcal{C}^{\mathrm{op}}$, corresponde a \( g \circ f \) en $\mathcal{C}$. Entonces, como \(F\) es un functor, se aplica
	$$ F(f \circ g) = F(f) \circ F(g), $$
	y al identificar \( f \circ g \) en $\mathcal{C}^{\mathrm{op}}$ con \( g \circ f \) en $\mathcal{C}$, se concluye
	$$F(g \circ f) = F(f) \circ F(g). $$
	
	\noindent Esto se ilustra en el siguiente diagrama conmutativo:
	
	$$
	{\Large
		\xymatrix{
			A \ar[rr]^f \ar[ddrr]_{g\circ f}& & B \ar[dd]^g &          &  & F(A)  & & F(B) \ar[ll]_{F(f)}  \\
			& &             & \ar[rr]^F &  &                                            & &\\
			& & C           &          &  &                                            & &  F(C)  \ar[uull]^{F(g\circ f)} \ar[uu]_{F(g)}
		}
	} 
	$$
	
\end{proof}

\begin{definicion}
	\upshape{
		Sean $\mathcal{C}$ y $\mathcal{D}$ dos categorías. Según el dominio de definición, distinguimos dos tipos de functores:
		\begin{enumerate}
			\item Un \textit{functor covariante} es un functor $F: \mathcal{C} \to \mathcal{D}$.
			\item Un \textit{functor contravariante} es un functor $F: \mathcal{C}^{\mathrm{op}} \to \mathcal{D}$, el cual se suele interpretar como una asignación
			$$
			F: \mathcal{C} \to \mathcal{D}
			$$
			que invierte la dirección de los morfismos.
		\end{enumerate}
	}
\end{definicion}





\begin{ejemplo}[Functor dual]\label{ejemplofunctor dual}
	\upshape{
		Sea $\mathbf{Vect}_{\mathbb{K}}$ la categoría de espacios vectoriales sobre un cuerpo $\mathbb{K}$.
		
		\begin{itemize}
			\item Para todo espacio vectorial $V$, se define su \textit{espacio dual} como $V^* := \mathrm{Hom}_{\mathbb{K}}(V, \mathbb{K}),$
			el conjunto de transformaciones lineales de $V$ en $\mathbb{K}$.
			
			\item Si $f: V \to W$ es una transformación lineal, se define su \textit{transformación dual} como $
			f^*: W^* \to V^*, \quad \phi \mapsto \phi \circ f,$ que es lineal.
		\end{itemize}
		Estas asignaciones definen el functor dual
		
		\begin{align*}
			F : \mathbf{Vect}_{\mathbb{K}}^{\mathrm{op}}& \longrightarrow \mathbf{Vect}_{\mathbb{K}}\\[5pt]
			V&\longmapsto V^\ast\\
			\left(f: V\to W \text{ en } \mathbf{Vect}_{\mathbb{K}}\right)&\longmapsto \left(f^\ast: W^\ast \to V^\ast\right).
		\end{align*}
		Es decir, a cada espacio $V$ le asigna su dual $V^*$, y a cada morfismo $f$ su dual $f^*$.\\
		
		
		\noindent Para verificar que $F$ define un functor $F: \textbf{Vect}^{\mathrm{op}} \to \textbf{Vect}$, seguimos los pasos de la definición:
		
		\begin{enumerate}
			\item Asignación sobre objetos:
			A cada objeto $V \in \mathrm{Ob}(\textbf{Vect}^{\mathrm{op}}) = \mathrm{Ob}(\textbf{Vect})$, el functor $F$ le asigna su dual $V^* = \mathrm{Hom}_{\textbf{Vect}}(V, \mathbb{K})$, que también es un espacio vectorial sobre $\mathbb{K}$. Por tanto, $F(V) \in \mathrm{Ob}(\textbf{Vect})$.
			
			\item Asignación sobre morfismos:
			A cada morfismo $f \in \mathrm{Hom}_{\textbf{Vect}^{\mathrm{op}}}(V, W)$, esto significa, $f: W \to V$ en $\textbf{Vect}$, el functor $F$ le asigna
			$$
			F(f) = f^* : V^* \to W^*, \quad \phi \mapsto \phi \circ f,
			$$
			que es un morfismo lineal entre espacios vectoriales.
			
			\item Preserva la identidad: Para todo espacio vectorial $V$,
			$$
			F(\mathrm{id}_V) = (\mathrm{id}_V)^* : V^* \to V^*, \quad \phi \mapsto \phi \circ \mathrm{id}_V = \phi,
			$$
			que es la identidad en $V^*$.
			
			\item Preserva la composición: Dados $f: V \to W$ y $g: W \to X$,
			$$
			F(g \circ f) = (g \circ f)^* : X^* \to V^*, \quad \phi \mapsto \phi \circ (g \circ f),
			$$
			mientras que
			$$
			F(f) \circ F(g) = f^* \circ g^* : X^* \to V^*, \quad \phi \mapsto f^*(g^*(\phi)) = f^*(\phi \circ g) = (\phi \circ g) \circ f = \phi \circ (g \circ f).
			$$
			Por lo tanto, $F(g \circ f) = F(f) \circ F(g)$.
		\end{enumerate}
		
		\noindent
		Con esto se concluye que $F$ es un functor de $\textbf{Vect}$ a sí misma.
		
	}
\end{ejemplo}

\begin{observacion}
	\upshape{
	En el Ejemplo \ref{ejemplofunctor dual} el functor dual es un functor contravariante.
	}
\end{observacion}
\begin{observacion}
	\upshape{
		Al definir un functor \( F \) desde \( \textbf{Vect}^{\mathrm{op}} \) en lugar de desde \( \textbf{Vect} \), se garantiza que la dirección de los morfismos sea compatible con la construcción dual. En efecto, una aplicación lineal \( f: V \to W \) induce naturalmente un morfismo \( f^*: W^* \to V^* \), que va en la dirección opuesta. Para que \( F \) sea un functor, se requiere que asigne a cada \( f \in \mathrm{Hom}_{\textbf{Vect}^{\mathrm{op}}}(V,W) \) un morfismo \( F(f) \in \mathrm{Hom}(F(V), F(W)) \), es decir, un morfismo de \( V^* \) en \( W^* \). Esto se logra al tomar como categoría de partida a \( \textbf{Vect}^{\mathrm{op}} \), de modo que el morfismo \( f: V \to W \) en \( \textbf{Vect} \) se interprete como un morfismo \( f: W \to V \) en \( \textbf{Vect}^{\mathrm{op}} \), y así \( f^*: V^* \to W^* \) encaje en la dirección correcta.
	}
\end{observacion}
  



\begin{ejemplo}[Functor bidual]\label{functorbidual}
	\upshape{
		Sea $\mathbf{Vect}_{\mathbb{K}}$ la categoría de espacios vectoriales sobre un cuerpo $\mathbb{K}$.
		
		\begin{itemize}
			\item Para cada espacio vectorial $V$, se define su \textit{bidual} como $V^{**} := (V^*)^*$, es decir, el dual del espacio dual de $V$.
			
			\item Dada una transformación lineal $f: V \to W$, su transformación bidual es
			$$
			f^{**} : V^{**} \to W^{**}, \quad \Psi \mapsto f^{**}(\Psi),
			$$
			donde $f^{**}(\Psi)$ es la aplicación lineal en $W^*$ dada por
			$$
			f^{**}(\Psi)(\phi) :=\left(\Psi\circ f^\ast\right)(\phi)= \Psi(\phi \circ f), \quad \text{para toda } \phi \in W^*.
			$$
			Esto es, $f^{**}$ se define como la dual de la aplicación dual \( f^* : W^* \to V^* \).
		\end{itemize}
		
		Estas asignaciones definen un functor
		\begin{align*}
			F : \mathbf{Vect}_{\mathbb{K}} &\longrightarrow \mathbf{Vect}_{\mathbb{K}}\\[5pt]
			V &\longmapsto V^{**}\\
			\left(f: V \to W\right) &\longmapsto \left(f^{**} : V^{**} \to W^{**}\right).
		\end{align*}
		
		\noindent Verifiquemos que $F$ define un functor:
		
		\begin{enumerate}
			\item Asignación sobre objetos: A cada espacio $V$, el functor $F$ le asigna su bidual $V^{**}$.
			
			\item Asignación sobre morfismos: A cada morfismo $f: V \to W$, se le asigna $f^{**} : V^{**} \to W^{**}$ como se definió arriba.
			
			\item Preserva la identidad: Para todo $V$,
			$$
			F(\mathrm{id}_V) = (\mathrm{id}_V)^{**} = \mathrm{id}_{V^{**}}.
			$$
			
			\item Preserva la composición: Dados $f: V \to W$ y $g: W \to X$, se tiene
			$$
			F(g \circ f) = (g \circ f)^{**} = g^{**} \circ f^{**} = F(g) \circ F(f).
			$$
		\end{enumerate}
		
		\noindent
		Por lo tanto, $F$ es un functor de $\mathbf{Vect}_{\mathbb{K}}$.
	}
\end{ejemplo}

\begin{observacion}
	\upshape{
		En el Ejemplo \ref{functorbidual} el functor dual es un functor covariante.
	}
\end{observacion}


\begin{ejemplo}[Functor olvidadizo]
	\upshape{
		En teoría de categorías, un \textit{functor olvidadizo} es un functor que “olvida” parte de la estructura o propiedades de los objetos y morfismos de una categoría para considerarlos en una categoría más simple. Por ejemplo, al pasar de espacios vectoriales a conjuntos, se olvida la estructura vectorial y se conserva únicamente la información del conjunto subyacente.\\
		
		\noindent Sea $\mathbf{Vect}_{\mathbb{K}}$ la categoría de espacios vectoriales sobre un cuerpo $\mathbb{K}$ y $\mathbf{Set}$ la categoría de conjuntos.
		
		\begin{itemize}
			\item A cada espacio vectorial $V$, el functor olvidadizo $F$ le asigna el conjunto subyacente de vectores de $V$.
			
			\item A cada transformación lineal $f: V \to W$, el functor le asigna la misma función $f$ vista como función entre conjuntos.
		\end{itemize} Así, el functor olvidadizo se define como
		\begin{align*}
			F : \mathbf{Vect}_{\mathbb{K}} &\longrightarrow \mathbf{Set}\\[5pt]
			V &\longmapsto F(V)\\
			(f: V \to W) &\longmapsto F(f).
		\end{align*}
		
		\noindent	Verificación de que $F$ es un functor:
		\begin{enumerate}
			\item Preserva la identidad: Para todo espacio vectorial $V$,
			$$
			F(\mathrm{id}_V) = \mathrm{id}_{F(V)}.
			$$
			
			\item Preserva la composición: Dados morfismos $f: V \to W$ y $g: W \to X$,
			$$
			F(g \circ f) = g \circ f = F(g) \circ F(f).
			$$
		\end{enumerate}
		
		Por lo tanto, $F$ es un functor de $\mathbf{Vect}_{\mathbb{K}}$ a $\mathbf{Set}$, llamado \textit{functor olvidadizo}.
	}
\end{ejemplo}



\begin{ejemplo}[Otros functores olvidadizos]
	\upshape{
		A continuación se muestran varios ejemplos de functores olvidadizos, indicando explícitamente su acción en objetos y morfismos mediante \texttt{align*}:
		
		\begin{itemize}
			\item De $\mathbf{Top}$ a $\mathbf{Set}$: \\
			El functor olvidadizo que asigna a cada espacio topológico su conjunto subyacente y a cada función continua la misma función vista como mapeo de conjuntos.
			\begin{align*}
				F : \mathbf{Top} &\longrightarrow \mathbf{Set}\\[5pt]
				(X, \tau) &\longmapsto X\\
				(f: X \to Y) &\longmapsto f: X \to Y.
			\end{align*}
			
			\item De $\mathbf{Vect}_{\mathbb{K}}^{\|\cdot\|}$ a $\mathbf{Top}$: \\
			El functor olvidadizo que a cada espacio vectorial normado le asigna el espacio topológico subyacente (con la topología inducida por la norma), y a cada transformación lineal continua la misma función.
			\begin{align*}
				F : \mathbf{Vect}_{\mathbb{K}}^{\|\cdot\|} &\longrightarrow \mathbf{Top}\\[5pt]
				(V, \|\cdot\|) &\longmapsto (V, \text{topología inducida por } \|\cdot\|)\\
				(f: V \to W,\ f\ \text{continua}) &\longmapsto f: V \to W.
			\end{align*}
			
			\item De $\mathbf{Vect}_{\mathbb{K}}^{\|\cdot\|}$ a $\mathbf{Vect}_{\mathbb{K}}$: \\
			El functor olvidadizo que “olvida” la norma, asignando a cada espacio vectorial normado su espacio vectorial subyacente y a cada transformación lineal continua la misma transformación lineal en $\mathbf{Vect}_{\mathbb{K}}$.
			\begin{align*}
				F : \mathbf{Vect}_{\mathbb{K}}^{\|\cdot\|} &\longrightarrow \mathbf{Vect}_{\mathbb{K}}\\[5pt]
				(V, \|\cdot\|) &\longmapsto V\\
				(f: V \to W,\ f\ \text{continua}) &\longmapsto f: V \to W.
			\end{align*}
			
			\item De $\mathbf{Grp}$ a $\mathbf{Set}$: \\
			El functor olvidadizo que asigna a cada grupo su conjunto subyacente y a cada homomorfismo de grupos la misma función vista como función de conjuntos.
			\begin{align*}
				F : \mathbf{Grp} &\longrightarrow \mathbf{Set}\\[5pt]
				G &\longmapsto G\\
				(\varphi: G \to H) &\longmapsto \varphi: G \to H.
			\end{align*}
			
			\item De $\mathbf{Ring}$ a $\mathbf{Ab}$: \\
			El functor olvidadizo que asigna a cada anillo su grupo aditivo subyacente y a cada homomorfismo de anillos la misma función como homomorfismo de grupos abelianos.
			\begin{align*}
				F : \mathbf{Ring} &\longrightarrow \mathbf{Ab}\\[5pt]
				(R, +, \cdot) &\longmapsto (R, +)\\
				(\psi: R \to S) &\longmapsto \psi: (R, +) \to (S, +).
			\end{align*}
		\end{itemize}
	}
\end{ejemplo}


\begin{ejemplo}[Functor Hom covariante]\label{functorhom}
	\upshape{
		Sean $\mathcal{C}$ una categoría y $A \in \mathrm{Ob}(\mathcal{C})$ un objeto fijo. Consideremos las siguientes construcciones:
		
		\begin{itemize}
			\item Para cada objeto $X \in \mathcal{C}$, se asocia el conjunto $\mathrm{Hom}_{\mathcal{C}}(A, X)$.
			
			\item Para cada morfismo $f : X \to Y$ en $\mathcal{C}$, se define la función
			\[
			\mathrm{Hom}_{\mathcal{C}}(A, f) : \mathrm{Hom}_{\mathcal{C}}(A, X) \to \mathrm{Hom}_{\mathcal{C}}(A, Y), \quad h \mapsto f \circ h.
			\]
		\end{itemize}
		
		Estas asignaciones definen un functor
		\begin{align*}
			\mathrm{Hom}_{\mathcal{C}}(A, -) : \mathcal{C} &\longrightarrow \mathbf{Set}\\[5pt]
			X &\longmapsto \mathrm{Hom}_{\mathcal{C}}(A, X)\\
			(f : X \to Y) &\longmapsto \left(\mathrm{Hom}_{\mathcal{C}}(A, f) : \mathrm{Hom}_{\mathcal{C}}(A, X) \to \mathrm{Hom}_{\mathcal{C}}(A, Y)\right).
		\end{align*}
		
		\noindent Verifiquemos que esta asignación define un functor:
		
		\begin{enumerate}
			\item Asignación sobre objetos: A cada objeto $X$ de $\mathcal{C}$, se le asigna el conjunto $\mathrm{Hom}_{\mathcal{C}}(A, X)$.
			
			\item Asignación sobre morfismos: A cada morfismo $f : X \to Y$, se le asigna la función $\mathrm{Hom}_{\mathcal{C}}(A, f)$ dada por la composición a derecha con $f$.
			
			\item Preserva la identidad: Para todo objeto $X$, se tiene
			\[
			\mathrm{Hom}_{\mathcal{C}}(A, \mathrm{id}_X)(h) = \mathrm{id}_X \circ h = h,
			\]
			para todo $h \in \mathrm{Hom}_{\mathcal{C}}(A, X)$, es decir,
			\[
			\mathrm{Hom}_{\mathcal{C}}(A, \mathrm{id}_X) = \mathrm{id}_{\mathrm{Hom}_{\mathcal{C}}(A, X)}.
			\]
			
			\item Preserva la composición: Dado $f : X \to Y$ y $g : Y \to Z$ en $\mathcal{C}$, se tiene que para todo $h \in \mathrm{Hom}_{\mathcal{C}}(A, X)$,
			\[
			\mathrm{Hom}_{\mathcal{C}}(A, g \circ f)(h) = (g \circ f) \circ h = g \circ (f \circ h) = \mathrm{Hom}_{\mathcal{C}}(A, g)(\mathrm{Hom}_{\mathcal{C}}(A, f)(h)).
			\]
			Así,
			\[
			\mathrm{Hom}_{\mathcal{C}}(A, g \circ f) = \mathrm{Hom}_{\mathcal{C}}(A, g) \circ \mathrm{Hom}_{\mathcal{C}}(A, f).
			\]
		\end{enumerate}
		
		\noindent
		Por lo tanto, $\mathrm{Hom}_{\mathcal{C}}(A, -)$ define un functor covariante de $\mathcal{C}$ a $\mathbf{Set}$.
	}
\end{ejemplo}
\begin{observacion}[Functor Hom contravariante]
	\upshape{
		De manera análoga al Ejemplo \ref{functorhom}, se puede fijar un objeto \( A \in \mathrm{Ob}(\mathcal{C}) \) y considerar la asignación
		\begin{align*}
			\mathrm{Hom}_{\mathcal{C}}(-, A) : \mathcal{C}^{\mathrm{op}} &\longrightarrow \mathbf{Set}\\[5pt]
			X &\longmapsto \mathrm{Hom}_{\mathcal{C}}(X, A)\\
			(f : X \to Y) &\longmapsto \left(\mathrm{Hom}_{\mathcal{C}}(f, A) : \mathrm{Hom}_{\mathcal{C}}(Y, A) \to \mathrm{Hom}_{\mathcal{C}}(X, A)\right),
		\end{align*}
		donde \( \mathrm{Hom}_{\mathcal{C}}(f, A)(h) = h \circ f \). \\
		
		\noindent Esta asignación define un functor contravariante \( \mathrm{Hom}_{\mathcal{C}}(-, A) \) de \( \mathcal{C} \) a \( \mathbf{Set} \), ya que:
		\begin{itemize}
			\item Preserva identidades: \( \mathrm{Hom}_{\mathcal{C}}(\mathrm{id}_X, A)(h) = h \circ \mathrm{id}_X = h \), para todo \( h \in \mathrm{Hom}_{\mathcal{C}}(X, A) \).
			
			\item Invierte la composición: dados \( f : X \to Y \) y \( g : Y \to Z \), se tiene
			\[
			\mathrm{Hom}_{\mathcal{C}}(g \circ f, A)(h) = h \circ (g \circ f) = (h \circ g) \circ f = \mathrm{Hom}_{\mathcal{C}}(f, A)(\mathrm{Hom}_{\mathcal{C}}(g, A)(h)),
			\]
			por lo que
			\[
			\mathrm{Hom}_{\mathcal{C}}(g \circ f, A) = \mathrm{Hom}_{\mathcal{C}}(f, A) \circ \mathrm{Hom}_{\mathcal{C}}(g, A).
			\]
		\end{itemize}
		
		\noindent
		Por lo tanto, \( \mathrm{Hom}_{\mathcal{C}}(-, A) \) es un functor contravariante de \( \mathcal{C} \) a \( \mathbf{Set} \), o equivalentemente, un functor covariante \( \mathcal{C}^{\mathrm{op}} \to \mathbf{Set} \).
	}
\end{observacion}


\section{Transformaciones naturales}

\begin{definicion}
	\upshape{
		Sean $\mathcal{C}$ y $\mathcal{D}$ dos categorías y $F, G: \mathcal{C} \to \mathcal{D}$ dos functores. Una \emph{transformación natural}
		$$
		\Phi: F \Longrightarrow G
		$$
		es una familia de morfismos 
		$$
		\{\Phi_A: F(A)\to G(A)\}_{\,A\in\mathrm{Ob}(\mathcal{C})\!}
		$$
		en $\mathcal{D}$ que satisface la siguiente condición: para todo morfismo $f: A\to B$ en $\mathcal{C}$, el diagrama
		
		$$
		{\Large 
			\xymatrix{
				F(A) \ar[rr]^{\Phi_A} \ar[dd]_{F(f)} && G(A) \ar[dd]^{G(f)} \\ \\
				F(B) \ar[rr]_{\Phi_B} && G(B)
			}
		}
		$$
		
	\noindent	conmuta, es decir,
		$G(f)\,\circ\,\Phi_A \;=\; \Phi_B\,\circ\,F(f).$\\
		
		\noindent Cada $\Phi_A$ se denomina \textit{componente} de $\Phi$ en $A$.
	}
\end{definicion}


\begin{ejemplo}\label{cascas}
	\upshape{
		Sean $\mathcal{C}$ y $\mathcal{D}$ dos categorías, y sea $F : \mathcal{C} \longrightarrow \mathcal{D}$ un functor. Definimos una transformación natural 
		$$
		\mathrm{id}_F : F \Longrightarrow F
		$$
		mediante la familia de morfismos 
		$$
		\bigl\{\, (\mathrm{id}_F)_A := \mathrm{id}_{F(A)} : F(A) \to F(A) \,\bigr\}_{\,A \in \mathrm{Ob}(\mathcal{C})}.
		$$
		Para verificar que $\mathrm{id}_F$ es natural, tomemos un morfismo arbitrario
		$f : A \longrightarrow B
		$ en $\mathcal{C}$. Debemos comprobar que el siguiente diagrama conmuta:
		$$
		{\large \xymatrix{
				F(A) \ar[rrr]^{(\mathrm{id}_F)_A} \ar[dd]_{F(f)} &&& F(A) \ar[dd]^{F(f)} \\
				&&& \\
				F(B) \ar[rrr]_{(\mathrm{id}_F)_B} &&& F(B)
		}}
		$$
		En efecto, se tiene:
		$$
		F(f) \circ (\mathrm{id}_F)_A = F(f) \circ \mathrm{id}_{F(A)} = F(f),
		$$
		y
		$$
		(\mathrm{id}_F)_B \circ F(f) = \mathrm{id}_{F(B)} \circ F(f) = F(f),
		$$
		por lo tanto,
		$$
		F(f) \circ (\mathrm{id}_F)_A = (\mathrm{id}_F)_B \circ F(f).
		$$
		Esto muestra que 
		$	\mathrm{id}_F : F \Longrightarrow F$ es una transformación natural.
	}
\end{ejemplo}




\begin{ejemplo}
	\upshape{
		En la categoría $\mathbf{Vect}_{\mathbb{K}}$, consideremos:
		\begin{itemize}
			\item El functor identidad 
			$$
			\mathrm{id}: \mathbf{Vect}_{\mathbb{K}}\longrightarrow \mathbf{Vect}_{\mathbb{K}},\quad V\longmapsto V,
			$$
			\item  El functor bidual 
			$$
			(-)^{**}: \mathbf{Vect}_{\mathbb{K}}\longrightarrow \mathbf{Vect}_{\mathbb{K}},\quad V\longmapsto V^{**}.
			$$
		\end{itemize}
		Definimos una transformación natural 
		$$
		\Phi: \mathrm{id} \Longrightarrow (-)^{**}
		$$
		mediante la familia de morfismos 
		$$
		\{\Phi_V : V \to V^{**}\}_{\,V\in\mathrm{Ob}(\mathbf{Vect}_{\mathbb{K}})} 
		$$
		donde, para cada espacio vectorial $V$, definimos el morfismo
		$$
		\Phi_V(v) \;:=\; \bigl[\;\phi \mapsto \phi(v)\bigr]\quad\text{para todo }v\in V,\;\phi\in V^*.
		$$
		
		\noindent Para verificar que $\Phi$ es natural, tomemos una transformación lineal 
		$$
		f: V \longrightarrow W,
		$$
		y debemos comprobar que el siguiente diagrama conmuta:
		$$
		{\Large \xymatrix{
			V \ar[rrr]^{\Phi_V} \ar[dd]_{f} && & V^{**} \ar[dd]^{f^{**}} \\
			&&&\\
			W \ar[rrr]_{\Phi_W} & && W^{**}
		}}
		$$
		En efecto, para cada $v\in V$ y cada $\psi\in W^*$, se tiene
		$$
		(f^{**}\circ \Phi_V)(v) \;=\; f^{**}(\Phi_V(v))
		\;=\; f^{**}\bigl[\;\phi\mapsto \phi(v)\bigr]
		\;=\; \bigl[\;\psi\mapsto \Phi_V(v)(\psi\circ f)\bigr]
		\;=\; \bigl[\;\psi\mapsto \psi(f(v))\bigr],
		$$
		y
		$$
		(\Phi_W\circ f)(v) \;=\; \Phi_W(f(v))
		\;=\; \bigl[\;\psi\mapsto \psi(f(v))\bigr].
		$$
		Por lo tanto,
		$$
		f^{**}\circ \Phi_V \;=\; \Phi_W\circ f,
		$$
		y el diagrama conmuta para todo $f$. Esto muestra que 
		$$
		\Phi: \mathrm{id} \Longrightarrow (-)^{**}
		$$
		es una transformación natural en $\mathbf{Vect}_{\mathbb{K}}$.
	}
\end{ejemplo}


\begin{ejemplo}
	\upshape{
		En la categoría \(\mathbf{Vect}_{\mathbb{K}}\), consideremos los siguientes functores:
		\begin{itemize}
			\item El functor identidad:
			\[
			\mathrm{id}: \mathbf{Vect}_{\mathbb{K}} \longrightarrow \mathbf{Vect}_{\mathbb{K}}, \quad V \longmapsto V,
			\]
			\item El functor bidual, revisar el Ejemplo \ref{functorbidual}:
			\[
			(-)^{**}: \mathbf{Vect}_{\mathbb{K}} \longrightarrow \mathbf{Vect}_{\mathbb{K}}, \quad V \longmapsto V^{**}.
			\]
		\end{itemize}
		Recordemos que todo vector \(v \in V\) define de manera natural un funcional sobre \(V^*\), dado por:
		$\phi \longmapsto \phi(v), \quad \text{para todo } \phi \in V^*$.
		Esto induce un morfismo lineal:
	
		\begin{align*}
			\Phi_V: V &\longrightarrow V^{**}\\
			v &\longmapsto 
			\left(
			\begin{array}{rcl}
				\Phi_V(v)\colon & V^* & \to \mathbb{K} \\
				& \phi & \mapsto \phi(v)
			\end{array}
			\right)
		\end{align*}
		De esta forma, definimos una transformación natural:
		$$
		\Phi: \mathrm{id} \Longrightarrow (-)^{**},
		$$
		dada por la familia de morfismos
		$$
		\left\{\, \Phi_V: V \to V^{**} \,\right\}_{V \in \mathrm{Ob}(\mathbf{Vect}_{\mathbb{K}})}.
		$$		
		Para verificar que \(\Phi\) es natural, tomemos una transformación lineal $
		f: V \longrightarrow W$,
		y comprobemos que el siguiente diagrama conmuta:
		$$
		{\Large \xymatrix{
				V \ar[rrr]^{\Phi_V} \ar[dd]_{f} &&& V^{**} \ar[dd]^{f^{**}} \\
				\\
				W \ar[rrr]_{\Phi_W} &&& W^{**}
		}}
		$$
		En efecto, para cada $v \in V$ y cada $\psi \in W^*$, se tiene: $$(f^{**} \circ \Phi_V)(v)
		= f^{\ast\ast}\left(\Phi_V(v)\right)=\Phi_V(v)\circ f^\ast ,$$ de donde se tiene $$(f^{**} \circ \Phi_V)(v)
		(\psi)=\Phi_V(v) \left(f^{\ast}(\psi)\right)=\Phi_V(v) \left(\psi \circ f\right)=\left(\psi \circ f\right)(v)$$
		y por otro lado:
		$$
		(\Phi_W \circ f)(v)(\psi) = \Phi_W(f(v))(\psi) = \psi(f(v)).
		$$
		Por lo tanto,
		\[
		f^{**} \circ \Phi_V = \Phi_W \circ f,
		\]
		y el diagrama conmuta para todo \(f\).
	}
\end{ejemplo}

\begin{definicion}
	\upshape{
		Sean \(\mathcal{C}\) y \(\mathcal{D}\) categorías, y \(F, G, H : \mathcal{C} \to \mathcal{D}\) tres functores. Y sean \(\Phi : F \Longrightarrow G\) y \(\Psi : G \Longrightarrow H\) transformaciones naturales. La \emph{composición} de \(\Phi\) con \(\Psi\), denotada por $\Psi \circ \Phi$ está definida para cada objeto \(A \in \mathcal{C}\) por
		$$
		(\Psi \circ \Phi)_A := \Psi_A \circ \Phi_A: F(A) \to H(A).
		$$ 		Así, la composición de transformaciones naturales se obtiene componiendo las componentes en cada objeto.
	}
\end{definicion}

\begin{proposicion}\label{lkjhgf}
	\upshape{
		La composición de transformaciones naturales es una transformación natural.
	}
\end{proposicion}

\begin{proof}
	Sea \(\mathcal{C}\) y \(\mathcal{D}\) dos categorías, y \(F, G, H : \mathcal{C} \to \mathcal{D}\) tres functores. Sean 
	\[
	\Phi : F \Longrightarrow G, \quad \Psi : G \Longrightarrow H
	\]
	transformaciones naturales con componentes \(\{\Phi_A\}_{A \in \mathrm{Ob}(\mathcal{C})}\) y \(\{\Psi_A\}_{A \in \mathrm{Ob}(\mathcal{C})}\).\\
	
	
	\noindent Queremos demostrar que \(\Psi \circ \Phi\) es una transformación natural, es decir, que para todo morfismo \(f : A \to B\) en \(\mathcal{C}\) el siguiente diagrama conmuta:
	\[
	{\Large
		\xymatrix{
			F(A) \ar[rr]^{(\Psi \circ \Phi)_A} \ar[dd]_{F(f)} && H(A) \ar[dd]^{H(f)} \\ \\
			F(B) \ar[rr]_{(\Psi \circ \Phi)_B} && H(B)
		}
	}
	\] Esto equivale a verificar que
	\[
	H(f) \circ (\Psi \circ \Phi)_A = (\Psi \circ \Phi)_B \circ F(f).
	\]	Por definición de composición,
	\[
	H(f) \circ \Psi_A \circ \Phi_A = \Psi_B \circ \Phi_B \circ F(f).
	\] Como \(\Psi\) es transformación natural,
	\[
	H(f) \circ \Psi_A = \Psi_B \circ G(f),
	\]
	y como \(\Phi\) es transformación natural,
	\[
	G(f) \circ \Phi_A = \Phi_B \circ F(f).
	\] Sustituyendo estas igualdades se obtiene:
	\[
	H(f) \circ \Psi_A \circ \Phi_A = \Psi_B \circ G(f) \circ \Phi_A = \Psi_B \circ \Phi_B \circ F(f).
	\] Por lo tanto, \(\Psi \circ \Phi\) es una transformación natural.
\end{proof}




\section{Functor de Yoneda}



\begin{proposicion}
	\upshape{
		Sean \( \mathcal{C} \) y \( \mathcal{D} \) categorías. Entonces, \( \mathcal{D}^\mathcal{C} \), cuya clase de objetos está formada por los functores \( F : \mathcal{C} \to \mathcal{D} \), y cuyos morfismos son las transformaciones naturales entre ellos, es una categoría.
	}
\end{proposicion}

\begin{proof}
	En efecto
	\begin{enumerate}
		
		\item \textbf{Objetos:} Son los functores \(F : \mathcal{C} \to \mathcal{D}\). Esto es $F\in \mathrm{Ob}\left(\mathcal{D}^\mathcal{C}\right)$.
		
		\item \textbf{Morfismos:} Son las transformaciones naturales, es decir, dados functores \(F, G : \mathcal{C} \to \mathcal{D}\), una transformación natural \(\Phi : F \Longrightarrow G\) es un morfismo, esto dignifica que, $\Phi \in \mathrm{Hom}_{\mathcal{D}^\mathcal{C}}(F,G)$
		
		\item \textbf{Composición:} Sean \(\Phi : F \Longrightarrow G\) y \(\Theta : G \Longrightarrow H\) transformaciones naturales. Por la Proposición \ref{lkjhgf} su composición \(\Theta \circ \Phi : F \Longrightarrow H\) es una transformación natural.
		
		\item \textbf{Identidad:} Para cada functor \(F\in \mathcal{D}^\mathcal{C}\), de Ejemlplo \ref{cascas}, existe la transformación natural identidad \(\mathrm{id}_F : F \Rightarrow F\). Verificamos cómo actúa respecto a la composición: \\
		
		\noindent Dado \(F \in \mathrm{Ob}(\mathcal{D}^\mathcal{C})\), sean \(G \in \mathrm{Ob}(\mathcal{D}^\mathcal{C})\), \(\Phi \in \mathrm{Hom}_{\mathcal{D}^\mathcal{C}}(F, G)\) y \(\Psi \in \mathrm{Hom}_{\mathcal{D}^\mathcal{C}}(G, F)\). Para cada objeto \(A \in \mathrm{Ob}(\mathcal{C})\), tenemos:
		\begin{align*}
			(\Phi \circ \mathrm{id}_F)_A &= \Phi_A \circ (\mathrm{id}_F)_A = \Phi_A \circ \mathrm{id}_{F(A)} = \Phi_A, \\
			(\mathrm{id}_F \circ \Psi)_A &= (\mathrm{id}_F)_A \circ \Psi_A = \mathrm{id}_{F(A)} \circ \Psi_A = \Psi_A.
		\end{align*}
		
		\noindent Por lo tanto,
		\[
		\Phi \circ \mathrm{id}_F = \Phi \quad \text{y} \quad \mathrm{id}_F \circ \Psi = \Psi.
		\]
	
		\item \textbf{Asociatividad:} Si \(\Phi : F \Rightarrow G\), \(\Theta : G \Rightarrow H\), \(\Psi : H \Rightarrow K\), entonces
		\[
		(\Psi \circ (\Theta \circ \Phi))_A = \Psi_A \circ (\Theta_A \circ \Phi_A) = (\Psi_A \circ \Theta_A) \circ \Phi_A = ((\Psi \circ \Theta) \circ \Phi)_A,
		\]
		por asociatividad de la composición en \(\mathcal{D}\). Luego,
		\[
		\Psi \circ (\Theta \circ \Phi) = (\Psi \circ \Theta) \circ \Phi.
		\]
	\end{enumerate}
\end{proof}
  
 
 

 
 
 
 
 
 
 \noindent Definimos 
 \begin{align*}
 	h : \mathcal{C} &\longrightarrow \mathbf{Set}^{\mathcal{C}^{\mathrm{op}}} \\
 	A &\longmapsto \mathrm{Hom}_{\mathcal{C}}(\,\_\,, A),\\
 	\left(f:A\to B\right)&\longmapsto \left( f^\ast: \mathrm{Hom}_{\mathcal{C}}(\_\,,A) \Longrightarrow \mathrm{Hom}_{\mathcal{C}}(\_\,,B)\right)
 \end{align*}
 donde 
 \begin{align*}
 	\mathrm{Hom}_{\mathcal{C}}(\,\_\,, A) : \mathcal{C}^{\mathrm{op}} &\longrightarrow \mathbf{Set} \\
 	X &\longmapsto \mathrm{Hom}_{\mathcal{C}}(X, A),
 \end{align*}
 y \(f^\ast\) está dada por la familia \(\left\{f^\ast_{X}: \mathrm{Hom}_{\mathcal{C}}(X,A)\to \mathrm{Hom}_{\mathcal{C}}(X,B)\right\}_{X \in \mathcal{C}}\), cuyas componentes se definen como
 \begin{align*}
 	f^\ast_X: \mathrm{Hom}_{\mathcal{C}}(X,A) &\longrightarrow \mathrm{Hom}_{\mathcal{C}}(X,B) \\
 	g &\longmapsto f \circ g.
 \end{align*}
 
 
 
 
 
 
 
 
 \noindent
 Para cada objeto \(A\) de \(\mathcal{C}\), denotamos por
 \[
 h_A := \mathrm{Hom}_{\mathcal{C}}(\_\,, A) : \mathcal{C}^{\mathrm{op}} \longrightarrow \mathbf{Set}.
 \]
 
 
 \begin{proposicion}
 	\upshape{
 		La asignación \(h : \mathcal{C} \to \mathbf{Set}^{\mathcal{C}^{\mathrm{op}}}\), dada por \(A \mapsto h_A\) y \(f \mapsto f^\ast\), define un functor covariante.
 	}
 \end{proposicion}
 
 
 
 
 \begin{proof}
 	Para demostrar que $h$ es un functor covariante, verificamos las siguientes propiedades:
 	\begin{enumerate}
 		
 		\item \textbf{Asignación a objetos:} Por definición, \( h(A) = h_A \in \mathbf{Set}^{\mathcal{C}^{\mathrm{op}}} \), ya que \( h_A \) es un functor contravariante de \( \mathcal{C} \) a \( \mathbf{Set} \).
 		
 		\item \textbf{Asignación a morfismos:} Dado \( f: A \to B \), definimos \( f^\ast: h_A \Longrightarrow h_B \) por \( f^\ast_X(g) := f \circ g \) para \( g \in \mathrm{Hom}_\mathcal{C}(X, A) \). Mostramos que esta es una transformación natural.\\
 		
 		\noindent Sea \( u: Y \to X \) en \( \mathcal{C}^{\mathrm{op}} \) (es decir, \( u: X \to Y \) en \( \mathcal{C} \)). Queremos probar que el siguiente diagrama conmuta:
 		
 		\[
 		\xymatrix @C=4.5em{
 			\mathrm{Hom}_\mathcal{C}(X, A) \ar[r]^{\displaystyle f^\ast_X} \ar[dd]_{\displaystyle h_A(u)} & \mathrm{Hom}_\mathcal{C}(X, B) \ar[dd]^{ \displaystyle h_B(u)} \\
 			&&\\
 			\mathrm{Hom}_\mathcal{C}(Y, A) \ar[r]_{ \displaystyle f^\ast_Y} & \mathrm{Hom}_\mathcal{C}(Y, B)
 		}
 		\]
 		
 		Para \( g \in \mathrm{Hom}_\mathcal{C}(X, A) \), tenemos:
 		\begin{align*}
 			(f^\ast_Y \circ h_A(u))(g) &= f^\ast_Y(g \circ u) && \text{(por definición de } h_A(u)\text{)} \\
 			&= f \circ (g \circ u) && \text{(por definición de } f^\ast_Y\text{)} \\
 			&= (f \circ g) \circ u && \text{(asociatividad)} \\
 			&= h_B(u)(f \circ g) && \text{(por definición de } h_B(u)\text{)} \\
 			&= (h_B(u) \circ f^\ast_X)(g) && \text{(por definición de } f^\ast_X\text{)}
 		\end{align*}
 		Así, \( f^\ast \) es una transformación natural.
 		
 		\item \textbf{Preservación de la composición:}
 		Sean $f: A \to B$ y $g: B \to C$ morfismos en $\mathcal{C}$. Debemos verificar que
 		\[
 		h(g \circ f) = h(g) \circ h(f),
 		\]
 		es decir, que
 		\[
 		(g \circ f)^\ast = g^\ast \circ f^\ast.
 		\]
 		Para cada objeto $X \in \mathcal{C}$ y cada $u \in \mathrm{Hom}_{\mathcal{C}}(X, A)$, tenemos:
 		\begin{align*}
 			((g \circ f)^\ast)_X(u) &= (g \circ f) \circ u && \text{(def. de } (g \circ f)^\ast_X) \\
 			&= g \circ (f \circ u) && \text{(asociatividad)} \\
 			&= g^\ast_X(f \circ u) && \text{(def. de } g^\ast_X) \\
 			&= g^\ast_X(f^\ast_X(u)) && \text{(def. de } f^\ast_X) \\
 			&= (g^\ast_X \circ f^\ast_X)(u) && \text{(composición de funciones)} \\
 			&= (g^\ast \circ f^\ast)_X(u) && \text{(def. de comp. de transf. nat.)}
 		\end{align*}
 		
 		\item \textbf{Preservación de identidades:}
 		Para cada objeto $A \in \mathcal{C}$, debemos verificar que
 		\[
 		h(\mathrm{id}_A) = \mathrm{id}_{h_A},
 		\]
 		es decir, que
 		\[
 		\mathrm{id}_A^\ast = \mathrm{id}_{h_A}.
 		\]
 		Para cada objeto $X \in \mathcal{C}$ y cada $g \in \mathrm{Hom}_{\mathcal{C}}(X, A)$:
 		\begin{align*}
 			(\mathrm{id}_A^\ast)_X(g) &= \mathrm{id}_A \circ g && \text{(por definición)} \\
 			&= g && \text{(identidad de la composición)} \\
 			&= (\mathrm{id}_{h_A})_X(g).
 		\end{align*}
 		Por lo tanto, $\mathrm{id}_A^\ast = \mathrm{id}_{h_A}$.
 		
 	\end{enumerate}
 	
 	Con esto, se concluye que \( h \) es un functor covariante.
 	
 
 \end{proof}
 
 
\begin{observacion}\upshape{
	
	Se utiliza la categoría opuesta \(\mathcal{C}^{\mathrm{op}}\) para definir \(h_A = \mathrm{Hom}_{\mathcal{C}}(\_, A)\) porque la variable donde actúa el morfismo en \(\mathrm{Hom}_{\mathcal{C}}(X, A)\) es contravariante. Es decir, un morfismo \(f: X \to Y\) induce una función \(\mathrm{Hom}_{\mathcal{C}}(Y, A) \to \mathrm{Hom}_{\mathcal{C}}(X, A)\) que va en dirección opuesta a \(f\). Por eso, para que \(h_A\) sea un functor covariante, consideramos los objetos y morfismos en \(\mathcal{C}^{\mathrm{op}}\), invirtiendo así las flechas y respetando la composición.
	}
\end{observacion}

 
\begin{definicion}
	\upshape{
		El \emph{functor de Yoneda} es el functor covariante
		\[
		h : \mathcal{C} \to \mathbf{Set}^{\mathcal{C}^\mathrm{op}},
		\]
		ya definido, que asigna a cada objeto \(A \in \mathcal{C}\) el functor  \(h_A = \mathrm{Hom}_{\mathcal{C}}(-, A)\) y a cada morfismo \(f: A \to B\) la transformación natural \(f^\ast : h_A \Longrightarrow h_B\).
	}
\end{definicion}


 
 
 
 
	\chapter{Objetos libres}
	\noindent Dado un conjunto base \(X\), una de las preguntas más fundamentales en distintas áreas del álgebra es cómo construir, a partir de sus elementos, una estructura algebraica con ciertas propiedades deseadas. Es decir, ¿cómo dotar a \(X\) de operaciones que satisfagan ciertos axiomas, de la manera más general posible? Por ejemplo:

\begin{itemize}
	\item ¿Cómo obtener el magma generado por \(X\)?
	\item ¿Cómo construir el semigrupo?
	\item ¿Cómo formar el monoide o el grupo sobre \(X\)?
	\item ¿Cómo generar, en general, otras estructuras algebraicas a partir de $X$?
\end{itemize}

\noindent Es importante resaltar que no se asume que \(X\) esté contenido en una estructura algebraica preexistente. \(X\) es simplemente un conjunto de símbolos formales sin relaciones entre ellos.\\

\noindent El objetivo es construir una estructura algebraica que:

\begin{itemize}
	\item contenga a \(X\), y
	\item esté generada por los elementos de \(X\) de manera puramente formal, sin imponer más relaciones que las necesarias para satisfacer los axiomas correspondientes.
\end{itemize}

\noindent A este tipo de construcción se le llama \emph{objeto libre}. Su rasgo distintivo es ser la estructura más general (o universal) posible generada por \(X\), en el sentido de que no introducen identificaciones ni relaciones adicionales entre sus elementos más allá de las que son estrictamente obligatorias por la definición de la estructura. Por esta razón, también se les considera el menor objeto que satisface esas propiedades a partir de \(X\).\\

\noindent Los objetos libres desempeñan un papel central en álgebra y teoría de categorías, al proporcionar un método sistemático para generar estructuras desde sus generadores, y entender cómo se extienden funciones desde conjuntos a morfismos algebraicos.






\section{Magma libre}

\noindent Sea \(X\) un conjunto. Construimos el \emph{magma libre generado por \(X\)} de la siguiente manera.\\

\noindent Definimos inductivamente una sucesión de conjuntos \(X_1, X_2, X_3, \dots\):

\begin{itemize}
	\item \(X_1 := X\),
	\item para \(n \geq 2\), 
	\[
	X_n := \bigsqcup_{\substack{i + j = n \\ i,j \geq 1}} X_i \times X_j.
	\]
\end{itemize}

\begin{observacion}
	\upshape
	Los primeros términos se expresan como:
\begin{itemize}
	\item \(X_2 \:\,= X_1 \times X_1 = X \times X\)
	\item \(X_3 \:\,= (X_1 \times X_2) \sqcup (X_2 \times X_1) = (X \times (X \times X)) \sqcup ((X \times X) \times X)\)
	\item \(X_4 \:\,= (X_1 \times X_3) \sqcup (X_2 \times X_2) \sqcup (X_3 \times X_1)\)
	
	\quad\ \, \(= \left[X \times ((X \times (X \times X)) \sqcup ((X \times X) \times X))\right]\)
	
	\quad	\quad\ \, \(\sqcup\ \left[(X \times X) \times (X \times X)\right]\)
	
	\quad\quad\ \, \(\sqcup\ \left[((X \times (X \times X)) \sqcup ((X \times X) \times X)) \times X\right]\).
\end{itemize}
\end{observacion}

\noindent Notemos que si \(a \in X_i\) y \(b \in X_j\), entonces \((a,b) \in X_i \times X_j \subseteq X_{i+j}\). \\

\noindent Definimos entonces
\[
M_X := \bigsqcup_{n \geq 1} X_n,
\]
y la operación interna llamada \emph{concatenación}
\[
M_X \times M_X \longrightarrow  M_X, \quad (a,b) \longmapsto (a,b).
\]

\begin{definicion}
	\upshape
	Dado un conjunto no vacío \(X\), el par \(\bigl(M_X, \text{concatenación}\bigr)\) se denomina el \emph{magma libre generado por \(X\)}. Además, es habitual escribir \(\mathrm{Mag}(X)\) para referirse a dicho magma libre.
\end{definicion}


\begin{observacion}
	\upshape
	El magma libre generado por \(X\) no es asociativo en general. Esto se debe a que, dados \(a,b,c \in M_X\), las dos expresiones \(((a,b),c)\) y \((a,(b,c))\) son elementos distintos de \(M_X\), ya que difieren en la forma de agrupar los paréntesis y, por tanto, no se identifican en la construcción. En otras palabras, la operación de concatenación definida como par ordenado no satisface la propiedad asociativa, porque
	\[
	((a,b),c) \neq (a,(b,c)).
	\]
\end{observacion}



\section{Semigrupo libre}


\noindent Los conjuntos \(X \times (Y \times Z)\) y \((X \times Y) \times Z\) son, por la naturaleza de los productos cartesianos, naturalmente isomorfos. Es decir, existe una biyección canónica
	\[
	(x, (y, z)) \mapsto ((x, y), z),
	\]
	que permite identificar de forma coherente elementos que difieren únicamente en la agrupación de paréntesis. Esta identificación es esencial para fundamentar la imposición de la asociatividad en la operación.\\
	
	\noindent A partir de esta identificación, definiremos una relación de equivalencia sobre el magma libre \(M_X\) que formaliza la equivalencia entre diferentes agrupaciones. Así, al considerar las clases de equivalencia, la operación se vuelve asociativa, lo que conduce a la construcción del semigrupo libre generado por \(X\).


\begin{definicion}
	\upshape
	Definimos una relación de equivalencia \(\sim\) en \(M_X = \bigsqcup_{n \geq 1} X_n\) por:
	\[
	a \sim b \iff \exists\, n \geq 1,\ a, b \in X_n,\ \text{y}\ b = \phi(a)
	\]
	donde \(\phi : X_n \to X_n\) es una biyección obtenida por reagrupación de paréntesis en el producto cartesiano.
\end{definicion}





Definimos ahora
$$
S_X := \frac{M_X}{\sim}.
$$
Para \(a \in M_X\), si por ejemplo \(a = \bigl((a_1, a_2), a_3\bigr) \in X_3\), escribiremos su clase de equivalencia como
$$
[a] = a_1 a_2 a_3.
$$
Entonces, la concatenación en \(M_X\) induce una operación interna en \(S_X\), también llamada \emph{concatenación}, definida por
$$
\begin{aligned}
	S_X \times S_X &\longrightarrow S_X \\
	(p, q) &\longmapsto p \bullet q,
\end{aligned}
$$
donde, por ejemplo, si \(p, q \in S_X\) con \(p = p_1 p_2 p_3 p_4\) y \(q = q_1 q_2\), se tiene
$$
p \bullet q = p_1 p_2 p_3 p_4 q_1 q_2.
$$

\begin{observacion}
	\upshape{
		La operación \(\bullet\) es asociativa, pues la relación de equivalencia \(\sim\) identifica las distintas maneras de agrupar la concatenación, eliminando así la ambigüedad en el orden de evaluación.
	}
\end{observacion}

\begin{observacion}
	\upshape{
		La operación \(\bullet\) es asociativa. Por ejemplo, consideremos \(p, q, r \in S_X\) con
		$$
		p = p_1 p_2 p_3 p_4 p_5, \quad q = q_1 q_2 q_3, \quad r = r_1 r_2.
		$$
		Entonces,
		$$
		(p \bullet q) \bullet r = (p_1 p_2 p_3 p_4 p_5 q_1 q_2 q_3) \bullet r = p_1 p_2 p_3 p_4 p_5 q_1 q_2 q_3 r_1 r_2,
		$$
		y de forma similar,
		$$
		p \bullet (q \bullet r) = p \bullet (q_1 q_2 q_3 r_1 r_2) = p_1 p_2 p_3 p_4 p_5 q_1 q_2 q_3 r_1 r_2.
		$$
		Ambas expresiones coinciden, lo que muestra que \(\bullet\) es asociativa.
	}
\end{observacion}


\begin{definicion}
	\upshape{
		Dado un conjunto no vacío \(X\), el par \(\left(S_X, \bullet\right)\) se denomina \emph{semigrupo libre generado por \(X\)}.
	}
\end{definicion}



\begin{definicion}
	\upshape
	Sea \(X\) un conjunto no vacío, llamado \emph{alfabeto} o \emph{conjunto generador}. El semigrupo libre \(S_X\) está generado por \(X\), y sus elementos se llaman \emph{palabras} sobre \(X\).\\
	
	\noindent Una palabra \(p \in S_X\) puede representarse como una concatenación linealizada
	$$
	p = p_1 p_2 \cdots p_n,
	$$
	donde cada \(p_i \in X\).
\end{definicion}

\begin{definicion}
	\upshape
	La \emph{longitud} de una palabra \(p = p_1 p_2 \cdots p_n \in S_X\) es el número \(n\) de generadores que la componen,
	$$
	|p| := n.
	$$
	Esta longitud está bien definida debido a que la relación de equivalencia \(\sim\) solo identifica distintas formas de agrupar la concatenación, sin alterar el número ni el orden de los elementos.
\end{definicion}


\begin{observacion}
	\upshape{
		Como $M_X := X \sqcup X_2 \sqcup X_3 \sqcup X_4 \sqcup \dotsb$, entonces
		$$
		S_X = \dfrac{M_X}{\sim} = X \sqcup X^2 \sqcup X^3 \sqcup X^4 \sqcup \dotsb,
		$$
		donde, para cada $n \geq 1$,
		$$
		X^n := \{\,p \in S_X : |p| = n\}.
		$$
	}
\end{observacion}













\section{Monoide libre}

\noindent El semigrupo libre $S_X$, construido a partir de las clases de equivalencia de tuplas no vacías de elementos de $X$, es inherentemente una estructura sin un elemento identidad con respecto a la concatenación. Esto significa que, por su definición, no existe en $S_X$ una palabra que, al concatenarse con cualquier otra, la deje inalterada. Para trascender esta limitación y construir un monoide libre generado por $X$, es esencial introducir artificialmente un elemento que cumpla el rol de identidad.\\

\begin{definicion}
	\upshape{
		La \textit{palabra vacía} es un elemento especial, denotado por $\mathbf{1}$ ($\mathrm{e}$,$\emptyset$), que no pertenece al semigrupo libre $S_X$. Se introduce formalmente como el elemento identidad para extender $S_X$ a un monoide. Representa la ``ausencia'' de elementos de $X$ y tiene la propiedad de que, para cualquier palabra $w \in S_X$, $w \bullet \mathbf{1} = \mathbf{1} \bullet w = w$.
	}
\end{definicion}
\begin{observacion}
	\upshape{
	La palabra vacía \(1\) tiene longitud cero, pues representa la ausencia de elementos.
	}
\end{observacion}

Definimos ahora  
$$
\mathrm{Mon}(X) := \{1\} \sqcup S_X = \{1\} \sqcup X \sqcup X^2 \sqcup X^3 \sqcup X^4 \sqcup \dotsb,
$$  
donde el símbolo \(1\) representa la \emph{palabra vacía}.\\

\noindent Entonces, la concatenación en \(S_X\) se extiende naturalmente a una operación interna en \(\mathrm{Mon}(X)\), también llamada \emph{concatenación}, definida por  
$$
\begin{aligned}
	\mathrm{Mon}(X) \times \mathrm{Mon}(X) &\longrightarrow \mathrm{Mon}(X) \\
	(p, q) &\longmapsto p \bullet q,
\end{aligned}
$$  
donde \(1 \bullet p = p \bullet 1 = p\) para todo \(p \in \mathrm{Mon}(X)\).

\begin{definicion}
	\upshape{
		El par \((\mathrm{Mon}(X), \bullet)\) se denomina \emph{monoide libre generado por \(X\)}.
	}
\end{definicion}
\begin{observacion}
	\upshape{
	Si $X$ es vació, entonces $\mathrm{Mon}(\emptyset)=\{1\}$.
	}
\end{observacion}




\section{Grupo libre}

\noindent Hemos visto que, dado un conjunto \(X\), el monoide libre \(\mathrm{Mon}(X)\) es la estructura más general con operación asociativa e identidad que puede construirse a partir de \(X\), sin imponer relaciones adicionales. Sus elementos son palabras formadas con símbolos de \(X\), incluida la palabra vacía.\\

\noindent Sin embargo, en un monoide no se exige la existencia de inversos. Por ejemplo, si \(x \in X \subseteq \mathrm{Mon}(X)\) es una palabra no vacía, en general no existe una palabra \(y \in \mathrm{Mon}(X)\) tal que \(x \cdot y = \mathbf{1}\).\\

\noindent Nuestro siguiente objetivo es construir un grupo a partir de \(X\), lo que requiere dotar a cada generador de un inverso. Para ello, introduciremos \emph{inversas formales} para cada elemento de \(X\) y definiremos una relación de equivalencia que exprese la cancelación entre elementos y sus inversos. Esta construcción nos llevará al \emph{grupo libre} generado por \(X\).


\begin{proposicion}\label{lkjhgkjhg}
	\upshape{
		Dado un conjunto no vacío \(X\), existe un conjunto \(X'\), también no vacío y biyectivo con \(X\), tal que \(X \cap X' = \emptyset\).
	}
\end{proposicion}


\begin{proof}
	Dado que \(X\) es no vacío, existe algún elemento \(x_0 \in X\). Definimos entonces
	\[
	X' := X \times \{x_0\}.
	\]
	Como el producto cartesiano con un conjunto unitario preserva la cardinalidad, el conjunto \(X'\) es no vacío y existe una biyectiva natural
	\[
	\begin{aligned}
		\varphi : X &\longrightarrow X' \\
		x &\longmapsto (x,x_0),
	\end{aligned}
	\]
	por lo que \(X'\) es biyectivo con \(X\).Finalmente, los elementos de \(X'\) tienen la forma \((x,x_0)\), mientras que los de \(X\) no son pares ordenados, por lo que \(X \cap X' = \emptyset\).
\end{proof}




\begin{definicion}
	\upshape{
		Dado un elemento \(a \in X\), definimos su \emph{inversa formal} como el elemento correspondiente \(a^{-1} := \varphi(a) \in X'\), donde \(\varphi : X \to X'\) es la biyección establecida anteriormente.
	}
\end{definicion}
 De esta manera, el conjunto
\[
X^{\pm 1} := X \cup X'
\]
denota el conjunto de generadores extendido con sus inversas formales, es decir, los símbolos que representan tanto los elementos de $X$ como sus inversos formales.

\begin{observacion}
	\upshape{
		Note que la palabra 
		\[
		abcdd^{-1}c^{-1}b^{-1}a^{-1}
		\]
		no es necesariamente la palabra vacía, pues en \(\mathrm{Mon}(X^{\pm 1})\), que es el monoide libre generado por \(X \cup X'\), sólo está definida la concatenación de palabras, sin ninguna regla de cancelación.
	}
\end{observacion}





\begin{definicion}
	\upshape{
		Definimos la relación de equivalencia \(\sim\) en \(\mathrm{Mon}(X^{\pm 1})\) como la relación que identifica dos palabras si una puede obtenerse de la otra eliminando una o varias veces subpalabras consecutivas de la forma \(a \cdot a^{-1}\) o \(a^{-1} \cdot a\), para algún \(a \in X\).\\
		
		\noindent Es decir, para toda palabra \(w_1, w_2 \in \mathrm{Mon}(X^{\pm 1})\) y todo \(a \in X\), se tiene
		\[
		w_1 \cdot a \cdot a^{-1} \cdot w_2 \sim w_1 \cdot w_2, \quad \text{y} \quad
		w_1 \cdot a^{-1} \cdot a \cdot w_2 \sim w_1 \cdot w_2.
		\] Esta relación formaliza la cancelación de pares de elementos inversos consecutivos en las palabras.
	}
\end{definicion}

\begin{definicion}\label{def:grupolibrerelacion}
	\upshape{
		Definimos la relación de equivalencia \(\sim\) en \(\mathrm{Mon}(X^{\pm 1})\) como la relación que identifica dos palabras \(w, w' \in \mathrm{Mon}(X^{\pm 1})\) si \(w'\) puede obtenerse a partir de \(w\) mediante unasecuencia finita de inserciones o eliminaciones de subpalabras consecutivas de la forma \(a \cdot a^{-1}\) o \(a^{-1} \cdot a\), con \(a \in X\).\\
		
		\noindent Dichas operaciones elementales son de la forma:
		\[
		w_1 \cdot a \cdot a^{-1} \cdot w_2 \longleftrightarrow w_1 \cdot w_2, \quad \text{y} \quad
		w_1 \cdot a^{-1} \cdot a \cdot w_2 \longleftrightarrow w_1 \cdot w_2,
		\]
		donde \(w_1, w_2 \in \mathrm{Mon}(X^{\pm 1})\) y \(a \in X\).\\
		
		\noindent Esta relación formaliza la cancelación de pares consecutivos de elementos inversos, permitiendo también su inserción. 
	}
\end{definicion}





\begin{proposicion}
	\upshape{
	La relación $\sim$ de la Definición \ref{def:grupolibrerelacion} es una relación de equivalencia.
	}
\end{proposicion}












\begin{proof}
	Verificamos que la relación \(\sim\) definida en \(\mathrm{Mon}(X^{\pm1})\) es una relación de equivalencia.
	
	\begin{itemize}
		\item \textbf{Reflexividad:} Sea \(w \in \mathrm{Mon}(X^{\pm1})\). Como no se requiere realizar ninguna operación para obtener \(w\) a partir de sí mismo, se tiene
		\[
		w \sim w.
		\]
		
		\item \textbf{Simetría:} Supongamos que \(w \sim v\). Entonces existe una secuencia finita de palabras
		\[
		w = w_0 \rightarrow w_1 \rightarrow \cdots \rightarrow w_k = v,
		\]
		donde cada paso corresponde a la inserción o eliminación de una subpalabra de la forma \(a a^{-1}\) o \(a^{-1} a\), con \(a \in X\).\\
		Invirtiendo la secuencia:
		\[
		v = w_k \rightarrow w_{k-1} \rightarrow \cdots \rightarrow w_0 = w,
		\]
		por lo tanto,
		\[
		v \sim w.
		\]
		
		\item \textbf{Transitividad:} Supongamos que \(w \sim v\) y \(v \sim u\). Entonces existen secuencias finitas:
		\[
		w = w_0 \rightarrow w_1 \rightarrow \cdots \rightarrow w_k = v, \quad
		v = v_0 \rightarrow v_1 \rightarrow \cdots \rightarrow v_m = u.
		\]
		Concatenando ambas secuencias, obtenemos:
		\[
		w = w_0 \rightarrow \cdots \rightarrow w_k = v = v_0 \rightarrow \cdots \rightarrow v_m = u,
		\]
		y por lo tanto,
		\[
		w \sim u.
		\]
	\end{itemize}
	
	Luego, \(\sim\) es una relación de equivalencia.
\end{proof}


\begin{observacion}
	\upshape{
		Dada una palabra \(u \in \mathrm{Mon}(X^{\pm 1})\), la palabra concatenada \(u \bullet u^{-1}\) es equivalente a la palabra vacía \(1_{\mathrm{Mon}(X^{\pm 1})}\) bajo la relación \(\sim\), es decir,
		\[
		u \bullet u^{-1} \sim 1_{\mathrm{Mon}(X^{\pm 1})}.
		\]
	}
\end{observacion}




Definimos  $$G_X:=\dfrac{\mathrm{Mon}(X^{\pm 1})}{\sim}$$

\begin{definicion}\label{def:palabrareducida}
	\upshape{
		Una palabra \( w \in \mathrm{Mon}(X^{\pm 1}) \) se dice \emph{reducida} si no contiene ninguna subpalabra consecutiva de la forma \(a a^{-1}\) o \(a^{-1} a\), para algún \(a \in X\).\\
			
		}
\end{definicion}

\begin{observacion}
	\upshape{
		Por definición, los elementos de \(G_X ) \) son clases de equivalencia de palabras en \(\mathrm{Mon}(X^{\pm 1})\). Sin embargo, debido a que cada clase contiene exactamente una palabra reducida (Definición \ref{def:palabrareducida}), se establece una identificación natural entre cada elemento \(g \in G_X\) y su representante reducido \(w \in \mathrm{Mon}(X^{\pm 1})\).\\
		
		\noindent Por lo tanto, aunque formalmente los elementos de \(G_X\) son clases de equivalencia, en la práctica los denominamos \emph{palabras reducidas} mediante esta identificación. Así, se puede escribir simplemente
		\[
		[g] = w,
		\]
		donde \(w\) es la palabra reducida representante de la clase \([g]\).
	}
\end{observacion}
La concatenación en \(\mathrm{Mon}(X^{\pm 1})\)
\[
\begin{aligned}
	\mathrm{Mon}(X^{\pm 1}) \times \mathrm{Mon}(X^{\pm 1}) &\longrightarrow \mathrm{Mon}(X^{\pm 1}) \\
	(p, q) &\longmapsto p \bullet q,
\end{aligned}
\]
induce una operación interna (también llamada concatenación) en el cociente $G_X $:
\[
\begin{aligned}
	G_X \times G_X &\longrightarrow G_X \\
	([p], [q]) &\longmapsto [p]\bullet[q]:=[p \bullet q].
\end{aligned}
\]



\begin{proposicion}
	\upshape{
		En \((G_X, \bullet)\) existe un elemento identidad \(1 \in G_X\), y todo elemento \( [p] \in G_X \) admite un inverso \( [p]^{-1} \in G_X \) tal que
		\[
		[p] \bullet [p]^{-1} = 1 \quad \text{y} \quad [p]^{-1} \bullet [p] = 1.
		\]
	}
\end{proposicion}


\begin{proof}
	En efecto.
	\begin{itemize}
		\item \textbf{Existencia de la identidad:} Definamos
		\[
		1 := [1_{\mathrm{Mon}(X^{\pm 1})}],
		\]
		donde \(1_{\mathrm{Mon}(X^{\pm 1})}\) es la palabra vacía en \(\mathrm{Mon}(X^{\pm 1})\). Para toda palabra reducida \(p \in \mathrm{Mon}(X^{\pm 1})\) se cumple
		\[
		[p] \bullet 1 = [p \bullet 1_{\mathrm{Mon}(X^{\pm 1})}] = [p], \quad 1 \bullet [p] = [1_{\mathrm{Mon}(X^{\pm 1})} \bullet p] = [p].
		\]
		Por lo tanto, \(1\) es el elemento identidad en \(G_X\).
		
		\item \textbf{Existencia de inversos:} Sea \([p] \in G_X\) con \(p = x_1 x_2 \cdots x_n\) palabra reducida. Definimos
		\[
		[p]^{-1} := [x_n^{-1} x_{n-1}^{-1} \cdots x_1^{-1}].
		\]
		Entonces se cumple
		\[
		[p] \bullet [p]^{-1} = [p \bullet (x_n^{-1} x_{n-1}^{-1} \cdots x_1^{-1})] = [1_{\mathrm{Mon}(X^{\pm 1})}] = 1,
		\]
		y análogamente,
		\[
		[p]^{-1} \bullet [p] = 1.
		\]
	\end{itemize}
\end{proof}



\begin{definicion}
	\upshape{
		Sea \(X\) un conjunto. El grupo \(G_X\) se denomina \emph{grupo libre generado por \(X\)}. Sus elementos se llaman \emph{palabras reducidas}.
	}
\end{definicion}

\begin{observacion}
	\upshape{
		Si \(X = \emptyset\), entonces el grupo libre generado por \(X\) es trivial: $$G_{\emptyset} = \{1\}.$$
	}
\end{observacion}


\section{Functor libre y propiedad universal de los objetos libres}


\noindent Recordemos que si \(\mathcal{C}\) es una categoría de cierta estructura algebraica, por ejemplo, magmas, semigrupos, monoides o grupos, entonces existe un funtor olvidadizo



\[
\begin{aligned}
	\mathcal{C} & \longrightarrow \mathbf{Set} \\
	(M, \cdot) & \longmapsto M \\
	\left(f : (M, \cdot) \to (N, \star)\right) & \longmapsto \left(f : M \to N\right)
\end{aligned}
\]que a cada objeto \((M, \cdot)\) de \(\mathcal{C}\) le asigna su conjunto subyacente \(M\), y a cada morfismo \(f\) en \(\mathcal{C}\) la función subyacente \(f\).\\

\noindent Ahora queremos hacer lo opuesto, es decir, queremos definir el \textit{funtor libre} 
\[
F : \mathbf{Set} \longrightarrow \mathcal{C}
\]que a cada conjunto \(X\) le asigna un objeto \(F(X)\) de \(\mathcal{C}\), llamado \emph{objeto libre}. Lo que queda pendiente es determinar cómo se asignan los morfismos. Para ello, analizaremos el caso en el que \(\mathcal{C}\) es la categoría de los grupos. En este contexto, utilizamos la propiedad universal del grupo libre.




\begin{definicion}
	\upshape{
		Sean \( M \) y \( N \) dos grupos. Un \textit{homomorfismo de grupos} es una función \( f : M \to N \) tal que 
		\[
		f(xy) = f(x) f(y),\quad \forall x, y \in M.
		\]
		Si además \( f \) es biyectiva, se dice que es un \textit{isomorfismo de grupos}.
	}
\end{definicion}

\begin{proposicion}
	\upshape{
		Sea \( f : M \to N \) un homomorfismo de grupos. Entonces:
		\begin{enumerate}
			\item \( f(1_M) = 1_N \).
			\item \( f(x^{-1}) = f(x)^{-1} \) para todo \( x \in M \).
		\end{enumerate}
	}
\end{proposicion}


\begin{proof}
	En efecto.
	\begin{enumerate}
		\item Como \( 1_M = 1_M \cdot 1_M \), se tiene
		\[
		f(1_M) = f(1_M \cdot 1_M) = f(1_M) \cdot f(1_M).
		\]	Multiplicando a ambos lados por \( f(1_M)^{-1} \), se obtiene
		\[
		1_N = f(1_M).
		\]
		
		\item Para todo \( x \in M \),
		\[
		1_N = f(1_M) = f(x x^{-1}) = f(x) f(x^{-1}),
		\]
		por lo que \( f(x^{-1}) = f(x)^{-1} \), ya que el inverso en \( N \) es único.
	\end{enumerate}
\end{proof}





\begin{proposicion}[Propiedad universal del grupo libre]
	\upshape{
		Sea \(X\) un conjunto y \(G_X\) el grupo libre generado por \(X\), con la función de inclusión o \textit{aplicación canónica}
		\[
		\iota : X \hookrightarrow G_X, \quad x \mapsto [x].
		\] Para todo grupo \(G\) y toda función \( j : X \to G \), existe un único homomorfismo de grupos \( f : G_X \to G \) que hace conmutar el siguiente diagrama:
		
		
		
		\[
		\xymatrix{
			X \ar[rr]^{\displaystyle j} \ar@{^(->}[dd]_{\displaystyle \iota } & & G \\
			&&\\
			G_X \ar@{-->}[rruu]_{\displaystyle \exists !f} &&
		}
		\] Es decir, \( f \circ \iota = j \).
	}
\end{proposicion}


\begin{proof}
	\upshape{
		Sea \( X \) un conjunto y sea \( G \) un grupo. Dada una función \( j : X \to G \), construiremos un único homomorfismo \( f : G_X \to G \) tal que \( f \circ \iota = j \).
		
		\begin{itemize}
			\item \textbf{Extensión de \(j\):}  
			Extendemos \(j\) a \(X^{\pm 1} := X \sqcup X^{-1}\), preservando la estructura inversa:
			
			
			\[
			\widetilde{j}(x) :=
			\begin{cases}
				j(x), & \text{si } x \in X, \\
				j(x')^{-1}, & \text{si } x = (x')^{-1} \in X^{-1}.
			\end{cases}
			\]
			
			
			Esta es la única forma de extender \( j \) respetando la inversibilidad en \( G \).
			
			\item \textbf{Definición de \(f\):}  
			Dada una clase \([x_1 \cdots x_n] \in G_X\), donde la palabra es reducida, definimos:
			
			
			\[
			f([x_1 \cdots x_n]) := \widetilde{j}(x_1) \cdots \widetilde{j}(x_n), \quad f(1_{G_X}) := 1_G.
			\]
			
			
			
			\item \textbf{Bien definida:}  
			La palabra reducida en \( G_X \) es única, lo que garantiza que \( f([w]) \) está bien definida.  
			Además, la única relación en \( G_X \) es la cancelación \( x x^{-1} \sim 1 \), y en \( G \) se cumple \( \widetilde{j}(x) \widetilde{j}(x^{-1}) = 1_G \), asegurando la coherencia entre \( G_X \) y \( G \).
			
			\item \textbf{Homomorfismo:} Sean \([p], [q] \in G_X\), con \(p = x_1 \cdots x_n\), \(q = y_1 \cdots y_m\). Entonces,
			\[
			f([p] \bullet [q]) = f([x_1 \cdots x_n y_1 \cdots y_m]) = \widetilde{j}(x_1) \cdots \widetilde{j}(x_n)\widetilde{j}(y_1) \cdots \widetilde{j}(y_m) = f([p]) \cdot f([q]),
			\]
			lo que garantiza que \( f \) respeta la estructura de grupo.
			
			\item \textbf{Conmutatividad:}  
			Para \( x \in X \), \(\iota(x) = [x]\), y por definición,
			
			
			\[
			f(\iota(x)) = f([x]) = \widetilde{j}(x) = j(x).
			\]
			
			
			Por lo tanto, \( f \circ \iota = j \).
			
			\item \textbf{Unicidad:}  
			Sea \( g : G_X \to G \) otro homomorfismo con \( g \circ \iota = j \). Como ambos asignan los mismos valores a los generadores y respetan la estructura de grupo,
			
			
			\[
			g([x_1 \cdots x_n]) = g([x_1]) \cdots g([x_n]) = \widetilde{j}(x_1) \cdots \widetilde{j}(x_n) = f([x_1 \cdots x_n]).
			\]
			
			
			Por lo tanto, \( g = f \) y la unicidad está probada.
		\end{itemize}
	}
\end{proof}

%\begin{proposicion}[Unicidad del grupo libre salvo isomorfismo]
%	\upshape{
	%		Dado un conjunto \(X\), el grupo libre \(G_X\) generado por \(X\) es \textbf{único salvo isomorfismo natural}.\\
	
	%		\noindent Es decir, si \( (H, k) \) es cualquier otro par que satisface la propiedad universal del grupo libre para el conjunto \( X \) (donde \( H \) es un grupo y \( k: X \to H \) es una función), entonces existe un \textbf{único isomorfismo de grupos} \( \varphi: G_X \to H \) tal que el siguiente diagrama conmuta:
	
	
	
	%		\[
	%		\xymatrix{
		%			X \ar[rrdd]_{\displaystyle \iota} \ar[rr]^{\displaystyle k} && H \\
		%			&&\\
		%			&& G_X \ar[uu]_{\displaystyle \exists ! \varphi}
		%		}
	%		\]
	
	
	%		Es decir, \( \varphi \circ \iota = k \).
	%	}
%\end{proposicion}


%\begin{proof}
%	\upshape{
	%		Sea \( (H, k) \) un par que satisface la propiedad universal del grupo libre para el conjunto \( X \). Como \( (G_X, \iota) \) también la satisface, aplicamos la propiedad universal de \( G_X \) con respecto a la función \( k : X \to H \), y obtenemos un único homomorfismo
	%		\[
	%		\varphi : G_X \to H \quad \text{tal que} \quad \varphi \circ \iota = k.
	%		\]
	
	%		Del mismo modo, aplicamos la propiedad universal de \( H \) con respecto a la función \( \iota : X \to G_X \), y obtenemos un único homomorfismo
	%		\[
	%		\psi : H \to G_X \quad \text{tal que} \quad \psi \circ k = \iota.
	%		\]
	
	%		Ahora probamos que \( \varphi \) y \( \psi \) son inversos uno del otro. Primero, para todo \( x \in X \),
	%		\[
	%		(\psi \circ \varphi)(\iota(x)) = \psi(k(x)) = \iota(x),
	%		\]
	%		y como \( \iota(X) \) genera \( G_X \), se deduce que \( \psi \circ \varphi = \mathrm{id}_{G_X} \).
	
	%		Análogamente, para todo \( x \in X \),
	%		\[
	%		(\varphi \circ \psi)(k(x)) = \varphi(\iota(x)) = k(x),
	%		\]
	%		y como \( k(X) \) genera \( H \), se deduce que \( \varphi \circ \psi = \mathrm{id}_H \).
	
	%		Por lo tanto, \( \varphi \) es un isomorfismo de grupos, y satisface \( \varphi \circ \iota = k \). La unicidad de \( \varphi \) se sigue de la unicidad en la propiedad universal del grupo libre.
	
	%		Así, el grupo libre sobre \( X \) es único salvo isomorfismo natural.
	%	}
%\end{proof}





Pues veamos ahora cómo es la asigmación de morfismos:  Sea $\varphi_ A\to B $, entonces $$  j:=A \xrightarrow{ \varphi} B \xhookrightarrow{\displaystyle \iota_B} G_B .$$ Para $A$ y $j: A\to G_B$ por la propiedad universal existe un único morfismo $G_A\to G_B$, es decir, $G(\varphi)=G_A \to G_B$.\\

\noindent Por lo que finalmente tenemos 


\[
\begin{aligned}
	G:\mathbf{Set} & \longrightarrow \mathcal{C}  \\
	X & \longmapsto G_X \\
	\left(\varphi : A\to B \right)& \longmapsto \left(G(\varphi) : G_A \to G_B\right) .
\end{aligned}
\]

\vspace{1cm}
La propiedad universal del grupo libre es un caso particular de una construcción más general en el contexto de las \emph{categorías algebraicas}. Este principio no se limita a los grupos, sino que se extiende a otras estructuras algebraicas como magmas, semigrupos, monoides, anillos y módulos. En una categoría algebraica \( \mathcal{C} \), el objeto libre generado por un conjunto \( X \) satisface una propiedad universal análoga.


\begin{proposicion}[Propiedad universal de los objetos libres en categorías algebraicas]
	\upshape{
		El par \( (X, \iota) \), donde \( X \) es un conjunto y \( \iota: X \hookrightarrow F(X) \) es la función de inclusión o apliación canónica en el objeto libre \( F(X) \), cumple la siguiente propiedad universal: \\ 
		
		\noindent Para todo objeto \( M \) en \( \mathcal{C} \) y toda función \( j : X \to M \), existe un único morfismo \( f : F(X) \to M \) en \( \mathcal{C} \) tal que el siguiente diagrama conmuta:
		
		
		\[
		\xymatrix{
			X \ar[rr]^{\displaystyle j} \ar@{^(->}[dd]_{\displaystyle \iota } & & M \\
			&&\\
			F(X) \ar@{-->}[rruu]_{\displaystyle \exists !f} &&
		}
		\] Es decir, \( f \circ \iota = j \).
	}
\end{proposicion}


En la construcción de un objeto libre, la inclusión o apliación canónica \( \iota : X \hookrightarrow F(X) \) surge de forma natural como parte intrínseca del proceso que inserta los generadores en la estructura algebraica, sin imponerse externamente. En los casos ya estudiados:
\begin{itemize}
	\item En el \textbf{magma libre} \( M_X \), \( \iota : X \hookrightarrow M_X \) introduce los generadores como símbolos sin restricciones.
	\item En el \textbf{semigrupo libre} \( S_X \), \( \iota : X \hookrightarrow S_X \) respeta la asociatividad mediante una relación de equivalencia.
	\item En el \textbf{monoide libre} \( \mathrm{Mon}(X) \), \( \iota : X \hookrightarrow \mathrm{Mon}(X) \) incorpora los generadores junto con la palabra vacía como identidad.
	\item En el \textbf{grupo libre} \( G_X \), \( \iota : X \hookrightarrow G_X \) da una representación que permite inversos formales.
\end{itemize} 

	\chapter{Grupo cociente y Teorema fundamental del homomorfismo de grupos}
	
\section{Subgrupo}
\begin{definicion}
	\upshape{
		Sea \( G \) un grupo. Un subconjunto no vacío \( H \subset G \) es un \textit{subgrupo} de \( G \) si \( H \) es un grupo con la operación inducida de \( G \). Y se denota como \(	H \leq G\).
	}
\end{definicion}



\begin{proposicion}\label{propsubgrupo}
	\upshape{
		Las siguientes condiciones son equivalentes:
		\begin{enumerate}
			\item \( H \leq G \), es decir, \( H \) es un subgrupo de \( G \).
			\item \( H \) satisface las siguientes propiedades:
			\begin{itemize}
				\item \( ab \in H \) para todo \( a, b \in H \) (cerradura bajo la operación).
				\item \( 1_G \in H \) (contiene la identidad).
				\item \( a^{-1} \in H \) para todo \( a \in H \) (cerradura bajo inversos).
			\end{itemize}
			\item \( H \) verifica: para todo \( a, b \in H \), se cumple \( ab^{-1} \in H \).
		\end{enumerate}
	}
\end{proposicion}

\begin{proof}
	{\color{white} Texto en azul.}\\
	
	\noindent [\(1 \Rightarrow 2\)] Se sigue de que \( H \) es un grupo con la operación inducida de \( G \), por lo tanto cumple las propiedades de cerradura, contiene la identidad y es cerrado bajo inversos.
	
	\vspace{1ex}
	\noindent [\(2 \Rightarrow 1\)] Tenemos:
	\begin{itemize}
		\item En \( H \) se induce la operación interna de \( G \).
		\item La operación en \( H \) es asociativa, ya que lo es en \( G \).
		\item \( 1_G \in H \) actúa como identidad.
		\item Todo \( a \in H \) tiene inverso en \( H \), pues por hipótesis \( a^{-1} \in H \) y \( a a^{-1} = a^{-1} a = 1_G \).
	\end{itemize}
	
	\vspace{1ex}
	\noindent [\(2 \Leftrightarrow 3\)] La condición \( ab^{-1} \in H \) implica simultáneamente la cerradura bajo productos e inversos. Es una forma abreviada de las tres propiedades mencionadas en (2), y viceversa.
\end{proof}


\begin{ejemplo}
	\upshape{
		Sea \( \mathbb{K} = \mathbb{R} \) o \( \mathbb{C} \). Consideramos el conjunto
		\[
		M_{n \times n}(\mathbb{K}) = \left\{ A = (a_{ij}) \mid a_{ij} \in \mathbb{K},\ 1 \leq i,j \leq n \right\},
		\]
		que es el conjunto de todas las matrices cuadradas de tamaño \( n \) con entradas en \( \mathbb{K} \).\\
		
		\noindent Definimos:
		\begin{itemize}
			\item El \textbf{grupo general lineal}:
			\[
			\mathrm{GL}_n(\mathbb{K}) = \left\{ A \in M_{n \times n}(\mathbb{K}) \mid \det(A) \neq 0 \right\},
			\]
			que consiste en las matrices invertibles con entradas en \( \mathbb{K} \).
			
			\item El \textbf{grupo especial lineal}:
			\[
			\mathrm{SL}_n(\mathbb{K}) = \left\{ A \in \mathrm{GL}_n(\mathbb{K}) \mid \det(A) = 1 \right\},
			\]
			formado por aquellas matrices invertibles cuyo determinante es exactamente uno.
		\end{itemize}
		
		\vspace{1ex}
		
		Por lo tanto, se tiene la inclusión de conjuntos (y de grupos donde corresponda):
		\[
		\mathrm{SL}_n(\mathbb{K}) \leq \mathrm{GL}_n(\mathbb{K}) \subseteq M_{n \times n}(\mathbb{K}).
		\]
	}
\end{ejemplo}






\begin{ejemplo}
	\upshape{
		Sea \( \mathbb{K} = \mathbb{R} \) o \( \mathbb{C} \). Consideramos el conjunto
		\[
		M_{n \times n}(\mathbb{K}) = \left\{ A = (a_{ij}) \mid a_{ij} \in \mathbb{K},\ 1 \leq i,j \leq n \right\},
		\]
		que es el conjunto de todas las matrices cuadradas de tamaño \( n \) con entradas en \( \mathbb{K} \), con la operación natural de suma y multiplicación de matrices.\\
		
		\noindent Definimos:
		\begin{itemize}
			\item El \textbf{grupo general lineal}:
			\[
			\mathrm{GL}_n(\mathbb{K}) = \left\{ A \in M_{n \times n}(\mathbb{K}) \mid \det(A) \neq 0 \right\},
			\]
			que consiste en las matrices invertibles con entradas en \( \mathbb{K} \). Este conjunto es un grupo con la operación de multiplicación de matrices.
			
			
			
			\item El \textbf{grupo especial lineal}:
			\[
			\mathrm{SL}_n(\mathbb{K}) = \left\{ A \in \mathrm{GL}_n(\mathbb{K}) \mid \det(A) = 1 \right\},
			\] También es un grupo con la multiplicación de matrices.
		\end{itemize}
		Por lo tanto, se tiene:
		\[
		\mathrm{SL}_n(\mathbb{K}) \leq \mathrm{GL}_n(\mathbb{K}).
		\]
	}
\end{ejemplo}


\begin{ejemplo}\label{klasj}
	\upshape{
	Sea $n\in \mathbb{N}$ fijo. Entonces $n\mathbb{Z} \leq (\mathbb{Z},+)$, donde $$n\mathbb{Z}:=\left\{na: a\in\mathbb{Z}\right\}.$$
	}
\end{ejemplo}
\begin{proof}{Justificación}
	Sea $na,nb\in n\mathbb{Z}$. Tenemos que $na+(-nb)=n(a-b)\in n\mathbb{Z}.$ De la Proposición \ref{propsubgrupo} se concluye que es un subgrupo.
\end{proof}

\begin{proposicion}
	\upshape{
	Los únicos subgrupos de $(\mathbb{Z},+)$ son de la forma $n\mathbb{Z}$.
	}
\end{proposicion}


\begin{proof}
	Dado \( n \in \mathbb{N}_0 \), por el Ejemplo~\ref{klasj}, se tiene que \( n\mathbb{Z} \leq \mathbb{Z} \). Demostraremos que si \( H \leq \mathbb{Z} \), entonces existe \( n \in \mathbb{Z} \) tal que \( H = n\mathbb{Z} \). En efecto:\\
	
	\noindent Si \( H = \{0\} \), entonces \( H = n\mathbb{Z} \) para $n=0$.\\
	
	\noindent Supongamos que \( H \neq \{0\} \). Entonces existen \( x \in H \) tal que \( x \neq 0 \), y como \( H \) es subgrupo, también \( -x \in H \). Por lo tanto, \( H \) contiene enteros positivos. Definimos:
	\[
	n := \min \{ h \in H : h > 0 \}.
	\]
	Veamos que \( H = n\mathbb{Z} \).
	
	\medskip
	\noindent \textbf{[$\subset$]} Sea \( h \in H \). Por el algoritmo de la división, existen \( q, r \in \mathbb{Z} \) tales que
	\[
	h = nq + r, \quad \text{con } 0 \leq r < n.
	\]
	Como \( h, nq \in H \), entonces \( r = h - nq \in H \). Si \( r > 0 \), esto contradice la minimalidad de \( n \), ya que \( r < n \). Por tanto, \( r = 0 \), y entonces \( h = nq \in n\mathbb{Z} \). Así, \( H \subset n\mathbb{Z} \).
	
	\medskip
	\noindent \textbf{[$\supset$]} Sea \( x \in n\mathbb{Z} \), o sea, \( x = na \) con \( a \in \mathbb{Z} \). Consideramos casos:
	\begin{itemize}
		\item Si \( a = 0 \), entonces \( x = 0 \in H \).
		\item Si \( a > 0 \), entonces \( x = \underbrace{n + \cdots + n}_{a\ \text{veces}} \in H \).
		\item Si \( a < 0 \), entonces \( x = \underbrace{(-n) + \cdots + (-n)}_{-a\ \text{veces}} \in H \).
	\end{itemize}
	
	\noindent En todos los casos, \( x \in H \), luego \( n\mathbb{Z} \subset H \).
	
	\medskip
	\noindent Por lo tanto, \( H = n\mathbb{Z} \), como queríamos.
\end{proof}


\begin{proposicion}
	\upshape{
	Sea $\left\{H_i\right\}_{i\in \Lambda}$ una familia de subgrupos de $G$. Entonces $H:=\bigcap_{i\in \Lambda}H_i$ es un subgrupo de $G$.
	}
\end{proposicion}

\begin{proof}
	\upshape{
	Sean $x,y\in H$, entonces $x,y\in H_i$ para todo $i\in \Lambda$. Así $xy^{-1}\in H_i$ para todo $i\in \Lambda$. Por lo que $xy^{-1}\in H$. De la Proposición \ref{propsubgrupo} se concluye que es un subgrupo.
	}
\end{proof}



\section{Subgrupo normal}

\begin{definicion}
	\upshape{
		Sea \( G \) un grupo, \( H \leq G \) un subgrupo, y \( a \in G \). Definimos:
		
		\begin{enumerate}
			\item La \textit{clase lateral izquierda} de \( H \) que contiene a \( a \) como
			\[
			aH := \{ ah : h \in H \}.
			\]
			
			\item La \textit{clase lateral derecha} de \( H \) que contiene a \( a \) como
			\[
			Ha := \{ ha : h \in H \}.
			\]
		\end{enumerate}
	}
\end{definicion}

\begin{definicion}\label{kjhg}
	\upshape{
	Sea $H\leq G$. Definimos las siguientes relaciones 
	\begin{enumerate}
		\item $a\sim b$ si y solo si $a^{-1}b\in H$.
		\item $a\sim' b$ si y solo si $ab^{-1}\in H$.
	\end{enumerate} 
	}
\end{definicion}


\begin{proposicion}
	\upshape{
		Sea \( H \leq G \). A partir de la Definición \ref{kjhg}, se tiene:
		\begin{enumerate}
			\item La relación \( \sim \) es una relación de equivalencia en \( G \), y \( [a] = aH \).
			\item La relación \( \sim' \) es una relación de equivalencia en \( G \), y \( [a]' = Ha \).
		\end{enumerate}
	}
\end{proposicion}

\begin{proof}
	Probaremos solo $[1]$ ya que $[2]$ es análogo. En efecto, \( \sim \) es relación de equivalencia:
	\begin{itemize}
		\item Reflexiva: \( a^{-1}a = 1 \in H \), pues \( H \) es subgrupo.
		\item Simétrica: Si \( a^{-1}b \in H \), entonces \( (a^{-1}b)^{-1} = b^{-1}a \in H \), y como \( H \) es cerrado bajo inversos, se concluye que \( b \sim a \).
		\item Transitiva: Si \( a \sim b \) y \( b \sim c \), entonces \( a^{-1}b \in H \) y \( b^{-1}c \in H \). Como \( H \) es cerrado bajo multiplicación, \( a^{-1}b \cdot b^{-1}c = a^{-1}c \in H \), luego \( a \sim c \).
	\end{itemize}
	Para cada \( a \in G \), se tiene
	\[
	[a] = \{ b \in G : a^{-1}b \in H \} = \{ ah : h \in H \} = aH.
	\]
\end{proof}


\begin{definicion}
	\upshape{
		Sea \( G \) un grupo y \( H \leq G \) un subgrupo. Decimos que \( H \) es un \textit{subgrupo normal} de \( G \) si para todo \( a \in G \) se cumple
		\[
		aH = Ha.
		\]	En este caso, escribimos \( H \trianglelefteq G \).
	}
\end{definicion}


\begin{proposicion}\label{lkjhg}
	\upshape{
		Sea \( H \leq G \). Las siguientes condiciones son equivalentes:
		\begin{enumerate}
			\item \( H \trianglelefteq G \), es decir, \( aH = Ha \) para todo \( a \in G \).
			\item Para todo \( a \in G \), se cumple \( aHa^{-1} \subseteq H \).
			\item Para todo \( a \in G \), se cumple \( aHa^{-1} = H \).
			\item Para todo \( a, b \in G \), se cumple \( (aH)(bH) = abH \).
		\end{enumerate}
	}
\end{proposicion}

\begin{proof}
	{\color{white} Texto en azul.}\\
	
	\noindent[$1 \Rightarrow 2$] Sea \( a \in G \), y sea \( h \in H \). Como \( aH = Ha \), existe \( h' \in H \) tal que \( ah = h'a \). Entonces
	\[
	a h a^{-1} = h' a a^{-1} = h' \in H,
	\]
	luego \( aHa^{-1} \subseteq H \)\\
	
	\noindent	\noindent[$2 \Rightarrow 3$] Sea \( a \in G \). Como \( aHa^{-1} \subseteq H \), al aplicar la misma propiedad a \( a^{-1} \) se tiene
	\[
	a^{-1}Ha \subseteq H \Rightarrow H \subseteq aHa^{-1},
	\]
	de modo que \( H = aHa^{-1} \).\\
	
	\noindent	\noindent[$3 \Rightarrow 1$] Sea \( a \in G \) y \( h \in H \). Como \( aHa^{-1} = H \), existe \( h' \in H \) tal que
	\[
	a h = h' a \Rightarrow aH \subseteq Ha.
	\]
	El mismo argumento para \( a^{-1} \) da \( Ha \subseteq aH \), luego \( aH = Ha \).\\

	\noindent[$3 \Rightarrow 4$] Sean \( a, b \in G \). Para cualquier \( h_1, h_2 \in H \),
	\[
	(a h_1)(b h_2) = a (h_1 b) h_2 = a b (b^{-1} h_1 b)h_2 .
	\]
	Como \( H \trianglelefteq G \), se tiene \( (b^{-1} h_1 b) h_2 \in H \), así que el producto está en \( abH \). Luego \( (aH)(bH) \subseteq abH \). La inclusión contraria se prueba similarmente, por lo que \[ (aH)(bH) = abH. \]
	
	\noindent[$4 \Rightarrow 3$] Sea \( a \in G \). Como por hipótesis se cumple
	\[
	(aH)(a^{-1}H) = aa^{-1}H = H,
	\]
	entonces cada elemento de \( H \) se puede escribir como \( (ah)(a^{-1}h') \) para algunos \( h, h' \in H \). En particular,
	\[
	(ah)(a^{-1}1) = aha^{-1} \in H.
	\]
	Esto muestra que \( aHa^{-1} \subseteq H \). Aplicando el mismo argumento a \( a^{-1} \), se obtiene \( a^{-1}Ha \subseteq H \), lo que implica que \( H \subseteq aHa^{-1} \). Por tanto, \( aHa^{-1} = H \).
	
	
	
\end{proof}

\begin{proposicion}
	\upshape{
		Todo subgrupo de un grupo abeliano es un subgrupo normal.
	}
\end{proposicion}

\begin{proof}
	Sea \( G \) un grupo abeliano y \( H \leq G \). Para todo \( a \in G \) y todo \( h \in H \), se tiene
	\[
	a h a^{-1} = a a^{-1} h = h,
	\]
	ya que \( G \) es abeliano. Entonces \( a h a^{-1} \in H \), lo que muestra que \( a H a^{-1} \subseteq H \). Luego por la Proposición \ref{lkjhg} se concluye que \( H \trianglelefteq G \).
\end{proof}






\section{Grupo cociente}

Dado \( H \leq G \), la relación de equivalencia definida por 
\[
a \sim b \quad \text{si y sólo si} \quad a^{-1}b \in H
\]
induce el conjunto cociente \( G/\sim \). Si queremos definir una operación en \( G/\sim \), naturalmente esperaríamos que 
\[
(aH)(bH) := abH,
\]
pero esto no siempre es posible. \\

\noindent Consideremos 
\[
A = \begin{pmatrix}
	-\frac{\sqrt{2}}{2} & -\frac{\sqrt{2}}{2} \\
	-\frac{\sqrt{2}}{2} & \frac{\sqrt{2}}{2}
\end{pmatrix} \in G := GL_2(\mathbb{R}).
\]

\vspace{0.5cm}

\noindent Notemos que 
\[
H := \{A, I_{2\times 2}\} \leq GL_2(\mathbb{R}).
\]Tomemos 
\[
a = \begin{pmatrix}
	-1 & 0 \\
	0 & 1
\end{pmatrix} \in G.
\] Entonces, como \( a^2 = I \), tenemos 
\[
a^2 H = H,
\]
pero también se verifica que 
\[
-a \in (aH)(aH),
\]
por lo que 
\[
(aH)(aH) \neq a^2 H.
\] Por lo tanto, vemos que la relación de equivalencia definida no siempre induce una estructura de grupo cociente. Para que esta operación esté bien definida y sea un grupo, es necesario que \( H \trianglelefteq G \), donde se cumple que \(
(aH)(bH) = abH.
\) Como se mostró en la Proposición \ref{lkjhg}.

\begin{definicion}
	\upshape{
		Sea \( G \) un grupo y \( H \trianglelefteq G \). La operación 
		\[
		(aH)(bH) := abH, \quad \forall\, a,b \in G
		\]
		define en \( G/\sim \) una estructura de grupo. Este grupo se denomina \emph{grupo cociente} y se denota por \( G/H \).
	}
\end{definicion}
\begin{observacion}
	\upshape{
		En \( G/H \):
		\begin{itemize}
			\item La operación interna está bien definida, es decir, no depende del representante de cada clase.
			\item El elemento identidad es la clase de \( 1 \in G \); es decir, \(1_{G/H}= 1_G H \).
			\item El inverso de \( aH \) es \( a^{-1}H \).
		\end{itemize}
	}
\end{observacion}



\section{Teorema fundamental de homomorfismo de grupos}


\begin{definicion}
	\upshape{
		Sea \( f: G \to H \) un homomorfismo de grupos. Definimos:
		\begin{enumerate}
			\item El \textit{núcleo} de \( f \) como
			\[
			\ker f := \{ a \in G : f(a) = 1_H \}.
			\]
			\item La \textit{imagen} de \( f \) como
			\[
			\operatorname{Im} f := \{ h \in H : h = f(a) \text{ para algún } a \in G \}.
			\]
			\item La \textit{proyección canónica} de un grupo sobre su cociente: Si \(N\) es un subgrupo normal de un grupo \(G\), la proyección canónica es la función
			\[
			\pi : G \to G/N, \quad a \mapsto aN.
			\]
			Esta función \(\pi\) es un homomorfismo de grupos sobreyectivo, cuyo núcleo es \(N\).
		\end{enumerate}
	}
\end{definicion}

\begin{observacion}\label{Obs: imagen y núcleo}
	\upshape{
		El núcleo y la imagen de un homomorfismo de grupos son subgrupos, es decir,  
		\[
		\ker f \leq G \quad \text{y} \quad \operatorname{Im} f \leq H.
		\]
	}
\end{observacion}


\begin{proposicion}\label{oikhugfd}
	\upshape{
		Sea \( f: G \to H \) un homomorfismo de grupos. Entonces:
		\begin{enumerate}
			\item \( \ker f \trianglelefteq G \).
			\item \( f \) es inyectivo si y solo si \( \ker f = \{1_G\} \).
		\end{enumerate}
	}
\end{proposicion}


\begin{proof}
	En efecto.
	\begin{enumerate}
		\item Sea \(x,y \in \ker f\), entonces \(f(x) = f(y) = 1_H\). Así,
		\[
		f(x y^{-1}) = f(x) f(y)^{-1} = 1_H \cdot 1_H^{-1} = 1_H.
		\]
		Por lo tanto, \(x y^{-1} \in \ker f\), luego \(\ker f \leq G\). Ahora verifiquemos que \(\ker f\) es normal. Sea \(a \in G\), \(k \in \ker f\), entonces
		\[
		f(a k a^{-1}) = f(a) f(k) f(a)^{-1} = f(a) \cdot 1_H \cdot f(a)^{-1} = 1_H,
		\]
		por lo que \(a (\ker f) a^{-1} \subseteq \ker f\). Así, \(\ker f \trianglelefteq G\).
		
		\item Veamos:
		
		\noindent \([ \Rightarrow ]\) Sea \(x \in \ker f\). Entonces
		\[
		f(x) = f(1_G) = 1_H.
		\]
		Como \(f\) es inyectiva, entonces \(x = 1_G\). Así, \(\ker f = \{1_G\}\).
		
		\noindent \([ \Leftarrow ]\) Sean \(x,y \in G\) tales que \(f(x) = f(y)\). Entonces,
		\[
		f(x y^{-1}) = f(x) f(y)^{-1} = 1_H.
		\]
		Por lo que \(x y^{-1} \in \ker f = \{1_G\}\), entonces \(x y^{-1} = 1_G\). Así, \(x = y\). Esto implica que \(f\) es inyectiva.
	\end{enumerate}
\end{proof}



\begin{teorema}[Propiedad universal del cociente]\label{Propiedad universal del cociente}
	\upshape{
		Sea \( f: G \to H \) un homomorfismo de grupos y sea \( N \trianglelefteq G \) tal que \( N \subseteq \ker f \). Entonces, existe un único homomorfismo de grupos \( \overline{f} : G / N \to H \) tal que el siguiente diagrama conmuta:
		
		
		
		\[
		\xymatrix{
			G \ar[rr]^{\displaystyle f} \ar[dd]_{\displaystyle \pi} && H \\
			&&\\
			G/N \ar@{-->}[uurr]_{\displaystyle \exists!\:\overline{f}} &&
		}
		\]
		
		
		Es decir, \( \overline{f} \circ \pi = f \), donde \( \pi : G \to G/N \) es la proyección canónica.
	}
\end{teorema}

\begin{proof}
	\upshape{
		Queremos definir un único homomorfismo \( \overline{f}: G/N \to H \) tal que \( \overline{f} \circ \pi = f \):
		
		\begin{itemize}
			\item \textbf{Definición:} Definimos $\overline{f}(aN) := f(a)$, para todo $ a \in G.$
			
			
			
			\item \textbf{Buena definición:} Si \( aN = bN \), entonces \( a^{-1}b \in N \subseteq \ker f \), así que \( f(a^{-1}b) = 1_H \), lo que implica \( f(a) = f(b) \). Por tanto, \( \overline{f}(aN) = \overline{f}(bN) \), y la aplicación está bien definida.
			
			\item \textbf{Homomorfismo:} Sean \( aN, bN \in G/N \). Entonces:
			\begin{align*}
				\overline{f}((aN)(bN)) &= \overline{f}(abN) = f(ab) = f(a)f(b) = \overline{f}(aN)\cdot \overline{f}(bN).
			\end{align*}
			Así, \( \overline{f} \) es homomorfismo.
			
			\item \textbf{Conmutatividad:} Para todo \( a \in G \),
			
			
			\[
			(\overline{f} \circ \pi)(a) = \overline{f}(aN) = f(a),
			\]
			
			
			por lo que \( \overline{f} \circ \pi = f \).
			
			\item \textbf{Unicidad:} Si \( \psi: G/N \to H \) es otro homomorfismo tal que \( \psi \circ \pi = f \), entonces para todo \( aN \in G/N \), se cumple
			
			
			\[
			\psi(aN) = \psi(\pi(a)) = f(a) = \overline{f}(aN),
			\]
			
			
			de donde \( \psi = \overline{f} \).
		\end{itemize}
	}
\end{proof}

\begin{corolario} \label{coikjhg}
	\upshape{
		En el Teorema \ref{Propiedad universal del cociente}. Si \( N = \ker f \), entonces \( \overline{f}:G/N \to H \) es inyectivo.
	}
\end{corolario}

\begin{proof}
	\upshape{
		Analizamos el núcleo de \( \overline{f} \):
		
		
		\[
		\ker \overline{f} = \{ aN \in G/N : \overline{f}(aN) = 1_H \} = \{ aN : f(a) = 1_H \} = \{aN: a\in \ker f \}.
		\]
		
		
		Como \( \ker f = N \), entonces
		
		
		\[
		\ker \overline{f} =\{aN: a\in N\}=\{N\},
		\]
		
		
		es decir, \( \overline{f} \) tiene núcleo trivial, por lo tanto es inyectivo.
		
		\vspace{-0.5em}
		\qedhere
	}
\end{proof}










\begin{teorema}[Teorema fundamental del homomorfismo de grupos]
	\upshape{
		Sean \( G \) y \( H \) grupos, y sea \( f: G \to H \) un homomorfismo de grupos. Entonces, \( f \) induce un isomorfismo de grupos
		
		
		\[
		\frac{G}{\ker(f)} \cong \operatorname{Im}(f).
		\]
		
		
	}
\end{teorema}

\begin{proof}
	En el Teorema \ref{Propiedad universal del cociente} consideremos \( N=\ker f \), por lo que existe el homomorfismo \( \overline{f}: G/\ker f \to H \). Así,  \( \overline{f}: G/\ker f \to \operatorname{Im}(f) \) es sobreyectivo. Además, \( \overline{f} \) es inyectivo por el Corolario \ref{coikjhg}. Por lo tanto, \( \overline{f} \) es un isomorfismo de grupos.
\end{proof}


\begin{ejemplo}
	\upshape{
		Considere el homomorfismo \( f: \left(\mathbb{Z},+\right) \to \left(\mathbb{Z}_n,+\right) \)
		dado por \( f(a) := [a] \) para todo \( a \in \mathbb{Z} \). Entonces, como \( f \) es sobreyectivo y	\[
		\ker f = \{ a \in \mathbb{Z} : [a] = [0] \} = \{ a \in \mathbb{Z} : n \mid a \} = n\mathbb{Z},
		\]por el Teorema fundamental del homomorfismo de grupos, tenemos que
		\[
		\dfrac{\mathbb{Z}}{n\mathbb{Z}} \cong \mathbb{Z}_n.
		\]
	}
\end{ejemplo}

\begin{ejemplo}
	\upshape{
		Considere el homomorfismo \( f: \left(GL_n(\mathbb{R}), \cdot \right) \to \left(\mathbb{R}^\ast, \cdot \right) \)
		definido por \( f(A) := \det(A) \) para toda matriz \( A \in GL_n(\mathbb{R}) \). Entonces, como \( f \) es sobreyectivo y  
		
		\[
		\ker f = \{ A \in GL_n(\mathbb{R}) \mid \det(A) = 1_{\mathbb{R}} \} = SL_n(\mathbb{R}),
		\]
		
		por el Teorema fundamental del homomorfismo de grupos, tenemos que
		
		
		\[
		\dfrac{GL_n(\mathbb{R})}{SL_n(\mathbb{R})} \cong \mathbb{R}^\ast.
		\]
		
		
	}
\end{ejemplo}

\begin{ejemplo}
	\upshape{
		Consideremos el grupo multiplicativo del círculo unitario 
		\[
		S^1 := \{ z \in \mathbb{C} : |z| = 1 \},
		\]
		que es un grupo bajo la multiplicación compleja.  \\
		
		\noindent Considere el homomorfismo \( f: \left(\mathbb{R}, + \right) \to \left(S^1, \cdot \right) \)
		dado por 	\[
		f(x) := e^{2\pi i x}, \quad \text{para todo } x \in \mathbb{R}.
		\] Entonces, como \( f \) es sobreyectivo y  
		\[
		\ker f = \{ x \in \mathbb{R} \mid e^{2\pi i x} = 1 \} = \mathbb{Z},
		\]por el Teorema fundamental del homomorfismo de grupos, tenemos que

		\[
		\dfrac{\mathbb{R}}{\mathbb{Z}} \cong S^1.
		\]

	}
\end{ejemplo}

\begin{observacion}
	\upshape{
		De manera similar, se tiene la siguiente isomorfía:		
		\[
		\dfrac{\mathbb{R}^n}{\mathbb{Z}^n} \cong \mathbb{T}^n := \underbrace{S^1 \times \cdots \times S^1}_{n\ \text{veces}}.
		\]
		
	}
\end{observacion}






	\chapter{Construcción de Grothendieck de un grupo a partir de un monoide conmutativo}
	
Sea \( M \) un monoide conmutativo. Buscamos construir un grupo conmutativo \( M^+ \) junto con un homomorfismo de monoides \( i: M \to M^+ \) que satisfaga una propiedad universal característica. Este grupo, conocido como \emph{grupo de Grothendieck} asociado a \( M \), permite dotar al monoide de una estructura completamente invertible extendiéndolo a un grupo que satisface una propiedad universal sin alterar sus propiedades fundamentales.

\begin{proposicion}[Propiedad universal del grupo de Grothendieck]
	\upshape{
		Para todo grupo conmutativo \( A \) y todo homomorfismo de monoides \( j : M \to A \), existe un único homomorfismo de grupos \( f : M^+ \to A \) tal que el siguiente diagrama conmuta:
		\[
		\xymatrix{
			M \ar[rr]^{\displaystyle j} \ar[dd]_{\displaystyle i } & & A \\
			&&\\
			M^+ \ar@{-->}[rruu]_{\displaystyle \exists !f} &&
		}
		\]Es decir, \( f \circ i = j \).
	}
\end{proposicion}

Demostraremos la propiedad universal después de definir \( M^+ \) y el homomorfismo \( i: M \to M^+ \). Por lo pronto, podemos establecer que cualquier otro grupo que satisfaga esta propiedad universal es isomorfo a \( M^+ \), garantizando así su unicidad.




\begin{corolario}
	\upshape{
		Sea \( (G, \varphi) \) otro par que cumple la propiedad universal del grupo de Grothendieck para \( M \). \\
		
		\noindent Entonces existe un único isomorfismo de grupos \( \psi : M^+ \to G \) tal que el siguiente diagrama conmuta:
		\[
		\xymatrix{
			&& M \ar[lld]_{\displaystyle i} \ar[rrd]^{\displaystyle \varphi} && \\
			M^+ \ar[rrrr]_{\displaystyle \exists!\,\psi}^{\displaystyle \cong} &&&& G
		}
		\]
	}
\end{corolario}


\begin{proof}
	Aplicamos la propiedad universal de \( M^+ \) al morfismo \( \varphi : M \to G \). Como \( G \) es un grupo conmutativo, existe un único homomorfismo de grupos \( \psi : M^+ \to G \) tal que el siguiente diagrama conmuta:
	\[
	\xymatrix{
		M \ar[rr]^{\displaystyle \varphi} \ar[dd]_{\displaystyle i} & & G \\
		&&\\
		M^+ \ar@{-->}[rruu]_{\displaystyle \exists! \psi} &&
	}
	\]	Recíprocamente, aplicamos la propiedad universal de \( G \) al morfismo \( i : M \to M^+ \), lo que garantiza la existencia de un único homomorfismo de grupos \( \theta : G \to M^+ \) tal que:
	\[
	\xymatrix{
		M \ar[rr]^{\displaystyle i} \ar[dd]_{\displaystyle \varphi} & & M^+ \\
		&&\\
		G \ar@{-->}[rruu]_{\displaystyle \exists! \theta} &&
	}
	\]
	Entonces,
	\[
	\theta \circ \psi \circ i = \theta \circ \varphi = i,
	\]
	Ahora, aplicamos la propiedad universal de \( M^+ \) al morfismo \( i : M \to M^+ \). Tanto \( \mathrm{id}_{M^+} \) como \( \theta \circ \psi \) son homomorfismos de grupos que hacen conmutar el siguiente diagrama:
	\[
	\xymatrix{
		M \ar[rrrr]^{\displaystyle i} \ar[dd]_{\displaystyle i} & &&& M^+ \\
		&&&&\\
		M^+ \ar@{-->}@/^1pc/[rrrruu]^{\displaystyle \mathrm{id}_{M^+}} \ar@{-->}@/_1pc/[rrrruu]_{\displaystyle \theta \circ \psi} &&&&
	}
	\] Por la unicidad del morfismo de grupos que hace conmutar el triángulo, se concluye que \( \theta \circ \psi = \mathrm{id}_{M^+} \).\\
	
	\noindent Análogamente, se prueba que \( \psi \circ \theta = \mathrm{id}_G \). Por lo tanto, \( \psi \) es un isomorfismo de grupos.
\end{proof}






\section{Construcción cuando \( M = (\mathbb{N}, +) \)}

\noindent Recordemos que \( \mathbb{N} := \{0; 1; 2; 3; \cdots\} \), y que en \( (\mathbb{N}, +) \) se cumple la \emph{propiedad cancelativa}: para todos \( a, b, c \in \mathbb{N} \), si \( a + c = b + c \), entonces \( a = b \).\\

\noindent Sobre \( \mathbb{N} \times \mathbb{N} \) definimos la relación:
\[
(a,b) \sim (c,d) \quad \text{si y solo si} \quad a + d = b + c.
\]Esta relación es una relación de equivalencia:
\begin{itemize}
	\item Reflexiva: \( a + b = b + a \), luego \( (a,b) \sim (a,b) \).
	\item Simétrica: Si \( (a,b) \sim (c,d) \), entonces \( a + d = b + c \), por lo tanto \( c + b = d + a \), y así \( (c,d) \sim (a,b) \).
	\item Transitiva: Si \( (a,b) \sim (c,d) \) y \( (c,d) \sim (e,f) \), entonces \( a + d = b + c \) y \( c + f = d + e \). Sumando y simplificando se obtiene \( a + f = b + e \). Así, \( (a,b) \sim (e,f) \).
\end{itemize}
\noindent Por lo tanto, podemos definir
\[
\mathbb{N}^{+} := \dfrac{\mathbb{N} \times \mathbb{N}}{\sim}.
\]
Y sobre \( \mathbb{N^+} \) definimos la operación
\[
[(a,b)] + [(c,d)] := [(a+c,\, b+d)].
\]
Esta operación está bien definida, pues si \( (a,b) \sim (a',b') \) y \( (c,d) \sim (c',d') \), entonces se tiene:
\[
a + b' = a' + b \quad \text{y} \quad c + d' = c' + d.
\]
Sumando ambas igualdades:
\[
(a + c) + (b' + d') = (a' + c') + (b + d),
\]
lo cual implica que \( (a + c,\, b + d) \sim (a' + c',\, b' + d') \), es decir:
\[
[(a+c,\, b+d)] = [(a'+c',\, b'+d')].
\]
Por lo tanto, la suma está bien definida sobre las clases de equivalencia.

\begin{proposicion}
	\upshape{
		\( (\mathbb{N^+}, +) \) es un grupo conmutativo.
	}
\end{proposicion}

\begin{proof}
	En efecto.
	\begin{itemize}
		\item La operación es una operación interna.
		\item La operación es conmutativa.
		\item El elemento identidad es \( [(0,0)] \).
		\item Dado \( [(a,b)] \in \mathbb{Z} \), su inverso es \( -[(a,b)] := [(b,a)] \).
	\end{itemize}
\end{proof}



\begin{proposicion}
	\upshape{
		Se cumple la isomorfía \( (\mathbb{Z}, +) \cong (\mathbb{N}^+, +) \).
	}
\end{proposicion}

\begin{proof}
	El grupo \((\mathbb{Z},+)\) junto con el morfismo \(i : \mathbb{N} \to \mathbb{Z}\), \(i(x) = x\), satisface la propiedad universal del grupo de Grothendieck del monoide \((\mathbb{N},+)\). En efecto, sea \(A\) un grupo abeliano y \(j : \mathbb{N} \to A\) un morfismo de monoides. Definimos
	\[
	f : \mathbb{Z} \to A,\quad
	f(n) :=
	\begin{cases}
		j(n), & \text{si } n \geq 0,\\
		-j(-n), & \text{si } n < 0.
	\end{cases}
	\] Probamos que \(f\) es un homomorfismo de grupos. Sea \(m,n \in \mathbb{Z}\). Consideramos los siguientes casos:
	
	\begin{itemize}
		\item \textbf{Caso 1: \(m,n \geq 0\)}. Entonces \(f(m+n) = j(m+n) = j(m) + j(n) = f(m) + f(n)\), pues \(j\) es morfismo de monoides.
		
		\item \textbf{Caso 2: \(m,n < 0\)}. Escribimos \(m = -a\), \(n = -b\) con \(a,b \in \mathbb{N}\). Entonces
		\[
		f(m+n) = f(-(a+b)) = -j(a+b) = -j(a) - j(b) = f(m) + f(n).
		\]
		
		\item \textbf{Caso 3: \(m \geq 0\), \(n < 0\)}. Escribimos \(n = -b\) con \(b \in \mathbb{N}\). Si \(m \geq b\), entonces \(m + n = m - b \geq 0\), y
		\[
		f(m + n) = j(m - b) = j(m) - j(b) = f(m) + f(n).
		\]
		Si \(m < b\), entonces \(m + n = -(b - m) < 0\), y
		\[
		f(m + n) = -j(b - m) = j(m) - j(b) = f(m) + f(n).
		\]
		
		\item \textbf{Caso 4: \(m < 0\), \(n \geq 0\)}. Similar al caso anterior. Se obtiene \(f(m+n) = f(m) + f(n)\).
	\end{itemize}
	
Por lo tanto, \(f\) es homomorfismo de grupos. Además, por definición,
\[
f \circ i(x) = f(x) = j(x),\quad \text{para todo } x \in \mathbb{N},
\]
es decir, \(f \circ i = j\).\\

\noindent \textit{Unicidad:} Supongamos que \(g : \mathbb{Z} \to A\) es un homomorfismo de grupos tal que \(g \circ i = j\), es decir, \(g(n) = j(n)\) para todo \(n \in \mathbb{N}\). Probamos que \(g(n) = f(n)\) para todo \(n \in \mathbb{Z}\), por casos:

\begin{itemize}
	\item Si \(n \geq 0\), entonces \(g(n) = j(n) = f(n)\).
	
	\item Si \(n < 0\). Como \(g\) es homomorfismo,
	\[
	0 = g(0) = g(n +(-n)) = g(n) + g(-n),
	\]
	por lo que \(g(n) = -g(-n) = -j(-n) = f(n)\).
\end{itemize} Así, \(g(n) = f(n)\) para todo \(n \in \mathbb{Z}\), y por lo tanto \(g = f\).\\

\noindent Concluimos que \(\mathbb{Z}\) satisface la propiedad universal del grupo de Grothendieck del monoide \(\mathbb{N}\), y por lo tanto, \(\mathbb{Z} \cong \mathbb{N}^+\).
\end{proof}




\section{Construcción general}

\noindent Definir en \( M \times M \) una relación de equivalencia como en el caso anterior no garantiza una relación de equivalencia, pues puede fallar la transitividad. Esto sucede porque el monoide \( M \) no necesariamente tiene la propiedad cancelativa.\\

\noindent Sea $M := \{0,1,2\} $
con la operación \(\min\). La operación \( \min \) no es cancelativa, pues \( \min(a, x) = \min(b, x) \) no implica \( a = b \), como se observa en \( \min(0,0) = \min(1,0) \). Definimos en \( M \times M \) la relación:
\[
(a,b) \sim (c,d) \quad \text{cuando} \quad \min(a,d) = \min(b,c).
\]Consideremos los pares:
\[
(0,1) \sim (0,0) \quad \text{ya que} \quad \min(0,0) = 0 = \min(1,0),
\]
\[
(0,0) \sim (1,0) \quad \text{ya que} \quad \min(0,0) = 0 = \min(0,1);
\]
pero
\[
(0,1) \not\sim (1,0) \quad \text{porque} \quad \min(0,1) = 0 \neq 1 = \min(1,0).
\]Por lo tanto, la relación no es transitiva y no es una relación de equivalencia en \( M \times M \).\\

Sea \(M\) un monoide conmutativo. Para ``forzar'' la transitividad definimos en \(M \times M\) la relación:
\[
(a,b) \sim (c,d) \quad \text{si y solo si} \quad \exists\, p \in M \text{ tal que } a + d + p = b + c + p.
\]

\begin{proposicion}
	\upshape{
		La relación \(\sim\) en \(M \times M\) es una relación de equivalencia.
	}
\end{proposicion}

\begin{proof}
	En efecto.
	\begin{itemize}
		\item Reflexiva: Para todo \((a,b) \in M \times M\), usando la conmutatividad de \(M\),
		\[
		a + b + p = b + a + p,
		\]
		para cualquier \(p \in M\). Por lo tanto, \((a,b) \sim (a,b)\).
		
		\item Simétrica: Si \((a,b) \sim (c,d)\), entonces existe \(p \in M\) tal que
		\[
		a + d + p = b + c + p.
		\]
		Por la conmutatividad de \(M\), esto implica
		\[
		c + b + p = d + a + p,
		\]
		y así \((c,d) \sim (a,b)\).
		
		\item \textbf{Transitiva:} Si \((a,b) \sim (c,d)\) y \((c,d) \sim (e,f)\), existen \(p,q \in M\) tales que
		\[
		a + d + p = b + c + p, \quad c + f + q = d + e + q.
		\]
		Sumando estas dos igualdades y reordenando términos, se obtiene
		\[
		a + f + (d + c + p + q) = b + e + (c + d + p + q),
		\]
		lo que implica \((a,b) \sim (e,f)\).
	\end{itemize}
\end{proof}

\begin{definicion}
	\upshape{
		Sea \(M\) un monoide conmutativo. Definimos el grupo conmutativo
	\[
	M^+ := (M \times M) / \sim,
	\]
	donde la relación de equivalencia \(\sim\) en \(M \times M\) está dada por
	\[
	(a,b) \sim (c,d) \quad \Longleftrightarrow \quad \exists\, p \in M \text{ tal que } a + d + p = b + c + p.
	\]
	El conjunto cociente \(M^+\) se llama el \emph{grupo de Grothendieck} asociado a \(M\). La operación en \(M^+\) se define por
	\[
	[(a,b)] + [(c,d)] := [(a+c, b+d)],
	\]
	y el elemento identidad está dado por  \([(e,e)]\), siendo \(e\) el elemento identidad de \(M\).
	}
\end{definicion}

\begin{observacion}
	\upshape{
		En notación multiplicativa, la relación se escribe \[ (a,b) \sim (c,d) \iff \exists\, p \in M \text{ tal que } a d p = b c p ,\] la operación es \( [(a,b)] \cdot [(c,d)] := [(a c, b d)] \).
	}
\end{observacion}










\section{Propiedad universal del grupo de Grothendieck}
\noindent Sea \( M \) un monoide conmutativo y \( i: M\to M^+ \) definido por \( i(a)=[(a,0)] \). Se cumple:\\

\noindent Para todo grupo conmutativo \( A \) y todo homomorfismo de monoides \( j : M \to A \), existe un único homomorfismo de grupos \( f : M^+ \to A \) tal que el siguiente diagrama conmuta:
\[
\xymatrix{
	M \ar[rr]^{\displaystyle j} \ar[dd]_{\displaystyle i } & & A \\
	&&\\
	M^+ \ar@{-->}[rruu]_{\displaystyle \exists !f} &&
}
\]
Es decir, \( f \circ i = j \).



\begin{proof}
	En efecto.
	\begin{itemize}
		\item Definimos \( f : M^+ \to A \) por
		\[
		f([(a,b)]) := j(a) - j(b).
		\]
		
		\item Comprobamos que \( f \) está bien definida: si \( (a,b) \sim (c,d) \), entonces existe \( p \in M \) tal que
		\[
		a + d + p = b + c + p.
		\]
		Aplicando \( j \) que es homomorfismo, se tiene
		\[
		j(a) + j(d) + j(p) = j(b) + j(c) + j(p),
		\]
		y al cancelar \( j(p) \), se obtiene \( j(a) + j(d) = j(b) + j(c) \), lo cual implica
		\[
		j(a) - j(b) = j(c) - j(d),
		\]
		así que \( f([(a,b)]) = f([(c,d)]) \).
		
		\item Verificamos que \( f \) es homomorfismo: dados \( [(a,b)], [(c,d)] \in M^+ \),
		\[
		f([(a,b)] + [(c,d)]) = f([(a+c, b+d)]) = j(a+c) - j(b+d) = j(a) + j(c) - j(b) - j(d),
		\]
		que coincide con
		\[
		f([(a,b)]) + f([(c,d)]).
		\]
		
		\item Verificamos que \( f \circ i = j \): para todo \( a \in M \),
		\[
		f(i(a)) = f([(a,0)]) = j(a) - j(0) = j(a),
		\]
		pues \( j(0) = 0 \) ya que \( j: M \to A \) es un homomorfismo de monoides y \( A \) es grupo. En efecto, al tener inversos, se tiene:
		\[
		j(0) + j(0) = j(0 + 0) = j(0) \quad \Rightarrow \quad j(0) = 0.
		\]
		
		
		\item Unicidad: si \( g : M^+ \to A \) es otro homomorfismo tal que \( g \circ i = j \), entonces para todo \( [(a,b)] \in M^+ \), se tiene
		\[
		[(a,b)] = [(a,0)] - [(b,0)] = i(a) - i(b),
		\]
		luego
		\[
		g([(a,b)]) = g(i(a)) - g(i(b)) = j(a) - j(b) = f([(a,b)]),
		\]
		por lo tanto \( g = f \).
	\end{itemize}
\end{proof}

\begin{observacion}
	\upshape{
		La demostración se hizo usando notación aditiva, pero es completamente análoga para la notación multiplicativa, cambiando suma por producto y el elemento neutro \(0\) por la unidad \(1\).
	}
\end{observacion}
	
\begin{proposicion}
	\upshape{
		Sea \(M\) un monoide conmutativo y \(i : M \to M^+\) el homomorfismo natural definido por \(i(a) = [(a,0)]\). Entonces, \(i\) es inyectivo si y solo si \(M\) es cancelativo, es decir,
		\[
		a + c = b + c \implies a = b \quad \text{para todo } a,b,c \in M.
		\]
	}
\end{proposicion}

\begin{proof}
	{\color{white} blanco.}\\
	
	\noindent \([ \Rightarrow ]\) Sea \(a + c = b + c\). Entonces como $j$ es un homomorfismo
	\[
	i(a) + i(c) = i(a + c) = i(b + c) = i(b) + i(c),
	\]
	luego coomo $M^{+}$ es grupo tenemos  \(i(a) = i(b)\), y por inyectividad de \(i\), se concluye \(a = b\).\\
	
	\noindent \([ \Leftarrow ]\) Sea \(i(a) = i(b)\), es decir, \([(a,0)] = [(b,0)]\), por lo que existe \(p \in M\) tal que \(a + p = b + p\). Por cancelatividad de \(M\), se deduce \(a = b\).
\end{proof}


\begin{observacion}
	\upshape{
		La construcción de Grothendieck permite ``forzar'' la existencia de inversos en un monoide conmutativo \(M\), generando así un grupo \(M^+\). Esta idea aparece en varias áreas del álgebra. Por ejemplo, en \emph{K-teoría}, se considera el monoide \( \mathrm{Vect}(X) \) formado por las clases de isomorfismo de fibrados vectoriales sobre un espacio topológico \( X \), con la operación suma directa. Su grupo de Grothendieck \( K_0(X) := \mathrm{Vect}(X)^+ \) recoge información topológica importante del espacio. Así, \( K_0(X) \) es una herramienta algebraica construida a partir de datos geométricos.
	}
\end{observacion}













\section{Construcción del grupo fundamental o universal de un monoide}

En esta construcción, no asumiremos que el monoide \( M \) sea conmutativo. Esta es una diferencia esencial respecto al grupo de Grothendieck, que solo se define para monoides conmutativos y produce un grupo abeliano. En particular, la definición de la relación de equivalencia utilizada en \( M^+ \) requiere conmutatividad, lo cual no será necesario aquí.

\medskip
\noindent El grupo resultante se llama \emph{grupo fundamental} o \emph{grupo universal} de un monoide. Recibe este nombre porque satisface una propiedad universal análoga a la del grupo de Grothendieck. Así, proporciona la extensión más general de \( M \) a un grupo, dotando a sus elementos de inversos sin imponer conmutatividad.

\medskip
\noindent Esta construcción no debe confundirse con el grupo fundamental de la topología. Aquí, ``fundamental'' se refiere a su carácter universal en álgebra, y no a consideraciones geométricas.\\






  
  
  
  Sea \( M \) un monoide. Consideramos \( M \) simplemente como un conjunto no vacío, olvidando momentáneamente su estructura de monoide.
  
  \medskip
  
  \noindent Sea \( G_M \) el grupo libre generado por \( M \) y \( \rho : M \to G_M \) la inclusión canónica de generadores. Para reflejar la operación del monoide \( M \) en \( G_M \), consideramos el subgrupo normal \( N \) de \( G_M \) generado por los elementos de la forma
  \[
  \rho(a) \bullet \rho(b) \bullet \rho(ab)^{-1}, \quad \text{para todo } a, b \in M.
  \]
  
  \noindent Definimos entonces el grupo cociente
  \[
  U(M) := G_M / N,
  \]
  y escribimos también \( \iota : M \to U(M) \) por  composición de la inclusión $\rho$ con el homomorfismo cociente de grupos $\pi$: 
  
  
  \[
  \xymatrix{
  	M \ar[rr]^{\displaystyle\rho} \ar@/_1.5pc/[rrrr]_{\displaystyle \iota:=\pi \circ \rho} && G_M \ar[rr]^{\displaystyle \pi} && U(M)
  }
  \]
 \noindent En efecto, veamos que \( \iota := \pi \circ \rho \) es un homomorfismo de monoides. Sean \( a, b \in M \). Como \( \rho \) es la inclusión de generadores en el grupo libre \( G_M \), y \( \pi \) respeta la relación
 \[
 \rho(a) \bullet \rho(b) \sim \rho(ab),
 \]
 se tiene:
 \[
 \iota(a \cdot b) = \pi(\rho(a \cdot b)) = \pi(\rho(a) \bullet \rho(b)) = \pi(\rho(a)) \cdot \pi(\rho(b)) = \iota(a) \cdot \iota(b).
 \]
 Por lo tanto, \( \iota \) es un homomorfismo de monoides.
  
  
  \begin{definicion}
  	\upshape{
  		El grupo \( U(M) \) se denomina \textit{grupo fundamental} o \textit{grupo universal} del monoide \( M \).
  	}
  \end{definicion}
  
  \noindent El grupo \( U(M) \) viene equipado con un morfismo natural de monoides \( \iota : M \to U(M) \), y satisface la siguiente propiedad universal:
  
  \begin{proposicion}[Propiedad universal del grupo universal de un monoide]
  	\upshape{
  		El par $(U(M),\iota)$ satisface lo siguiente:
  		
  		\medskip
  		\noindent Para todo grupo \( G \) y todo morfismo de monoides \( j : M \to G \), existe un único morfismo de grupos \( f : U(M)\to G \) tal que el siguiente diagrama conmuta:
  		
  		
  		\[
  		\xymatrix{
  			M \ar[rr]^{\displaystyle j} \ar[dd]_{\displaystyle \iota} & & G \\
  			&&\\
  			U(M) \ar@{-->}[rruu]_{\displaystyle \exists !f} &&
  		}
  		\]Es decir, \( f \circ \iota = j \).
  	}
  \end{proposicion}
  

\begin{proof}
	En efecto:
	\begin{itemize}
		\item Sea \( j : M \to G \) un morfismo de monoides hacia un grupo \( G \). Como \( G_M \) es el grupo libre generado por el conjunto \( M \), existe un único morfismo de grupos
		\[
		\tilde{f} : G_M \to G
		\]
		tal que \( \tilde{f} \circ \rho = j \), donde \( \rho : M \to G_M \) es la inclusión de generadores.
		
		\item Mostramos que \( \tilde{f} \) anula al subgrupo normal \( N \trianglelefteq G_M \) generado por los elementos
		\[
		\rho(a)\bullet \rho(b)\bullet \rho(ab)^{-1}, \quad \text{para } a, b \in M.
		\]
		En efecto, para todo \( a, b \in M \), se tiene:
		\[
		\tilde{f}(\rho(a)\bullet \rho(b)\bullet \rho(ab)^{-1}) = j(a) \cdot j(b) \cdot j(ab)^{-1}.
		\]
		Pero como \( j \) es un morfismo de monoides, se cumple \( j(ab) = j(a)j(b) \), así que:
		\[
		j(a)j(b)j(ab)^{-1} = j(ab)j(ab)^{-1} = 1_G.
		\]
		Por lo tanto, \( \tilde{f} \) se anula en los generadores de \( N \), y como es homomorfismo de grupos, también en todo \( N \). Así, \( N \subseteq \ker(\tilde{f}) \), y \( \tilde{f} \) induce un morfismo cociente.
		
		\item Por la propiedad universal del cociente, existe un único morfismo de grupos
		\[
		f : U(M) = G_M/N \to G
		\]
		tal que \( f \circ \pi = \tilde{f} \), donde \( \pi : G_M \to U(M) \) es el morfismo cociente.
		
		\item Verificamos que \( f \circ \iota = j \): recordemos que \( \iota := \pi \circ \rho \). Entonces:
		\[
		f \circ \iota = f \circ \pi \circ \rho = \tilde{f} \circ \rho = j.
		\]
		
		\item Unicidad: supongamos que \( g : U(M) \to G \) es otro morfismo de grupos tal que \( g \circ \iota = j \). Entonces:
		\[
		g \circ \pi \circ \rho = g \circ \iota = j = \tilde{f} \circ \rho.
		\]
		Notemos que \( \rho(M) \) genera \( G_M \) y todos los mapas involucrados son homomorfismos de grupos. Para un elemento \( a \in G_M \), como los elementos de \( G_M \) son concatenaciones de generadores y sus inversos, podemos suponer que \( a = \rho(x)\bullet\rho(y)^{-1} \) (caso representativo). Entonces:		
		\begin{align*}
			(g\circ\pi)(a) 
			&= (g\circ\pi)(\rho(x)\bullet\rho(y)^{-1}) \\
			&= g\left(\pi(\rho(x)) \cdot \pi(\rho(y))^{-1}\right) \\
			&= g\left(\iota(x) \cdot \iota(y)^{-1}\right) \\
			&= g(\iota(x)) \cdot \left(g(\iota(y))\right)^{-1} \\
			&= j(x) \cdot j(y)^{-1} \\
			&= \tilde{f}(\rho(x)) \cdot \left(\tilde{f}(\rho(y))\right)^{-1} \\
			&= \tilde{f}(\rho(x)\bullet\rho(y)^{-1}) = \tilde{f}(a)
		\end{align*}
		
		Por lo tanto, \( g \circ \pi = \tilde{f} \).  Y por otro lado, también \( f \circ \pi = \tilde{f} \), así que:
		\[
		g \circ \pi = f \circ \pi.
		\]
		Como \( \pi : G_M \to U(M) \) es sobreyectivo, se deduce que \( g = f \). En efecto, dado \( b \in U(M) \), existe \( z \in G_M \) tal que \( \pi(z) = b \). Entonces:
		\[
		g(b) = g(\pi(z)) = (g \circ \pi)(z) = (f \circ \pi)(z) = f(\pi(z)) = f(b).
		\]
		Por lo tanto, \( f \) es el único morfismo de grupos que hace conmutar el diagrama.
	\end{itemize}
\end{proof}

 En la prueba de la propiedad universal del grupo universal de un monoide, se hizo uso del hecho de que un conjunto de elementos puede generar el subgrupo normal más pequeño que los contiene. Estudiaremos esta construcción con mayor profundidad en el próximo capítulo.


\begin{corolario}
	\upshape{
		Sea \( (G, \varphi) \) otro par que satisface la propiedad universal del grupo universal de un monoide para \( M \), es decir, para todo grupo \( H \) y todo morfismo de monoides \( j : M \to H \), existe un único morfismo de grupos \( f : G \to H \) tal que \( f \circ \varphi = j \).
		
		\medskip
		
		\noindent Entonces existe un único isomorfismo de grupos \( \psi : U(M) \to G \) tal que el siguiente diagrama conmuta:
		\[
		\xymatrix{
			&& M \ar[lld]_{\displaystyle \iota} \ar[rrd]^{\displaystyle \varphi} && \\
			U(M) \ar[rrrr]_{\displaystyle \exists!\,\psi}^{\displaystyle \cong} &&&& G
		}
		\]
		Es decir, \( \psi \circ \iota = \varphi \).
	}
\end{corolario}


\begin{proof}
	Aplicamos la propiedad universal de \( U(M) \) al morfismo \( \varphi : M \to G \). Como \( G \) es un grupo, existe un único homomorfismo de grupos \( \psi : U(M) \to G \) tal que el siguiente diagrama conmuta:
	\[
	\xymatrix{
		M \ar[rr]^{\displaystyle \varphi} \ar[dd]_{\displaystyle \iota} & & G \\
		&&\\
		U(M) \ar@{-->}[rruu]_{\displaystyle \exists! \psi} &&
	}
	\]
	Recíprocamente, aplicamos la propiedad universal de \( G \) al morfismo \( \iota : M \to U(M) \), lo que garantiza la existencia de un único homomorfismo de grupos \( \theta : G \to U(M) \) tal que:
	\[
	\xymatrix{
		M \ar[rr]^{\displaystyle \iota} \ar[dd]_{\displaystyle \varphi} & & U(M) \\
		&&\\
		G \ar@{-->}[rruu]_{\displaystyle \exists! \theta} &&
	}
	\] Entonces,
	\[
	\theta \circ \psi \circ \iota = \theta \circ \varphi = \iota.
	\] Ahora, aplicamos la propiedad universal de \( U(M) \) al morfismo \( \iota : M \to U(M) \). Tanto \( \mathrm{id}_{U(M)} \) como \( \theta \circ \psi \) son homomorfismos de grupos que hacen conmutar el siguiente diagrama:
	\[
	\xymatrix{
		M \ar[rrrr]^{\displaystyle \iota} \ar[dd]_{\displaystyle \iota} & &&& U(M) \\
		&&&&\\
		U(M) \ar@{-->}@/^1pc/[rrrruu]^{\displaystyle \mathrm{id}_{U(M)}} \ar@{-->}@/_1pc/[rrrruu]_{\displaystyle \theta \circ \psi} &&&&
	}
	\]
	Por la unicidad del morfismo de grupos que hace conmutar el triángulo, se concluye que \( \theta \circ \psi = \mathrm{id}_{U(M)} \).
	
	\medskip
	
	\noindent Análogamente, se prueba que \( \psi \circ \theta = \mathrm{id}_G \). Por lo tanto, \( \psi \) es un isomorfismo de grupos.
\end{proof}

\begin{ejemplo}\label{kjhgcsa}
	\upshape{
		Sea \( q \in \mathbb{N} \), y sean \( x \) y \( y \) variables formales. Definimos el conjunto
		\[
		M_q := \{0\} \cup \{mx : m \in \mathbb{N}\} \cup \{ny : n \in \mathbb{N}\},
		\]
		con una operación binaria \( + \) definida por las siguientes reglas:
		\begin{itemize}
			\item \( 0 + c = c \) para todo \( c \in M_q \),
			\item \( mx + px = (m + p)x \),
			\item \( my + py = (m + p)y \),
			\item \( mx + ny = (qm + n)y \).
		\end{itemize} Estas reglas implican, por ejemplo, que \( x + y = (q + 1)y \), y más generalmente:
		\[
		mx + ny = (qm + n)y.
		\] Entonces \( M_q \) es un monoide. Afirmamos que su grupo universal es isomorfo a \( \mathbb{Z} \), es decir:
		\[
		U(M_q) \cong \mathbb{Z}.
		\]	Para justificarlo, definimos un morfismo de monoides \( \iota : M_q \to \mathbb{Z} \) dado por:
		\[
		\iota(0) = 0, \quad \iota(ny) = n, \quad \iota(mx) = qm.
		\]Este morfismo induce, por la propiedad universal de \( U(M_q) \), un morfismo de grupos
		\[
		f : U(M_q) \to \mathbb{Z}.
		\]Recíprocamente, dado un grupo \( H \) y un morfismo de monoides \( j : M_q \to H \), definimos:
		\[
		\tilde{j} : \mathbb{Z} \to H, \quad \tilde{j}(n) := j(ny).
		\]
		Esta definición extiende a un morfismo de grupos tal que \( \tilde{j} \circ \iota = j \), lo que muestra que \( \mathbb{Z} \) también cumple la propiedad universal para \( M_q \).
		
		\medskip 
		\noindent Por la unicidad del grupo universal (salvo isomorfismo), se concluye que \( U(M_q) \cong \mathbb{Z} \).
	}
\end{ejemplo}



\begin{observacion}
	\upshape{
		Aunque el monoide \( M \) no sea conmutativo, su grupo fundamental \( U(M) \) puede serlo, como en el Ejemplo \ref{kjhgcsa}. Sin embargo, esto no ocurre en general y no debe asumirse.
	}
\end{observacion}



	\chapter{Subestructuras generadas por un cojunto en una estructura algebraica y Abelianización de un grupo}
	\section{Subestructuras generadas}

\noindent Sea \( M \) un grupo (respectivamente, semigrupo o monoide). Sea \( X \subset M \) un subconjunto.

\begin{definicion}
	\upshape{
		El \textit{subgrupo generado} por \( X \) (respectivamente, semigrupo o monoide) es el menor subgrupo (respectivamente, semigrupo o monoide) que contiene a \( X \), y se denota por \( \langle X \rangle \).
	}
\end{definicion}

\begin{observacion}
	\upshape{
		Si \( X = \{x_1, x_2, \dots, x_n\} \), entonces se escribe \( \langle x_1, x_2, \dots, x_n \rangle := \langle X \rangle \).
	}
\end{observacion}


\noindent Nuestro objetivo es saber cómo es \( \langle X \rangle \); pero antes, tengamos en cuenta la siguiente definición.

\begin{definicion}
	\upshape{
		Sea \( (G, \ast) \) un grupo, \( e \) su elemento neutro, y \( n \in \mathbb{Z} \). Para todo \( a \in G \), definimos:
		\[
		a^n :=
		\begin{cases}
			\underbrace{a \ast \cdots \ast a}_{n\ \text{veces}} & \text{si } n \geq 1, \\
		\:	e & \text{si } n = 0, \\
			\underbrace{a^{-1} \ast \cdots \ast a^{-1}}_{-n\ \text{veces}} & \text{si } n \leq -1.
		\end{cases}
		\]
	}
\end{definicion}
\begin{proposicion}
	\upshape{
		Sean \( a, x, y \) elementos de un grupo \( G \) y \( n \in \mathbb{Z} \). Se verifican las siguientes igualdades:
		\begin{enumerate}
			\item  \( (xy)^{-1} = y^{-1} x^{-1} \).
			\item  \( a^{-n} = (a^{-1})^n \).
			\item  \( (a^{-1})^n = (a^n)^{-1} \).
		\end{enumerate}
	}
\end{proposicion}

\begin{proof}
	{\color{white}saa}
	
	\medskip
	\noindent [$1$] Para todo \( x, y \in G \), se tiene:
	\[
	(xy)(y^{-1}x^{-1}) = x(yy^{-1})x^{-1} = xex^{-1} = xx^{-1} = e,
	\]
	y análogamente,
	\[
	(y^{-1}x^{-1})(xy) = y^{-1}(x^{-1}x)y = y^{-1}ey = y^{-1}y = e.
	\]
	Por unicidad del inverso, \( (xy)^{-1} = y^{-1}x^{-1} \).\\
	
	\noindent [$2$] Demostraremos inicialmente por inducción que \( a^{-n} = (a^{-1})^n \) para  \( n \in \mathbb{N}\).
	\begin{itemize}
		\item  Para \( n = 0 \): Por definición,
		\[
		a^0 = e = (a^{-1})^0.
		\]
		\item  Hipótesis inductiva: Supongamos que \( a^{-k} = (a^{-1})^k \) para algún \( k \geq 0 \).
		\item Para $k+1$ tenemos:
		\[
		a^{-(k+1)} = a^{-k} a^{-1} = (a^{-1})^k a^{-1} = (a^{-1})^{k+1}.
		\]
	\end{itemize}
	\noindent Por inducción, la igualdad es válida para todo \( n \geq 0 \).
	
	\medskip
	\noindent Para \( n < 0 \), escribimos \( n = -m \) con \( m > 0 \). Entonces:
	\[
	a^{-n} = a^m = \left((a^{-1})^{-1}\right)^m = (a^{-1})^{-m} = (a^{-1})^n.
	\]
	
	\medskip
	\noindent [$3$] Aplicamos el caso anterior al elemento \( a^{-1} \), obteniendo:
	\[
	(a^{-1})^n = ((a^{-1})^{-1})^{-n} = (a^1)^{-n} = (a^n)^{-1}.
	\]
	Esto completa la demostración.
\end{proof}

\vspace{0.8cm}
Las siguientes tablas nos describen cómo son los elementos de $\langle X \rangle$.\\


\begin{table}[H]
	\centering
	{\large
		\renewcommand{\arraystretch}{0.8}
		\resizebox{160mm}{!}{
			\begin{tabular}{
					>{\centering\arraybackslash}m{6cm}
					>{\centering\arraybackslash}m{12cm}
				}
				\toprule
				\textbf{Estructura} &  \(\langle X \rangle\)  \\
				\midrule
				Semigrupos & 
				\[
				\left\{ x_1^{k_1} x_2^{k_2} \cdots x_n^{k_n} \mid x_i \in X, n \geq 1, k_i \in \mathbb{N}^+ \right\}
				\] \\
				Monoides & 
				\[
				\left\{ x_1^{k_1} x_2^{k_2} \cdots x_n^{k_n} \mid x_i \in X, n \geq 1, k_i \in \mathbb{N} \right\}
				\] \\
				Grupos & 
				\[
				\left\{ x_1^{k_1} x_2^{k_2} \cdots x_n^{k_n} \mid x_i \in X, n \geq 1, k_i \in \mathbb{Z} \right\}
				\] \\
				\bottomrule
			\end{tabular}
		}
	}
	\caption{Descripción de \(\langle X \rangle\) con notación multiplicativa}
\end{table}

\vspace{0.8cm}
\noindent En el caso de ser el grupo un grupo aditivo tenemos:\\

\begin{table}[H]
	\centering
	{\large
		\renewcommand{\arraystretch}{0.8}
		\resizebox{160mm}{!}{
			\begin{tabular}{
					>{\centering\arraybackslash}m{6cm}
					>{\centering\arraybackslash}m{12cm}
				}
				\toprule
				\textbf{Estructura} &  \(\langle X \rangle\)  \\
				\midrule
				Semigrupos & 
				\[
				\left\{ k_1 x_1 + k_2 x_2 + \cdots + k_n x_n \mid x_i \in X, n \geq 1, k_i \in \mathbb{N}^+ \right\}
				\] \\
				Monoides & 
				\[
				\left\{ k_1 x_1 + k_2 x_2 + \cdots + k_n x_n \mid x_i \in X, n \geq 1, k_i \in \mathbb{N} \right\}
				\] \\
				Grupos & 
				\[
				\left\{ k_1 x_1 + k_2 x_2 + \cdots + k_n x_n \mid x_i \in X, n \geq 1, k_i \in \mathbb{Z} \right\}
				\] \\
				\bottomrule
			\end{tabular}
		}
	}
	\caption{Descripción de \(\langle X \rangle\) con notación aditiva}
\end{table}

\medskip
 \noindent Entonces formalicemos la que describen las tablas. 

\begin{proposicion}
	\upshape{
		Sea el conjunto
		\[
		H_X := \left\{ k_1 x_1 + k_2 x_2 + \cdots + k_n x_n \mid x_i \in X,\ n \geq 1,\ k_i \in \mathbb{Z} \right\}.
		\]
		Entonces \( H_X \) es un subgrupo de \( M \) y, además, se cumple que \( \langle X \rangle = H_X \).
	}
\end{proposicion}

\begin{proof}
	{\color{white}asas}
	
	\medskip
	\noindent [\(H_X\) es un subgrupo] Sea \( a, b \in H_X \). Entonces, existen \( x_1, \ldots, x_n \in X \), \( y_1, \ldots, y_m \in X \), \( k_1, \ldots, k_n \in \mathbb{Z} \), \( \ell_1, \ldots, \ell_m \in \mathbb{Z} \) tales que:
	\[
	a = x_1^{k_1} x_2^{k_2} \cdots x_n^{k_n}, \quad b = y_1^{\ell_1} y_2^{\ell_2} \cdots y_m^{\ell_m}.
	\]
	Como \( b^{-1} = y_m^{-\ell_m} \cdots y_1^{-\ell_1} \) (invirtiendo el orden y los exponentes), se tiene:
	\[
	ab^{-1} = x_1^{k_1} \cdots x_n^{k_n} \cdot y_m^{-\ell_m} \cdots y_1^{-\ell_1} \in H_X,
	\]
	pues es una combinación finita de potencias de elementos de \( X \) con exponentes en \( \mathbb{Z} \). Así, \( H_X \) es un subgrupo.\\
	

	\noindent[\(H_X = \langle X \rangle\)] Verifiques ahora que es el menor subgrupo que contiene a $X$, esto es 
	\begin{itemize}
		\item \( X \subset H_X \): En efecto, para todo \( x \in X \), se tiene \( x = x^1 \in H_X \), ya que \( 1 \in \mathbb{Z} \).
		
		\item Sea \( H \leq M \) tal que \( X \subset H \). Como \( H \) es un grupo, contiene los inversos y es cerrado bajo productos, por lo que toda combinación de la forma
		\[
		x_1^{k_1} x_2^{k_2} \cdots x_n^{k_n} \in H
		\]
		pertenece a \( H \). Por tanto, \( H_X \subset H \). 
	\end{itemize}
	Se concluye que \( \langle X \rangle = H_X \).
\end{proof}

\vspace{0.5cm}
Ahora nos interesa describir en particular el caso en que \( X = \{x\} \), es decir, cuando el conjunto generador está formado por un solo elemento. Este caso es fundamental, pues muchos subgrupos importantes en álgebra son generados por un solo elemento.

\begin{definicion}
	\upshape{
		Dado \( x \in M \), el subgrupo de la forma \( \langle x \rangle \) se denomina \textit{subgrupo cíclico} o, simplemente, \textit{grupo cíclico}.
	}
\end{definicion}

\vspace{0.5cm}
 La siguiente tabla describe cómo son los elementos de $\langle x\rangle$.


 
\begin{table}[H]
	\centering
	{\large
		\renewcommand{\arraystretch}{1.2}
		\resizebox{160mm}{!}{
			\begin{tabular}{
					>{\centering\arraybackslash}m{6cm}
					>{\raggedright\arraybackslash}m{12cm}
				}
				\toprule
				\textbf{Estructura} &  \(\langle x \rangle\)  \\ 
				\midrule
				Semigrupos & 
				\[
				\{ x, x^2, x^3, x^4, \ldots \}
				\] \\ 
				Monoides & 
				\[
				\{ e, x, x^2, x^3, x^4, \ldots \}
				\] \\ 
				Grupos & 
				\[
				\{ \ldots, x^{-3}, x^{-2}, x^{-1}, e, x, x^{2}, x^{3}, \ldots \}
				\] \\ 
				\bottomrule
			\end{tabular}
		}
	}
	\caption{Elementos por extensión del conjunto \(\langle x \rangle\) generados por un solo elemento, en notación multiplicativa}
\end{table}
Dependiendo del comportamiento del elemento \( x \), el subgrupo cíclico \( \langle x \rangle \) puede ser finito o infinito. A continuación se presentan ejemplos en ambos casos.
\begin{ejemplo}
	\leavevmode
	\upshape{
	
	\begin{itemize}
		\item Caso infinito: Sea \( G = (\mathbb{Z}, +) \), el grupo aditivo de los enteros. Toma \( x = 1 \in \mathbb{Z} \). Entonces:
		\[
		\langle x \rangle = \langle 1 \rangle = \{ n \cdot 1 \mid n \in \mathbb{Z} \} = \mathbb{Z},
		\]
		es decir, \( \langle x \rangle \) es un grupo cíclico infinito.
		
		\item Caso finito: Sea \( G = (\mathbb{Z}_6, +) \), el grupo aditivo de los enteros módulo 6. Toma \( x = \overline{2} \in \mathbb{Z}_6 \). Entonces:
		\[
		\langle \overline{2} \rangle = \{ \overline{0}, \overline{2}, \overline{4} \},
		\]
		pues:
		\[
		\overline{2} + \overline{2} = \overline{4}, \quad
		\overline{2} + \overline{2} + \overline{2} = \overline{0}.
		\]
		Así, \( \langle x \rangle \) es un grupo cíclico finito de orden 3.
	\end{itemize}
	}
\end{ejemplo}


\begin{ejemplo}
	\upshape{
		Sea \( X \) un espacio topológico y sea \( f \in \mathrm{Homeo}(X) \), es decir, un homeomorfismo de \( X \) en sí mismo. Como \( \mathrm{Homeo}(X) \) es un grupo bajo la composición de funciones, el elemento \( f \) genera:
		
		\begin{itemize}
			\item El submonoide cíclico generado por \( f \):
			\[
			\langle f \rangle := \{ f^{\circ n} : n \in \mathbb{N} \},
			\]
			que consiste en las iteradas no negativas de \( f \). 
			\begin{center}
				\begin{tikzpicture}[scale=1.2, >=Stealth, thin]
					
					% Coordenadas de los nodos
					\coordinate (p0) at (180:0.4);
					\coordinate (p1) at (135:0.6);
					\coordinate (p2) at (90:0.8);
					\coordinate (p3) at (45:1.0);
					\coordinate (p4) at (0:1.2);
					\coordinate (p5) at (-45:1.4);
					\coordinate (p6) at (-90:1.6);
					\coordinate (p7) at (-135:1.8);
					\coordinate (endArrow) at (-150:3.2); % Flecha un poco más larga que antes
					
					% Nodos con etiqueta refinada
					\node[circle, fill=black, inner sep=1.2pt, label=below left:{\( x \)}] at (p0) {};
					\node[circle, fill=black, inner sep=1.2pt, label=above left:{\( f(x) \)}] at (p1) {};
					\node[circle, fill=black, inner sep=1.2pt, label=above:{\( f^{\circ 2}(x) \)}] at (p2) {};
					\node[circle, fill=black, inner sep=1.2pt, label=above right:{\( f^{\circ 3}(x) \)}] at (p3) {};
					\node[circle, fill=black, inner sep=1.2pt, label=right:{\( f^{\circ 4}(x) \)}] at (p4) {};
					\node[circle, fill=black, inner sep=1.2pt, label=below right:{\( f^{\circ 5}(x) \)}] at (p5) {};
					\node[circle, fill=black, inner sep=1.2pt, label=below:{\( f^{\circ 6}(x) \)}] at (p6) {};
					\node[circle, fill=black, inner sep=1.2pt, label=above left:{\( f^{\circ 7}(x) \)}] at (p7) {};
					
					% Flechas curvas más delgadas usando bend left/right
					\draw[->, thin, bend left] (p0) to (p1);
					\draw[->, thin, bend left] (p1) to (p2);
					\draw[->, thin, bend left] (p2) to (p3);
					\draw[->, thin, bend left] (p3) to (p4);
					\draw[->, thin, bend right] (p4) to (p5);
					\draw[->, thin, bend right] (p5) to (p6);
					\draw[->, thin, bend right] (p6) to (p7);
					
					% Flecha curva adicional después de f^7 pero sin nodo final
					\draw[->, thin, bend left] (p7) to (endArrow);
					
				\end{tikzpicture}
			\end{center}
			
			\item El subgrupo cíclico generado por \( f \):
			\[
			\langle f \rangle := \{ f^{\circ n} : n \in \mathbb{Z} \},
			\]
			que incluye también las iteradas inversas \( f^{-1}, f^{-2}, \dots \), ya que \( f \) es invertible.
		\end{itemize} Ambas estructuras son relevantes en el estudio de dinámicas topológicas, ya que permiten analizar las trayectorias de puntos bajo iteraciones del mismo homeomorfismo.
	}
\end{ejemplo}


\vspace{0.5cm}
Hasta ahora sabemos cómo generar subgrupos a partir de un subconjunto de un grupo. Sin embargo, no podemos esperar que estos subgrupos sean normales en general.



\begin{ejemplo}
	\upshape{
		Sea \( G \) el grupo de matrices invertibles \( 2 \times 2 \) con entradas reales, es decir, \( G = \mathrm{GL}_2(\mathbb{R}) \). Consideremos la matriz
		\[
		A = \begin{pmatrix}
			1 & 1 \\
			0 & 1
		\end{pmatrix}.
		\]
		El subgrupo generado por \( A \) está dado por
		\[
		\langle A \rangle = \left\{ A^n \mid n \in \mathbb{Z} \right\} = \left\{ \begin{pmatrix}
			1 & n \\
			0 & 1
		\end{pmatrix} \mid n \in \mathbb{Z} \right\}.
		\] Ahora, tomemos la matriz
		\[
		P = \begin{pmatrix}
			0 & 1 \\
			1 & 0
		\end{pmatrix}.
		\]
		Calculamos el conjugado de \( A \) por \( P \):
		\[
		P A P^{-1} = \begin{pmatrix}
			1 & 0 \\
			1 & 1
		\end{pmatrix}.
		\]
		Para que \( P A P^{-1} \) esté en \( \langle A \rangle \), debería existir algún \( m \in \mathbb{Z} \) tal que
		\[
		A^m = \begin{pmatrix}
			1 & m \\
			0 & 1
		\end{pmatrix} = \begin{pmatrix}
			1 & 0 \\
			1 & 1
		\end{pmatrix},
		\]
		lo cual es imposible pues los elementos en la posición (2,1) son distintos (0 en \( A^m \) y 1 en \( P A P^{-1} \)). 	Por tanto, \( P A P^{-1} \notin \langle A \rangle \) y el subgrupo \( \langle A \rangle \) no es normal en \( G \).
	}
\end{ejemplo}


\vspace{0.4cm}
Por ello, nos interesa describir cómo deben ser los elementos de un grupo para generar subgrupos normales.

\begin{definicion}
	\upshape{
		El \textit{subgrupo normal generado} por \( X \), denotado por \( \langle X \rangle^{\trianglelefteq} \), es el menor subgrupo normal de \( M \) que contiene a \( X \).
	}
\end{definicion}

\begin{proposicion}
	\upshape{
		Sea \( M \) un grupo y \( X \subseteq M \) un subconjunto no vacío. Definimos el conjunto
		\[
		H_N := \langle \{ a x^k a^{-1} \mid a \in M,\, k \in \mathbb{Z} \} \rangle.
		\] Entonces, \( H_N \) es un subgrupo normal de \( M \) y además
		\[
		\langle x \rangle^{\trianglelefteq} =  H_N.
		\]
	}
\end{proposicion}




\begin{proof}
	Es claro que \( H_N \) es un subgrupo de \( M \), ya que está definido como el subgrupo generado por un conjunto. Veamos ahora que \( H_N \) es normal en \( M \).
	
	\medskip	
	\noindent Para todo \( m \in M \) y para todo generador \( a x^k a^{-1} \in H_N \) (con \( a \in M \), \( x \in X \), \( k \in \mathbb{Z} \)), se tiene:
	\[
	m (a x^k a^{-1}) m^{-1} = (m a) x^k (m a)^{-1} \in H_N,
	\]
	ya que \( m a \in M \), \( x \in X \), y \( k \in \mathbb{Z} \). Como \( H_N \) contiene a todos sus conjugados, es normal en \( M \).\\
	
	\medskip 
	\noindent Ahora mostraremos que \( \langle X \rangle^{\trianglelefteq} = H_N \). Para ello:
	\begin{itemize}
		\item Como \( x = e x e^{-1} \in H_N \) para todo \( x \in X \), se tiene \( X \subseteq H_N \).
		\item Sea \( H \) un subgrupo normal de \( M \) que contiene a \( X \). Entonces, para todo \( a \in M \), \( x \in X \), \( k \in \mathbb{Z} \), se tiene:
		\[
		x^k \in H \quad \text{(porque \( H \) es subgrupo)}, \quad \text{y} \quad a x^k a^{-1} \in H \quad \text{(porque \( H \) es normal)}.
		\]
		Por tanto, todos los generadores de \( H_N \) pertenecen a \( H \), y luego \( H_N \subseteq H \).
	\end{itemize}
	Esto prueba que \( H_N \) es el menor subgrupo normal de \( M \) que contiene a \( X \), es decir, \( \langle X \rangle^{\trianglelefteq} = H_N \).
\end{proof}

\vspace{0.4cm}

\begin{observacion}
	\upshape{
		La caracterización mediante elementos del subgrupo normal generado por un único elemento \( x \in M \) es:
		\[
		\langle x \rangle^{\trianglelefteq} = 
		\left\{ a_1 x^{k_1} a_1^{-1} \cdots a_n x^{k_n} a_n^{-1} \;\middle|\; 
		n \in \mathbb{N},\, a_i \in M,\, k_i \in \mathbb{Z} \right\}.
		\]
		Es decir, está formado por todos los productos finitos de conjugados de potencias enteras de \( x \). Así, por ejemplo, algunos de sus elementos son:
		\[
		e,\quad x,\quad x^{-1},\quad
		a x a^{-1},\quad a x^{-1} a^{-1},\quad
		b x^{2} b^{-1},\quad
		c x^{-3} c^{-1},\quad
		(a x a^{-1})(b x b^{-1}).
		\]
	}
\end{observacion}

\vspace{0.4cm}

\begin{proposicion}
	\upshape{
		Sea \( M \) un grupo y \( X \subseteq M \) un subconjunto no vacío. Definimos el conjunto
		\[
		H_N := \left\langle \left\{ a x^k a^{-1} \mid a \in M,\, x \in X,\, k \in \mathbb{Z} \right\} \right\rangle.
		\]
		Entonces, \( H_N \) es un subgrupo normal de \( M \) y además
		\[
		\langle X \rangle^{\trianglelefteq} = H_N.
		\]
	}
\end{proposicion}


\begin{proof}
	Es claro que \( H_N \) es un subgrupo de \( M \), ya que está definido como el subgrupo generado por un conjunto. Veamos ahora que \( H_N \) es normal en \( M \).
	
	\medskip
	\noindent Para todo \( m \in M \) y para todo generador \( a x^k a^{-1} \in H_N \) (con \( a \in M \), \( x \in X \), \( k \in \mathbb{Z} \)), se tiene:
	\[
	m (a x^k a^{-1}) m^{-1} = (m a) x^k (m a)^{-1} \in H_N,
	\]
	ya que \( m a \in M \), \( x \in X \), y \( k \in \mathbb{Z} \). Por tanto, \( H_N \) es normal en \( M \).
	
	\medskip
	\noindent Veamos ahora que \( \langle X \rangle^{\trianglelefteq} = H_N \). Para ello:
	\begin{itemize}
		\item Como para cada \( x \in X \), se tiene \( x = e x e^{-1} \in H_N \), entonces \( X \subseteq H_N \).
		\item Sea \( H \) un subgrupo normal de \( M \) tal que \( X \subseteq H \). Entonces:
		\begin{itemize}
			\item Para todo \( x \in X \) y \( k \in \mathbb{Z} \), se tiene \( x^k \in H \) (porque \( H \) es subgrupo),
			\item Para todo \( a \in M \), \( a x^k a^{-1} \in H \) (porque \( H \) es normal).
		\end{itemize}
		Así, todos los generadores de \( H_N \) están en \( H \), por lo que \( H_N \subseteq H \).
	\end{itemize}
	Esto prueba que \( H_N \) es el menor subgrupo normal de \( M \) que contiene a \( X \), es decir,
	\[
	\langle X \rangle^{\trianglelefteq} = H_N. \qedhere
	\]
\end{proof}



\section{Clasificación de grupo cíclicos}




\begin{definicion}
	\upshape{
		Sea $G$ un grupo. El \textit{orden de $G$}, denotado por $|G|$, es el cardinal del conjunto subyacente $G$. Si $|G|$ es finito, decimos que $G$ es un \textit{grupo finito}.
	}
\end{definicion}

\begin{proposicion}\label{lkjhgfñ-lkj}
	\upshape
	Sea \( G \) un grupo finito y \( x \in G \). Entonces, existe un entero positivo \( n \) tal que \( x^n = e \). Además, el subgrupo generado por \( x \) es
	\[
	\langle x \rangle = \{ e,\ x,\ x^2,\ \dots,\ x^{n-1} \}.
	\]
\end{proposicion}

\begin{proof}
	\noindent\textit{Existencia de \( n \):} Como \( G \) es finito, la secuencia \( x, x^2, x^3, \dots \) contiene infinitos términos pero toma valores en un conjunto finito. Por el principio del palomar, existen enteros \( k > \ell \geq 1 \) tales que \( x^k = x^\ell \). Entonces, $ x^{k - \ell} = e$. Sea \( n := k - \ell \), que es un entero positivo tal que \( x^n = e \).\\
	

	
	\noindent\textit{Estructura de \( \langle x \rangle \):} El subgrupo generado por \( x \) está formado por todas sus potencias:
	\[
	\langle x \rangle = \{ x^m : m \in \mathbb{Z} \}.
	\]
	Sea \( m \in \mathbb{Z} \), al escribir \( m = qn + r \) con \( 0 \leq r < n \), se tiene:
	\[
	x^m = x^{qn + r} = (x^n)^q x^r = e^q x^r = x^r.
	\]
	Por tanto, todo \( x^m \) coincide con alguna de las potencias \( x^r \) para \( 0 \leq r < n \), y entonces:
	\[
	\langle x \rangle = \{ e, x, x^2, \dots, x^{n-1} \}.
	\]
\end{proof}

\begin{definicion}
	\upshape{
		Sea $G$ un grupo y $x \in G$. El \textit{orden de $x$}, denotado por $\mathrm{ord}(x)$, es el menor entero positivo $n$ tal que $x^n = e$. Si no existe tal $n$, decimos que $x$ tiene \textit{orden infinito}.
	}
\end{definicion}

\begin{proposicion}
	\upshape
	Sea \( G \) un grupo finito y \( x \in G \). Entonces, el orden de \( x \) es el menor entero positivo \( k \) tal que \( x^k = e \).
\end{proposicion}

\begin{proof}
	Por la Proposición \ref{lkjhgfñ-lkj}, existe un entero positivo \( n \) tal que \( x^n = e \). Definamos
	\[
	k := \min \left\{ n \in \mathbb{N} \setminus \{0\} \mid x^n = e \right\}.
	\]
	Además, por la misma proposición, tenemos que
	\[
	\langle x \rangle = \{ e, x, x^2, \dots, x^{k-1} \}.
	\]
	Que \( k \) sea el menor entero positivo tal que \( x^k = e \) implica que todos los elementos de ese conjunto son distintos entre sí. Supongamos que existen \( 1 \leq i, j < k \) con \( i \neq j \) tales que \( x^i = x^j \). Sin pérdida de generalidad, supongamos \( i > j \), entonces \( x^{i - j} = e \). Así,  $k':=i-j$, \( 0 < k '< k \) contradice la minimalidad de \( k \). 
	
	\medskip
	\noindent Por lo tanto, \( \operatorname{Ord}(x) = k \), y este es el menor entero positivo tal que \( x^k = e \).
\end{proof}

\begin{definicion}
	\upshape{
		Un grupo $G$ es \textit{cíclico} si $G = \langle g \rangle$ para algún $g \in G$.
	}
\end{definicion}

\begin{proposicion}\label{lkjhghjklñ}
	\upshape
	Sea \( G \) un grupo cíclico. Entonces se cumple:
	\begin{enumerate}
		\item Si \( G \) tiene orden infinito, entonces \( G \cong \mathbb{Z} \).
		\item Si \( G \) tiene orden \( n \), entonces \( G \cong \mathbb{Z}/n\mathbb{Z} \).
	\end{enumerate}
\end{proposicion}

\begin{proof}
	Como \( G \) es cíclico, existe un elemento \( a \in G \) tal que \( G = \langle a \rangle = \{ a^m : m \in \mathbb{Z} \} \).
	
	\medskip
	\noindent$[1)]$ Supongamos que \( G \) tiene orden infinito. Definimos el homomorfismo:
		\[
		\begin{aligned}
			f : \mathbb{Z} &\longrightarrow G \\
			m &\longmapsto a^m.
		\end{aligned}
		\]
		Veamos que \( f \) es un isomorfismo.
		\begin{itemize}
			\item Homomorfismo:
			\[
			f(m + n) = a^{m+n} = a^m a^n = f(m)f(n).
			\]
			
			\item Sobreyectividad: Sea $g\in G$, entonces \(g= a^m \) para algún \( m \in \mathbb{Z} \). Así $f(m)=g$.
			
			\item Inyectividad:
			\[
			\ker f = \{ m \in \mathbb{Z} : f(m) = e \} = \{ m \in \mathbb{Z} : a^m = e \}.
			\]
			Como \( a \) tiene orden infinito, no existe ningún entero positivo \( m \) tal que \( a^m = e \), excepto \( m = 0 \). Por lo tanto, \( \ker f = \{0\} \), es decir, \( f \) es inyectivo.
		\end{itemize}
		
		Como \( f \) es un homomorfismo biyectivo, es un isomorfismo. Por lo tanto, \( G \cong \mathbb{Z} \).\\
		
		\noindent $[2)]$ Supongamos ahora que \( G =\langle a\rangle\) tiene orden \( n \in \mathbb{N} \). Entonces \( a^n = e \) y \( n \) es el menor entero positivo con esta propiedad. Consideramos el mismo homomorfismo:
		\[
		\begin{aligned}
			f : \mathbb{Z} &\longrightarrow G \\
			m &\longmapsto a^m.
		\end{aligned}
		\]
		Como antes, \( f \) es un homomorfismo y es sobreyectivo. Calculemos su núcleo:
		\[
		\ker f = \{ m \in \mathbb{Z} : a^m = e \} = \{ m \in \mathbb{Z} : n \text{ divide a } m \} = n\mathbb{Z}.
		\]
		Entonces por el teorema fundamental del homomorfismo de grupos:
		\[ \mathbb{Z} / n\mathbb{Z} \cong G.
		\]

\end{proof}


\begin{proposicion}
	\upshape{
		Todo subgrupo de un grupo cíclico es cíclico.
	}
\end{proposicion}

\begin{proof}
	Como \( G \) es cíclico, existe un elemento \( a \in G \) tal que \( G = \langle a \rangle \). Sea \( H \leq G \) un subgrupo. Queremos probar que \( H \) también es cíclico.
	
	\medskip
	\noindent
	Si \( H = \{e\} \), entonces \( H = \langle e \rangle \), que es cíclico.
	
	\medskip
	\noindent
	Supongamos ahora que \( H \neq \{e\} \). Entonces existe un entero \( n \neq 0 \) tal que \( a^n \in H \). Sea
	\[
	k := \min \left\{ n \in \mathbb{N} \setminus \{0\} \mid a^n \in H \right\}.
	\]
	Es decir, \( a^k \in H \) y \( k \) es el menor entero positivo con esta propiedad. Afirmamos que
	\[
	H = \langle a^k \rangle.
	\]
	\begin{itemize}
		\item \( \langle a^k \rangle \subseteq H \): Como \( a^k \in H \) y \( H \) es un subgrupo, entonces todas las potencias de \( a^k \) también están en \( H \).
		
		\item  \( H \subseteq \langle a^k \rangle \): Sea \( h \in H \). Como \( G = \langle a \rangle \), existe \( m \in \mathbb{Z} \) tal que \( h = a^m \). Escribamos \( m = qk + r \) con \( 0 \leq r < k \). Entonces,
		\[
		h = a^m = a^{qk + r} = (a^k)^q a^r.
		\]
		Como \( (a^k)^q \in \langle a^k \rangle \subseteq H \), y \( h = (a^k)^q a^r \in H \), se deduce que \( a^r = (a^k)^{-q} h \in H \). Pero por la definición de \( k \), ningún entero positivo menor que \( k \) satisface que \( a^r \in H \), a menos que \( r = 0 \). Así, \( h = (a^k)^q \in \langle a^k \rangle \).
	\end{itemize}
	\medskip
	\noindent
	Por tanto, \( H \subseteq \langle a^k \rangle \), y concluimos que \( H = \langle a^k \rangle \), lo que demuestra que \( H \) es cíclico.
\end{proof}











\section{Abelianización de un grupo}\label{kjhgf}

\begin{definicion}
	\upshape{
		Sea \( G \) un grupo. El \textit{conmutador} de dos elementos \( a, b \in G \) se define como
		\[
		[a,b] := aba^{-1}b^{-1}.
		\]
	}
\end{definicion}

\begin{observacion}
	\upshape{
		Si \( [a, b] = e \), entonces \( ab = ba \), pues la igualdad \( aba^{-1}b^{-1} = e \) implica \( aba^{-1} = b \), lo que equivale a \( ab = ba \).
		
		\medskip
		\noindent En resumen, dos elementos conmutan si y solo si su conmutador es la identidad.
	}
\end{observacion}
\begin{observacion}
	\upshape{
		Sea \( f: G \to H \) un homomorfismo de grupos. Entonces, para cualesquiera \( a, b \in G \), se cumple
		\[
		f([a,b]) = [f(a), f(b)].
		\]
		Pues como \( f \) es homomorfismo, se tiene:
		\[
		f([a,b]) = f(aba^{-1}b^{-1}) = f(a)f(b)f(a)^{-1}f(b)^{-1} = [f(a), f(b)].
		\]
	}
\end{observacion}

\begin{definicion}
	\upshape{
		El \textit{grupo conmutador} o \textit{grupo derivado} de \( G \) es el subgrupo
		\[
		[G, G] := \langle [a, b] : a, b \in G \rangle.
		\]
		Es decir, el subgrupo normal generado por todos los conmutadores de elementos de \( G \). Además, \( [G, G] \trianglelefteq G \).
	}
\end{definicion}

\begin{proposicion}
	\upshape{
		El grupo conmutador \( [G, G] \) es un subgrupo normal de \( G \).
	}
\end{proposicion}

\begin{proof}
	Sea \( x \in G \). Como \( [G, G] \) está generado por conmutadores \( [a,b] = aba^{-1}b^{-1} \), basta probar que para todo conmutador \( [a,b] \) se cumple que \( x [a,b] x^{-1} \in [G,G] \). Calculemos:
	\[
	x [a,b] x^{-1} = x (aba^{-1}b^{-1}) x^{-1} = (x a x^{-1})(x b x^{-1})(x a x^{-1})^{-1}(x b x^{-1})^{-1} \in [G,G].
	\] Por tanto, el conjugado de cualquier elemento generado por conmutadores pertenece a \( [G,G] \). Así, \( [G,G] \) es normal en \( G \).
\end{proof}

\vspace{0.4cm}
Por lo tanto, el cociente
\[
\dfrac{G}{[G,G]}
\]
tiene una estructura de grupo con la operación inducida por \( G \), es decir, para \( a, b \in G \) se define
\[
(a [G,G]) \cdot (b [G,G]) := (ab) [G,G].
\]
\begin{observacion}
	\upshape{
		Sea \( f : G \to H \) un homomorfismo de grupos. Entonces,
		\[
		f([G,G]) \trianglelefteq f(G).
		\]
		En general, \( f([G,G]) \) no es normal en \( H \) ni en \( [H,H] \).
	}
\end{observacion}

\begin{definicion}
	\upshape{
		Sea \( G \) un grupo. La \textit{abelianización} de \( G \) es el grupo
		\[
		G^{\mathrm{ab}} := \frac{G}{[G,G]}.
		\]
	}
\end{definicion}





\medskip

\noindent La abelianización de un grupo proporciona una forma canónica de obtener un grupo conmutativo a partir de cualquier grupo \( G \). Esta construcción es fundamental en teoría de grupos y permite estudiar \( G \) desde una perspectiva más manejable, reduciendo su estructura a la información esencial sobre su conmutatividad. 

\begin{ejemplo}
	\upshape{
	Sea \( G = S_3 \), el grupo simétrico de las permutaciones de tres elementos. Sabemos que \( S_3 \) no es abeliano, pero su grupo conmutador \( [S_3, S_3] \) es el subgrupo alternante \( A_3 \), que es un grupo normal de orden 3. Por lo tanto, la abelianización de \( S_3 \) es el cociente
	\[
	S_3^{\mathrm{ab}} = \frac{S_3}{[S_3, S_3]} = \frac{S_3}{A_3} \cong \mathbb{Z}_2,
	\]
	que es un grupo abeliano de dos elementos.
	
	\medskip
	\noindent Esto muestra que \( S_3 \) se “reduce” a un grupo abeliano simple, \(\mathbb{Z}_2\), cuando consideramos solo su parte conmutativa.
	}
\end{ejemplo}











\section{El functor de abelianización}

\noindent Para estudiar la relación entre los grupos arbitrarios y los grupos abelianos, es útil considerar la categoría en la que residen. La categoría \( \mathbf{Grp} \) tiene como objetos los grupos y como morfismos los homomorfismos de grupos.

\medskip
\noindent Dentro de \( \mathbf{Grp} \) encontramos la subcategoría \( \mathbf{Ab} \), que contiene todos los grupos abelianos junto con sus respectivos morfismos. Esto nos permite definir el functor de abelianización, denotado por \( (-)^{\operatorname{ab}} \).

\medskip
\noindent Primero veamos la asignación de objetos. De la Sección \ref{kjhgf} podemos hacer 
\begin{align*}
	(-)^{\operatorname{ab}}: \mathbf{Grp}&\longrightarrow  \mathbf{Ab}\\
	G&\longmapsto G^{\operatorname{ab}}.
\end{align*}
Por lo que ahora debemos ver cómo se asignan los morfismos. Sea \(f: G\to H\) en \(\mathbf{Grp}\). Entonces definimos
\begin{align*}
	f^{\operatorname{ab}}:G^{\operatorname{ab}}&\longrightarrow H^{\operatorname{ab}}\\
	[a]&\longmapsto [f(a)],
\end{align*}
esta asignación tiene sentido porque, como recordamos en la Sección \ref{kjhgf}, tanto \(G^{\operatorname{ab}}\) como \(H^{\operatorname{ab}}\) son grupos cocientes y, por un abuso de notación, usamos el mismo símbolo para sus clases. Veamos entonces que \(f^{\operatorname{ab}}\) está bien definido. Recordemos que 
\[
G^{\operatorname{ab}} = G / [G, G], \quad H^{\operatorname{ab}} = H / [H, H],
\]
donde \([G,G]\) y \([H,H]\) son los subgrupos conmutadores de \(G\) y \(H\) respectivamente.

\medskip
\noindent Para que \(f^{\operatorname{ab}}\) esté bien definido, debemos comprobar que si dos elementos \(a, a' \in G\) pertenecen a la misma clase en \(G^{\operatorname{ab}}\), entonces sus imágenes bajo \(f^{\operatorname{ab}}\) también pertenezcan a la misma clase en \(H^{\operatorname{ab}}\). Es decir, si 
\[
[a] = [a'] \implies a' = a \cdot c,
\]
para algún \(c \in [G, G]\) por ser $[G,G]\trianglelefteq G$. Como \(f\) es un homomorfismo de grupos, tenemos
\[
f(a') = f(a \cdot c) = f(a) \cdot f(c).
\]Además, como \(c \in [G,G]\), entonces \(f(c) \in [H,H]\) porque los homomorfismos de grupos 
\[
[f(a')] = [f(a) \cdot f(c)] = [f(a)] \cdot [f(c)] = [f(a)],
\] ya que $[f(c)]$ es una unidad en $H^{\operatorname{ab}}$. Veamos ahora que \(f^{\operatorname{ab}}\) es un homomorfismo de grupos. Sean \( [a], [b] \in G^{\operatorname{ab}} \), es decir, \( a, b \in G \). Entonces,
\[
f^{\operatorname{ab}}([a] \cdot [b]) = f^{\operatorname{ab}}([ab]) = [f(ab)].
\]
Como \(f\) es homomorfismo, se tiene que \(f(ab) = f(a)f(b)\), y por tanto
\[
f^{\operatorname{ab}}([a] \cdot [b]) = [f(ab)] = [f(a)f(b)] = [f(a)] \cdot [f(b)] = f^{\operatorname{ab}}([a]) \cdot f^{\operatorname{ab}}([b]).
\]
Por lo tanto, \(f^{\operatorname{ab}}\) respeta la operación de grupo y es un homomorfismo.

\begin{proposicion}
	\upshape{
		Consideremos la asignación
		\begin{align*}
			(-)^{\operatorname{ab}}: \mathbf{Grp} &\longrightarrow \mathbf{Ab} \\
			G &\longmapsto G^{\operatorname{ab}}, \\
			\left(f : G \to H\right) &\longmapsto \left(f^{\operatorname{ab}} : G^{\operatorname{ab}} \to H^{\operatorname{ab}}\right).
		\end{align*}
		Entonces \( (-)^{\operatorname{ab}} \) es un functor.
	}
\end{proposicion}

\begin{proof}
	\upshape{
		Para ello verificamos las condiciones de la Definición \ref{ppoiujhgf}:
		
		\begin{enumerate}
			\item Asignación sobre objetos: A cada grupo \(G \in \mathbf{Grp}\), se le asigna el grupo abeliano \(G^{\operatorname{ab}} = G / [G, G]\). Esta construcción se ha justificado previamente.
			
			\item Asignación sobre morfismos: Sea \(f: G \to H\) un homomorfismo de grupos. Definimos \(f^{\operatorname{ab}}: G^{\operatorname{ab}} \to H^{\operatorname{ab}}\) por
			\[
			f^{\operatorname{ab}}([a]) := [f(a)],
			\]
			que, como vimos antes, es un homomorfismo de grupos.
			
			\item Preservación de identidades: Sea \(G \in \mathbf{Grp}\). Entonces,
			\[
			(\mathrm{id}_G)^{\operatorname{ab}}([a]) = [\mathrm{id}_G(a)] = [a],
			\]
			para todo \(a \in G\). Así, \((\mathrm{id}_G)^{\operatorname{ab}} = \mathrm{id}_{G^{\operatorname{ab}}}\).
			
			\item Preservación de la composición: Sean \(f: G \to H\) y \(g: H \to K\) morfismos en \( \mathbf{Grp} \). Entonces, para todo \( [a] \in G^{\operatorname{ab}} \),
			\[
			(g \circ f)^{\operatorname{ab}}([a]) = [(g \circ f)(a)] = [g(f(a))] = g^{\operatorname{ab}}([f(a)]) = g^{\operatorname{ab}}(f^{\operatorname{ab}}([a])) = \left( g^{\operatorname{ab}} \circ f^{\operatorname{ab}} \right)([a]).
			\]
			Por tanto,
			\[
			(g \circ f)^{\operatorname{ab}} = g^{\operatorname{ab}} \circ f^{\operatorname{ab}}.
			\]
		\end{enumerate}
		
		\noindent
		Concluimos que \( (-)^{\operatorname{ab}} \) es un funtor bien definifo de \( \mathbf{Grp} \) a \( \mathbf{Ab} \).
	}
\end{proof}




\begin{definicion}
	\upshape{
		El functor
		$(-)^{\operatorname{ab}}$
		de la categoría de grupos a la categoría de grupos abelianos está definido por
		\begin{align*}
			(-)^{\operatorname{ab}} : \mathbf{Grp} &\longrightarrow \mathbf{Ab}, \\
			G &\longmapsto G^{\operatorname{ab}}, \\
			\left(f : G \to H\right) &\longmapsto \left(f^{\operatorname{ab}} : G^{\operatorname{ab}} \to H^{\operatorname{ab}}\right).
		\end{align*}
		Se denomina \emph{functor de abelianización}.
	}
\end{definicion}





	\chapter{Presentación de un grupo, Grupo diedral, Teorema de Cayley, y Acción de un grupo}
	\section{Presentación de un grupo}

En teoría de grupos, una \textbf{presentación de grupo} es una forma de describir un grupo mediante un conjunto de \textbf{generadores}, que producen todos los elementos del grupo, y un conjunto de \textbf{relaciones}, que imponen restricciones sobre cómo estos generadores interactúan. Se expresa simbólicamente como:
\[
\langle \text{generadores} \mid \text{relaciones} \rangle
\]
Esta notación proporciona una descripción compacta y útil para estudiar la estructura de los grupos, sin necesidad de enumerar todos sus elementos.ensin necesidad de describir explícitamente todos sus elementos.\\



Dado un conjunto \( X \), se denota por \( G_X \) al grupo libre generado por \( X \) y con $e$ denotaremos su elemento identidad. 




\begin{definicion}
	\upshape{
		Sean \( X \) un conjunto y \( R \subseteq G_X \) un subconjunto. Decimos que un grupo \( G \) está definido por los generadores en \( X \) y las relaciones \( w = e \),\( w \in R \), si
		\[
		G \cong G_X / N_R,
		\]
		donde \( N_R \) es el subgrupo normal de \( G_X \) generado por \( R \).
		
		\medskip
		
		\noindent En este caso, denotamos
		\[
		 G_X / N_R:= \langle X : w = e,\ w \in R \rangle \quad \text{o bien} \quad \langle X : R \rangle.
		\]
		Y se le denomina \emph{presentación} de \( G \).
	}
\end{definicion}


\begin{observacion}
	\upshape{
		Las relaciones \( w = e \), \( w \in R \), consisten en que para todo \( w \in R \) se tiene que \( w N_R = N_R = e N_R \); es decir, indican que los elementos \( w \in R \), considerados como clases, corresponden a la identidad en \( \langle X : R \rangle \).
	}
\end{observacion}

\begin{observacion}
	\upshape{
		Si \( X = \varnothing \), entonces el grupo libre generado por \( X \) es el grupo trivial:
		\[
		G_X = \{e\}.
		\]
		En consecuencia, para cualquier conjunto de relaciones \( R \subseteq G_X \), se tiene que \( N_R = \{e\} \), y por tanto
		\[
		G_X / N_R = \{eN_R\} = \{N_R\},
		\]
		lo cual significa que se trata de un grupo trivial.
		
		\medskip
		
		\noindent Por tanto, \( \langle \emptyset : R \rangle \) es la presentación de todo grupo trivial.
	}
\end{observacion}

Habiendo establecido la definición y notación de presentación de un grupo, la pregunta natural es cómo construirlas. El método principal para obtener presentaciones consiste en usar la propiedad universal del grupo libre junto con el teorema fundamental del homomorfismo de grupos. A continuación, el Ejemplo \ref{poiuyt} ilustrará este proceso con detalle.

\begin{ejemplo}\label{poiuyt}
	\upshape{
	
		El grupo ciclico $$\dfrac{\mathbb{Z}}{n\mathbb{Z}}=\{\:\overline{0},\overline{1},\overline{2},\cdots, \overline{n-1}\:\}$$ se puede escribir como $\langle x: x^n=e\rangle$.
	}
\end{ejemplo}


\begin{proof}[Justificación]
	Consideremos \( X := \{x\} \) y la función 
	\begin{align*}
		j : X &\longrightarrow \dfrac{\mathbb{Z}}{n\mathbb{Z}}\\[4pt]
		x &\longmapsto \overline{1}.
	\end{align*}
	Por la propiedad universal del grupo libre, existe un único homomorfismo de grupos \( f : G_X \to \mathbb{Z}/n\mathbb{Z} \) tal que \( f \circ \iota = j \):
	\[
	\xymatrix{
		X \ar[rr]^{\displaystyle j} \ar@{^(->}[dd]_{\displaystyle \iota } & & \dfrac{\mathbb{Z}}{n\mathbb{Z}} \\
		&&\\
		G_X \ar@{-->}[rruu]_{\displaystyle \exists !f} &&
	}
	\]  
	De donde se tiene que 
	\begin{align*}
		f : G_X &\longrightarrow \dfrac{\mathbb{Z}}{n\mathbb{Z}} \\[4pt]
		[x]^m &\longmapsto m\;\overline{1}.
	\end{align*}
	Ya que como \( f \) es un homomorfismo de grupos, se cumple que 
	\[
	f\left([x]^m\right) = mf\left([x]\right)= m\left((f\circ \iota)(x)\right) = mj(x) = m\;\overline{1}.
	\] Además es fácil ver que es sobreyectivo. Luego, por el teorema fundamental del homomorfismo de grupos, se tiene que 
	\[
	\dfrac{G_{X}}{\ker f} \cong \dfrac{\mathbb{Z}}{n\mathbb{Z}},
	\]
	por lo que nos interesa ahora saber cómo es el núcleo de \( f \). 
	\begin{align*}
		\ker f &= \left\{\:[x]^m : f([x]^m) = \overline{0}\:\right\}\\[4pt]
		&= \{\:[x]^m : m\;\overline{1} = \overline{0}\:\}\\[4pt]
		&= \{\:[x]^m : m \text{ es múltiplo de } n\:\}\\[4pt]
		&= \langle [x]^n \rangle.
	\end{align*} 
	Así, tenemos \( X = \{x\} \) y consideramos \( R = \{[x]^n\} \subset G_X \); podemos verificar facilmente que, \( N_R = \langle [x]^n\rangle\), y por definición se tiene que  
	\[
	\dfrac{G_{X}}{\ker f} = \dfrac{G_{\{x\}}}{\langle [x]^n \rangle} =: \langle x : [x]^n = e \rangle.
	\] 	Por otro lado, notemos que \( [x] = x \), pues \( x \) es una palabra reducida (no se puede reducir más). Entonces se tiene, por igualdad de conjuntos, que 
	\[
	\langle x : [x]^n = e \rangle = \langle\, x : x^n = e\, \rangle.
	\] 	Por lo tanto, finalmente obtenemos
	\[
	\dfrac{\mathbb{Z}}{n\mathbb{Z}} \cong \langle \,x : x^n = e \,\rangle
	\]
\end{proof}




\begin{ejemplo}
	\upshape{
		El grupo de los números enteros $\mathbb{Z}$ se puede escribir como $\langle x : \emptyset \rangle$.
	}
\end{ejemplo}

\begin{proof}[Justificación]
	Consideremos \( X := \{x\} \) y la función 
	\begin{align*}
		j : X &\longrightarrow \mathbb{Z}\\[4pt]
		x &\longmapsto 1.
	\end{align*}
	Por la propiedad universal del grupo libre, existe un único homomorfismo de grupos \( f : G_X \to \mathbb{Z} \) tal que \( f \circ \iota = j \):
	\[
	\xymatrix{
		X \ar[rr]^{\displaystyle j} \ar@{^(->}[dd]_{\displaystyle \iota } & & \mathbb{Z} \\
		&&\\
		G_X \ar@{-->}[rruu]_{\displaystyle \exists !f} &&
	}
	\]	
	De donde se tiene que 
	\begin{align*}
		f : G_X &\longrightarrow \mathbb{Z} \\[4pt]
		[x]^m &\longmapsto m \cdot 1 = m.
	\end{align*}	
	Ya que como \( f \) es un homomorfismo de grupos, se cumple que 
	\[
	f\left([x]^m\right) = m f\left([x]\right) = m\left((f\circ \iota)(x)\right) = m j(x) =  m1 = m.
	\]Además es sobreyectivo. Luego, por el teorema fundamental del homomorfismo de grupos, se tiene que 
	\[
	\dfrac{G_{X}}{\ker f} \cong \mathbb{Z},
	\]
	por lo que nos interesa ahora saber cómo es el núcleo de \( f \). 
	\begin{align*}
		\ker f &= \left\{\:[x]^m : f([x]^m) = 0\:\right\}\\[4pt]
		&= \{\:[x]^m : m = 0\:\}\\[4pt]
		&= \{\:[x]^0 \:\}\\[4pt]
		&=\{[e]\}=\{e\}\\[4pt]
		&= \langle e \rangle.
	\end{align*}	
	Así, tenemos \( X = \{x\} \) y consideramos \( R = \emptyset \subset G_X \); esto es, \( N_R = \ker f \), y por definición se tiene que 	
	\[
	\dfrac{G_{X}}{\ker f} = \dfrac{G_{\{x\}}}{\langle e \rangle} =:\langle x : w=e,w\in R \rangle= \langle x : \emptyset \rangle.
	\] 	Por lo tanto, obtenemos
	\[
	\mathbb{Z} \cong \langle \,x : \emptyset \,\rangle.
	\]
\end{proof}


\section{Grupo simétrico}


Una permutación es, en esencia, una reordenación de los elementos de un conjunto. Podemos entenderla formalmente como una función biyectiva.

\begin{definicion}
	\upshape{
		Sea \( X \) un conjunto no vacío. Una función \(\sigma : X \to X\) biyectiva se
		denomina permutación del conjunto \( X \).
	}
\end{definicion}

La composición de biyecciones del conjunto \( X \) vuelve a ser una biyección, y la función identidad es una biyección. Además, toda biyección tiene inversa que también es biyección. Así, se cumple que:

\begin{itemize}
	\item La composición de permutaciones es una operación interna.
	\item La composición es asociativa.
	\item La función identidad es la permutación identidad.
	\item Dada \(\sigma\) una permutación, entonces la función inversa \(\sigma^{-1}\) es la permutación inversa.
\end{itemize}

Por lo que tenemos una estructura de grupo.

\begin{definicion}
	\upshape{
		El conjunto de todas las permutaciones (biyecciones) de un conjunto \( X \) con la composición de funciones forma un grupo llamado \textit{grupo de permutaciones} (grupo de biyecciones), denotado por
		\[
		\operatorname{Biyec}(X).
		\]
	}
\end{definicion}

\begin{definicion}
	\upshape{
		Si el conjunto \( X \) es un subconjunto de los números naturales, es decir,
		\[
		X = \{1, 2, \ldots, n\},
		\]
		entonces su correspondiente grupo de permutaciones \(\operatorname{Biyec}(X)\) se denomina \emph{grupo simétrico} de \(n\) elementos y se denota usualmente como \(S_n\) o también como \(\Sigma_n\).
		
		\medskip
		\noindent
		A los elementos del grupo simétrico se les denomina \emph{permutaciones de \(n\) letras}.
	}
\end{definicion}

\begin{observacion}
	\upshape
	El grupo simétrico \(\Sigma_n\) tiene  orden \(n!\).
\end{observacion}



\begin{ejemplo}
	\upshape
	Sea \(X = \{1; 2\}\). Entonces, las únicas permutaciones de \(X\) son:
	\[
	\begin{array}{rcl}
		\sigma_1: & X \to X & \\
		& 1 \mapsto 1, & \\
		& 2 \mapsto 2 &
	\end{array}
	\qquad \text{y} \qquad
	\begin{array}{rcl}
		\sigma_2: & X \to X & \\
		& 1 \mapsto 2, & \\
		& 2 \mapsto 1 &
	\end{array}
	\]
	Con lo cual, \(\Sigma_2 = \{\sigma_1, \sigma_2\}\). Más aún, un elemento \(\sigma \in \Sigma_2\) queda determinado por su imagen, y podemos escribir:
	\[
	\sigma = \left(\begin{array}{cc}
		1 & 2 \\
		\sigma(1) & \sigma(2)
	\end{array}\right).
	\]
	En este caso:
	\[
	\sigma_1 = \left(\begin{array}{cc}
		1 & 2 \\
		1 & 2
	\end{array}\right)
	\qquad \text{y} \qquad
	\sigma_2 = \left(\begin{array}{cc}
		1 & 2 \\
		2 & 1
	\end{array}\right). 
	\]
\end{ejemplo}

\medskip
\begin{observacion}
	\upshape{
		Note que $\sigma_1\sigma_2=\sigma_2\sigma_1$ por lo que $\Sigma_2$ es un grupo abeliano.
	}
\end{observacion}

\medskip
\begin{ejemplo}
	\upshape{
		Sea \(X = \{1, 2, 3\}\). Entonces, \(\Sigma_3\) tiene \(3! = 6\) elementos:
		
		
		\[
		I = \left(\begin{array}{ccc}
			1 & 2 & 3 \\
			1 & 2 & 3
		\end{array}\right), \quad
		\alpha = \left(\begin{array}{ccc}
			1 & 2 & 3 \\
			2 & 3 & 1
		\end{array}\right), \quad
		\alpha' = \left(\begin{array}{ccc}
			1 & 2 & 3 \\
			3 & 1 & 2
		\end{array}\right),
		\]
		\[
		\beta = \left(\begin{array}{ccc}
			1 & 2 & 3 \\
			3 & 2 & 1
		\end{array}\right), \quad
		\beta' = \left(\begin{array}{ccc}
			1 & 2 & 3 \\
			1 & 3 & 2
		\end{array}\right), \quad
		\beta'' = \left(\begin{array}{ccc}
			1 & 2 & 3 \\
			2 & 1 & 3
		\end{array}\right).
		\]
		
		\medskip
		\noindent Además, observamos que:
		\begin{align*}
			\alpha^2 &= \alpha \alpha = \left(\begin{array}{ccc}
				1 & 2 & 3 \\
				3 & 1 & 2
			\end{array}\right) = \alpha', \\[1ex]
			\alpha^3 &= \alpha \alpha^2 = \left(\begin{array}{ccc}
				1 & 2 & 3 \\
				1 & 2 & 3
			\end{array}\right) = I, \\[1ex]
			\alpha \beta &= \left(\begin{array}{ccc}
				1 & 2 & 3 \\
				1 & 3 & 2
			\end{array}\right) = \beta', \\[1ex]
			\alpha^2 \beta &= \left(\begin{array}{ccc}
				1 & 2 & 3 \\
				2 & 1 & 3
			\end{array}\right) = \beta''.
		\end{align*}
		Por lo tanto, 
		\[
		\Sigma_3 = \{I, \alpha, \alpha^2, \beta, \alpha \beta, \alpha^2 \beta\}.
		\]
	}
\end{ejemplo}

\begin{observacion}
	\upshape{
		Podemos obtener la siguiente tabla de operaciones de los elementos del grupo simétrico $\Sigma_3$.
		\begin{table}[H]
			\centering
			{\large
				\renewcommand{\arraystretch}{2.2}% Aumenta el espacio entre filas
				\resizebox{160mm}{!}{
					\begin{tabular}{
							>{\centering\arraybackslash}p{3cm}
							>{\centering\arraybackslash}p{3cm}
							>{\centering\arraybackslash}p{3cm}
							>{\centering\arraybackslash}p{3cm}
							>{\centering\arraybackslash}p{3cm}
							>{\centering\arraybackslash}p{3cm}
							>{\centering\arraybackslash}p{3cm}
						}
						\toprule
						$\scalebox{1.5}{$\circ$}$ & $\scalebox{1.5}{$I$}$ & $\scalebox{1.5}{$\alpha$}$ & $\scalebox{1.5}{$\alpha^2$}$ & $\scalebox{1.5}{$\beta$}$ & $\scalebox{1.5}{$\alpha\beta$}$ & $\scalebox{1.5}{$\alpha^2\beta$}$ \\
						\midrule
						$\scalebox{1.5}{$I$}$ & $\scalebox{1.5}{$I$}$ & $\scalebox{1.5}{$\alpha$}$ & $\scalebox{1.5}{$\alpha^2$}$ & $\scalebox{1.5}{$\beta$}$ & $\scalebox{1.5}{$\alpha\beta$}$ & $\scalebox{1.5}{$\alpha^2\beta$}$ \\
						$\scalebox{1.5}{$\alpha$}$ & $\scalebox{1.5}{$\alpha$}$ & $\scalebox{1.5}{$\alpha^2$}$ & $\scalebox{1.5}{$I$}$ & $\scalebox{1.5}{$\alpha\beta$}$ & $\scalebox{1.5}{$\alpha^2\beta$}$ & $\scalebox{1.5}{$\beta$}$ \\
						$\scalebox{1.5}{$\alpha^2$}$ & $\scalebox{1.5}{$\alpha^2$}$ & $\scalebox{1.5}{$I$}$ & $\scalebox{1.5}{$\alpha$}$ & $\scalebox{1.5}{$\alpha^2\beta$}$ & $\scalebox{1.5}{$\beta$}$ & $\scalebox{1.5}{$\alpha\beta$}$ \\
						$\scalebox{1.5}{$\beta$}$ & $\scalebox{1.5}{$\beta$}$ & $\scalebox{1.5}{$\alpha^2\beta$}$ & $\scalebox{1.5}{$\alpha\beta$}$ & $\scalebox{1.5}{$I$}$ & $\scalebox{1.5}{$\alpha^2$}$ & $\scalebox{1.5}{$\alpha$}$ \\
						$\scalebox{1.5}{$\alpha\beta$}$ & $\scalebox{1.5}{$\alpha\beta$}$ & $\scalebox{1.5}{$\beta$}$ & $\scalebox{1.5}{$\alpha^2\beta$}$ & $\scalebox{1.5}{$\alpha$}$ & $\scalebox{1.5}{$I$}$ & $\scalebox{1.5}{$\alpha^2$}$ \\
						$\scalebox{1.5}{$\alpha^2\beta$}$ & $\scalebox{1.5}{$\alpha^2\beta$}$ & $\scalebox{1.5}{$\alpha\beta$}$ & $\scalebox{1.5}{$\beta$}$ & $\scalebox{1.5}{$\alpha^2$}$ & $\scalebox{1.5}{$\alpha$}$ & $\scalebox{1.5}{$I$}$ \\
						\bottomrule
					\end{tabular}
			}}
			\caption{Tabla de operaciones del grupo simétrico \(\Sigma_3\).}
		\end{table}
		\noindent De esta tabla podemos ver que $\Sigma_3$ no es un grupo abeliano, pues existe al menos un par de elementos que no conmutan:
		
		\[
		\alpha \beta =
		\left(
		\begin{array}{ccc}
			1 & 2 & 3 \\ 
			3 & 1 & 2
		\end{array}
		\right)
		\neq 
		\left(
		\begin{array}{ccc}
			1 & 2 & 3 \\ 
			2 & 3 & 1
		\end{array}
		\right)
		= \beta \alpha.
		\]
		
	}
\end{observacion}

\medskip
\begin{proposicion}
	\upshape{
		El grupo simétrico \(\Sigma_n\) es abeliano si \(n = 2\) y  es no abeliano si \(n \geq 3\).
	}
\end{proposicion}

\begin{proof}
	Si \(n = 2\), entonces \(\Sigma_2\) contiene solo dos permutaciones:
	\[
	I = \left(\begin{array}{cc}
		1 & 2 \\
		1 & 2
	\end{array}\right),
	\qquad
	\sigma = \left(\begin{array}{cc}
		1 & 2 \\
		2 & 1
	\end{array}\right).
	\]
	Como \(I \sigma = \sigma I = \sigma\), se concluye que \(\Sigma_2\) es abeliano.
	
	\medskip
	\noindent Supongamos ahora que \(n \geq 3\). Definimos dos permutaciones en \(\Sigma_n\):
	\[
	\sigma = \left(\begin{array}{ccccccc}
		1 & 2 & 3 & 4 & \cdots & n \\
		2 & 3 & 1 & 4 & \cdots & n
	\end{array}\right),
	\qquad
	\rho = \left(\begin{array}{ccccccc}
		1 & 2 & 3 & 4 & \cdots & n \\
		1 & 3 & 2 & 4 & \cdots & n
	\end{array}\right).
	\]
	Estas permutaciones actúan no trivialmente solo sobre las posiciones \(1\), \(2\) y \(3\), dejando fijas las demás.
	Entonces:
	\[
	\sigma\rho = \left(\begin{array}{ccccccc}
		1 & 2 & 3 & 4 & \cdots & n \\
		2 & 1 & 3 & 4 & \cdots & n
	\end{array}\right)\quad\text{y}\quad
	\rho\sigma = \left(\begin{array}{ccccccc}
		1 & 2 & 3 & 4 & \cdots & n \\
		3 & 2 & 1 & 4 & \cdots & n
	\end{array}\right).
	\]
	
	\medskip
	\noindent Como
	\[
	\sigma\rho \neq \rho\sigma,
	\]
	se concluye que \(\Sigma_n\) no es abeliano para \(n \geq 3\).
\end{proof}

Escribir una permutación \(\sigma \in \Sigma_n\) como  
\[
\sigma = \left(\begin{array}{ccccccc}
	1 & 2 & 3 & 4 & \cdots & n \\
	\sigma(1) & \sigma(2) & \sigma(3) & \sigma(4) & \cdots & \sigma(n)
\end{array}\right)
\]
resulta redundante, ya que siempre se repite la primera fila de manera fija. Por esta razón, se introduce una notación más eficiente y concisa conocida como \textit{notación cíclica}.

\medskip
\noindent Para comprender esta nueva forma de denotar permutaciones, consideremos el siguiente ejemplo:
\[
\sigma = \left(\begin{array}{cccccccc}
	1 & 2 & 3 & 4 & 5 & 6 & 7 & 8 \\
	4 & 2 & 7 & 5 & 1 & 8 & 6 & 3
\end{array}\right)\in \Sigma_8.
\]

\noindent Las imágenes sucesivas de \(1\) mediante \(\sigma\) son:
\[
\sigma^0(1) = 1,\quad \sigma(1) = 4,\quad \sigma^2(1) = 5,\quad \sigma^3(1) = 1,\quad \text{y} \quad \sigma^j(1) = \sigma^{j-3}(1)\text{ para }j \geq 4.
\]
Entonces, estas imágenes sucesivas de \(1\) mediante \(\sigma\) se escriben en forma cíclica como:
\[
\sigma_1 := (\:1\quad 4\quad 5\:).
\]

\noindent Tomamos ahora el menor elemento de \(\{1,2,3,\dots,n\}\), con \(n = 8\), que no esté en el conjunto \(\{1,4,5\}\). Ese elemento es \(2\), y sus imágenes mediante \(\sigma\) son:
\[
\sigma^0(2) = 2,\quad \sigma(2) = 2,
\]
por lo tanto, escribimos:
\[
\sigma_2 := (2).
\]

\noindent De modo similar, tomamos el menor elemento que no está en \(\{1,2,4,5\}\), es decir, \(3\), y calculamos:
\[
\sigma^0(3) = 3,\quad \sigma(3) = 7,\quad \sigma^2(3) = 6,\quad \sigma^3(3) = 8,\quad \sigma^4(3) = 3,
\]
lo que se denota cíclicamente como:
\[
\sigma_3 := (\:3\quad 7\quad 6\quad 8\:).
\]

\noindent Finalmente, se intentaría continuar con el siguiente menor elemento que no está en \(\{1,2,3,4,5,6,7,8\}\), pero ya no quedan más elementos. Por lo tanto, hemos terminado el proceso y podemos escribir:
\[
\sigma = \sigma_3 \sigma_2 \sigma_1.
\]
\noindent A cada \(\sigma_j\) se le denomina \emph{ciclo}, y la función \(k \mapsto \sigma_j(k)\) se define así: si el ciclo es
\[
\sigma_j = (\:a_1\quad a_2\quad a_3\quad\cdots\quad a_r\:), \quad r \leq n,
\]
entonces, para \(k \in \{1,2,\dots,n\}\), se define:
\[
\sigma_j(k) := \begin{cases}
	a_{s+1} & \text{si } k = a_s \in \{a_1, a_2, \dots, a_{r-1}\}, \\
	a_1     & \text{si } k = a_r, \\
	k       & \text{si } k \notin \{a_1, a_2, \dots, a_r\}.
\end{cases}
\]

\noindent Así, por ejemplo:
\begin{itemize}
	\item \(\sigma(1) = \sigma_3\sigma_2\sigma_1(1) = 4\), porque en el ciclo \(\sigma_1 = (1\ 4\ 5)\), el número que sigue a \(1\) es \(4\); además, \(\sigma_2(4) = 4\) y \(\sigma_3(4) = 4\).
	\item \(\sigma(4) = \sigma_3\sigma_2\sigma_1(4) = 5\), porque en el ciclo \((1\ 4\ 5)\), el número que sigue a \(4\) es \(5\).
	\item \(\sigma(5) = \sigma_3\sigma_2\sigma_1(5) = 1\), porque en el ciclo \((1\ 4\ 5)\), el número que sigue a \(5\) es \(1\).
	\item \(\sigma(2) = \sigma_3\sigma_2\sigma_1(2) = 2\), porque el ciclo \(\sigma_2 = (2)\) deja a \(2\) fijo.
	\item \(\sigma(3) = 7\), \(\sigma(7) = 6\), \(\sigma(6) = 8\) y \(\sigma(8) = 3\), por lectura directa del ciclo \((3\ 7\ 6\ 8)\).
\end{itemize}

y recupereamos \[
\sigma =\sigma_3\sigma_2\sigma_1= \left(\begin{array}{cccccccc}
	1 & 2 & 3 & 4 & 5 & 6 & 7 & 8 \\
	4 & 2 & 7 & 5 & 1 & 8 & 6 & 3
\end{array}\right)\in \Sigma_8.
\]


\begin{observacion}
	\upshape{
		Cuando un ciclo es trivial, como en el caso de \(\sigma_2 = (2)\), se suele omitir en la notación cíclica. Por tanto, simplemente escribiremos:
		\[
		\sigma = \sigma_3 \sigma_2 \sigma_1 = \sigma_3 \sigma_1.
		\]
		O más simple aún, podemos escribir directamente:
		\[
		\sigma = (\:3\quad 7\quad 6\quad 8\:)(\:1\quad 4\quad 5\:).
		\]
	}
\end{observacion}

\begin{ejemplo}
	\upshape{
		La permutación \[
		\alpha = \left(\begin{array}{cccccccc}
			1 & 2 & 3 & 4 & 5 & 6 \\
			3 & 1 & 4 & 2 & 6 & 5 
		\end{array}\right)\in \Sigma_6
		\] se escribe en notación ciclica de la forma $\alpha=(\:5\quad 6\:)(\:1\quad 3\quad 4\quad 2\:)$
	}
\end{ejemplo}

\begin{observacion}
	\upshape{
		En general dos ciclos no conmutan. Por ejemplo, en $\Sigma_6$ consideremos los ciclos $\beta_1=(\:1\quad 3\quad 5\:)$ y $\beta_2=(\:3\quad 5\quad 6\:)$ y se tiene que $\beta_2\beta_1\neq \beta_1\beta_2$ ya que
		\[
		\beta_2\beta_1 = \left(\begin{array}{cccccccc}
			1 & 2 & 3 & 4 & 5 & 6  \\
			5 & 2 & 6 & 4 & 1 & 3 
		\end{array}\right)=(\:3\quad 6\:)(\:1\quad 5\:)
		\]
		
		\[
		\beta_1\beta_2 = \left(\begin{array}{cccccccc}
			1 & 2 & 3 & 4 & 5 & 6  \\
			3 & 2 & 1 & 4 & 6 & 5 
		\end{array}\right)=(\:5\quad 6\:)(\:1\quad 3\:).
		\]
	}
\end{observacion}

\medskip
Una pregunta natural que surge al trabajar con permutaciones es: ¿cuándo dos ciclos conmutan, es decir, cuándo su composición da el mismo resultado independientemente del modo en que se combinen? La Proposición \ref{propojhgf} proporciona una condición sencilla y útil para saber cuándo dos ciclos pueden conmutar.


\begin{definicion}
	\upshape{
		Dos ciclos $\rho = (a_1\ a_2\ \cdots\ a_p)$ y $\gamma = (b_1\ b_2\ \cdots\ b_q)$ se denominan \textit{disjuntos} si no tienen ningún elemento en común, es decir, si
		\[
		\{a_1, a_2, \dots, a_p\} \cap \{b_1, b_2, \dots, b_q\} = \emptyset.
		\]
	}
\end{definicion}

\begin{proposicion}\label{propojhgf}
	\upshape{
		Si $\rho$ y $\gamma$ son ciclos disjuntos en $\Sigma_n$, entonces conmutan: 
		\[
		\rho\gamma = \gamma\rho.
		\]
	}
\end{proposicion}

\begin{proof}
	Sea \( x \in \{1, 2, \dots, n\} \). Examinamos los posibles casos:
	
	\begin{itemize}
		\item Si \( x \) pertenece al ciclo \( \rho \) (es decir, \( x \in \{a_1, \dots, a_p\} \)), entonces \( \gamma(x) = x \) porque \( \rho \) y \( \gamma \) son disjuntos. Así,
		\[
		\rho\gamma(x) = \rho(x) = \gamma\rho(x).
		\]
		
		\item Si \( x \in \{b_1, \dots, b_q\} \), entonces \( \rho(x) = x \), y del mismo modo:
		\[
		\rho\gamma(x) = \gamma(x) = \gamma\rho(x).
		\]
		
		\item Si \( x \notin \{a_1, \dots, a_p, b_1, \dots, b_q\} \), entonces tanto \( \rho(x) = x \) como \( \gamma(x) = x \), y por tanto:
		\[
		\rho\gamma(x) = x = \gamma\rho(x).
		\]
	\end{itemize}
	En todos los casos se tiene que \( \rho\gamma(x) = \gamma\rho(x) \), por lo tanto, \( \rho\gamma = \gamma\rho \).
\end{proof}
\begin{observacion}
	\upshape{
		La recíproca de la Proposición \ref{propojhgf} no es válida. Es decir, aunque dos ciclos no vacíos conmutan, no necesariamente deben ser disjuntos. Esto ocurre, por ejemplo, cuando ambos ciclos describen la misma permutación. Por ejemplo, los ciclos
		\[
		\rho = (1\ 4\ 5) \quad \text{y} \quad \gamma = (4\ 5\ 1)
		\]
		no son disjuntos (pues comparten todos sus elementos), pero satisfacen \(\rho\gamma = \gamma\rho\), ya que producen el mismo efecto al aplicarse. En estos casos, la conmutatividad se debe a que ambos ciclos representan la misma permutación, y no a que sean disjuntos.
	}
\end{observacion}


\begin{proposicion}\label{poiuyhgf}
	\upshape{
		Toda permutación \(\sigma \in \Sigma_n\) puede escribirse como composición de ciclos disjuntos:
		\[
		\sigma = \sigma_1 \circ \sigma_2 \circ \cdots \circ \sigma_k,
		\]
		donde cada \(\sigma_i\) es un ciclo y los ciclos son disjuntos entre sí. Esta descomposición es única salvo la secuencia en que se componen los ciclos y la elección del primer elemento en cada ciclo.
	}
\end{proposicion}
\begin{proof}
	Para \(\sigma \in \Sigma_n\), consideremos el conjunto \(\{1, 2, \ldots, n\}\). Empezamos con el menor elemento \(a_1 = 1\) y formamos su ciclo:
	\[
	\sigma_1 = (a_1 \quad \sigma(a_1) \quad \sigma^2(a_1) \quad \cdots \quad \sigma^{m_1 - 1}(a_1)),
	\]
	donde \(m_1\) es el menor entero positivo tal que \(\sigma^{m_1}(a_1) = a_1\). Este \(m_1\) existe, pues \(\Sigma_n\) es finito y \(\sigma: \Sigma_n \to \Sigma_n\).
	
	\medskip
	\noindent Luego, tomamos el menor elemento \(a_2\) que no pertenece al conjunto \(\{a_1, \sigma(a_1), \ldots, \sigma^{m_1 - 1}(a_1)\}\) y definimos su ciclo
	\[
	\sigma_2 = (a_2 \quad \sigma(a_2) \quad \cdots \quad \sigma^{m_2 - 1}(a_2)),
	\]
	con \(m_2\) definido análogamente.
	
	\medskip
	\noindent
	Repetimos este proceso, que es finito, hasta cubrir todos los elementos de \(\{1, \ldots, n\}\). Por construcción, los ciclos \(\sigma_i\) son disjuntos, pues no comparten elementos, y si \(\sigma_k\) es el último ciclo obtenido, tenemos:
	\[
	\sigma = \sigma_1 \circ \sigma_2 \circ \cdots \circ \sigma_k.
	\]
	
	\medskip
	\noindent
	\vspace{1ex}
	Unicidad: Supongamos que existe otra descomposición en ciclos disjuntos de \(\sigma\):
	\[
	\sigma = \tau_1 \circ \tau_2 \circ \cdots \circ \tau_p,
	\]
	donde cada \(\tau_j\) es un ciclo disjunto.
	
	\medskip
	\noindent
	Sea \(a \in \{1,2,\dots,n\}\) y definamos la órbita de \(a\) bajo \(\sigma\) como
	\[
	O(a) = \{a, \sigma(a), \sigma^2(a), \dots, \sigma^{m_a -1}(a)\},
	\]
	donde \(m_a\) es el menor entero positivo tal que \(\sigma^{m_a}(a) = a\).
	
	\medskip
	\noindent
	Dado que los ciclos en ambas descomposiciones son disjuntos, el único ciclo en cada descomposición que mueve a \(a\) debe contener exactamente los elementos de \(O(a)\). Por tanto, existe un ciclo \(\sigma_i\) en la primera descomposición y un ciclo \(\tau_j\) en la segunda, tales que
	\[
	\sigma_i = (a \quad \sigma(a) \quad \sigma^2(a) \quad \dots \quad \sigma^{m_a -1}(a)) \quad \text{y} \quad \tau_j = (a \quad \sigma(a) \quad \sigma^2(a) \quad \dots \quad \sigma^{m_a -1}(a)).
	\]
	
	Como esto ocurre para todo \(a \in \{1,2,\dots,n\}\), las dos descomposiciones coinciden salvo en la forma en que se presentan, es decir, en la secuencia en la que aparecen los ciclos y en la elección del elemento inicial de cada ciclo.
	
	\medskip
	\noindent 
	En otros términos, la unicidad se debe a que los ciclos son disjuntos y, por la Proposición \ref{propojhgf}, la composición de ciclos disjuntos conmutan, lo que garantiza que la secuencia en la que se componen los ciclos no afecta el resultado final. Finalmente, la elección del primer elemento de cada ciclo no afecta la permutación que representa dicho ciclo, pues cualquier ciclo puede escribirse comenzando por cualquiera de sus elementos sin cambiar la permutación.
	
\end{proof}




\begin{definicion}
	\upshape{
		Una permutación \(\sigma \in \Sigma_n\) es un \(r\)-ciclo si se puede escribir como un ciclo de la forma
		\[
		\sigma = (a_1 \ a_2 \ \dots \ a_r),
		\]
		donde los elementos \(a_1, a_2, \dots, a_r\) son distintos entre sí. Es decir, \(\sigma(a_i) = a_{i+1}\) para \(i = 1, \dots, r-1\), y \(\sigma(a_r) = a_1\), mientras que \(\sigma(x) = x\) para todo \(x \notin \{a_1, \dots, a_r\}\).
	}
\end{definicion}


\begin{ejemplo}
	\upshape{
		En \(\Sigma_n\), \((\:1\quad 2\quad 3\:)\) es un \(3\)-ciclo. Además, notemos que:
		
		\[
		\begin{aligned}
			(\:1\quad 2\quad 3\:)&= 
			\left(\begin{array}{ccccccc}
				1 & 2 & 3 & 4 & \cdots & n \\
				2 & 3 & 1 & 4 & \cdots & n
			\end{array}\right)
			\\[2ex]
			(\:1\quad 2\quad 3\:)^{\circ 2} &= 
			\left(\begin{array}{ccccccc}
				1 & 2 & 3 & 4 & \cdots & n \\
				3 & 1 & 2 & 4 & \cdots & n
			\end{array}\right)
			\\[2ex]
			(\:1\quad 2\quad 3\:)^{\circ 3} &= 
			\left(\begin{array}{ccccccc}
				1 & 2 & 3 & 4 & \cdots & n \\
				1 & 2 & 3 & 4 & \cdots & n
			\end{array}\right) = I
		\end{aligned}
		\]
	}
\end{ejemplo}




\begin{proposicion}\label{proppaosjha}
	\upshape{
		Si \(\sigma\) es un \(r\)-ciclo, entonces \(\sigma^{\circ r} = I\), donde \(I\) es la permutación identidad. Además, \(\sigma^{\circ (r-1)} \neq I\).
	}
\end{proposicion}


\begin{proof}
	Supongamos que \(\sigma = (a_1\ a_2\ \dots\ a_r)\) es un \(r\)-ciclo en \(\Sigma_n\). Por definición de ciclo:
	\[
	\sigma(a_1) = a_2,\quad \sigma(a_2) = a_3,\quad \dots,\quad \sigma(a_{r-1}) = a_r,\quad \sigma(a_r) = a_1,
	\]y \(\sigma(x) = x\) para todo \(x \in \{1, 2, \dots, n\}\) tal que \(x \notin \{a_1, \dots, a_r\}\).
	
	\medskip
	\noindent
	Apliquemos \(\sigma\) varias veces a un elemento \(a_i\):
	
	
	\[
	\sigma^{\circ k}(a_i) = a_{i+k},
	\] donde los subíndices se toman módulo \(r\), entendiéndolo como residuo euclideano no negativo; es decir,
	\[
	a_{i+k} := a_{((i+k-1) \bmod r) + 1}.
	\]
	
	\medskip\noindent
	Por ejemplo, tomando \(6 \bmod 4\) (entendido como residuo euclideano no negativo, el cual pertenece a \(\{0,1,2,3\}\)) tenemos que \(6 \bmod 4 = 2\).
	
	\medskip
	\noindent
	Entonces, dado que \((i+r-1) \bmod r = i-1\), al aplicar \(\sigma\) consigo misma \(r\) veces obtenemos:
	\[
	\sigma^{\circ r}(a_i) = a_{i+r} = a_{(i-1)+1}=a_i,
	\]
	lo cual vale para todo \(i = 1, \dots, r\), y, además, claramente:
	\[
	\sigma^{\circ r}(x) = x \quad \text{para todo } x \notin \{a_1, \dots, a_r\}.
	\]
	Por tanto, \(\sigma^{\circ r} = I\).
	
	\medskip
	\noindent
	Para ver que \(\sigma^{\circ (r-1)} \ne I\), basta notar que:
	\[
	\sigma^{\circ (r-1)}(a_1) = a_r \ne a_1,
	\]	pues los elementos \(a_1, \dots, a_r\) son distintos. Por lo tanto, \(\sigma^{\circ (r-1)} \ne I\).
\end{proof}


\begin{definicion}\label{jhgsd}
	\upshape{
		Sea \(\sigma \in \Sigma_n\). Si \(\sigma\) es un \(2\)-ciclo, entonces \(\sigma\) se denomina\textit{ transposición}.
	}
\end{definicion}

\begin{observacion}
	\upshape{
		Podríamos haber definido la transposición simplemente diciendo que es la permutación que mueve exactamente dos elementos distintos \(i \neq j\). Sin embargo, dar la Definición \ref{jhgsd} resulta útil en las demostraciones, mientras que entenderla como se menciona al inicio es práctico en aplicaciones.
	}
\end{observacion}
\begin{observacion}\label{hghssd}
	\upshape{
		La inversa de una transposición es también una transposición. En efecto, si \( \tau = (i\ j) \) con \( i \neq j \), entonces \( \tau \circ \tau = \text{id} \), lo cual implica que \( \tau^{-1} = \tau \), y por tanto es nuevamente una transposición.
	}
\end{observacion}

\begin{proposicion}\label{kjhgfhj}
	\upshape{
		Todo ciclo se puede escribir como composición de transposiciones.
	}
\end{proposicion}

\begin{proof}
	Sea \(\sigma = (a_1\ a_2\ \dots\ a_r)\) un \(r\)-ciclo en \(\Sigma_n\), con \(r \geq 2\). Entonces, \(\sigma\) puede escribirse como la siguiente composición de \(r-1\) transposiciones:
	\[
	\sigma = (a_1\ a_r) \circ (a_1\ a_{r-1}) \circ \cdots \circ (a_1\ a_2).
	\]
	Verifiquemos que esta expresión es correcta. Al aplicar las composiciones, la permutación resultante es
	\[ (\:a_2\quad a_3\ \dots\ a_{r-1}\quad a_r\quad a_1\:)
	\]
	que coincide con \(\sigma\) pues ambos ciclos nos dan la misma permutación.\\
	
	\noindent Si $\sigma =(a_j)$, que podemos suponer $j\neq1 $, entonces podemos hacer $$\sigma=(a_1\ a_j)\circ(a_j\ a_1)$$ 
	
	\noindent En todo los casos podemos ver que el ciclo se puede expresar como composición de transposiciones. 
\end{proof}




\begin{proposicion}
	\upshape{
		Toda permutación se puede escribir como composición de transposiciones:
		\[
		\sigma = \tau_r \circ \tau_{r-1} \circ \cdots \circ \tau_2 \circ \tau_1,
		\]
		donde \( r \) es único salvo paridad.
	}
\end{proposicion}

\begin{proof}
	Por la Proposición \ref{poiuyhgf}, toda permutación se puede escribir como composición de ciclos disjuntos, y por la Proposición \ref{kjhgfhj}, cada uno de esos ciclos se puede escribir como composición de transposiciones. Por tanto, toda permutación puede expresarse como composición de transposiciones.
	
	\medskip
	\noindent
	Resta demostrar que si una permutación se escribe como composición de \( r \) transposiciones y también como composición de \( s \) transposiciones, entonces \( r \equiv s \pmod{2} \); es decir, que la paridad del número de transposiciones es única, aunque la cantidad exacta pueda variar.
	
	\medskip
	\noindent
	Supongamos que \( \sigma = \tau_1 \circ \cdots \circ \tau_r = \tau'_1 \circ \cdots \circ \tau'_s \) son dos descomposiciones de \( \sigma \) en transposiciones. Entonces,
	\[
	\tau_r^{-1} \circ \cdots \circ \tau_1^{-1} \circ \tau'_1 \circ \cdots \circ \tau'_s = \text{id},
	\]
	y por la Observación \ref{hghssd} tenemos 
	\[
	\tau_r \circ \cdots \circ \tau_1 \circ \tau'_1 \circ \cdots \circ \tau'_s = \text{id},
	\]
	es decir, la composición de \( r + s \) transposiciones es la permutación identidad.
	
	\medskip
	\noindent
	Finalmente, notemos que identidad no puede escribirse como composición de un número impar de transposiciones. Por tanto, \( r + s \) debe ser par, lo que implica que \( r \equiv s \pmod{2} \), como queríamos.
\end{proof}

\begin{ejemplo}
	\upshape{
		Sea 
		\[
		\sigma = \left(\begin{array}{ccccc}
			1 & 2 & 3 & 4 & 5 \\
			2 & 4 & 5 & 3 & 1
		\end{array}\right) \in \Sigma_5.
		\]
		Determinemos su signo.
		
		\medskip
		\noindent
		Primero escribimos \( \sigma \) como composición de ciclos disjuntos. Observamos que:
		\[
		\sigma = (1\ 2\ 4\ 3\ 5),
		\]
		que es un ciclo de longitud \(5\).
		
		\medskip
		\noindent
		Luego, este ciclo puede escribirse como una composición de \(4\) transposiciones:
		\[
		(1\ 2\ 4\ 3\ 5) = (1\ 5)\circ(1\ 3)\circ(1\ 4)\circ(1\ 2).
		\]
		Por tanto, el número de transposiciones necesarias es \(r=4\), que es par. Así que:
		\[
		\operatorname{sgn}(\sigma) = (-1)^4 = 1.
		\]
	}
\end{ejemplo}



\begin{ejemplo}[Matriz de permutación]
	\upshape{
		Sea \(\mathbb{K} = \mathbb{R}\) o \(\mathbb{C}\). Sea \(\sigma \in \Sigma_n\), y sea \(E \in \mathbb{K}^{n \times n}\) la matriz identidad, que puede escribirse como:
		\[
		E = \left[\:e_1\quad e_2\quad \cdots\quad e_n\right],
		\]
		donde cada \(e_i \in \mathbb{K}^n\) es la \(i\)-ésima columna canónica, es decir, el vector con \(1\) en la posición \(i\) y ceros en las demás. Explícitamente:
		\[
		e_i = \left(\begin{array}{c}
			0\\
			\vdots\\
			0\\
			1\\
			0\\
			\vdots\\
			0
		\end{array}\right) \in \mathbb{K}^n,
		\]
		con el \(1\) en la \(i\)-ésima posición.
		
		\medskip
		\noindent
		La \textit{matriz de permutación} asociada a \(\sigma\), denotada por \(M_\sigma\), es la matriz obtenida al permutar las columnas de \(E\) según \(\sigma\), es decir:
		\[
		M_\sigma := \left[\:e_{\sigma(1)}\quad e_{\sigma(2)}\quad \cdots\quad e_{\sigma(n)}\right].
		\]
		
		\medskip
		\noindent
		Por ejemplo, sea \(n=3\) y sea \(\sigma = (1\ 3\ 2)\), es decir,
		\[
		\sigma(1) = 3, \quad \sigma(2) = 1, \quad \sigma(3) = 2.
		\]
		Entonces, la matriz de permutación \(M_\sigma \in \mathbb{K}^{3 \times 3}\) es:
		\[
		M_\sigma = \begin{pmatrix}
			0 & 0 & 1 \\
			1 & 0 & 0 \\
			0 & 1 & 0
		\end{pmatrix}.
		\]
		Si aplicamos esta matriz a un vector columna \(x = (x_1, x_2, x_3)^T \in \mathbb{K}^3\), el resultado es:
		\[
		M_\sigma x = \begin{pmatrix} x_3 \\ x_1 \\ x_2 \end{pmatrix}.
		\]
		
		\medskip
		\noindent
		Notemos que \(M_\sigma\), la matriz de permutación asociada a \(\sigma \in \Sigma_n\), reorganiza las entradas de un vector columna de acuerdo al inverso de la permutación \(\sigma\). Esto es importante tenerlo en cuenta, ya que puede resultar confuso al inicio.
	}
\end{ejemplo}


\begin{observacion}
	\upshape{
		Las matrices de permutación no siempre son simétricas. Una matriz de permutación \(M_\sigma\) es simétrica si y solo si la permutación \(\sigma\) es una involución (es decir, \(\sigma^2 = I\)).
		
		\medskip
		\noindent Por ejemplo, para \(\sigma = (1\ 3\ 2)\), la matriz de permutación \(M_\sigma\) no es simétrica:
		\[
		\begin{pmatrix}
			0 & 0 & 1 \\
			1 & 0 & 0 \\
			0 & 1 & 0
		\end{pmatrix}
		\neq
		\begin{pmatrix}
			0 & 1 & 0 \\
			0 & 0 & 1 \\
			1 & 0 & 0
		\end{pmatrix}=\begin{pmatrix}
			0 & 0 & 1 \\
			1 & 0 & 0 \\
			0 & 1 & 0
		\end{pmatrix}^{T}.
		\]
		Es importante destacar que toda matriz de permutación es una matriz ortogonal, lo que significa que su transpuesta es igual a su inversa ($M_\sigma^T = M_\sigma^{-1}$).
	}
\end{observacion}




\begin{proposicion}
	\upshape{
		Sea \(A = [V_1 \quad V_2 \quad \cdots \quad V_n] \in \mathbb{K}^{n \times n}\), y sea \(M_\sigma\) la matriz de permutación asociada a \(\sigma \in \Sigma_n\). Entonces:
		\[
		AM_\sigma = [V_{\sigma(1)} \quad V_{\sigma(2)} \quad \cdots \quad V_{\sigma(n)}].
		\]
	}
\end{proposicion}

\begin{proof}
	Por definición, \(M_\sigma = [e_{\sigma(1)} \quad e_{\sigma(2)} \quad \cdots \quad e_{\sigma(n)}]\), donde \(e_j\) es la \(j\)-ésima columna de la matriz identidad \(E\).
	
	\medskip
	\noindent
	Recordemos que al multiplicar una matriz \(A\) por una columna \(e_k\), se obtiene la \(k\)-ésima columna de \(A\), es decir, \(Ae_k = V_k\). Así, la \(j\)-ésima columna de \(AM_\sigma\) es:
	\[
	AM_\sigma^{(j)} = A e_{\sigma(j)} = V_{\sigma(j)}.
	\]
	Esto es válido para todo \(j = 1, \dots, n\), por lo tanto:
	\[
	AM_\sigma = [V_{\sigma(1)} \quad V_{\sigma(2)} \quad \cdots \quad V_{\sigma(n)}],
	\]
	como se quería demostrar.
\end{proof}

\begin{proposicion}
	\upshape{
		Sea \(\sigma \in \Sigma_n\) una permutación, y sea \(M_\sigma\) la matriz de permutación asociada. Entonces:
		\[
		\det\left(M_\sigma\right) = \operatorname{sgn}(\sigma).
		\]
	}
\end{proposicion}

\begin{proof}
	Recordemos que la matriz de permutación \(M_\sigma \in \mathbb{K}^{n \times n}\) se obtiene al permutar las columnas de la matriz identidad \(E = [e_1 \ e_2 \ \cdots \ e_n]\) según la permutación \(\sigma\), es decir,
	\[
	M_\sigma = [e_{\sigma(1)} \quad e_{\sigma(2)} \quad \cdots \quad e_{\sigma(n)}].
	\]
	
	\medskip
	\noindent
	El intercambiar dos columnas cambia el signo del determinante. Como toda permutación puede escribirse como composición de transposiciones, y cada transposición corresponde a un intercambio de dos columnas, basta notar que cada transposición cambia el signo del determinante.
	
	\medskip
	\noindent
	Así, si \(\sigma = \tau_1 \circ \cdots \circ \tau_r\) es una descomposición de \(\sigma\) en \(r\) transposiciones, se tiene:
	\[
	\det(M_\sigma) = (-1)^r.
	\]
	
	\noindent
	Por definición del signo de una permutación, \(\operatorname{sgn}(\sigma) := (-1)^r\). Como la paridad de \(r\) es independiente de la descomposición elegida, se concluye que:
	\[
	\det(M_\sigma) = \operatorname{sgn}(\sigma).
	\]
\end{proof}













\section{Teorema de Cayley}

Sea $G$ un grupo, entonces existe un homomorfismo de grupos inyectivo 
\begin{align*}
	\phi: G &\longrightarrow \operatorname{Biyec}(G) \\
	a &\longmapsto \left(  
	\begin{array}{rcl}
		\phi_a:G &\to& G \\
		x &\mapsto& ax
	\end{array} \right)
\end{align*}




\noindent En efecto. Fijado $a \in G$, debemos mostrar primero que
\[
\begin{array}{rcl}
	\phi_a: G &\to& G \\
	x &\mapsto& ax
\end{array}
\]
es una biyección que, por definición, es una permutación:
\begin{itemize}
	\item Inyectividad: Sean $x, y \in G$ tales que $\phi_a(x) = \phi_a(y)$, entonces $ax = ay$. Y como $G$ es grupo, se sigue que $x = y$.
	
	\item Sobreyectividad: Sea $z \in G$, entonces existe $x = a^{-1}z$ tal que $\phi_a(x) = z$.
\end{itemize}

\noindent Ahora debemos ver que  
\begin{align*}
	\phi: G &\longrightarrow \operatorname{Biyec}(G) \\
	a &\longmapsto \phi(a) := \phi_a
\end{align*}
es un homomorfismo inyectivo. Sean $a, b \in G$:
\begin{itemize}
	\item Homomorfismo: Sea $x \in G$, entonces
	\[
	\phi_{ab}(x) = (ab)x = a(bx) = \phi_a(bx) = \phi_a(\phi_b(x)) = (\phi_a \circ \phi_b)(x).
	\]
	Así, por la definición, se tiene que
	\[
	\phi(ab) = \phi(a) \circ \phi(b).
	\]
	
	\item Inyectividad: Sea $\phi(a) = \phi(b)$, entonces por definición $\phi_a(x) = \phi_b(x)$ para todo $x \in G$. Luego $ax = bx$ y, como $G$ es grupo, se concluye que $a = b$.
\end{itemize}

\noindent Con lo desarrollado, tenemos el teorema de Cayley:

\begin{teorema}
	\upshape{
		Todo grupo $G$ es isomorfo a un subgrupo del grupo de permutaciones de $G$. En particular, si $G$ es finito, entonces $G$ es isomorfo a un subgrupo del grupo simétrico de $G$.
	}
\end{teorema}

\begin{proof}
	De lo desarrollado basta notar que como $\phi: G \to \operatorname{Biyec}(G)$ es un homomorfismo inyectivo de grupos, entonces $\phi(G)$ es un subgrupo de $\operatorname{Biyec}(G)$ que es isomorfo al grupo $G$.
\end{proof}






\section{Grupo diedral}


%\begin{definicion}
%	\upshape{
	%		El \emph{grupo diedral} de orden \( 2n \), denotado \( D_n \), es el subgrupo de \( \operatorname{Sym}(n) \) generado por un \( n \)-ciclo \( r := (1\ 2\ \cdots\ n) \), y una transposici\'on/reflexi\'on
	%		\[
	%		s := (1\ n)(2\ n{-}1)\cdots.
	%		\]
	%		Este grupo modela las simetr\'ias (rotaciones y reflexiones) de un pol\'igono regular de \( n \) lados.
	%	}
%\end{definicion}

\begin{definicion}\label{deiugv}
	\upshape{
		El \textit{grupo diedral} \( D_n \) de grado \( n \geq 3 \) se define como el subgrupo de \( \Sigma_n \) generado por el \( n \)-ciclo
		\[
		a = (\:1\quad 2\quad 3\quad \cdots\quad n\:),
		\]
		y por la permutación
		\[
		b = \left(\begin{array}{cccccc}
			1 & 2 & 3 & \cdots & n-1 & n \\
			1 & n & n-1 & \cdots & 3 & 2
		\end{array}\right).
		\]
	}
\end{definicion}

\medskip
\begin{ejemplo}
	\upshape{
		El grupo diedral de grado \(3\) es el subgrupo de $\Sigma_3$ dado por 
		\[
		D_3 := \left\langle (\:1\quad 2\quad 3\:),\quad  
		\left(\begin{array}{ccc}
			1 & 2 & 3 \\
			1 & 3 & 2 
		\end{array}\right) \right\rangle.
		\]
		Otra forma de  obtener o construir este grupo diedral es mediante movimientos de rotación y reflexión de un triángulo equilátero en el plano.
		
		\medskip
		
		%% Primer diagrama: Rotación con flecha más grande y aún más separación
		\begin{center}
			\begin{tikzpicture}
				% Triángulo original
				\draw[thick] (0,0) -- (2,0) -- (1,1.732) -- cycle;
				\node[below left] at (0,0) {1};
				\node[below right] at (2,0) {3};
				\node[above] at (1,1.732) {2};
				
				% Flecha indicando rotación (más grande y centrada)
				\draw[->, thick] (3,0.866) -- (7,0.866) node[midway, above] {Rotación};
				
				% Triángulo con numeración rotada (2,3,1) más separado
				\draw[thick] (8,0) -- (10,0) -- (9,1.732) -- cycle;
				\node[below left] at (8,0) {2};
				\node[below right] at (10,0) {1};
				\node[above] at (9,1.732) {3};
			\end{tikzpicture}
		\end{center}
		
		\medskip
		
		% Segundo diagrama: Reflexión con eje de simetría
		\begin{center}
			\begin{tikzpicture}
				% Triángulo original
				\draw[thick] (0,0) -- (2,0) -- (1,1.732) -- cycle;
				\node[below left] at (0,0) {1};
				\node[below right] at (2,0) {3};
				\node[above] at (1,1.732) {2};
				
				% Flecha indicando reflexión
				\draw[->, thick] (3,0.866) -- (7,0.866) node[midway, above] {Reflexión};
				
				% Triángulo reflejado con numeración ajustada
				\draw[thick] (8,0) -- (10,0) -- (9,1.732) -- cycle;
				\node[above left] at (8,0) {1};
				\node[below right] at (10,0) {2};
				\node[above] at (9,1.732) {3};
				
				% Eje de simetría (mediana desde vértice 1, extendida simétricamente)
				\draw[dashed, thick] (7.25,-0.433) -- (10.3,1.326) 
				node[pos=0.468, below right, sloped, text depth=0.1ex] 
				{\scriptsize eje de simetría};
			\end{tikzpicture}
		\end{center}
	}
\end{ejemplo}

\vspace{0.5cm}




\begin{proposicion}\label{lkjhgjkl}
	\upshape{
	De la Definción \ref{deiugv} se tiene que 
	\begin{enumerate}
		\item El orden de $\langle a\rangle$ es $n$ y el orden de $\langle b\rangle$ es $2$.
		\item  $b^ra^{s}=a^{-s}b^r$ para todo $r,s\in \mathbb{Z}$ con $r\neq 0$.
		\item $\langle a\rangle$  es un subgrupo normal de $D_n$.
		\item $\langle a \rangle\langle b \rangle=D_n$ y $\langle a \rangle\cap\langle b \rangle=\{I\}$.
	\end{enumerate}
	}
\end{proposicion}

\begin{proof}
	\textcolor{white}{text}
	
	\medskip
	\noindent $[1]$ Como $a$ es un $n$-ciclo, por la Proposición \ref{proppaosjha} tenemos que $a^{\circ n}=I$ y que $a^{\circ(n-1)}\neq I$. Entonces el orden de $a$ es $n$. Así, el orden de $\langle a\rangle$ es $n$. 
	
	\medskip
	\noindent Por otro lado, claramente $b\neq I$ y $b^{\circ 2}=I$. Entonces el orden del subgrupo  $b$ es $2$. Así, el orden de $\langle b\rangle$ es $2$.         \\
	
	\noindent $[2]$ Veamos inicialmente que $ba=a^{-1}b$. En efecto,  
	\[
	ba = \begin{pmatrix}
		1 & 2 & 3 & \cdots & n \\
		1 & n & n-1 & \cdots & 2
	\end{pmatrix}
	(1\ 2\ 3\ \cdots\ n) 
	= \begin{pmatrix}
		1 & 2 & 3 & 4 & \cdots & n-1 & n \\
		n & 1 & 2 & 3 & \cdots & n-2 & n-1
	\end{pmatrix}
	\]     
	
	\[
	a^{-1}b = (n\ n-1\ \cdots\ 2\ 1)\begin{pmatrix}
		1 & 2 & 3 & \cdots & n \\
		1 & n & n-1 & \cdots & 2
	\end{pmatrix} = \begin{pmatrix}
		1 & 2 & 3 & 4 & \cdots & n-1 & n \\
		n & 1 & 2 & 3 & \cdots & n-2 & n-1
	\end{pmatrix}
	\] Así, claramente tenemos que $ba=a^{-1}b$.
	
	\medskip
	\noindent Veamos ahora que $ba^{s}=a^{-s}b$ para todo $s\in \mathbb{Z}$. En efecto:
	\begin{itemize}
		\item Para $s=0$ está claro que es válido.
		\item Es válido para $s\geq 1$: Para $s=1$ hemos visto que es válido. Supongamos que $ba^s=a^{-s}b$ para $s\in \mathbb{N}$ con $s>1$. Entonces $ba^{s+1}=ba^sa=a^{-s}ba=a^{-s}a^{-1}b=a^{-(s+1)}b$.
		\item Es válido para $s\leq -1$: Como $-s\geq 1$ tenemos entonces que $ba^{-s}=a^{-(-s)}b=a^sb$. Entonces $$ba^s=b(a^sb)b=b(ba^{-s})b=a^{-s}b.$$
	\end{itemize}
	\noindent Veamos ahora que dado $s\in \mathbb{Z}$ se cumple que $b^ra^{s}=a^{-s}b^r$ para todo $r\in \mathbb{Z}\setminus\{0\}$. En efecto:
	\begin{itemize}
		\item Es válido para $r\geq 1$: Para $s=1$ hemos visto que es válido. Supongamos que $b^ra^s=a^{-s}b^r$ para $r\in \mathbb{N}$ con $r>1$. Entonces $b^{r+1}a^{s}=bb^ra^s=ba^{-s}b^s=a^{-s}bb^r=a^{-s}b^{r+1}$.
		\item Es válido para $r\leq -1$: Como $-r\geq 1$ tenemos entonces que $b^{-r}a^{s}=a^{-s}b^{-r}$. Entonces \[
		b^r a^{-s} = b^r \left( a^{-s} b^{-r} \right) b^r = b^r \left( b^{-r} a^{s}\right) b^r=a^sb^r.
		\]
		Y como $s\in \mathbb{Z}$ podemos reescribir 
		\[
		b^r a^{s} =a^{-s}b^r.
		\]
	\end{itemize}
	
	
	\noindent $[3]$ Probaremos que $d\langle a\rangle d^{-1}\subset \langle a\rangle$ para todo $d\in D_n$. En efecto:
	\begin{itemize}
		\item Si $d=a^mb^n$:   	
		\begin{align*}
			da^k d^{-1} 
			&= (a^m b^n) a^k (b^{-n} a^{-m}) \\
			&= a^m (b^n a^k) (a^{m} b^{-n}) \\
			&= a^m a^{-k} (b^n a^{m}) b^{-n} \\
			&= a^m a^{-k} a^{-m} b^n b^{-n} \\
			&= a^{-k} \in \langle a \rangle.
		\end{align*}
		
		\item Si $d=b^na^m$: 
		\begin{align*}
			da^k d^{-1} 
			&= (b^n a^m) a^k (a^{-m} b^{-n}) \\
			&= b^na^{k}b^{-n} \\
			&= a^{-k} b^n b^{-n} \\
			&= a^{-k} \in \langle a \rangle.
		\end{align*}
	\end{itemize}    
	Así, tenemos que $\langle a\rangle \trianglelefteq D_n$.\\ 
	
	\noindent $[4]$ En efecto, $\langle a\rangle \langle b \rangle=\{a^{\circ m}b^{\circ n}:m,n\in\mathbb{Z}\}=\langle a,b\rangle=D_n$.
	
	\medskip
	\noindent Por otro lado. Demostremos que $b \notin \langle a \rangle$. Supongamos por contradicción que existe $k \in \mathbb{Z}$ tal que $b = a^k$. Analizando
	\begin{itemize}
		\item Por definición, $b(1) = 1$
		\item $a^k(1) = 1$ solo si:
		\begin{itemize}
			\item Cuando $k = 0$: $a^0 = I$, pero $b \neq I$ pues $b(2) = n \neq 2$
			\item Cuando $k > 0$: $a^k(1) = 1$ implica que $k$ es múltiplo de $n$ (pues $a$ es un $n$-ciclo), luego $a^k = I$, contradicción
			\item Cuando $k < 0$: similarmente lleva a $a^k = I$, contradicción
		\end{itemize}
	\end{itemize}
	Por tanto, no existe tal $k$ y $b \notin \langle a \rangle$. Así $\langle a\rangle \cap \langle b\rangle=\{I\}$.
	
	
\end{proof}

\begin{observacion}
	\upshape{
		De la Proposición \ref{lkjhgjkl} tenemos que $$a^{\circ n}=I ,\qquad b^{\circ 2}=I,\qquad bab=a^{-1}$$
	}
\end{observacion}
Esta observación nos permite dar la presentación del grupo diedral $D_n$.  
\begin{proposicion}
	\upshape{
		El grupo diedral puede escribirse como
		\[
		D_n = \langle x, y : x^n = e,\ y^2 = e,\ yxy = x^{-1} \rangle.
		\]
	}
\end{proposicion}

\begin{proof}[Justificación]
	Consideremos \( X := \{x, y\} \) y definimos una función
	\begin{align*}
		j : X &\longrightarrow D_n \\[4pt]
		x &\longmapsto a \\[2pt]
		y &\longmapsto b
	\end{align*}
	donde \( D_n = \langle a, b \rangle \) que satisfacen:
	\[
	a^n = I,\quad b^2 = I,\quad baba= I.
	\]
	
	\medskip
	
	\noindent Por la propiedad universal del grupo libre, existe un único homomorfismo de grupos
	$f : G_X \longrightarrow D_n$ tal que $f \circ \iota = j$
	\[
	\xymatrix{
		X \ar[rr]^{\displaystyle j} \ar@{^(->}[dd]_{\displaystyle \iota } & & D_n \\ 
		&&\\
		G_X \ar@{-->}[rruu]_{\displaystyle \exists !f} &&
	}
	\] Entonces, el homomorfismo \( f \) se determina en los generadores por:
	\begin{align*}
		f([x]) &= (f\circ\iota)(x)=j(x)=a \\[4pt]
		f([y]) &= (f\circ \iota)(y)=j(y)=b
	\end{align*}
	y, por ser homomorfismo de grupos, se tiene que
	\begin{align*}
		f([x]^n) &= f([x])^n = a^n = I, \\[4pt]
		f([y]^2) &= f([y])^2 = b^2 = I, \\[4pt]
		f([y][x][y][x]) &=f([x])f([y])f([x])f([y]) = abab = I.
	\end{align*}
	Además es un homomorfismo sobreyectivo pues dado $z=a^\alpha b^\beta\in D_n$ existe $[x]^\alpha[y]^\beta\in G_X$ de modo que $f\left([x]^\alpha[y]^\beta\right)=f([x])^\alpha f([y])^\beta=a^\alpha b^\beta=z$.
	
	\medskip
	\noindent
	Definimos el conjunto:
	\[
	R := \{ [x]^n,\ [y]^2,\ [y][x][y][x] \} \subset G_X.
	\] Sea \( N_R \) el subgrupo normal de \( G_X \) generado por \( R \).	Por el teorema fundamental del homomorfismo de grupos, se tiene:
	\[
	G_X / \ker f \cong D_n,
	\]
	Nos interesa saber cómo es el núcleo de $f$. Notemos que $$H:= \langle [x]^n ,\ [y]^2 ,\ [y][x][y][x] \rangle =\ker f.$$ En efecto
	\begin{itemize}
		\item Como los generadores de $H$ se anulan mediante $f$, entonces $H\subset \ker f$.
		\item Notemos primero que como $bab=a^{-1}$ y $b^2=I$, entonces $ba=a^{-1}b$. Sea $w\in \ker f$, entonces existe $k\in \mathbb{N} $ tal que
		\[
		w = [x]^{\alpha_1}[y]^{\beta_1}[x]^{\alpha_2}[y]^{\beta_2} \cdots [x]^{\alpha_k}[y]^{\beta_k}, 
		\quad \text{donde } \alpha_i,\ \beta_i \in \mathbb{Z},\ \text{para } i = 1, \cdots, k
		\] de modo que \[
		f(w) = a^{\alpha_1} b^{\beta_1} a^{\alpha_2} b^{\beta_2} \cdots a^{\alpha_k} b^{\beta_k}=I
		\] Ahora como $ba=a^{-1}b$, entonces $ba^{r}=a^{-r}b$ para todo $r\in \mathbb{Z}$ con lo cual $$b^sa^{r}=a^{-r}b^s,\quad \text{para todo }s,r\in \mathbb{Z}.$$ Entonces
		\[
		f(w) = a^{\alpha_1-\alpha_2-\alpha_3-\cdots-\alpha_k} b^{\beta_1+\beta_2+\cdots+\beta_k}=I
		\]
		Así, \begin{itemize}
			\item $\beta_1+\beta_2+\cdots+\beta_k$ es múltiplo  dos.
			\item $\alpha_1-\alpha_2-\alpha_3-\cdots-\alpha_k$ es múltiplo de $n$. 
			\item 
		\end{itemize}
		
		
	\end{itemize}
	
	
	
	
	
	
	qui
	Debemos probar que $H=\ker f$. La primera inclusión es trivial por lo visto arriba. Sea $w\in \ker(f)$ debemos ver que $w\in H$. En efecto 
	
	
	es un subgrupo normal y por tanto $N_R=H$. En efecto
	
	
	
	
	
	
	
	
	
	
	
	\begin{align*}
		\ker f&=\{\:\xi\in G_X: f\left(\xi\right)=I\:\}\\[4pt]
	\end{align*}
	
	\[
	D_n \cong G_X / N_R = \langle x, y : [x]^n = e,\ [y]^2 = e,\ [y][x][y][x] = e \rangle.
	\]
	
	Finalmente, como las palabras \( [x] \), \( [y] \), \( [x]^n \), etc., están en forma reducida, podemos omitir los corchetes:
	\[
	\langle x, y : [x]^n = e,\ [y]^2 = e,\ [y][x][y][x] = e \rangle = \langle x, y : x^n = e,\ y^2 = e,\ yxy = x^{-1} \rangle.
	\]
\end{proof}









\section{Acción de un grupo}

Sea \( X \) un conjunto no vacío y \( G \) un grupo.

\begin{definicion}
	\upshape{
	Se dice que \( G \) actúa sobre \( X \) si existe una función, denominada \textit{acción de grupo},
	\[
	\varphi : G \times X \to X
	\]
	tal que, para todo \( a, b \in G \) y todo \( x \in X \), se cumplen:
	\begin{enumerate}
		\item Identidad: \( \varphi(e, x) = x \), donde \( e \) es el elemento neutro de \( G \).
		\item Compatibilidad: \( \varphi(ab, x) = \varphi(a, \varphi(b, x)) \).
	\end{enumerate}
	}
\end{definicion}



\begin{observacion}
	\upshape{
		La compatibilidad por sí sola no garantiza que \( \varphi \) sea una acción de grupo: también se requiere la neutralidad.
		
		\medskip
		\noindent Por ejemplo. Sea \( G = (\mathbb{Z}, +) \) y sea \( X = \{a, b\} \). Definimos la aplicación \( \varphi: \mathbb{Z} \times X \to X \) como sigue:
		\[
		\varphi(n, a) = b, \quad \varphi(n, b) = a, \quad \text{para todo } n \in \mathbb{Z}.
		\]
		Es decir, \( \varphi(n, x) \) intercambia \( a \) y \( b \), sin importar el valor de \( n \). Y tenemos
		
		\begin{itemize}
			\item Compatibilidad: Para todo \( n, m \in \mathbb{Z} \) y \( x \in X \), se cumple:			 
			\[
			\varphi(n, \varphi(m, x)) = \varphi(n, y) = x=\varphi(n + m, x),
			\]
			
			\item No se cumple la indentidad. La identidad de \( G \) es \( 0 \), pero:
			\[
			\varphi(0, a) = b \neq a, \quad\text{ y }\quad \varphi(0, b) = a \neq b.
			\]
			Por tanto, la condición $\varphi(e,x)=x$ para todo $x\in X$ falla.
		\end{itemize}
		Concluimos que \( \varphi \) no define una acción de grupo, a pesar de que satisface la compatibilidad.
	}
\end{observacion}




\begin{observacion}\label{hjgfhjkl}
	\upshape{
		Si \( G \) actúa sobre \( X \) mediante \( \varphi \), entonces, para todo \( a \in G \) y \( x \in X \), se cumple:
		\[
		\varphi(a, \varphi(a^{-1}, x)) = x 
		\quad \text{y} \quad 
		\varphi(a^{-1}, \varphi(a, x)) = x.
		\]
		Esto significa que cada aplicación inducida \( \varphi_a : X \to X \), definida por \( \varphi_a(x) := \varphi(a, x) \), es una biyección. Además,
		\[
		\varphi_a^{-1} = \varphi_{a^{-1}}.
		\]
	}
\end{observacion}



\begin{ejemplo}[Acción regular izquierda]
	\upshape{
		Sea \( X = G \). La aplicación \( \varphi(a, x) = ax \) define una acción de \( G \) sobre \( G \).
	}
\end{ejemplo}

\begin{ejemplo}[Acción regular derecha]
	\upshape{
		Sea \( X = G \). No se puede definir la acción como \( \varphi(a, x) = xa \), ya que no cumple la compatibilidad. En su lugar, se define:
		\[
		\varphi(a, x) = x a^{-1},
		\]
		que sí verifica las condiciones de una acción de grupo.
	}
\end{ejemplo}

\begin{ejemplo}[Acción por conjugación]
	\upshape{
		Sea \( X = G \). La aplicación \( \varphi(a, x) = axa^{-1} \) define una acción de \( G \) sobre \( G \) por conjugación.
	}
\end{ejemplo}


\begin{ejemplo}[Acción de traslación]
	\upshape{
		Sea \( G \) un grupo y \( A \) un conjunto. Consideramos el conjunto \( A^G \) de todas las funciones de \( G \) en \( A \). Definimos la aplicación
			\begin{align*}
			\varphi: G\times A^G &\longrightarrow A^G \\[4pt]
			(b,f) &\longmapsto \left(  
			\begin{array}{rcl}
				\varphi(b,f) : G &\to& A \\[2pt]
				x &\mapsto& f(b^{-1}x)
			\end{array} \right)
		\end{align*}
		Esta aplicación define una acción de \( G \) sobre \( A^G \), llamada \emph{acción de traslación (izquierda)}.
	}
\end{ejemplo}



\begin{proposicion}
	\upshape{
		Las siguientes afirmaciones son equivalentes:
		\begin{enumerate}
			\item Toda acción \( \varphi: G \times X \to X \) induce un homomorfismo de grupos \( G \to \operatorname{Biyec}(X) \). Es decir, la aplicación
			\begin{align*}
				\rho: G &\longrightarrow \operatorname{Biyec}(X) \\[4pt]
				a &\longmapsto \left(  
				\begin{array}{rcl}
					\varphi_a : X &\to& X \\[2pt]
					x &\mapsto& \varphi(a, x)
				\end{array} \right)
			\end{align*}
			es un homomorfismo de grupos.
			
			\item Todo homomorfismo de grupos \( \rho: G \to \operatorname{Biyec}(X) \) induce una acción. Es decir, la aplicación 
			\begin{align*}
				\varphi: G \times X &\longrightarrow X\\
				(a,x) &\longmapsto \rho_a(x),
			\end{align*}
			donde \( \rho(a) := \rho_a \in \operatorname{Biyec}(X) \).
		\end{enumerate}
	}
\end{proposicion}


\begin{proof}
	\textcolor{white}{text}
	
	\medskip
	\noindent
	\([1 \Rightarrow 2]\) Supongamos que \( \varphi(ab, x) = \varphi(a, \varphi(b,x)) \) para todo \( a,b \in G \), \( x \in X \).
	
	\medskip
	\noindent 
	Queremos ver que \( \varphi_a \in \operatorname{Biyec}(X) \) para cada \( a \in G \), y que la aplicación
	\[
	\rho : G \to \operatorname{Biyec}(X),\quad a \mapsto \varphi_a,
	\]
	es un homomorfismo de grupos.
	
	\begin{itemize}
		\item Biyectividad: Por la Observación~\ref{hjgfhjkl}, cada \( \varphi_a \) es una biyección con inversa dada por \( \varphi_{a^{-1}} \).
		
		\item Homomorfismo: Para todo \( a,b \in G \) y todo \( x \in X \), se cumple:
		\[
		\varphi_{ab}(x) = \varphi(ab, x) = \varphi(a, \varphi(b, x)) = \varphi_a(\varphi_b(x)) = (\varphi_a \circ \varphi_b)(x).
		\]
		Entonces \( \rho(ab) = \rho(a)\circ\rho(b) \), por lo tanto \( \rho \) es un homomorfismo de grupos.
	\end{itemize}
	
	\medskip
	
	\noindent\([2 \Rightarrow 1]\) Supongamos que \( \rho : G \to \operatorname{Biyec}(X) \), \( a \mapsto \rho_a \), es un homomorfismo de grupos. Entonces, para todo \( a,b \in G \) y \( x \in X \), se cumple:
	\[
	\varphi(ab, x) := \rho_{ab}(x) = (\rho_a \circ \rho_b)(x) = \rho_a(\rho_b(x)) = \varphi(a, \varphi(b, x)),
	\]
	lo que demuestra la compatibilidad. Finalmente, veamos que \( \varphi(e,x) = x \). En efecto, como \( \rho \) es un homomorfismo de grupos, se tiene \( \rho_e = \operatorname{id}_X \), por lo que:
	\[
	\varphi(e,x) = \rho_e(x) = \operatorname{id}_X(x) = x.
	\]
\end{proof}


\begin{observacion}
	\upshape{
		Es usual omitir la notación explícita de la acción \( \varphi: G \times X \to X \) y, en su lugar, denotar simplemente que \( G \) actúa sobre \( X \) escribiendo \( G : X \to X \), y usar la notación \( a \cdot x \) o simplemente \( ax \) para representar la acción de \( a \in G \) sobre \( x \in X \). Con esta convención, las propiedades de la acción se expresan como:
		\[
		a(bx) = (ab)x \quad \text{y} \quad ex = x.
		\]
	}
\end{observacion}
\vspace{0.5cm}

 Cuando \( G \) es un monoide, la acción ya no induce un homomorfismo hacia \( \operatorname{Biyec}(X) \), pues los elementos de \( G \) pueden no ser invertibles. En su lugar, la acción define una aplicación en el monoide de funciones:\[
\operatorname{Fun}(X) := \{ f: X \to X \mid f \text{ es una función} \}.
\]
Con la composición como operación, este conjunto forma un monoide donde las transformaciones no requieren ser biyectivas ni tener inversa.
		
\begin{proposicion}
	\upshape{
		Sea \( M \) un monoide. Son equivalentes:
		\begin{enumerate}
			\item \( M \) actúa sobre \( X \) si existe una función $\varphi : M \times X \to X$ tal que, para todo \( a, b \in M \) y todo \( x \in X \), se cumplen:
			\[
			\varphi(ab, x) = \varphi(a, \varphi(b, x)) \quad\text{ y }\quad \varphi(e, x) = x.
			\]
		
			
			\item \( M \) actúa sobre \( X \) si existe un homomorfismo de monoides $M \to \operatorname{Fun}(X)$.			
			
		\end{enumerate}
	}
\end{proposicion}





\begin{ejemplo}[Aplicación a sistemas dinámicos continuos: Problema de Cauchy]\label{enkbhjgf}
	\upshape{
		Sea \( E \) un espacio de Banach. Sea \( A: E \to E \) un operador cerrado, es decir, una aplicación que lleva conjuntos cerrados en conjuntos cerrados. Considere el problema de Cauchy
		\begin{equation*}
			\left\{
			\begin{aligned}
				u'(t) &= A \cdot u(t), \\
				u(0) &= u_0.
			\end{aligned}
			\right.
		\end{equation*}
		La existencia y unicidad de la solución (global) implica la existencia de una acción del monoide \( (\mathbb{R}_{\geq 0}, +) \), dada por
		\begin{align*}
			\varphi: \mathbb{R}_{\geq 0} \times E &\longrightarrow E \\[4pt]
			(t, u_0) &\longmapsto e^{tA} u_0.
		\end{align*}
	}
\end{ejemplo}


\begin{definicion}
	\upshape{
		Si \( \varphi : G \times X \to X \) es una acción, entonces la \textbf{órbita} de \( x \in X \) es el conjunto
		\[
		G \cdot x := \{ \varphi(a, x) \::\: a \in G \},
		\]
		también denotado por \( \operatorname{Orb}(x) \).
	}
\end{definicion}

\begin{observacion}
	\upshape{
		Obtenemos la misma definición si \( G \) se reemplaza por un monoide \( M \), es decir, si \( \varphi : M \times X \to X \) es una acción de monoide, entonces la órbita de \( x \in X \) se define igual como
		\[
		M \cdot x := \{ \varphi(m, x) \::\: m \in M \},
		\]
		y también se denota por \( \operatorname{Orb}(x) \).
	}
\end{observacion}


\begin{ejemplo}
	\upshape{
		Del Ejemplo \ref{enkbhjgf} (problema de Cauchy), se tiene que
		\[
		\operatorname{Orb}(u_0) = \{ e^{tA} u_0 \::\: t \geq 0 \},
		\]
		esto es, el conjunto de soluciones del sistema dinámico que pasan por \( u_0 \) en \( t = 0 \).
		
		
		
	\begin{center}
		\begin{tikzpicture}[>=Stealth, scale=1.5]
			% Axes
			\draw[->] (-0.5,0) -- (7.5,0) node[right]{$t$};  % Extended x-axis
			\draw[->] (0,-0.5) -- (0,3.5) node[above]{$E$};   % Extended y-axis
			
			% 5 parallel curves (orbits)
			\foreach \y in {0.5,1,1.5,2,2.5} {
				\draw[thick] (0,\y) .. controls (2,\y+0.7) and (5,\y+0.3) .. (7,\y+1);
			}
			
			% Highlighted orbit (u_0)
			\draw[ultra thick] (0,1.5) .. controls (2,2.2) and (5,1.8) .. (7,2.5) 
			node[midway, above left]{$\operatorname{Orb}(u_0)$};
			
			% Initial point
			\filldraw (0,1.5) circle (1.5pt) node[below left]{$u_0$};
		\end{tikzpicture}
	\end{center}
	}
\end{ejemplo}
\begin{ejemplo}[Aplicación a sistemas dinámicos discretos]
	\upshape{
		Sea \( X \) un espacio topológico y definamos
		\[
		\operatorname{Homeo}(X) := \{f : X \to X \mid f \text{ es un homeomorfismo} \}.
		\]
		Sea \( f \in \operatorname{Homeo}(X) \) fijado. Entonces, la aplicación
		\begin{align*}
			\mathbb{Z} &\longrightarrow \operatorname{Homeo}(X) \\[4pt]
			n &\longmapsto f^{\circ n}
		\end{align*}
		define un sistema dinámico discreto con tiempo pasado y futuro. 
		
		\medskip
		\noindent Por otro lado, definamos
		\[
		C^{0}(X) := \{f : X \to X \mid f \text{ es continua} \},
		\]
		entonces la aplicación
		\begin{align*}
			\mathbb{N} &\longrightarrow C^{0}(X) \\[4pt]
			n &\longmapsto f^{\circ n}
		\end{align*}
		define un sistema dinámico discreto con tiempo solo hacia el futuro.
		
		\medskip
		\noindent
		En particular, dado \( z_0 \in \mathbb{C} \), para 
		\begin{align*}
			f : \mathbb{C} &\longrightarrow \mathbb{C} \\[4pt]
			z &\longmapsto z^2 + 1,
		\end{align*}
		tenemos que
		\[
		\operatorname{Orb}(z_0) = \{z_0,\ f(z_0),\ f^{\circ 2}(z_0),\ f^{\circ 3}(z_0),\ f^{\circ 4}(z_0),\ \dots \}.
		\]
		Por ejemplo:
		\begin{equation*}
			\begin{aligned}
				\operatorname{Orb}(i) &= \{ i,\ 0,\ 1,\ 2,\ 5,\ \dots \}, \\[5pt]
				\operatorname{Orb}(i+1) &= \{ i+1,\ 1+2i,\ -2+4i,\ -11-16i,\ \dots \}.
			\end{aligned}
		\end{equation*}
		
	}
\end{ejemplo}

\begin{definicion}
	\upshape{
		Si \( \varphi : G \times X \to X \) es una acción, entonces el \textit{estabilizador} de \( x \in X \) es el conjunto
		\[
		\operatorname{Stab}(x) := \{ a \in G : \varphi(a, x) = x \}.
		\]
	}
\end{definicion}

\begin{proposicion}
	\upshape{
		Para cualquier \( x \in X \), el conjunto \( \operatorname{Stab}(x) \) es un subgrupo de \( G \).
	}
\end{proposicion}

\begin{proof}
	Dado $x\in X$. Sean $a,b\in \operatorname{Stab}(x)$ entonces $\varphi(a,x)=x$ y $\varphi(b,x)=x$. Entonces $$\varphi(ab^{-1},x)=\varphi\left(a,\varphi(b^{-1},x)\right)=\varphi\left(a,\varphi(b^{-1},\varphi(b,x))\right)=\varphi\left(a,\varphi(e,x)\right)=\varphi(a,x)=x.$$Por tanto, \( ab^{-1} \in \operatorname{Stab}(x) \), lo que demuestra que \( \operatorname{Stab}(x) \) es un subgrupo de \( G \).
\end{proof}




	\chapter{Producto directo, y Producto semidirecto }
	

\section{Producto directo}

Sean \( G \) y \( H \) dos grupos. En \( G \times H \) se define la siguiente operación:


\[
(g,h)(g',h') := (gg', hh'), \quad \text{para todo } g, g' \in G; \, h, h' \in H.
\]


Esta operación define una operación interna en \( G \times H \). Además:

\begin{itemize}
	\item Es asociativa: Para todo \( g, g', g'' \in G \) y \( h, h', h'' \in H \), se cumple:
	
	
	\[
	(g,h)((g',h')(g'',h'')) = ((g,h)(g',h'))(g'',h'').
	\]
	
	
	
	\item \( (e_G, e_H) \) es la identidad en \( G \times H \), pues para todo \( (g,h) \in G \times H \), se cumple:
	
	
	\[
	(e_G, e_H) (g,h) = (g,h) \quad \text{y} \quad (g,h) (e_G, e_H) = (g,h).
	\]
	
	
	
	\item Todo elemento \( (g,h) \in G \times H \) admite un inverso. Existe \( (g^{-1}, h^{-1}) \in G \times H \) tal que:
	\[
	(g,h)(g^{-1}, h^{-1}) = (e_G, e_H).
	\]
	Por lo que el inverso de \( (g,h) \) está dado por:
	\[
	(g,h)^{-1} := (g^{-1}, h^{-1}).
	\]	
\end{itemize}
Por lo tanto, \( G \times H \) tiene una estructura de grupo.

\begin{definicion}
	\upshape{
		Dado \( G \) y \( H \) dos grupos, el producto directo de \( G \) y \( H \) es el grupo \( G \times H \), con la operación definida por:
		
		
		\[
		(g,h)(g',h') := (gg', hh').
		\]
		
		
	}
\end{definicion}


\begin{ejemplo}
	\upshape{
		Sean los grupos \( (\mathbb{Z}/3\mathbb{Z}, +) \) y \( (\langle i \rangle, \cdot) \). Entonces el producto directo es
		\[
		\begin{array}{rllll}
			\mathbb{Z}/3\mathbb{Z} \times \langle i \rangle = \Bigl\{ & (\overline{0}, 1), & (\overline{1}, 1), & (\overline{2}, 1),&(\overline{0}, i),\\[4pt]
			& (\overline{1}, i), & (\overline{2}, i), &(\overline{0}, -1), & (\overline{1}, -1),\\[4pt]
			& (\overline{2}, -1), &(\overline{0}, -i), & (\overline{1}, -i), & (\overline{2}, -i)\:\: \Bigr\}
		\end{array}
		\]	
		Además, este grupo es cíclico, ya que puede ser generado por el elemento \( (\,\overline{1}, i) \):
		
		\[
		\mathbb{Z}/3\mathbb{Z} \times \langle i \rangle = \left\langle (\overline{1}, i) \right\rangle.
		\]
		Pero en general, el producto directo de dos grupos cíclicos no es cíclico.
	}
\end{ejemplo}

\begin{observacion}
	\upshape{
		Solo cuando \( G \) y \( H \) son grupos abelianos aditivos se suele escribir \( G \oplus H := G \times H \) y lo denominamos \textit{suma directa} de \( G \) y \( H \).
	}
\end{observacion}







\section{Producto semidirecto}

Sea \( G \) un grupo, entonces  
\[
\operatorname{Aut}(G) := \bigl\{\,f: G \to G \mid f \text{ es un isomorfismo de grupos}\bigr\},
\]  
con la operación de composición, es un grupo llamado \emph{grupo de automorfismos} de \(G\).

\medskip

Sea \(\varphi: H \to \operatorname{Aut}(G)\). Definimos una operación interna en \(G \times H\) (visto como conjunto, no como producto directo) por
\[
(g,h)\,\rtimes_{\varphi}\,(g',h')
: = \bigl(g\,\varphi_h(g'),\,h h'\bigr),
\quad
\text{para todo }g,g'\in G,\;h,h'\in H,
\]
donde \(\varphi_h := \varphi(h)\in \operatorname{Aut}(G)\). Además 














\begin{itemize}
	\item Es asociativa: Dado \((g,h),(g',h'),(g'',h'')\in G\times H\),  
	\begin{align*}
		\bigl((g,h)\!\rtimes_{\varphi}\!(g',h')\bigr)\!\rtimes_{\varphi}\!(g'',h'')
		&=
		\bigl(g\,\varphi_h(g'),\,h\,h'\bigr)\!\rtimes_{\varphi}\!(g'',h'')
		\\
		&=
		\bigl(\,g\,\varphi_h(g')\;\varphi_{h h'}(g''),\;(h\,h')\,h''\bigr),
		\\
		(g,h)\!\rtimes_{\varphi}\!\bigl((g',h')\!\rtimes_{\varphi}\!(g'',h'')\bigr)
		&=
		(g,h)\!\rtimes_{\varphi}\!\bigl(g'\,\varphi_{h'}(g''),\,h'\,h''\bigr)
		\\
		&=
		\bigl(\,g\;\varphi_h\bigl(g'\,\varphi_{h'}(g'')\bigr),\;h\,(h'\,h'')\bigr).
	\end{align*}
	Por la propiedad de que \(\varphi_h\) es un homomorfismo y \(\varphi_{h h'}=\varphi_h\circ\varphi_{h'}\), ambas expresiones coinciden, garantizando la asociatividad.
	
	\item Elemento identidad: Sea \((e_G,e_H)\). Para cualquier \((g,h)\),  
	\begin{align*}
		(e_G,e_H)\!\rtimes_{\varphi}\!(g,h)
		&=
		\bigl(e_G\,\varphi_{e_H}(g),\,e_H\,h\bigr)
		=
		(e_G\,g,\,h)
		=
		(g,h),
		\\
		(g,h)\!\rtimes_{\varphi}\!(e_G,e_H)
		&=
		\bigl(g\,\varphi_h(e_G),\,h\,e_H\bigr)
		=
		(g\,e_G,\,h)
		=
		(g,h),
	\end{align*}
	puesto que \(\varphi_{e_H}=\mathrm{id}_G\) y \(\varphi_h(e_G)=e_G\).
	
	\item Inverso: Dados \((g,h)\), definimos
	\[
	(g,h)^{-1}
	:=
	\bigl(\varphi_{h^{-1}}(g^{-1}),\,h^{-1}\bigr).
	\]
	Verifiquemos:
	\begin{align*}
		(g,h)\!\rtimes_{\varphi}\!(g,h)^{-1}
		&=
		\bigl(g\,\varphi_h\bigl(\varphi_{h^{-1}}(g^{-1})\bigr),\,h\,h^{-1}\bigr)
		=
		\bigl(g\,(g^{-1}),\,e_H\bigr)
		=
		(e_G,e_H),
		\\
		(g,h)^{-1}\!\rtimes_{\varphi}\!(g,h)
		&=
		\bigl(\varphi_{h^{-1}}(g^{-1})\,\varphi_{h^{-1}}(g),\,h^{-1}\,h\bigr)
		=
		\bigl(\varphi_{h^{-1}}(g^{-1}g),\,e_H\bigr)
		=
		(e_G,e_H),
	\end{align*}
	puesto que \(\varphi_{h^{-1}}\) es un isomorfismo de grupos.
\end{itemize}
Por lo tanto, \(\bigl(G\times H,\rtimes_{\varphi}\bigr)\) es un grupo.


\begin{definicion}
	\upshape{
		Dados \(G\) y \(H\) grupos y un homomorfismo \(\varphi\colon H\to\operatorname{Aut}(G)\), el \textit{producto semidirecto} de $G$ y $H$ via $\varphi$ es el grupo  \(\left(G\times H, \rtimes_{\varphi}\right)\). Se denota 
		\[
		G\rtimes_{\varphi} H.
		\]
	}
\end{definicion}


\begin{observacion}
	\upshape{
		A diferencia del producto directo, en el producto semidirecto uno de los grupos actúa sobre el otro. En efecto, si $G$ y $H$ son grupos y $\varphi: H \to \operatorname{Aut}(G)$ es un homomorfismo, entonces $H$ actúa  sobre $G$ por medio de automorfismos (inducen una acción), y esta interacción se refleja en la operación de $G \rtimes_\varphi H$. De este modo, el producto semidirecto incorpora no solo la estructura de $G$ y $H$, sino también cómo $H$ transforma a $G$.
	}
\end{observacion}



\begin{proposicion}\label{ñlkjhgfhjkwwwwww}
	\upshape{
		Sean \( G \) un grupo, \( N \trianglelefteq G \) y \( H \leq G \) dos subgrupos. Supongamos que:
		\begin{enumerate}
			\item Para todo \( g \in G \) existen únicos \( n' \in N \) y \( h' \in H \) tales que \( g = n'h' \), \label{lkjhgf.kjhg}
			\item Existe un homomorfismo \( \varphi: H \to \operatorname{Aut}(N) \) tal que \( \varphi_h(n) = hnh^{-1} \) para todo \( n \in N \) y \( h \in H \), donde \( \varphi_h := \varphi(h) \in \operatorname{Aut}(N) \).
		\end{enumerate}
		Entonces 
		\[
		G \cong N \rtimes_{\varphi} H.
		\]
	}
\end{proposicion}

\medskip
\begin{proof}
	Definamos la aplicación
	\begin{align*}
			\theta: N \rtimes_{\varphi} H &\longrightarrow G \\
			 (n, h) &\longmapsto n h
	\end{align*}
	
	\smallskip
	\noindent Verificamos que \( \theta \) es un homomorfismo. Sean \( (n, h), (n', h') \in N \rtimes_{\varphi} H \). Entonces:
	\begin{align*}
		\theta\big((n,h)\rtimes_{\varphi}(n',h')\big) &= \theta\big(n \varphi_h(n'), hh'\big) \\
		&= n \varphi_h(n') \cdot hh' \\
		&= n \cdot h n' h^{-1} \cdot hh' \\
		&= n h n' h' \\
		&= (nh)(n'h') = \theta(n,h)\cdot \theta(n',h').
	\end{align*}
	Por lo tanto, \( \theta \) es un homomorfismo de grupos.
	
	\smallskip
	\noindent Ahora mostramos que \( \theta \) es biyectiva:	
	\begin{itemize}
		\item Inyectividad: Supongamos que \( \theta(n,h) = \theta(n',h') \), es decir, \( nh = n'h' \). Como por hipótesis la descomposición \( g = n''h'' \) es única para cada elemento de \( G \), se concluye que \( n = n' \) y \( h = h' \), por lo tanto \( (n,h) = (n',h') \).
		
		\item Sobreyectividad: Por hipótesis, todo \( g \in G \) se puede escribir como \( g = n h \) con \( n \in N \), \( h \in H \), lo cual implica que \( g = \theta(n,h) \) para algún \( (n,h) \in N \rtimes_{\varphi} H \). Así, \( \theta \) es sobreyectiva.
	\end{itemize}
	Concluimos que \( \theta \) es un isomorfismo de grupos, es decir,
	\[
	G \cong N \rtimes_{\varphi} H.
	\]
\end{proof}

\begin{observacion}\label{equiv-condicion-descomposicion}
	\upshape
	En la Proposici\'on \ref{ñlkjhgfhjkwwwwww}, la condici\'on \ref{lkjhgf.kjhg} (descomposici\'on \'unica) puede sustituirse de manera equivalente por:  
	\[  
	N \cap H = \{e\} \quad \text{y} \quad NH = G.  
	\]  
	Es decir, $N$ y $H$ tienen intersecci\'on trivial y generan $G$.  
\end{observacion}

\begin{definicion}\label{kljhgf}
	\upshape{
		Sean \( A \), \( B \) y \( C \) tres grupos. Una \emph{sucesión exacta corta de grupos} es una secuencia de morfismos de grupos
		\[
		1 \longrightarrow A \xrightarrow{\:\:f\:} B \xrightarrow{\:\:g\:} C \longrightarrow 1
		\]
		tales que:
		\begin{enumerate}
			\item \( f \) es inyectivo,
			\item \( g \) es suprayectivo,
			\item \( \operatorname{Im}(f) = \ker(g) \).
		\end{enumerate}
		Aquí, \( 1 \) denota a cualquier grupo trivial.
	}
\end{definicion}

\begin{observacion}
	\upshape{
		De la Definición \ref{kljhgf},
		\begin{enumerate}
			\item  \( A \) se identifica con un subgrupo normal de \( B \), es decir, $A\cong f(A)\trianglelefteq B$.
			\item  \( C \) es isomorfo al cociente \( B / f(A) \).
		\end{enumerate}
	}
\end{observacion}
\begin{proof}[Justificación]
	\textcolor{white}{text}
	
	\noindent $[1]$ Por la Proposición \ref{oikhugfd} tenemos que $\ker g \trianglelefteq B$. Entonces $$f(A)=\operatorname{Im}f=\ker g \trianglelefteq B.$$
	\noindent $[2]$ Aplicando el Teorema Fundamnetal de Homomorfismo de grupos para $g: B\to C$ tenemos que $$C\cong \dfrac{B}{\ker g}=\dfrac{B}{f(A)}.$$
\end{proof}






\begin{proposicion}
	\upshape{
		Sean \( G \) un grupo, \( N \leq G \) y \( H \leq G \) dos subgrupos. Son equivalentes:
		\begin{enumerate}
			\item \( G \cong N \rtimes_{\varphi} H \) para algún homomorfismo \(\varphi: H \to \operatorname{Aut}(N)\).
			\item Existe una sucesión exacta corta escindida:
			\[
			1 \to N \xrightarrow{\:\:\beta\:} G \xrightarrow{\:\:\alpha\:} H \to 1,
			\]
			con \(\alpha \circ \gamma = \text{id}_H\) para algún \(\gamma: H \to G\).
		\end{enumerate}
	}
\end{proposicion}







\begin{proof}
	\textcolor{white}{text} % Espacio invisible para alineación
	
	\medskip
	\noindent
	$[1 \Rightarrow 2]$ Consideremos 
	$$\begin{array}{rl}
		\beta': N&\longrightarrow N\rtimes_{\varphi}H\\
		n&\longmapsto (n,e_H)
	\end{array}\quad\text{y}\quad\begin{array}{rl}
		\alpha': N\rtimes_{\varphi}H&\longrightarrow H\\
		(n,h)&\longmapsto h
	\end{array}$$
	Entonces
	\begin{itemize}
		\item $\beta'$ es inyectiva: pues $\ker \beta'=\{e_N\}$.
		\item $\alpha'$ es sobreyectiva: dado $h\in H$ existe $(e_N,h) \in N\rtimes_{\varphi}H$ tal que $\alpha'(e_N,h)=h$.
		\item $\ker \alpha'=\{(n,e_H):n\in N\}=\operatorname{Im}\beta'$
	\end{itemize}
	Sea ahora $\rho: G \to N \rtimes_\varphi H$ un isomorfismo (dado por la hipótesis). Entonces:
	\[
	\beta := \rho^{-1} \circ \beta', \quad \alpha := \alpha' \circ \rho,
	\]
	Verifiquemos que $\ker \alpha = \operatorname{Im} \beta$:
	\begin{align*}
		\ker \alpha &= \{ g \in G \mid \alpha(g) = e_H \} \\
		&= \{ g \in G \mid \alpha'(\rho(g)) = e_H \} \\
		&= \{ g \in G \mid \rho(g) \in \ker \alpha' \} \\
		&= \{ g \in G \mid \rho(g) \in \operatorname{Im} \beta' \} \\
		&= \{ g \in G \mid \exists n \in N \text{ tal que } \rho(g) = \beta'(n) = (n, e_H) \} \\
		&= \{ g \in G \mid \exists n \in N \text{ con } g = \rho^{-1}(n, e_H) = \beta(n) \} \\
		&= \operatorname{Im} \beta.
	\end{align*}
	Así, obtenemos la sucesión exacta corta:
	\[
	1 \to N \xrightarrow{\:\:\beta\:} G \xrightarrow{\:\:\alpha\:} H \to 1,
	\]
	Definimos además:
	\[
	\gamma': H \to N \rtimes_\varphi H,\quad h \mapsto (e_N, h), \quad \text{y} \quad \gamma := \rho^{-1} \circ \gamma'.
	\]
	Entonces $\alpha \circ \gamma = \operatorname{id}_H$, como se requería.\\
	
	\noindent
	$[2 \Rightarrow 1]$ Sea $g\in G$, entonces $\alpha(g)\in H$. Así, $\alpha(g)=h$ para algún $h\in H$. Consideremos $g\gamma(h^{-1})\in G$, entonces
	$$\alpha\left(g\gamma(h^{-1})\right)=\alpha(g)\alpha\left(\gamma(h^{-1})\right)=h\operatorname{id}_H(h^{-1})=hh^{-1}=e_H.$$ 
	Por lo que $g\gamma(h^{-1})\in \ker \alpha=\operatorname{Im}\beta$, entonces $g\gamma(h^{-1})=\beta(n)$ para algún $n\in N$. Así, $$g=\beta(n)\gamma(h).$$ 
	Supongamos que existe $n'\in N$ y $h'\in H$ tales que $$g=\beta(n)\gamma(h)=\beta(n')\gamma(h'),$$ luego $\gamma(hh'^{-1})=\beta(n'n^{-1})$, y como $\ker \alpha =\operatorname{Im}\beta$ tenemos que $$hh'^{-1}=\alpha\left(\gamma\left(hh'^{-1}\right)\right)=\alpha\left(\beta(n'n^{-1})\right)=e_H,$$ por lo que $$h=h'.$$ 
	Esto implica que $\beta(n)=\beta(n')$, y dado que $\beta$ es inyectiva, entonces $n=n'$.
	
	\medskip
	\noindent Entonces tenemos que cada $g\in G$ se escribe como 
	\begin{equation}\label{poiuhygf}
		g=\beta(n)\gamma(h),\quad \text{con }n\in N\text{ único y }h\in H\text{ único.}
	\end{equation}
	
	\medskip
	\noindent
	Definamos
	\begin{align*}
		\varphi: H &\longrightarrow \operatorname{Aut}(N) \\[4pt]
		h &\longmapsto \left(  
		\begin{array}{rcl}
			\varphi_h : N &\to& N \\[2pt]
			n &\mapsto& \varphi_h(n)
		\end{array} \right),
	\end{align*}
	donde $\varphi_h(n):=\beta^{-1}\left(\gamma(h)\beta(n)\gamma(h^{-1})\right)$ y el homomorfismo inverso $\beta^{-1}: \beta(N)\to N$.
	
	\medskip
	\noindent Esta aplicación está bien definida pues $\operatorname{Im}\beta=\ker \alpha \trianglelefteq G$, con lo que $$\gamma(h)\beta(n)\gamma(h^{-1})\in \operatorname{Im}\beta.$$ 
	Y para cada $h\in H$ se tiene que $\varphi_h: N\to N$ es un isomorfismo. En efecto,
	
	\begin{itemize}
		\item Homomorfismo:
		\begin{align*}
			\varphi_h(nn')    & =\beta^{-1}\left(\gamma(h)\beta(nn')\gamma(h^{-1})\right)\\
			& =\beta^{-1}\left(\gamma(h)\beta(n)\gamma(h^{-1})\gamma(h)\beta(n')\gamma(h^{-1})\right) \\
			&= \beta^{-1}\left(\gamma(h)\beta(n)\gamma(h^{-1})\right) \beta^{-1}\left(\gamma(h)\beta(n')\gamma(h^{-1})\right)\\
			&=\varphi_h(n)\varphi_h(n')
		\end{align*}
		
		\item Biyectiva: Considere $\varphi_{h^{-1}}$. Entonces 
		\begin{align*}
			\varphi_{h^{-1}}\left(\varphi_h(n)\right)&=\varphi_{h^{-1}} \left(\beta^{-1}\left(\gamma(h)\beta(n)\gamma(h^{-1})\right)\right)\\[4pt]
			&=\beta^{-1}\left(\gamma(h^{-1})\beta\left(\beta^{-1}\left(\gamma(h)\beta(n)\gamma(h^{-1})\right)\right)\gamma(h)\right)\\[4pt]
			&=\beta^{-1}\left(\gamma(h^{-1}h)\beta(n)\gamma(h^{-1}h)\right)\\[4pt]
			&=\beta^{-1}(\beta(n))=n\\[4pt]
			&=\operatorname{id}_N(n).
		\end{align*}
		De manera análoga tenemos que $\varphi_h\left(\varphi_{h^{-1}}(n)\right)=\operatorname{id}_N(n)$ para todo $n\in N$. Con lo cual $\varphi_h$ es biyectiva.
	\end{itemize}
	Así, para todo $h\in H$ tenemos que $\varphi_h \in\operatorname{Aut}(N)$.
	
	\medskip
	\noindent
	Ahora debemos ver que $\varphi: H\to \operatorname{Aut}(N)$ es un homomorfismo de grupos. En efecto,
	\begin{align*}
		\varphi_{h_1 h_2}(n) 
		&= \beta^{-1}\left( \gamma(h_1 h_2)\, \beta(n)\, \gamma\left((h_1 h_2)^{-1}\right) \right) \\[4pt]
		&= \beta^{-1}\left( \gamma(h_1)\gamma(h_2)\, \beta(n)\, \gamma\left(h_2^{-1}\right)\gamma\left(h_1^{-1}\right) \right) \\[4pt]
		&= \beta^{-1}\left( \gamma(h_1)\, \left( \gamma(h_2)\, \beta(n)\, \gamma\left(h_2^{-1}\right) \right)\, \gamma\left(h_1^{-1}\right) \right) \\[4pt]
		&= \beta^{-1}\left( \gamma(h_1)\, \beta\left( \varphi_{h_2}(n) \right)\, \gamma\left(h_1^{-1}\right) \right) \\[4pt]
		&= \varphi_{h_1}\left( \varphi_{h_2}(n) \right) \\[4pt]
		&= \left( \varphi_{h_1} \circ \varphi_{h_2} \right)(n).
	\end{align*}
	
	\noindent Ahora estamos en condiciones de finalmente mostrar que $$G\cong N\rtimes_\varphi H.$$
	
	\noindent
	Definamos la aplicación
	\begin{align*}
		\theta: N \rtimes_{\varphi} H &\longrightarrow G \\
		(n, h) &\longmapsto \beta(n)\gamma(h)
	\end{align*}
	
	\smallskip
	\noindent
	Verificamos que $\theta$ es un homomorfismo. Sean $(n, h), (n', h') \in N \rtimes_{\varphi} H$. Entonces:
	\begin{align*}
		\theta\left( (n,h)\rtimes_\varphi (n',h') \right) 
		&= \theta\left( \left( n \cdot \varphi_h(n'),\, hh' \right) \right) \\[4pt]
		&= \beta\left( n \cdot \varphi_h(n') \right)\gamma(hh') \\[4pt]
		&= \beta(n)\beta\left( \varphi_h(n') \right)\gamma(h)\gamma(h') \\[4pt]
		&= \beta(n)\gamma(h)\beta(n')\gamma(h^{-1})\gamma(h)\gamma(h') \\[4pt]
		&= \beta(n)\gamma(h)\beta(n')\gamma(h') \\[4pt]
		&= \theta(n,h)\cdot\theta(n',h').
	\end{align*}
	
	\noindent
	Por lo tanto, $\theta$ es un homomorfismo de grupos.\\
	
	\noindent
	Ahora mostramos que $\theta$ es biyectiva:    
	\begin{itemize}
		\item Inyectividad: Supongamos que $\theta(n,h) = \theta(n',h')$, es decir,
		\[
		\beta(n)\gamma(h) = \beta(n')\gamma(h').
		\]
		Por la unicidad de la descomposición en \eqref{poiuhygf} se tiene que $n = n'$ y $h = h'$, por lo tanto $(n,h) = (n',h')$.
		
		\item Sobreyectividad: Por \eqref{poiuhygf}, todo $g \in G$ se puede escribir de manera única como $g = \beta(n)\gamma(h)$ con $n \in N$, $h \in H$. Esto implica que $g = \theta(n,h)$, y por tanto $\theta$ es sobreyectiva.
	\end{itemize}
	
	\noindent
	Concluimos que $\theta$ es un isomorfismo de grupos, es decir,
	\[
	G \cong N \rtimes_{\varphi} H.
	\]
\end{proof}


\begin{proposicion}\label{ñlkjhgfhjk}
	\upshape
	Sea \( n \geq 3 \). Consideremos el producto semidirecto \( \mathbb{Z}/n\mathbb{Z} \rtimes_{\varphi} \mathbb{Z}/2\mathbb{Z} \), para algún homomorfismo \( \varphi: \mathbb{Z}/2\mathbb{Z} \to \operatorname{Aut}(\mathbb{Z}/n\mathbb{Z}) \), y consideremos al grupo diedral $D_n$. Entonces,
	\[
	D_n \cong \mathbb{Z}/n\mathbb{Z} \rtimes_{\varphi} \mathbb{Z}/2\mathbb{Z}.
	\]
\end{proposicion}



\begin{proof}
	Recordemos de la Definición \ref{deiugv} y la Proposición \ref{lkjhgjkl} que  \(D_n=\langle a,b\rangle\), donde \(a\) es un \(n\)-ciclo tal que \(a^n=I\), y \(b\) es una permutación tal que \(b^2=I\). Consideremos \(N:=\langle a\rangle\), un grupo cíclico de orden \(n\), y \(H:=\langle b\rangle\), un grupo cíclico de orden \(2\). Entonces \(N\trianglelefteq D_n\), \(N\cap H = \{I\}\), y \(NH = D_n\).
	
	\medskip
	\noindent Definimos la acción
	\[
	\rho: \langle b \rangle \longrightarrow \operatorname{Aut}(\langle a\rangle), \quad h \longmapsto \left(  
	\begin{array}{rcl}
		\rho_h : \langle a\rangle &\to& \langle a\rangle \\
		n &\mapsto& hnh^{-1}
	\end{array} \right).
	\]
	Esta aplicación está bien definida porque \(N \trianglelefteq D_n\). Así, la conjugación por elementos de \(H\) induce automorfismos de \(N\). Además, \(\rho\) es un homomorfismo de grupos.
	
	\medskip
	\noindent
	Entonces, por la Proposición \ref{ñlkjhgfhjkwwwwww} y la Observación \ref{equiv-condicion-descomposicion}, tenemos que
	\[
	D_n \cong \langle a\rangle\rtimes_{\rho} \langle b \rangle.
	\]
	Por otro lado, sabemos por la Proposición \ref{lkjhghjklñ} que
	\[
	\langle a\rangle \cong \mathbb{Z}/n\mathbb{Z} \quad \text{y} \quad \langle b\rangle \cong \mathbb{Z}/2\mathbb{Z}.
	\]
	Definimos ahora 	
	\begin{align*}
		\varphi: \dfrac{\mathbb{Z}}{2\mathbb{Z}} &\longrightarrow  \operatorname{Aut}\left(\dfrac{\mathbb{Z}}{n\mathbb{Z}}\right)\\[5pt]
		\overline{O}& \longmapsto \varphi_{\overline{0}}: =\operatorname{id}_{\mathbb{Z}/n\mathbb{Z}}\\[5pt]
		\overline{1}& \longmapsto  \left(  
		\begin{array}{rcl}
			\varphi_{\overline{1}} : \mathbb{Z}/n\mathbb{Z} &\to&\mathbb{Z}/n\mathbb{Z} \\[2pt]
			[r] &\mapsto& -[r]
		\end{array} \right)
	\end{align*}
	Así para cada $t \in \mathbb{Z}/2\mathbb{Z}$ tenemos que $\varphi_t\in \operatorname{Aut}\left(\mathbb{Z}/n\mathbb{Z}\right)$ y que $\varphi$ es un homomorfismo de grupos ya que $$\varphi(\overline{0}+\overline{1})=\varphi_{\overline{1}}=\varphi_{\overline{0}}\circ\varphi_{\overline{1}}=\varphi\left(\overline{0}\right)\circ \varphi\left(\overline{1}\right).$$
	
	$$\varphi(\overline{1}+\overline{0})=\varphi_{\overline{1}}=\varphi_{\overline{1}}\circ \varphi_{\overline{0}}=\varphi\left(\overline{1}\right)\circ \varphi\left(\overline{0}\right).$$
	
	
	\medskip
	\noindent
	Entonces formamos el producto semidirecto
	\[
	\mathbb{Z}/n\mathbb{Z} \rtimes_{\varphi} \mathbb{Z}/2\mathbb{Z}.
	\]
	Sean ahora los isomorfismos \(f: \langle a \rangle \to \mathbb{Z}/n\mathbb{Z}\), y \(g: \langle b \rangle \to \mathbb{Z}/2\mathbb{Z}\). Definimos
	
	\begin{align*}
		\theta: \langle a\rangle\rtimes_{\rho} \langle b \rangle&\longrightarrow \dfrac{\mathbb{Z}}{n\mathbb{Z}} \rtimes_{\varphi} \dfrac{\mathbb{Z}}{2\mathbb{Z}} \\[5pt]
		\quad (x, y) &\longmapsto \big(f(x), g(y)\big).
	\end{align*}
	Verificamos que \(\theta\) es un homomorfismo. Tomemos \((x,y), (x', y')\) en \(\langle a\rangle \rtimes_{\rho} \langle b\rangle\). Entonces:
	\begin{align*}
		\theta((x,y)\rtimes_{\rho}(x',y')) 
		&= \theta\left(x  \rho_{y}(x'),\; yy'\right) \\
		&= \left(f(x  \rho_{y}(x')),\; g(yy')\right) \\
		&= \left(f(x) f(\rho_{y}(x')),\; g(y) g(y')\right).
	\end{align*}
	Como \(\rho\) y \(\varphi\) están relacionadas por \(f\) y \(g\) vía:
	\[
	f(\rho_h(z)) = \varphi_{g(h)}(f(z)),\quad h\in \langle b\rangle; \: z\in \langle a\rangle.
	\]
En efecto:
	\begin{itemize}
		\item Si \(h = I\), entonces \(g(h) = \overline{0}\), y como \(\rho_h = \operatorname{id}\), se tiene:
		\[
		f(\rho_h(z)) = f(z) = \varphi_{\overline{0}}(f(z)).
		\]
		\item Si \(h = b\), entonces \(g(h) = \overline{1}\), y \(\rho_b(z) = bzb^{-1} = z^{-1}\), dado que  \(ba = a^{-1}b\). Así,
		\[
		f(\rho_b(z)) = f(z^{-1}) = -f(z) = \varphi_{\overline{1}}(f(z)).
		\]
		
		\noindent
		\textbf{Nota.} Aunque \( \rho_h(z) = hzh^{-1} \), no podemos separar \( f(hzh^{-1}) \) como \( f(h) + f(z) + f(h^{-1}) \), ya que \( f \) sólo está definido en \( \langle a \rangle \), y \( h \notin \langle a \rangle \).
	\end{itemize}
	Entonces,
	\begin{align*}
		\theta((x,y)\rtimes_{\rho}(x',y')) 
		&= \left(f(x) f(\rho_{y}(x')),\; g(y) g(y')\right)\\
		&=\left(f(x) \varphi_{g(y)}(f(x')),\; g(y) g(y')\right) \\
		&= \left(f(x), g_{g(y)} \left(f(x')\right),\:(f(x'), g(y')\right) \\
		&=\left(f(x),g(y)\right)\rtimes_{\varphi}\left(f(x'),g(x')\right)\\
		&= \theta(x, y) \rtimes_{\varphi} \theta(x', y').
	\end{align*}
	Así, \(\theta\) es un homomorfismo de grupos. Como \(f\) y \(g\) son isomorfismos, \(\theta\) también lo es. En efecto
	\begin{itemize}
		\item Inyectividad: Supongamos que \( \theta(z_1, h_1) = \theta(z_2, h_2) \). Entonces,
		\[
		(f(z_1), g(h_1)) = (f(z_2), g(h_2)).
		\]
		Como \( f \) y \( g \) son inyectivos, se concluye que \( z_1 = z_2 \) y \( h_1 = h_2 \), por lo que \( \theta \) es inyectiva.
		
		\item Sobreyectividad: Sea cualquier par \( ([r], [t]) \in \mathbb{Z}/n\mathbb{Z} \rtimes_\varphi \mathbb{Z}/2\mathbb{Z} \). Como \( f \) y \( g \) son isomorfismos, existen \( z \in \langle a \rangle \) y \( h \in \langle b \rangle \) tales que \( f(z) = [r] \) y \( g(h) = [t] \). Entonces \( \theta(z, h) = ([r], [t]) \), y por tanto \( \theta \) es sobreyectiva.
	\end{itemize}
	Como consecuencia tenemos que 
	\[
	\langle a\rangle \rtimes_{\rho} \langle b \rangle \cong \mathbb{Z}/n\mathbb{Z} \rtimes_{\varphi} \mathbb{Z}/2\mathbb{Z}.
	\]
	Finalmente, como ya se había establecido que \(D_n \cong \langle a\rangle \rtimes_{\rho} \langle b \rangle\), se concluye:
	\[
	D_n \cong \mathbb{Z}/n\mathbb{Z} \rtimes_{\varphi} \mathbb{Z}/2\mathbb{Z}. \qedhere
	\]
\end{proof}




\medskip

Así, como puede notarse, la demostración de la Proposición \ref{ñlkjhgfhjk} ha resultado extensa. A continuación presentaremos una demostración más directa. No obstante, esta primera demostración cumple un propósito importante: nos ha permitido observar en detalle cómo se relacionan dos productos semidirectos y bajo qué condiciones pueden ser isomorfos. Esto será analizado en la Proposición \ref{prop:isomorfismo_semi}.\\




\noindent \textbf{Proposicion}
	\upshape
	Sea \( n \geq 3 \). Consideremos el producto semidirecto \( \mathbb{Z}/n\mathbb{Z} \rtimes_{\varphi} \mathbb{Z}/2\mathbb{Z} \), para algún homomorfismo \( \varphi: \mathbb{Z}/2\mathbb{Z} \to \operatorname{Aut}(\mathbb{Z}/n\mathbb{Z}) \), y consideremos al grupo diedral $D_n$. Entonces,
	\[
	D_n \cong \mathbb{Z}/n\mathbb{Z} \rtimes_{\varphi} \mathbb{Z}/2\mathbb{Z}.
	\]


\begin{proof}
	Sea \(G := \mathbb{Z}/n\mathbb{Z} \rtimes_\varphi \mathbb{Z}/2\mathbb{Z}\), donde la acción \(\varphi: \mathbb{Z}/2\mathbb{Z} \to \operatorname{Aut}(\mathbb{Z}/n\mathbb{Z})\) está dada por
	\[
	\varphi_{\overline{0}} = \operatorname{id}, \qquad \varphi_{\overline{1}}([r]) = [-r].
	\]
	Denotemos
	\[
	p := ([1], \overline{0}), \qquad q := ([0], \overline{1}).
	\]
	Veamos que estos elementos satisfacen las relaciones del grupo diedral \(D_n\):
	
	\begin{itemize}
		\item \(p^n = ([1], \overline{0})^n = ([n], \overline{0}) = ([0], \overline{0})\).
		
		\item \(q^2 = ([0], \overline{1})^2 = ([0] + \varphi_{\overline{1}}([0]), \overline{1} + \overline{1}) = ([0], \overline{0})\).
		
		\item \(q p q = ([0], \overline{1})([1], \overline{0})([0], \overline{1})\). Primero,
		\[
		([0], \overline{1})([1], \overline{0}) = ([0] + \varphi_{\overline{1}}([1]), \overline{1}) = ([-1], \overline{1}),
		\]
		y luego
		\[
		([-1], \overline{1})([0], \overline{1}) = ([-1] + \varphi_{\overline{1}}([0]), \overline{0}) = ([-1], \overline{0}) = p^{-1}.
		\]
		Así, \(q p q = p^{-1}\).
	\end{itemize}
	Además, \(G = \langle p, q \rangle\). En efecto, notemos que todo elemento de \(G = \mathbb{Z}/n\mathbb{Z} \rtimes_\varphi \mathbb{Z}/2\mathbb{Z}\) es un par \(([\ell], \overline{t})\), donde \([\ell] \in \mathbb{Z}/n\mathbb{Z}\) y \(\overline{t} \in \mathbb{Z}/2\mathbb{Z}\). Entonces, hay dos casos:
	
	\begin{itemize}
		\item Si \(\overline{t} = \overline{0}\), entonces \(([\ell], \overline{0}) = ([1], \overline{0})^\ell = p^\ell\).
		
		\item Si \(\overline{t} = \overline{1}\), entonces
		\[
		([\ell], \overline{1}) = ([\ell], \overline{0})([0], \overline{1}) = ([1], \overline{0})^\ell \cdot q = p^\ell q.
		\]
	\end{itemize}
	Por lo tanto, cualquier elemento de \(G\) puede expresarse como un producto de potencias de \(p\) y \(q\), lo cual prueba que \(G = \langle p, q \rangle\).
	
	
	\medskip
	\noindent Así,  tanto $D_n$ como $G$ tienen la misma presentación de grupos $$\langle x,y: x^n=e,y^2=2,yxy=x^{-1}\rangle$$ de donde 
	\[
	D_n \cong \langle x,y: x^n=e,y^2=2,yxy=x^{-1}\rangle \cong\mathbb{Z}/n\mathbb{Z} \rtimes_\varphi \mathbb{Z}/2\mathbb{Z}. 
	\]
\end{proof}


%\begin{proposicion}\label{dasdasdasdasad}\upshape
%	Sean \( A \), \( B \), \( A' \) y \( B' \) grupos, y sean \( \varphi \colon B \to \operatorname{Aut}(A) \) y \( \varphi' \colon B' \to \operatorname{Aut}(A') \) homomorfismos de grupos. Entonces, son equivalentes las siguientes condiciones:
%	\begin{enumerate}
%		\item Existe un isomorfismo de grupos
%		\[
%		A \rtimes_{\varphi} B \cong A' \rtimes_{\varphi'} B'.
%		\]
		
%		\item Existen isomorfismos \( f \colon A \to A' \) y \( g \colon B \to B' \) tales que, para todo \( b \in B \), se cumple:
%		\[
%		f \circ \varphi_{b} = \varphi'_{g(b)} \circ f.
%		\]
%	\end{enumerate}
%\end{proposicion}


\begin{proposicion}\label{prop:isomorfismo_semi}
	\upshape
	Sean \( A \), \( B \), \( A' \) y \( B' \) grupos, y sean $$ \varphi \colon B \to \operatorname{Aut}(A) \qquad \text{y}\qquad\varphi' \colon B' \to \operatorname{Aut}(A') $$ homomorfismos de grupos. Supongamos que existen isomorfismos de grupos \( f \colon A \to A' \) y \( g \colon B \to B' \) tales que, para todo \( b \in B \), se cumple:
	\[
	f \circ \varphi_b = \varphi'_{g(b)} \circ f.
	\]
	Entonces, el grupo producto semidirecto \( A \rtimes_\varphi B \) es isomorfo a \( A' \rtimes_{\varphi'} B' \).
\end{proposicion}

\begin{proof}
	\textcolor{white}{text}
	
	\medskip
	\noindent 
	Definamos
	\[
	\theta: A \rtimes_{\varphi} B \longrightarrow A' \rtimes_{\varphi'} B', \qquad (a,b) \longmapsto \left(f(a), g(b)\right).
	\]
	Veamos que \( \theta \) es homomorfismo. Tomemos \( (a,b), (a',b') \in A \rtimes_{\varphi} B \). Entonces:
	\begin{align*}
		\theta\big((a,b) \rtimes_{\varphi} (a',b')\big) 
		&= \theta\big(a  \varphi_b(a'),\; bb'\big)\\
		&= \left(f\left(a  \varphi_b(a')\right),\; g(bb')\right)\\
		&= \left(f(a)  f(\varphi_b(a')),\; g(b)g(b')\right)\\
		&= \left(f(a)  \varphi'_{g(b)}(f(a')),\; g(b)g(b')\right)\\
		&= \left(f(a),g(b)\right) \rtimes_{\varphi'} \left(f(a'),g(b')\right)\\
		&= \theta(a,b) \rtimes_{\varphi'} \theta(a',b').
	\end{align*}
Como \(f\) y \(g\) son isomorfismos, \(\theta\) también lo es. En efecto,
\begin{itemize}
	\item Inyectividad: Supongamos que \( \theta(a_1, b_1) = \theta(a_2, b_2) \). Entonces,
	\[
	(f(a_1), g(b_1)) = (f(a_2), g(b_2)).
	\]
	Como \( f \) y \( g \) son inyectivos, se concluye que \( a_1 = a_2 \) y \( b_1 = b_2 \), por lo que \( \theta \) es inyectiva.
	
	\item Sobreyectividad: Sea cualquier par \( (a', b') \in A' \rtimes_{\psi} B' \). Como \( f \) y \( g \) son isomorfismos, existen \( a \in A \) y \( b \in B \) tales que \( f(a) = a' \) y \( g(b) = b' \). Entonces \( \theta(a, b) = (a', b') \), y por tanto \( \theta \) es sobreyectiva.
\end{itemize}
Concluimos que
\[
A \rtimes_{\varphi} B \cong A' \rtimes_{\psi} B'. \qedhere
\]

	
\end{proof}

\begin{proposicion}
	\upshape{
	Sea \( \mathbb{K} \) un cuerpo. Consideremos los siguientes subconjuntos de \( \operatorname{GL}_n(\mathbb{K}) \), el grupo de matrices invertibles \( n \times n \) con entradas en \( \mathbb{K} \), con la operación de multiplicación de matrices:
	\begin{align*}
		\mathbb{T}_n &:= \{ A \in \operatorname{GL}_n(\mathbb{K}) \mid A \text{ es triangular superior} \}, \\[5pt]
		\mathbb{D}_n &:= \{ A \in \operatorname{GL}_n(\mathbb{K}) \mid A \text{ es diagonal} \}, \\[5pt]
		\mathbb{U}_n &:= \{ A \in \mathbb{T}_n \mid \text{todas las entradas diagonales de } A \text{ son iguales a } 1 \}.
	\end{align*}
	Entonces:
	\begin{itemize}
		\item \( \mathbb{T}_n \), \( \mathbb{D}_n \) y \( \mathbb{U}_n \) son subgrupos de \( \operatorname{GL}_n(\mathbb{K}) \) bajo la multiplicación de matrices.
		\item Se tiene \( \mathbb{D}_n \leq \mathbb{T}_n \), \( \mathbb{U}_n \trianglelefteq \mathbb{T}_n \), y
		\[
		\mathbb{T}_n \cong \mathbb{U}_n \rtimes_{\varphi} \mathbb{D}_n,
		\]
		para algún homomorfismo \( \varphi \colon \mathbb{D}_n \to \operatorname{Aut}(\mathbb{U}_n) \).
	\end{itemize}
	}
\end{proposicion}







\begin{proof}
	\textcolor{white}{text}
	
	\medskip
	\noindent Nos interesa ver que  $\mathbb{T}_n \cong \mathbb{U}_n \rtimes_{\varphi} \mathbb{D}_n,$ por lo que probaremos inicialmente que $ \mathbb{U}_n \trianglelefteq \mathbb{T}_n $. En efecto, sean
\[
T =
\begin{pmatrix}
	d_1 & t_{12} & \cdots & t_{1n} \\
	0   & d_2    & \cdots & t_{2n} \\
	\vdots & \ddots & \ddots & \vdots \\
	0 & \cdots & 0 & d_n
\end{pmatrix} \in \mathbb{T}_n, \quad
U =
\begin{pmatrix}
	1 & u_{12} & \cdots & u_{1n} \\
	0 & 1 & \cdots & u_{2n} \\
	\vdots & \ddots & \ddots & \vdots \\
	0 & \cdots & 0 & 1
\end{pmatrix} \in \mathbb{U}_n.
\]
Entonces \( T^{-1} \) también es triangular superior e invertible, y tiene la forma:
\[
T^{-1} =
\begin{pmatrix}
	d_1^{-1} & s_{12} & \cdots & s_{1n} \\
	0   & d_2^{-1} & \cdots & s_{2n} \\
	\vdots & \ddots & \ddots & \vdots \\
	0 & \cdots & 0 & d_n^{-1}
\end{pmatrix},
\quad \text{con ciertos } s_{ij} \in \mathbb{K}.
\]
Así, la conjugación
\[
T U T^{-1} =
\begin{pmatrix}
	1 & * & \cdots & * \\
	0 & 1 & \cdots & * \\
	\vdots & \ddots & \ddots & \vdots \\
	0 & \cdots & 0 & 1
\end{pmatrix} \in \mathbb{U}_n,
\]
Así, tenemos que  \( \mathbb{U}_n \trianglelefteq \mathbb{T}_n \).

\medskip
\noindent Como siguienye paso definamos 
	\[
	\varphi \colon \mathbb{D}_n \to \operatorname{Aut}(\mathbb{U}_n), \quad \varphi_D(U) := D U D^{-1}.
	\]
	Veamos que esto está bien definido:
	
	\begin{itemize}
		\item Dado \( D \in \mathbb{D}_n \), se tiene que \( \varphi_D \) es un automorfismo de \( \mathbb{U}_n \):
		\begin{itemize}
			\item Homomorfismo: Para \( U_1, U_2 \in \mathbb{U}_n \),
			\[
			\varphi_D(U_1 U_2) = D U_1 U_2 D^{-1} = D U_1 D^{-1} D U_2 D^{-1} = \varphi_D(U_1) \varphi_D(U_2).
			\]
			\item Biyectiva: La inversa de \( \varphi_D \) es \( \varphi_{D^{-1}} \), pues
			\[
			\varphi_{D^{-1}}(\varphi_D(U)) = D^{-1} (D U D^{-1}) D = U=\operatorname{id}_{\mathbb{U}_n}(U).
			\]
			\[
			\varphi_{D^{-1}}(\varphi_D(U)) = D^{-1}(D U D^{-1})D = U = \operatorname{id}_{\mathbb{U}_n}(U).
			\]
			
			
			
			
		\end{itemize}
		\item \( \varphi \) es un homomorfismo: Para \( D_1, D_2 \in \mathbb{D}_n \) y \( U \in \mathbb{U}_n \),
		\[
		\varphi_{D_1 D_2}(U) = \left(D_1 D_2\right) U \left(D_1 D_2\right)^{-1} = D_1 \left(D_2 U D_2^{-1}\right) D_1^{-1} = \varphi_{D_1}\left(\varphi_{D_2}(U)\right).
		\]
		
	\end{itemize}
	Debemos ver $\mathbb{T}=\mathbb{U}_n\mathbb{D}_n$, es decir, toda matriz \( T \in \mathbb{T}_n \) puede escribirse como \( T = U D \), con \( U \in \mathbb{U}_n \), \( D \in \mathbb{D}_n \). En efecto, sea
	\[
	T = (t_{ij}) \in \mathbb{T}_n,
	\]
	definimos
	\[
	D := \operatorname{diag}(t_{11}, t_{22}, \dots, t_{nn}) \in \mathbb{D}_n,
	\]
	y luego como \[
	D^{-1} = \operatorname{diag}\left(t_{11}^{-1},\, t_{22}^{-1},\, \dots,\, t_{nn}^{-1}\right)\in \mathbb{D}_n
	\]
	 es tal que 	
	\[
	T D^{-1} = 
	\begin{pmatrix}
		t_{11} & t_{12} & \cdots & t_{1n} \\
		0 & t_{22} & \cdots & t_{2n} \\
		\vdots & \vdots & \ddots & \vdots \\
		0 & 0 & \cdots & t_{nn}
	\end{pmatrix}
	\begin{pmatrix}
		t_{11}^{-1} & 0 & \cdots & 0 \\
		0 & t_{22}^{-1} & \cdots & 0 \\
		\vdots & \vdots & \ddots & \vdots \\
		0 & 0 & \cdots & t_{nn}^{-1}
	\end{pmatrix}
	=
	\begin{pmatrix}
		1 & \frac{t_{12}}{t_{22}} & \cdots & \frac{t_{1n}}{t_{nn}} \\
		0 & 1 & \cdots & \frac{t_{2n}}{t_{nn}} \\
		\vdots & \vdots & \ddots & \vdots \\
		0 & 0 & \cdots & 1
	\end{pmatrix}
	\in \mathbb{U}_n
	\] definimos  \( U:=TD^{-1} \in \mathbb{U}_n \), y se tiene \( T = U D \) como queríamos.
	
	\medskip
	\noindent
	También debemos ver que \( \mathbb{U}_n \cap \mathbb{D}_n = \{I\} \). En efecto,
	si \( A \in \mathbb{U}_n \cap \mathbb{D}_n \), entonces \( A \) es diagonal e igual a la identidad en la diagonal, por lo que
	\[
	A = 
	\begin{pmatrix}
		1 & 0 & \cdots & 0 \\
		0 & 1 & \cdots & 0 \\
		\vdots & \vdots & \ddots & \vdots \\
		0 & 0 & \cdots & 1
	\end{pmatrix}
	= I.
	\]
	Así, \( \mathbb{U}_n \cap \mathbb{D}_n = \{I\} \).
	
	\medskip
	\noindent	
	Finalmente, por la ver Proposición~\ref{ñlkjhgfhjkwwwwww} y Observación~\ref{equiv-condicion-descomposicion} concluimos que
	\[
	\mathbb{T}_n \cong \mathbb{U}_n \rtimes_\varphi \mathbb{D}_n. \qedhere
	\]
\end{proof}



	
	\chapter{Anillos y Cuerpos}
	\section{Anillos}
\begin{definicion}
	\upshape{
		Un \textit{anillo} es un conjunto \( R \) provisto de dos operaciones internas
		\[
		+: R \times R \to R, \qquad \cdot: R \times R \to R,
		\]
		que satisfacen las siguientes condiciones:
		\begin{enumerate}
			\item \( (R, +) \) es un grupo abeliano.
			\item \( (R, \cdot) \) es un semigrupo.
			\item Propiedad distributiva: para todo \( a, b, c \in R \),
			\[
			a(b + c) = ab + ac \quad \text{y} \quad (a + b)c = ac + bc.
			\]
		\end{enumerate}
		De esta manera, denotamos por \( (R, +, \cdot) \) al anillo que se define.
	}
\end{definicion}
\begin{observacion}
	\upshape{
		\textcolor{white}{text}
		\medskip
		\begin{enumerate}
			\item Llamaremos \textit{elemento neutro} a la identidad de \( (R, +) \), y lo denotaremos por \( 0_R \).
			\item Si \( (R, \cdot) \) es un monoide, entonces diremos que \( (R, +, \cdot) \) es un \textit{anillo con unidad}, y denotaremos a dicha identidad por \( 1_R \), la cual también será llamada \textit{identidad del anillo}.
			\item Si \( (R, \cdot) \) es conmutativo, entonces \( (R, +, \cdot) \) se denomina un \textit{anillo conmutativo}.
		\end{enumerate}
	}
\end{observacion}


\begin{ejemplo}
	\upshape{
		Sea \( X \) un espacio topológico. Entonces
		\[
		C^0(X) := \{\, f : X \to \mathbb{R} \mid f \text{ es continua} \,\}
		\]
		con la suma definida punto a punto \((f + g)(x) := f(x) + g(x)\) y el producto definido  \((fg)(x) := f(x)g(x)\), es un anillo conmutativo con unidad.
	}
\end{ejemplo}



\begin{ejemplo}
	\upshape{
		Sea \( f \colon \mathbb{R}^n \to \mathbb{R} \) una función continua. Decimos que \( f \) tiene \textit{soporte compacto} si
		\[
		\operatorname{supp}(f) := \overline{\{ x \in \mathbb{R}^n \mid f(x) \neq 0 \}},
		\]
		es un subconjunto compacto de \( \mathbb{R}^n \).
		
		\medskip
		\noindent
		El conjunto 
		\[
		C_c(\mathbb{R}^n) := \{\, f \in C^0(\mathbb{R}^n) \mid \operatorname{supp}(f) \text{ es compacto} \,\}.
		\]
		con la suma y el producto definidos punto a punto es un anillo conmutativo, pero no posee unidad.
	}
\end{ejemplo}

\begin{proof}
	\upshape{
		Basta probar que si \( f, g \) tienen soporte compacto, entonces \( f + g \) y \( fg \) también lo tienen. En efecto:
		\begin{itemize}
			\item \textbf{\( f + g \):} Como \( \operatorname{supp}(f) \) y \( \operatorname{supp}(g) \) son compactos, su unión \( \operatorname{supp}(f) \cup \operatorname{supp}(g) \) también lo es. Y como
			\[
			\{ x \in \mathbb{R}^n \mid (f + g)(x) \ne 0 \} \subseteq \{ x \in \mathbb{R}^n \mid f(x) \ne 0 \} \cup \{ x \in \mathbb{R}^n \mid g(x) \ne 0 \},
			\]
			se concluye que
			\[
			\operatorname{supp}(f + g) = \overline{\{ x \in \mathbb{R}^n \mid (f + g)(x) \ne 0 \}} \subseteq \operatorname{supp}(f) \cup \operatorname{supp}(g),
			\]
			y por tanto \( \operatorname{supp}(f + g) \) es compacto.
			
			\item \textbf{\( fg \):} Análogamente,
			\[
			\{ x \in \mathbb{R}^n \mid (fg)(x) \ne 0 \} \subseteq \{ x \in \mathbb{R}^n \mid f(x) \ne 0 \} \cap \{ x \in \mathbb{R}^n \mid g(x) \ne 0 \},
			\]
			por lo que
			\[
			\operatorname{supp}(fg) = \overline{\{ x \in \mathbb{R}^n \mid (fg)(x) \ne 0 \}} \subseteq \operatorname{supp}(f) \cap \operatorname{supp}(g),
			\]
			que también es compacto, pues la intersección de dos compactos es cerrada y contenida en un compacto.
		\end{itemize}
		Finalmente, si \( C_c(\mathbb{R}^n) \) tuviera unidad, esta sería la función constante \( 1 \colon \mathbb{R}^n \to \mathbb{R} \); pero
		\[
		\operatorname{supp}(1) = \overline{\{ x \in \mathbb{R}^n \mid 1(x) \ne 0 \}} = \overline{\mathbb{R}^n} = \mathbb{R}^n,
		\]
		que no es compacto. Por tanto, \( 1 \notin C_c(\mathbb{R}^n) \), y el anillo no tiene unidad.
	}
\end{proof}
\begin{definicion}
	\upshape{
		En el anillo \( R \), para todo \( a, b \in R \), se define
		\[
		a - b := a + (-b),
		\]
		donde \( -b \) denota el inverso aditivo de \( b \) en \( R \).
	}
\end{definicion}



\begin{proposicion}
	\upshape{
		Sea \( R \) un anillo. Para todo \( a, b, c \in R \), se cumple:
		\begin{enumerate}
			\item \( 0a = a0 = 0 \).
			\item \( (-a)b = a(-b) = -(ab) \).
			\item \( (-a)(-b) = ab \).
			\item \( a(b - c) = ab - ac \) y \( (b - c)a = ba - ca \).
				\item Para todo \( a_1, \dots, a_m, b_1, \dots, b_n \in R \), se tiene
			\[
			\left( \sum_{i=1}^{m} a_i \right) \left( \sum_{j=1}^{n} b_j \right) = \sum_{i=1}^{m} \sum_{j=1}^{n} a_i b_j.
			\]
		\end{enumerate}
		
	}
\end{proposicion}

\begin{proof}
	\textcolor{white}{text}
	
	\medskip
	\noindent $[1]$ Sea \( a \in R \). Usamos la distributividad:
	\[
	a \cdot 0 = a(0 + 0) = a \cdot 0 + a \cdot 0.
	\]
	Restando \( a \cdot 0 \) a ambos lados:
	\[
	a \cdot 0 = 0.
	\]
	Análogamente, \( 0 \cdot a = 0 \).\\
	

	\noindent $[2]$ Sea \( a, b \in R \). Como \( -a + a = 0 \), entonces:
	\[
	(-a)b + ab = (-a + a)b = 0b = 0,
	\]
	así que \( (-a)b = -(ab) \). Similarmente,
	\[
	a(-b) + ab = a(-b + b) = a \cdot 0 = 0,
	\]
	por lo que \( a(-b) = -(ab) \). \\
	

	
	\noindent $[3]$ Aplicando el punto anterior:
	\[
	(-a)(-b) = -\left( a(-b) \right) = -\left( -(ab) \right) = ab.
	\]
	
	
	\noindent $[4]$ Sea \( a, b, c \in R \). Usamos la distributividad:
	\[
	a(b - c) = a(b + (-c)) = ab + a(-c) = ab - ac,
	\]
	y
	\[
	(b - c)a = (b + (-c))a = ba + (-c)a = ba - ca.
	\]
	
	
	\noindent $[5]$ Procedemos por distributividad iterada. Sea \( a_1, \dots, a_m, b_1, \dots, b_n \in R \). Aplicando distributividad en etapas:
	\[
	\left( \sum_{i=1}^m a_i \right) \left( \sum_{j=1}^n b_j \right)
	= \sum_{i=1}^m \left( a_i \sum_{j=1}^n b_j \right)= \sum_{i=1}^m \left(  \sum_{j=1}^n a_ib_j \right)
	= \sum_{i=1}^m \sum_{j=1}^n a_i b_j.
	\]
	
\end{proof}

\begin{proposicion}
	\upshape{
		Sea $R$ un anillo conmutativo. Entonces, para cualesquiera $a, b \in R$ y para todo entero $n \geq 0$, se cumple
		\[
		(a + b)^n = \sum_{k = 0}^{n} \binom{n}{k} a^k b^{n - k}.
		\]
	}
\end{proposicion}
\begin{proof}
	La demostración se realiza por inducción sobre $n$ y sigue el mismo esquema que la demostración estándar del binomio de Newton en $\mathbb{R}$.
\end{proof}


\begin{definicion}
	\upshape{
		Un elemento \( a \in R \setminus \{0\} \) se llama \textit{divisor de cero} si existe \( b \in R \setminus \{0\} \) tal que $$ ab = 0  \quad\text{
		o} \quad ba = 0 .$$
	}
\end{definicion}

\medskip
\begin{ejemplo}
	\upshape{
		Sea \( \mathbb{K} = \mathbb{R} \) o \( \mathbb{C} \), y sea \( R = M_{2 \times 2}(\mathbb{K}) \). Entonces \( R \) tiene divisores de cero.\\
		
		
		\noindent
		Si \( \mathbb{K} = \mathbb{R} \), considérese, por ejemplo, las matrices
		\[
		A =
		\begin{pmatrix}
			1 & 0 \\
			0 & 0
		\end{pmatrix},
		\qquad
		B =
		\begin{pmatrix}
			0 & 1 \\
			0 & 0
		\end{pmatrix}.
		\]
		Entonces
		\[
		AB =
		\begin{pmatrix}
			0 & 1 \\
			0 & 0
		\end{pmatrix},
		\qquad
		BA =
		\begin{pmatrix}
			0 & 0 \\
			0 & 0
		\end{pmatrix}.
		\]
		Como \( A \ne 0 \), \( B \ne 0 \) y \( BA = 0 \), se concluye que \( A \) es un divisor de cero por la derecha y \( B \) por la izquierda.\\
		
	
		
		\noindent
		Si \( \mathbb{K} = \mathbb{C} \), consideremos las matrices
		\[
		M =
		\begin{pmatrix}
			1 & i \\
			-i & 1
		\end{pmatrix},
		\qquad
		N =
		\begin{pmatrix}
			1 & -i \\
			i & 1
		\end{pmatrix}.
		\]
		Entonces
		\[
		MN =
		\begin{pmatrix}
			1 + i^2 & -i + i \\
			-i + i & i^2 + 1
		\end{pmatrix}
		=
		\begin{pmatrix}
			0 & 0 \\
			0 & 0
		\end{pmatrix}.
		\]
		Como \( M \ne 0 \), \( N \ne 0 \), y \( MN = 0 \), se concluye que \( M \) y \( N \) son divisores de cero en \( M_{2 \times 2}(\mathbb{C}) \).
	}
\end{ejemplo}


\begin{definicion}
	\upshape{
		Sea \( R \) un anillo. Entonces:
		\begin{enumerate}
			\item Decimos que \( R \) es un \textit{dominio} si es un anillo sin divisores de cero.
			\item Si \( R \) es un dominio y además es conmutativo, se dice que \( R \) es un \textit{dominio de integridad}.
		\end{enumerate}
	}
\end{definicion}
\begin{ejemplo}
	\upshape{
		El anillo \( \mathbb{Z}/n\mathbb{Z} \) es un dominio de integridad si y solo si \( n \) es primo.
	}
\end{ejemplo}

\begin{proof}
	\upshape{
		\textcolor{white}{text}
		
		\medskip
		\noindent $[\Rightarrow]$ Supongamos primero que \( n \) no es primo. Entonces existen \( a, b \in \mathbb{Z} \) tales   con \( 1 < a < n \), \( 1 < b < n \), tales que \( ab \equiv 0 \pmod{n} \). Esto implica que las clases \( \overline{a}, \overline{b} \in \mathbb{Z}/n\mathbb{Z} \) son no nulas, pero \( \overline{a} \cdot \overline{b} = \overline{0} \), así que hay divisores de cero y \( \mathbb{Z}/n\mathbb{Z} \) no es dominio de integridad.\\
		
		\noindent $[\Leftarrow]$ Recíprocamente, supongamos que \( n \) es primo. Si \( \overline{a}, \overline{b} \in \mathbb{Z}/n\mathbb{Z} \) satisfacen \( \overline{a} \cdot \overline{b} = \overline{0} \), entonces \( ab \equiv 0 \pmod{n} \). Como \( n \) es primo, esto implica que \( n \mid a \) o \( n \mid b \), es decir, que \( \overline{a} = \overline{0} \) o \( \overline{b} = \overline{0} \). Por tanto, no hay divisores de cero, y \( \mathbb{Z}/n\mathbb{Z} \) es un dominio de integridad.
	}
\end{proof}




\begin{ejemplo}
	\upshape{
		El anillo \( C^0(\mathbb{R}) \), con la suma y el producto definidos punto a punto, no es un dominio.
		
		\medskip
		
		\noindent Sean
		\[
		f(x) := \begin{cases}
			e^{-1/(1 - x^2)} & \text{si } |x| < 1, \\
			0 & \text{si } |x| \geq 1,
		\end{cases}
		\quad
		g(x) := \begin{cases}
			e^{-1/(1 - (x - 2)^2)} & \text{si } |x - 2| < 1, \\
			0 & \text{si } |x - 2| \geq 1.
		\end{cases}
		\]
		Entonces \( f, g \in C^0(\mathbb{R}) \), \( f \ne 0 \), \( g \ne 0 \), pero se cumple que \( fg = 0 \), ya que \( f(x)g(x) = 0 \) para todo \( x \in \mathbb{R} \).
		
		\medskip
		
		\noindent Por tanto, \( f \) y \( g \) son divisores de cero, y se concluye que \( C^0(\mathbb{R}) \) no es un dominio.
	}
\end{ejemplo}






\begin{observacion}
	\upshape{
		Se llama \textit{función bump} a toda función \( \varphi \colon \mathbb{R}^n \to \mathbb{R} \) que cumple:
		\begin{itemize}
			\item \( \varphi \in C^\infty(\mathbb{R}^n) \), es decir, es infinitamente diferenciable en todo \( \mathbb{R}^n \);
			\item \( \operatorname{supp}(\varphi) \subseteq K \) para algún compacto \( K \subseteq \mathbb{R}^n \), es decir, tiene soporte compacto.
		\end{itemize}
		
		\medskip
		\noindent Una construcción clásica de funciones bump es la siguiente. Sea
		\[
		\varphi(x) := 
		\begin{cases}
			\exp\left( -\dfrac{1}{1 - \|x\|^2} \right) & \text{si } \|x\| < 1, \\[6pt]
			0 & \text{si } \|x\| \geq 1,
		\end{cases}
		\]donde \( \|x\| \) denota la norma euclídea. Esta función pertenece a \( C^\infty(\mathbb{R}^n) \) y su soporte está contenido en la bola cerrada \( \overline{B(0,1)} \).
		
		\medskip
		
		\noindent De forma más general, dados \( a \in \mathbb{R}^n \) y \( r > 0 \), se define
		
		
		\[
		\varphi_{a,r}(x) := 
		\begin{cases}
			\exp\left( -\dfrac{1}{1 - \left( \frac{\|x - a\|}{r} \right)^2} \right) & \text{si } \|x - a\| < r, \\[6pt]
			0 & \text{si } \|x - a\| \geq r.
		\end{cases}
		\]
		Entonces \( \varphi_{a,r} \in C^\infty(\mathbb{R}^n) \), su soporte está contenido en \( \overline{B(a,r)} \), y todas sus derivadas se anulan sobre la frontera \( \partial B(a,r) \).
		
		\medskip
		
		\noindent Este tipo de funciones es fundamental en análisis y geometría, pues permite construir funciones suaves con comportamiento controlado en regiones localizadas del espacio.
	}
\end{observacion}

\begin{ejemplo}
	\upshape{
		Denotamos por \( C^\infty(\mathbb{R}) \) al conjunto de todas las funciones \( f \colon \mathbb{R} \to \mathbb{R} \) que son derivables infinitas veces. Con la suma y el producto definidos punto a punto, este conjunto forma un anillo conmutativo con unidad.
		
		\medskip
		
		\noindent Este anillo no es un dominio. En efecto, consideremos las funciones
		\[
		f(x) := \begin{cases}
			e^{-1/(1 - x^2)} & \text{si } |x| < 1, \\
			0 & \text{si } |x| \geq 1,
		\end{cases}
		\quad
		g(x) := \begin{cases}
			e^{-1/(1 - (x - 2)^2)} & \text{si } |x - 2| < 1, \\
			0 & \text{si } |x - 2| \geq 1.
		\end{cases}
		\]
		Ambas funciones pertenecen a \( C^\infty(\mathbb{R}) \). Así, \( f \ne 0 \), \( g \ne 0 \), pero \( fg = 0 \), lo que muestra que \( f \) y \( g \) son divisores de cero. En consecuencia, el anillo \( C^\infty(\mathbb{R}) \) no es un dominio.
	}
\end{ejemplo}

\begin{observacion}
	\upshape{
		El anillo \( C^\infty_c(\mathbb{R}^n) \) no es un dominio. Basta notar que existen funciones bump \( f, g \in C^\infty_c(\mathbb{R}^n) \), con soportes disjuntos, tales que \( f \ne 0 \), \( g \ne 0 \), pero \( fg = 0 \).
	}
\end{observacion}










\begin{definicion}
	\upshape{
		Sea \( R \) un anillo no trivial (es decir, \( R \neq \{0\} \)). Decimos que \( R \) cumple la \textit{ley cancelativa} si para cualesquiera \( a, b, c \in R \) con $a\neq 0$ se tiene:
		\[
		ba = ca \;\:\:\text{entonces}\:\; b = c,
		\qquad\text{y}\qquad
		ab = ac \;\:\:\text{entonces}\:\; b = c.
		\]
	}
\end{definicion}


\begin{proposicion}
	\upshape{
		Sea \( R \) un anillo. Son equivalentes:
		\begin{enumerate}
			\item \( R \) es un dominio.
			\item \( R \) cumple la ley cancelativa.
		\end{enumerate}
	}
\end{proposicion}

\begin{proof}
	\textcolor{white}{text}
	
	\medskip
	\noindent $[1 \Rightarrow 2]$ Sean \( a, b, c \in R \) con \( a \ne 0 \) y \( ab = ac \). Entonces
	\[
	ab - ac = a(b - c) = 0.
	\]
	Como \( R \) no tiene divisores de cero y \( a \ne 0 \), se concluye que \( b - c = 0 \), es decir, \( b = c \). El argumento para \( ba = ca \) es análogo.
	
	\medskip
	\noindent $[2 \Rightarrow 1]$ Sean \( a, b \in R \) tales que \( ab = 0 \). Entonces:
	\begin{itemize}
		\item Si \( a \ne 0 \), entonces
		\[
		ab = 0 = a \cdot 0,
		\]
		y por la ley cancelativa, \( b = 0 \).
		
		\item Si \( b \ne 0 \), entonces
		\[
		ab = 0 = 0 \cdot b,
		\]
		y por la ley cancelativa, \( a = 0 \).
	\end{itemize}
	Por tanto, \( R \) no tiene divisores de cero y es un dominio.
\end{proof}












\section{Anillo monoidal}

Sea \( (M, \star) \) un monoide conmutativo con unidad.

\begin{definicion}\label{lkjhgxdcsd}
	\upshape{
		Sea \( (R, +, \bullet) \) un anillo conmutativo con unidad. Definimos el \textit{anillo monoidal generado por \( M \)} como
		\[
		R[M] := \{ f \colon M \to R \mid f \text{ tiene soporte finito} \}.
		\]
	}
\end{definicion}

\begin{observacion}
	\upshape{
		Diremos que una función \( f \colon M \to R \) tiene \textit{soporte finito} si existe un subconjunto finito \( S \subseteq M \) tal que
		\[
		f(m) = 0 \quad \text{para todo } m \in M \setminus S.
		\]
	}
\end{observacion}


En la Definición \ref{lkjhgxdcsd}, se introdujo el conjunto \( R[M] \) como el conjunto de funciones de soporte finito de \( M \) en \( R \). Para dotar a \( R[M] \) de estructura de anillo, se definen dos operaciones: la suma se realiza punto a punto, mientras que el producto se define mediante convolución, utilizando la operación del monoide \( M \).

\begin{align*}
	\text{Suma:} \quad & (f + g)(a) := f(a) + g(a), \quad \text{para todo } a \in M, \\[8pt]
	\text{Producto (convolución):} \quad & (f\cdot g)(a) := \sum_{m \star n = a} f(m) \bullet g(n), \quad \text{para todo } a \in M.
\end{align*}
\begin{observacion}
	\upshape{
		En la fórmula del producto, \( m \star n = a \) significa que la suma recorre todos los pares \( (m, n) \in M \times M \) tales que \( m \star n = a \), es decir, que combinan mediante la operación del monoide para dar \( a \).
	}
\end{observacion}

Probemos que \( R[M] \) es un anillo. En efecto:



\begin{itemize}
	
	\item \( (R[M], +) \) es un grupo
	\begin{itemize}
		\item \( + \) es un producto interno: si \( f, g \in R[M] \), entonces \( f + g \in R[M] \), ya que existen subconjuntos finitos \( S_f, S_g \subseteq M \) tales que
		\[
		f(m) = 0 \quad \text{para todo } m \in M \setminus S_f,
		\qquad
		g(m) = 0 \quad \text{para todo } m \in M \setminus S_g.
		\]
		Sea \( S := S_f \cup S_g \), que también es finito. Entonces, para todo \( m \in M \setminus S \), se tiene
		\[
		(f + g)(m) = f(m) + g(m) = 0 + 0 = 0,
		\]
		por lo tanto, \( f + g \in R[M] \).
		
	
		
		\item Asociatividad: para todo \( a \in M \),
		\begin{align*}
			((f + g) + h)(a) 
			&= (f + g)(a) + h(a) \\
			&= f(a) + g(a) + h(a) \\
			&= f(a) + (g(a) + h(a)) \\
			&= (f + (g + h))(a). 
		\end{align*}

		\item Conmutatividad: como \( (R, +) \) es abeliano, para todo \( a \in M \),
		\[
		(f + g)(a) = f(a) + g(a) = g(a) + f(a) = (g + f)(a).
		\]
		
		\item Elemento neutro: la función nula \( 0 \in R[M] \), dada por \( 0(a) := 0_R \) para todo \( a \in M \), cumple
		\[
		(f + 0)(a) = f(a) + 0_R = f(a).
		\]
		
		\item Elemento inverso: si \( f \in R[M] \), entonces \( -f \in R[M] \), definido por \( (-f)(a) := -f(a) \), cumple
		\[
		(f + (-f))(a) = f(a) + (-f(a)) = 0_R.
		\]
	\end{itemize}
	
	\item \( (R[M], \cdot) \) es un semigrupo
	\begin{itemize}
		\item El producto es un producto interno: si \( f, g \in R[M] \), entonces \( fg \in R[M] \), ya que existen subconjuntos finitos \( S_f, S_g \subseteq M \) tales que
		\[
		f(m) = 0 \quad \text{para todo } m \in M \setminus S_f,
		\qquad
		g(n) = 0 \quad \text{para todo } n \in M \setminus S_g.
		\]
		Sea
		\[
		S := \{ m \star n \mid m \in S_f,\ n \in S_g \} \subseteq M.
		\]
		Como \( S_f \) y \( S_g \) son finitos, y \( \star \) es una operación binaria, el conjunto \( S \) también es finito. Si \( a \in M \setminus S \), entonces no existen \( (m, n) \in S_f \times S_g \) tales que \( m \star n = a \). Por tanto, si \( m \star n = a \), se tiene que \( m \notin S_f \) o \( n \notin S_g \), y en consecuencia \( f(m) = 0 \) o \( g(n) = 0 \), lo que implica \( f(m) \bullet g(n) = 0 \). Así,
		\[
		(fg)(a) = \sum_{m \star n = a} f(m) \bullet g(n) = 0,
		\]
		lo cual muestra que \( fg \) se anula fuera de un subconjunto finito \( S \), es decir, \( fg \in R[M] \).
		
	
		
		\item Asociatividad: para todo \( a \in M \),
		\begin{align*}
			((fg)h)(a) 
			&= \sum_{x \star y = a} (fg)(x) \bullet h(y) \\[5pt]
			&= \sum_{x \star y = a} \left( \sum_{m \star n = x} f(m) \bullet g(n) \right) \bullet h(y) \\[5pt]
			&= \sum_{\substack{m, n, y \in M \\ (m \star n) \star y = a}} \left(f(m) \bullet g(n)\right) \bullet h(y) \\[5pt]
			&= \sum_{\substack{m, n, y \in M \\ m \star (n \star y) = a}} f(m) \bullet\left( g(n) \bullet h(y)\right) \\[5pt]
			&= \sum_{m \star p = a} f(m) \bullet \left( \sum_{n \star y = p} g(n) \bullet h(y) \right) \\[5pt]
			&= \sum_{m \star p = a} f(m) \bullet (gh)(p) \\[5pt]
			&= (f(gh))(a). \\[5pt]
		\end{align*}
	\end{itemize}
	
	\item Distributividad
	\begin{itemize}
		\item Para todo \( a \in M \),
		\begin{align*}
			f(g + h)(a) 
			&= \sum_{m \star n = a} f(m) \bullet (g + h)(n) \\[5pt]
			&= \sum_{m \star n = a} f(m) \bullet (g(n) + h(n)) \\[5pt]
			&= \sum_{m \star n = a} f(m) \bullet g(n) + f(m) \bullet h(n) \\[5pt]
			&= \sum_{m \star n = a} f(m) \bullet g(n) + \sum_{m \star n = a} f(m) \bullet h(n) \\[5pt]
			&= (fg)(a) + (fh)(a). \\[5pt]
		\end{align*}
		\item Análogamente se prueba que \( (g + h)f = gf + hf \).
	\end{itemize}
	
\end{itemize}

\begin{proposicion}
	\upshape{
		\( R[M] \) es un anillo conmutativo con unidad.
	}
\end{proposicion}


\begin{proof}
	Ya se ha verificado que \( R[M] \) es un anillo. Solo falta demostrar que la convolución es conmutativa y que existe la identidad multiplicativa.
	
	\medskip
	
	\noindent Conmutatividad: Sea \( f, g \in R[M] \) y \( a \in M \). 	Dado que \( M \) es conmutativo, el conjunto de pares \( (m,n) \in M \times M \) con \( m \star n = a \) es el mismo que el conjunto de pares \( (n,m) \) con \( n \star m = a \). Como además \( R \) es conmutativo, se tiene \( f(m) \bullet g(n) = g(n) \bullet f(m) \). Por lo tanto:
	\[
	(f \cdot g)(a) = \sum_{m \star n = a} f(m) \bullet g(n) = \sum_{n \star m = a} g(n) \bullet f(m) = (g \cdot f)(a),
	\]
	y por ende \( f \cdot g = g \cdot f \), es decir, el producto es conmutativo.\\
	
	\medskip
	
	\noindent Para la existencia de identidad, definimos \( \delta \in R[M] \) como
	\[
	\delta(m) := 
	\begin{cases}
		1_R & \text{si } m = e_M, \\[4pt]
		0   & \text{en otro caso}.
	\end{cases}
	\]
	Entonces, para todo \( f \in R[M] \) y \( a \in M \), se tiene
	\[
	(\delta \cdot f)(a) = \sum_{m \star n = a} \delta(m) \bullet f(n).
	\]
	El único \( m \) con \( \delta(m) \ne 0 \) es \( m = e_M \), en cuyo caso \( n = a \), por lo que:
	\[
	(\delta \cdot f)(a) = \delta(e_M) \bullet f(a) = 1_R \bullet f(a) = f(a).
	\]
	De modo análogo, \( (f \cdot \delta)(a) = f(a) \). Por lo tanto, \( \delta \) actúa como la identidad en \( R[M] \).\\
	

	
	\noindent Concluimos que \( R[M] \) es un anillo conmutativo con unidad.
\end{proof}


\medskip
\begin{ejemplo}\label{ekjasckañsc}
	\upshape{
		Sea \( R = \mathbb{Z} \) y \( M = \mathbb{N}\) con la suma como operación. Consideramos las funciones \( f, g \in \mathbb{Z}[M] \) definidas por:
		\[
		f(n) := \begin{cases}
			1 & \text{si } n = 0, \\
			3 & \text{si } n = 1, \\
			0 & \text{si } n \geq 2,
		\end{cases}
		\quad\quad
		g(n) := \begin{cases}
			2 & \text{si } n = 0, \\
			5 & \text{si } n = 1, \\
			4 & \text{si } n = 2, \\
			0 & \text{si } n \geq 3.
		\end{cases}
		\]
		
		\medskip
		
		\noindent Calculemos el producto \( fg \) en \( \mathbb{Z}[M] \), usando la convolución:
		\begin{align*}
			(fg)(0) &= f(0) \cdot g(0) 
			= 1 \cdot 2 = 2, \\[5pt]
			(fg)(1) &= f(0) \cdot g(1) 
			+ f(1) \cdot g(0) \\ 
			&= 1 \cdot 5 + 3 \cdot 2 = 5 + 6 = 11, \\[5pt]
			(fg)(2) &= f(0) \cdot g(2) 
			+ f(1) \cdot g(1) 
			+ f(2) \cdot g(0) \\
			&= 1 \cdot 4 + 3 \cdot 5 + 0 \cdot 2 = 4 + 15 + 0 = 19, \\[5pt]
			(fg)(3) &= f(0) \cdot g(3) 
			+ f(1) \cdot g(2) 
			+ f(2) \cdot g(1) 
			+ f(3) \cdot g(0) \\
			&= 1 \cdot 0 + 3 \cdot 4 + 0 \cdot 5 + 0 \cdot 2 = 0 + 12 + 0 + 0 = 12, \\[5pt]
			(fg)(n) &= 0 \quad \text{para todo } n \geq 4.
		\end{align*}
		
		
		
		
		\medskip
		
		\noindent Luego, el producto \( fg \in \mathbb{Z}[M] \) es la función:
		\[
		(fg)(n) := \begin{cases}
			2 & \text{si } n = 0, \\
			11 & \text{si } n = 1, \\
			19 & \text{si } n = 2, \\
			12 & \text{si } n = 3, \\
			0 & \text{si } n \geq 4,
		\end{cases}
		\]
		Estas funciones corresponden a los polinomios \( f(x) = 3x + 1 \) y \( g(x) = 4x^2 + 5x + 2 \), y su producto usual es
		\[
		fg(x) = (3x + 1)(4x^2 + 5x + 2) = 12x^3 + 19x^2 + 11x + 2.
		\]
	}
\end{ejemplo}

\vspace{0.5cm}

Además, en el Ejemplo \ref{ekjasckañsc}, notemos que
\[
g(x) = g(0)x^0 + g(1)x^1 + g(2)x^2 + \cdots = \sum_{n \in M} g(n)x^n.
\]
En el contexto del anillo \( R[M] \), es común identificar formalmente el símbolo \( x^n \) con el elemento \( n \in M \). Por tanto, podemos escribir a \( g \) como una combinación formal:
\[
g = \sum_{n \in M} g(n)  n,
\]
donde \( g(n) \in R \) y \( n \in M \), y se entiende que \( g(n)n \) representa la función \( \delta_n \colon M \to R \) definida por
\[
\delta_n(m) = \begin{cases}
	g(n) & \text{si } m = n, \\
	0 & \text{en otro caso}.
\end{cases}
\]
De este modo, \( g \) es la suma (finita) de tales funciones. \\

\begin{observacion}
	\upshape{
		En general, para \( f \in R[M] \), se suele usar la notación
		\[
		f = \sum_{n \in M} f(n) n,
		\]
		lo cual tiene sentido ya que, al tener \( f \) soporte finito, esta suma es finita.
	}
\end{observacion}

En el Ejemplo \ref{ekjasckañsc} vimos que un polinomio puede identificarse con un elemento de $\mathbb{Z}[\mathbb{N}]$. Esto se debe a que el conjunto de polinomios en una indeterminada $x$ con coeficientes en $\mathbb{Z}$, dotado de las operaciones usuales de suma y multiplicación, forma un anillo $\mathbb{Z}[x]$ que es isomorfo al anillo monoidal $\mathbb{Z}[\mathbb{N}]$.

\medskip
\noindent
Esta identificación proporciona una definición rigurosa del anillo de polinomios como un caso particular de un anillo monoidal, donde el monoide $\mathbb{N}$ (con la operación de suma) indexa los exponentes de la indeterminada $x$. Esta construcción se ilustrará con mayor detalle en los ejemplos siguientes.

\begin{ejemplo}[Anillo de polinomios en una variable]
	\upshape{
		Sea \( (R,+,\cdot) \) un anillo conmutativo con unidad, y consideremos el monoide conmutativo con unidad \( (\mathbb{N}, +) \). Entonces, el\textit{ anillo de polinomios en una variable} con coeficientes en \( R \) se define como
		\[
		R[X] := R[\mathbb{N}],
		\]
		donde la estructura de anillo se da mediante las siguientes operaciones:
		
		\begin{align*}
			\text{Suma:} \quad & (f + g)(a) := f(a) + g(a), \quad \text{para todo } a \in \mathbb{N}, \\[8pt]
			\text{Producto (convolución):} \quad & (f \cdot g)(a) := \sum_{m + n = a} f(m) \cdot g(n), \quad \text{para todo } a \in \mathbb{N}.
		\end{align*}
		
	}
\end{ejemplo}

\begin{observacion}
	\upshape{
		Para cada \( n \in \mathbb{N} \), definimos el monomio \( x^n \in R[\mathbb{N}] \) como la función \( f_n \colon \mathbb{N} \to R \) dada por
		\[
		f_n(k) := \begin{cases}
			1_R & \text{si } k = n, \\[4pt]
			0 & \text{en otro caso}.
		\end{cases}
		\]
		Usando la notación \( f = \sum_{k \in \mathbb{N}} f(k) \cdot k \), tenemos:
		\[
		x^n := f_n = f_n(n) \cdot n = 1_R \cdot n.
		\]
		
		\medskip
		
		\noindent Entonces, para cualquier \( r \in R \), definimos el monomio \( r x^n \in R[\mathbb{N}] \) como la función \( f \colon \mathbb{N} \to R \) dada por
		\[
		f(k) := r \cdot f_n(k) = \begin{cases}
			r & \text{si } k = n, \\[4pt]
			0 & \text{en otro caso}.
		\end{cases}
		\]
		Así,
		\[
		r x^n = r \cdot f_n = r \cdot n.
		\]
		
		\medskip
		
		\noindent Por ejemplo,
		\[
		x^0 := 1_R \cdot 0, \quad
		x := x^1 = 1_R \cdot 1, \quad
		x^2 := 1_R \cdot 2, \quad
		x^3 := 1_R \cdot 3, \quad \dots
		\]
		
		\medskip
		
		\noindent Y combinaciones como \( 4x^5 - 3x + 7 \) se escriben como
		\[
		4 \cdot 5 + (-3) \cdot 1 + 7 \cdot 0 \in R[\mathbb{N}].
		\]
		
		\medskip
		
		\noindent Más generalmente, un polinomio
		\[
		a_0 + a_1 x + a_2 x^2 + \dots + a_n x^n
		\]
		se representa como
		\[
		a_0 \cdot 0 + a_1 \cdot 1 + a_2 \cdot 2 + \dots + a_n \cdot n = \sum_{k = 0}^{n} a_k \cdot k \in R[\mathbb{N}].
		\]
	}
\end{observacion}



\begin{ejemplo}[Anillo de polinomios en \( n \) variables]
	\upshape{
		Sea \( (R,+,\cdot) \) un anillo conmutativo con unidad, y consideremos el monoide conmutativo con unidad \( (\mathbb{N}^n,+).\) Entonces, el \textit{anillo de polinomios en \( n \) variables} con coeficientes en \( R \) se define como
		\[
		R[x_1,x_2\dots,x_n] := R[\mathbb{N}^n],
		\]
		donde la estructura de anillo se da mediante las siguientes operaciones:
		
		\begin{align*}
			\text{Suma:} \quad & (f + g)(\alpha) := f(\alpha) + g(\alpha), \quad \text{para todo } \alpha \in \mathbb{N}^n, \\[8pt]
			\text{Producto (convolución):} \quad & (f \cdot g)(\alpha) := \sum_{\beta + \gamma = \alpha} f(\beta) \cdot g(\gamma), \quad \text{para todo } \alpha \in \mathbb{N}^n.
		\end{align*}
		Donde los elementos de \( \mathbb{N}^n \) se interpretan como multiíndices \( \alpha = (\alpha_1,\dots,\alpha_n) \).
	}
\end{ejemplo}




\begin{observacion}
	\upshape{
		Para cada \( \alpha = (\alpha_1, \dots, \alpha_n) \in \mathbb{N}^n \), definimos el monomio $$x^\alpha := x_1^{\alpha_1} \cdots x_n^{\alpha_n} \in R[\mathbb{N}^n] $$ como la función \( f_\alpha \colon \mathbb{N}^n \to R \) dada por
		\[
		f_\alpha(\beta) := \begin{cases}
			1_R & \text{si } \beta = \alpha, \\[4pt]
			0 & \text{en otro caso}.
		\end{cases}
		\]
		
		\medskip
		
		\noindent Usando la notación \( f = \sum_{\beta \in \mathbb{N}^n} f(\beta) \cdot \beta \), tenemos:
		\[
		x^\alpha := f_\alpha = f_\alpha(\alpha) \cdot \alpha = 1_R \cdot \alpha.
		\]
		
		\medskip
		
		\noindent Entonces, para cualquier \( r \in R \), definimos el monomio \( r x^\alpha \in R[\mathbb{N}^n] \) como la función \( f \colon \mathbb{N}^n \to R \) dada por
		\[
		f(\beta) := r \cdot f_\alpha(\beta) = \begin{cases}
			r & \text{si } \beta = \alpha, \\[4pt]
			0 & \text{en otro caso}.
		\end{cases}
		\]
		Así,
		\[
		r x^\alpha = r \cdot f_\alpha = r \cdot \alpha.
		\]
		
		\medskip
		
		\noindent Por ejemplo, en \( R[\mathbb{N}^3] \), denotando por \( x := x_1, y := x_2, z := x_3 \), tenemos:
		\[
		x := x^{(1,0,0)} = 1_R \cdot (1,0,0), \quad
		y := x^{(0,1,0)} = 1_R \cdot (0,1,0), \quad
		xyz := x^{(1,1,1)} = 1_R \cdot (1,1,1), \quad \dots
		\]
		
		\medskip
		
		\noindent Y combinaciones como \( 2x^2z + 5yz^3 - 4 \) se escriben como
		\[
		2 \cdot (2,0,1) + 5 \cdot (0,1,3) + (-4) \cdot (0,0,0) \in R[\mathbb{N}^3].
		\]
		
		\medskip
		
		\noindent Más generalmente, un polinomio
		\[
		\sum_{\alpha \in \mathbb{N}^n} a_\alpha x^\alpha
		\]
		se representa como
		\[
		\sum_{\alpha \in \mathbb{N}^n} a_\alpha \cdot \alpha \in R[\mathbb{N}^n].
		\]
	}
\end{observacion}





\section{Anillo de polinomios}



		Definir el anillo de polinomios como un caso particular de anillo monoidal ofrece una formulación general y rigurosa. No obstante, el enfoque clásico mediante expresiones algebraicas resulta igualmente valioso, ya que permite una manipulación directa de los polinomios en diversos contextos.\medskip




\begin{definicion}\label{ñlkjpoiuyy}
	\upshape{
		Sea \( R \) un anillo conmutativo con unidad y \( x \) un símbolo que no pertenece a \( R \). Un \textit{polinomio en \( x \)} con coeficientes en \( R \) es una expresión formal
		\[
		p(x) = a_0 + a_1x + a_2x^^2 + \dots + a_nx^n,
		\]
		donde \( n \in \mathbb{N} \), cada \( a_k \in R \), y \( a_n \) puede ser cero (en cuyo caso el término \( x^n \) puede omitirse). El conjunto de todos estos polinomios se denota por
		\[
		R[x] := \left\{ \sum_{k=0}^n a_k x^k : n \in \mathbb{N}, \; a_k \in R \text{ para } 0 \leq k \leq n \right\}.
		\]
		Un elemento \( p \in R[x] \) suele escribirse como \( p = p(x) \), enfatizando su dependencia formal de \( x \).
	}
\end{definicion}



\vspace{0.5cm}
En \( R[x] \), las operaciones de \textit{suma} y \textit{producto} se definen de la siguiente manera. Sean
\[
p(x) = \sum_{k = 0}^n a_k x^k, \quad q(x) = \sum_{k = 0}^m b_k x^k,
\]
con \( a_k, b_k \in R \). Entonces:
\begin{align*}
	\text{Suma:} \quad (p + q)(x) &:= \sum_{k = 0}^{\max(n,m)} c_k x^k, \quad \text{donde } c_k :=
	\begin{cases}
		a_k + b_k & \text{si } k \leq \min(n,m), \\[4pt]
		a_k & \text{si } n > m \text{ y } k > m, \\[4pt]
		b_k & \text{si } m > n \text{ y } k > n,
	\end{cases} \\[10pt]
	\text{Producto:} \quad (p \cdot q)(x) &:= \sum_{k = 0}^{n+m} \left( \sum_{\substack{0 \leq i \leq n,\: 0 \leq j \leq m \\ i + j = k}} a_i b_j \right) x^k.
\end{align*}

\noindent Es decir, la suma se obtiene sumando término a término los coeficientes de igual grado, y el producto resulta de aplicar la distributividad junto con la regla \( x^i \cdot x^j = x^{i+j} \).\\

\noindent Estas operaciones hacen de \( R[x] \) un anillo conmutativo con unidad, donde el neutro aditivo es el polinomio nulo \( 0(x) := 0 \), y el elemento identidad para el producto es el polinomio constante \( 1(x) := 1 \in R \).

\begin{proposicion}
	\upshape{
		\( R[x] \) es un anillo conmutativo con unidad.
	}
\end{proposicion}




\begin{definicion}
	\upshape{
		Sea \( R \) un anillo conmutativo con unidad. Se define \( R[x,y] \) como el conjunto de todas las expresiones finitas de la forma
		\[
		\sum_{i = 0}^n \sum_{j = 0}^m a_{i,j} x^i y^j,
		\]
		con \( a_{i,j} \in R \). La suma y el producto se definen término a término y distributivamente, extendiendo la regla $$ x^i y^j \cdot x^{i'} y^{j'} = x^{i+i'} y^{j+j'}.$$
	}
\end{definicion}

\vspace{0.3cm}

\begin{definicion}
	\upshape{
		Análogamente, al considerar el anillo \( (R[x])[y] \), sus elementos son polinomios en \( y \) cuyos coeficientes están en \( R[x] \), es decir, expresiones finitas de la forma
		\[
		\sum_{j = 0}^m f_j(x) y^j,
		\]
		donde cada \( f_j(x) \in R[x] \), y a su vez cada uno puede escribirse como
		\[
		f_j(x) = \sum_{i = 0}^{n_j} a_{i,j} x^i, \quad \text{con } a_{i,j} \in R.
		\]
		Por tanto, todo elemento de \( (R[x])[y] \) puede escribirse como una suma finita de términos \( a_{i,j} x^i y^j \).
		
		\medskip
		
		\noindent La suma y el producto en \( (R[x])[y] \) se definen siguiendo las mismas reglas de la  Definición \ref{ñlkjpoiuyy}).
	}
\end{definicion}


\vspace{0.3cm}
\begin{proposicion}
	\upshape{
		Los conjuntos \( (R[x])[y] \) y \( R[x,y] \) son iguales, ya que en ambos casos sus elementos son combinaciones finitas de términos de la forma \( a_{i,j} x^i y^j \) con \( a_{i,j} \in R \), \( i, j \in \mathbb{N} \).
	}
\end{proposicion}


\begin{proof}
	Por doble inclusión.
	
	\medskip
	
	\noindent \( (R[x])[y] \subseteq R[x,y] \): 
	Sea \( f \in (R[x])[y] \). Entonces, por definición, existe una expresión finita
	\[
	f = \sum_{j=0}^m f_j(x) y^j,
	\]
	donde cada \( f_j(x) \in R[x] \). Como \( f_j(x) = \sum_{i=0}^{n_j} a_{i,j} x^i \) con \( a_{i,j} \in R \), se tiene:
	\[
	f = \sum_{j=0}^m \left( \sum_{i=0}^{n_j} a_{i,j} x^i \right) y^j = \sum_{j=0}^m \sum_{i=0}^{n_j} a_{i,j} x^i y^j.
	\]
	Como la suma es finita, esto es un elemento de \( R[x,y] \), es decir, un polinomio en dos variables con coeficientes en \( R \). Por tanto, \( f \in R[x,y] \).\\
	
	\noindent \( R[x,y] \subseteq (R[x])[y] \):
	Sea \( f \in R[x,y] \). Por definición, se puede escribir como una suma finita
	\[
	f = \sum_{(i,j) \in F} a_{i,j} x^i y^j,
	\]
	con \( F \subseteq \mathbb{N}^2 \) finito y \( a_{i,j} \in R \). Reuniendo los términos con la misma potencia de \( y \), escribimos:
	\[
	f = \sum_{j=0}^m \left( \sum_{i=0}^{n_j} a_{i,j} x^i \right) y^j.
	\]
	Cada coeficiente entre paréntesis es un elemento de \( R[x] \), por lo que \( f \in (R[x])[y] \).
	
	\medskip
	
	\noindent Como ambas inclusiones se cumplen, concluimos que \( (R[x])[y] = R[x,y] \) como conjuntos.
\end{proof}

\begin{observacion}
	\upshape{
		A pesar de que \( (R[x])[y] = R[x,y] \) como conjuntos, no podemos afirmar que son el mismo anillo, ya que las operaciones de suma y producto en cada caso se definen inicialmente de manera distinta, al menos en su construcción formal. Sin embargo, más adelante veremos que, a pesar de estas diferencias en la definición, ambos anillos resultan ser isomorfos, es decir, existe una correspondencia que respeta las operaciones y preserva la estructura algebraica.
	}
\end{observacion}


\vspace{0.5cm}
De manera general, se define el anillo de polinomios en varias variables como una extensión natural del caso de una sola variable, preservando su estructura algebraica.

\begin{definicion}[Anillo de polinomios en \( n \) variables]
	\upshape{
		Sea \( R \) un anillo conmutativo con unidad. Se define \( R[x_1, x_2, \dots, x_n] \) como el conjunto de todas las expresiones finitas de la forma
		\[
		\sum_{\alpha \in \mathbb{N}^n} a_{\alpha} x^{\alpha},
		\]
		donde \( a_{\alpha} \in R \) y, para cada multiíndice \( \alpha = (\alpha_1, \dots, \alpha_n) \in \mathbb{N}^n \), se define
		\[
		x^{\alpha} := x_1^{\alpha_1} x_2^{\alpha_2} \cdots x_n^{\alpha_n}.
		\]
		La suma y el producto se definen término a término y mediante linealidad, usando la regla
		\[
		x^{\alpha} \cdot x^{\beta} = x^{\alpha + \beta},
		\]
		donde \( \alpha + \beta := (\alpha_1 + \beta_1, \dots, \alpha_n + \beta_n) \). Estas operaciones hacen de \( R[x_1, x_2, \dots, x_n] \) un anillo conmutativo con unidad.
	}
\end{definicion}



\begin{definicion}
	\upshape{
		Sea \( p(x) = \sum_{k = 0}^n a_k x^k \in R[x] \), con \( a_k \in R \). Definimos:
		\begin{enumerate}
			\item Si \( p(x) \neq 0 \), el \textit{grado} de \( p(x) \), denotado por \( \deg(p) \), es el mayor entero \( k \) tal que \( a_k \neq 0 \).
			
			\item Por convenio, si \( p(x) = 0 \), definimos \( \deg(0) := -\infty \).
			
			\item El coeficiente \( a_k \) correspondiente al término de mayor grado se denomina \textit{coeficiente principal} de \( p(x) \).
			
			\item Si el coeficiente principal o líder de \( p(x) \) es igual a \( 1_R \), decimos que \( p(x) \) es un \textit{polinomio monico}.
			
			\item Decimos que \( p(x) \) es un \textit{monomio} si contiene un único término no nulo, es decir, si existe \( k \in \mathbb{N} \) tal que \( p(x) = a_k x^k \) con \( a_k \neq 0 \).
			
		\end{enumerate}
	}
\end{definicion}

\begin{definicion}
	\upshape{
		Sea \( p(x_1, x_2, \dots, x_n) \in R[x_1, x_2, \dots, x_n] \) un polinomio de la forma
		\[
		p(x_1, x_2, \dots, x_n) = \sum_{\alpha \in \mathbb{N}^n} a_{\alpha} x^{\alpha} = \sum_{\alpha \in \mathbb{N}^n} a_{\alpha} x_1^{\alpha_1} x_2^{\alpha_2} \cdots x_n^{\alpha_n},
		\]
		con \( a_{\alpha} \in R \) y solo una cantidad finita de coeficientes no nulos.
		
		\medskip
		
		\noindent Si \( p(x_1, \dots, x_n) \neq 0 \), definimos el \textit{grado total} de \( p \), denotado por \( \deg(p) \), como:
		\[
		\deg(p) := \max \left\{ \alpha_1 + \alpha_2 + \dots + \alpha_n \, : \, a_{\alpha} \neq 0 \right\}.
		\]
		Por convenio, se define \( \deg(0) := -\infty \).
	}
\end{definicion}






\begin{proposicion}
	\upshape{
		Sea \( R \) un dominio de integridad. Entonces, el anillo de polinomios \( R[x_1, x_2, \dots, x_n] \) también es un dominio de integridad.
	}
\end{proposicion}




\begin{proof}
	Sean \( f, g \in R[x_1, \dots, x_n] \) tales que \( fg = 0 \). Supongamos que \( f \neq 0 \) y \( g \neq 0 \). Entonces podemos tomar \( ax^{\alpha} \) el término principal de \( f \), con \( a \neq 0_R \), y \( bx^{\beta} \) el término principal de \( g \), con \( b \neq 0_R \).
	
		\medskip
	\noindent
	Así, el producto \( fg \) contiene el término \( abx^{\alpha + \beta} \), cuyo coeficiente es \( ab \neq 0_R \), ya que \( R \) es un dominio de integridad.
	
		\medskip
	\noindent
	Además, todos los demás términos de \( fg \) tienen grado total menorque grado total de \( abx^{\alpha + \beta} \), por lo que dicho término no puede anularse al sumarse con los demás. Es decir, $fg\neq 0$. Esto contradice que \( fg = 0 \). Por tanto, si \( fg = 0 \), necesariamente \( f = 0 \) o \( g = 0 \), lo que prueba que \( R[x_1, \dots, x_n] \) es un dominio de integridad.
\end{proof}

\vspace{0.5cm}
Como comentario final sobre los polinomios, podemos mencionar su relevancia en diversas ramas de las matemáticas y sus aplicaciones:

\begin{itemize}
	\item En análisis, los polinomios permiten aproximar funciones continuas en intervalos cerrados, según establece el Teorema de Stone-Weierstrass. Además, toda función diferenciable admite una aproximación local mediante su polinomio de Taylor.
	
	\item En análisis complejo, el teorema de Cauchy implica que toda función holomorfa es analítica, es decir, puede representarse localmente mediante una serie de potencias, lo que refuerza la presencia de los polinomios en el estudio de funciones complejas.
	
	
		\item Las funciones trigonométricas satisfacen ecuaciones polinomiales como $c^2+s^2=1$:
	\[
	\cos^2(\theta) + \sen^2(\theta) = 1,
	\]
	lo que evidencia la conexión entre polinomios y relaciones funcionales básicas.
	
	
	
	\item En aplicaciones prácticas, los polinomios aparecen en problemas de optimización, modelado estadístico mediante regresión polinomial, y en teoría de control, donde se utilizan para describir el comportamiento de sistemas dinámicos.
	
	
	\item En geometría algebraica, los polinomios constituyen el lenguaje fundamental para describir variedades algebraicas, generalizando así la geometría analítica clásica. Un ejemplo notable son las curvas cúbicas en dos variables, conocidas como curvas elípticas, que tienen aplicaciones relevantes en teoría de números, como la conjetura de Birch y Swinnerton-Dyer, uno de los problemas del milenio, y en criptografía, especialmente sobre cuerpos finitos.

	\begin{table}[H]
		\centering
		{\large
			\renewcommand{\arraystretch}{1.6}
			\resizebox{160mm}{!}{
				\begin{tabular}{
						>{\raggedright\arraybackslash}p{4.5cm}
						>{\raggedright\arraybackslash}p{5.5cm}
						>{\raggedright\arraybackslash}p{5.5cm}
					}
					\toprule
					\multicolumn{1}{c}{\textbf{Tipo de Ecuación}} & 
					\multicolumn{1}{c}{\textbf{Ejemplo}} & 
					\multicolumn{1}{c}{\textbf{Gráfico Asociado}} \\
					\midrule
					Lineal en dos variables & $ax + by + c = 0$ & Recta en el plano \\
					Lineal en tres variables & $ax + by + cz + d = 0$ & Plano en el espacio \\
					Cuadrática en dos variables & $ax^2 + by^2 + cxy + dx + ey + f = 0$ & Cónica: circunferencia, parábola, elipse o hipérbola \\
					Cuadrática en tres variables & $ax^2 + by^2 + cz^2 + dxy + exz + fyz + gx + hy + iz + j = 0$ & Superficie cuadrática: esfera, elipsoide, hiperboloide, paraboloide, etc. \\
					Cúbica en dos variables & $y^2 = x^3 + ax + b$ & Curva elíptica \\
					\bottomrule
				\end{tabular}
		}}
		\caption{Ecuaciones polinomiales y sus gráficos asociados.}
	\end{table}
	
	
	

\end{itemize}

Estos ejemplos reflejan cómo los polinomios constituyen una herramienta fundamental en la comprensión y el modelado de fenómenos tanto en matemáticas puras como aplicadas.









\section{Cuerpos y Clausura algebraica}

Dado un anillo \( (R,+,\cdot) \) con unidad.

\begin{definicion}
	\upshape{
		Un elemento \( a \in R \) es \textit{inversible} si es inversible con respecto a \( (R,\cdot) \), es decir, existe \( b \in R \) tal que \( ab = 1_R = ba \).
	}
\end{definicion}

\begin{observacion}
	\upshape{
		El neutro aditivo del anillo, \( 0 \in R \), no es inversible. En efecto, si existiera \( b \in R \) tal que \( 0 \cdot b = 1_R \), se tendría \( 0 = 1_R \), lo cual es absurdo, pues salvo en el anillo trivial, donde \( 1_R = 0 \), en general estos elementos son distintos.
	}
\end{observacion}
 
 
 El conjunto de todos los elementos invertibles del anillo \( R \) se denota por \( R^\times \), es decir,
 \[
 R^\times := \{ \, a \in R \mid a \text{ es invertible} \, \}.
 \]
 Este conjunto se denomina \textit{grupo multiplicativo} del anillo \( R \), o también \textit{grupo de los elementos invertibles} de \( R \). La razón es la siguiente proposición.
 
 \begin{proposicion}
 	\upshape{
 		\(\left( R^\times, \cdot \right)\) es un grupo.
 	}
 \end{proposicion}
 
\begin{proof}
	En efecto,
	\begin{itemize}
		\item Operación interna: Si \( a, b \in R^\times \), existen \( a^{-1}, b^{-1} \in R \) tales que \( a a^{-1} = b b^{-1} = 1_R \). Además, \( ab \) es invertible con inverso \( b^{-1} a^{-1} \), ya que:
		\[
		(ab)(b^{-1} a^{-1}) = a (b b^{-1}) a^{-1} = a (1_R) a^{-1} = a a^{-1} = 1_R,
		\]
		y de modo similar, \( (b^{-1} a^{-1})(ab) = 1_R \). Por tanto, \( ab \in R^\times \).
		
		\item Asociatividad: Heredada de la operación en \( R \), pues \( \cdot \) es asociativa en \( R \).
		
		\item Elemento identidad: La unidad \( 1_R \) pertenece a \( R^\times \) y actúa como identidad.
		
		\item Elemento inverso: Por definición, todo elemento de \( R^\times \) es invertible en \( R \).
	\end{itemize}
	Por tanto, \( \left( R^\times, \cdot \right) \) es un grupo.
\end{proof}

\begin{ejemplo}
	\upshape{
		El grupo multiplicativo de \( \mathbb{Z} \) es
		\[
		\mathbb{Z}^{\times} = \{ -1,\ +1 \}.
		\]
	}
\end{ejemplo}

\begin{ejemplo}
	\upshape{
		Sea \( \mathbb{Z}/n\mathbb{Z} \) el anillo de clases módulo \( n \), es decir, \( \mathbb{Z}/n\mathbb{Z} = \{ [m] : m \in \mathbb{Z} \} \). Entonces, su grupo multiplicativo es
		\[
		\left( \dfrac{\mathbb{Z}}{n\mathbb{Z}} \right)^{\times} = \{ [m] \in \mathbb{Z}/n\mathbb{Z} : \, \mathrm{mcd}(m,n) = 1 \}.
		\]
	}
\end{ejemplo}
\begin{proof}[Justificación]
	Un elemento \( [m] \in \mathbb{Z}/n\mathbb{Z} \) es invertible si y solo si existe \( [s] \in \mathbb{Z}/n\mathbb{Z} \) tal que
	\[
	[m][s] = [1].
	\]
	Esto equivale a decir que \( ms \equiv 1 \pmod{n} \), lo cual sucede si y solo si \( m \) y \( n \) son coprimos, es decir, \( \mathrm{mcd}(m,n) = 1 \).
	
	\medskip
	\noindent Por tanto, los elementos invertibles de \( \mathbb{Z}/n\mathbb{Z} \) son exactamente aquellos \( [m] \) tales que \( \mathrm{mcd}(m,n) = 1 \), como se indicó.
\end{proof}





\begin{definicion}
	\upshape{
		Sea \( R \) un anillo. Un elemento \( a \in R \) se dice \textit{nilpotente} si existe un entero positivo \( k \) tal que
		\[
		a^k = 0.
		\]
	}
\end{definicion}

\begin{observacion}\label{ñlkjhgjklasx}
	\upshape{
		Todo elemento nilpotente no nulo de un anillo es un divisor de cero.
	}
\end{observacion}
\begin{proof}[Justificación]
	\upshape{
	En efecto, sea \( a \in R \) un elemento nilpotente, es decir, existe un entero positivo \( k \) tal que \( a^k = 0 \). Si además \( a \neq 0 \), tomamos el menor de esos enteros \( k \), el cual necesariamente cumple \( k \geq 2 \). Entonces:
	\[
	a^{k-1} \neq 0 \quad \text{y} \quad a \cdot a^{k-1} = a^k = 0,
	\]
	por lo que \( a \) es divisor de cero.
	}
\end{proof}

\medskip
\begin{observacion}
	\upshape{
		Si \( R \) es un dominio, entonces no existen elementos nilpotentes distintos de \( 0_R \). Pues si existiera, por la Observación \ref{ñlkjhgjklasx}, \( R \) tendría divisores de cero.
	}
\end{observacion}

\medskip

\begin{proposicion}\label{lñkjhgsc}
	\upshape{
		En un anillo conmutativo, la suma de elementos nilpotentes es nilpotente.
	}
\end{proposicion}

\begin{proof}
	Sean \( a, b \in R \) nilpotentes, es decir, existen enteros positivos \( m, n \) tales que \( a^m = 0 \) y \( b^n = 0 \). Consideramos el binomio \( a + b \) y calculamos:
	\[
	(a + b)^{m+n} = \sum_{k=0}^{m+n} \binom{m+n}{k} a^k b^{m+n - k}.
	\]
	En cada término, si \( k \geq m \), entonces \( a^k = 0 \); y si \( k < m \), entonces \( m+n - k > n \), por lo que \( b^{m+n - k} = 0 \).
	
	\medskip
	\noindent Por lo tanto, todos los términos de la expansión son nulos y se tiene:
	\[
	(a + b)^{m+n} = 0.
	\]
	Por inducción se extiende a sumas finitas de nilpotentes.
\end{proof}

\medskip
\begin{observacion}
	\upshape{
		Si \( R \) no es conmutativo, la Proposición \ref{lñkjhgsc} es falsa. En efecto, consideremos el anillo \( M_{2\times 2}(\mathbb{\mathbb{R}}) \). Definimos:
		\[
		A = \begin{pmatrix} 0 & 1 \\ 0 & 0 \end{pmatrix}, \quad B = \begin{pmatrix} 0 & 0 \\ 1 & 0 \end{pmatrix}.
		\]
		Ambas matrices son nilpotentes, ya que:
		\[
		A^2 = B^2 = \begin{pmatrix} 0 & 0 \\ 0 & 0 \end{pmatrix}.
		\]
		Sin embargo, su suma es:
		\[
		A + B = \begin{pmatrix} 0 & 1 \\ 1 & 0 \end{pmatrix},
		\]
		y se verifica que:
		\[
		(A + B)^2 = \begin{pmatrix} 1 & 0 \\ 0 & 1 \end{pmatrix} = I_2,
		\]
		la matriz identidad. Como \( (A + B)^2 = I_2 \neq 0 \), concluimos que \( A + B \) no es nilpotente. Por lo tanto, en anillos no conmutativos, la suma de nilpotentes puede no ser nilpotente.
	}
\end{observacion}

\vspace{0.5cm}
\begin{proposicion}\label{lkjhgjnm}
	\upshape{
		Sea \( R \) un anillo conmutativo. Si \( a \in R \) es inversible y \( b \in R \) es nilpotente, entonces \( a + b \) es inversible.
	}
\end{proposicion}

\begin{proof}
	Como \( a \) es inversible, existe \( a^{-1} \in R \) tal que \( a a^{-1} = 1 \). Podemos escribir
	\[
	a + b = a \left( 1 + a^{-1} b \right).
	\]
	El elemento \( a^{-1} b \) es nilpotente, ya que \( b \) lo es y \( R \) es conmutativo. Sea \( k \) el menor entero positivo tal que \( (a^{-1} b)^{k} = 0 \). Entonces,  el inverso de \( 1 + a^{-1} b \) es 
	\[
	1 - a^{-1} b + (a^{-1} b)^2 - \cdots + (-1)^{k - 1} (a^{-1} b)^{k - 1}.
	\]
	En efecto, al multiplicar:
	\[
	\left( 1 + a^{-1} b \right) \left( 1 - a^{-1} b + (a^{-1} b)^2 - \cdots + (-1)^{k - 1} (a^{-1} b)^{k - 1} \right) = 1,
	\]
	
	\noindent Por lo tanto, \( 1 + a^{-1} b \) es inversible, y como \( a \) también lo es, concluimos que \( a + b \) es inversible.
\end{proof}



\begin{observacion}
	\upshape{
		Para hallar el inverso de \( 1 + a^{-1} b \) en la Proposición \ref{lkjhgjnm}, se utilizó la expresión de la serie geométrica finita. Recordemos primero la forma general:
		
		\[
		1 + r + r^2 + \cdots + r^k = \frac{1 - r^{k + 1}}{1 - r},
		\]
		Que para el caso de anillos se tiene que exigir que \( 1 - r \) sea inversible. Entonces la expresión válida es 		
		\[
		1 + r + r^2 + \cdots + r^k = \left( 1 - r^{k + 1} \right)(1 - r)^{-1}.
		\]
		
		\medskip
		
		\noindent En nuestro caso particular, tomando \( r = - a^{-1} b \), se deduce que:
		
		\[
		1 + (- a^{-1} b) + (- a^{-1} b)^2 + \cdots + (- a^{-1} b)^{k-1} = \left( 1 - (- a^{-1} b)^{k } \right) (1 + a^{-1} b)^{-1}.
		\] Reescribiendo  y como \( (a^{-1} b)^{k} = 0 \) tenemos 
		\[
	(1 + a^{-1} b)\left( 1 - a^{-1} b + (a^{-1} b)^2 - \cdots + (-1)^{k - 1} (a^{-1} b)^{k - 1} \right) =  1 - (-1)^k( a^{-1} b)^{k }.
		\]
		
	}
\end{observacion}



















\begin{proposicion}
	\upshape{
		Sea \( R \) un anillo conmutativo con identidad. Son equivalentes
	
	\begin{enumerate}
		\item Un polinomio \( p(x) = a_nx^n + \cdots + a_1x + a_0 \in R[x] \) es inversible.
		\item Para $p(x)$ se cumplen 
		\begin{itemize}
			\item \( a_0 \) es un elemento inversible de \( R \),
			\item Los coeficientes \( a_1, \dots, a_n \) son elementos nilpotentes de \( R \).
		\end{itemize}
	\end{enumerate}
		
	}
\end{proposicion}

\begin{proof}
	\text{white}
	
	\medskip
	\noindent $[ 1 \Rightarrow 2]$
	Supongamos que \( p(x) \) es inversible en \( R[x] \). Sea \( q(x) = b_m x^m + \cdots + b_0 \in R[x] \) tal que:
	\[
	p(x) q(x) = 1.
	\]
	Examinando el término constante:
	\[
	a_0 b_0 = 1,
	\]
	por lo que \( a_0 \) es inversible en \( R \)\\
	
	\noindent Para los coeficientes \( a_1, \dots, a_n \), se procede por inducción descendente.
	
	\medskip
	
	\noindent Caso base: Consideremos el término de mayor grado. Sea \( a_n \) el coeficiente principal de \( p(x) \) y \( b_m \) el de \( q(x) \). El coeficiente principal del producto \( p(x)q(x) \) es
	\[
	a_n b_m x^{n + m} = 0,
	\]
	ya que el producto total es 1, que es de grado 0. Por lo tanto:
	\[
	a_n b_m = 0.
	\]
	Multiplicando ambos lados de la igualdad \( p(x)q(x) = 1 \) por \( a_n^m \), se obtiene:
	\[
	a_n^m p(x) q(x) = a_n^m.
	\]
	El término líder del lado izquierdo tiene grado estrictamente menor que \( m \), pues \( a_n b_m = 0 \). Al iterar este razonamiento se concluye que \( a_n \) es nilpotente.
	
	\medskip
	
	\noindent \textbf{Paso inductivo:} Suponiendo que \( a_n, a_{n-1}, \dots, a_{k+1} \) son nilpotentes, definimos:
	\[
	p_1(x) = p(x) - \sum_{i = k+1}^n a_i x^i.
	\]
	El polinomio \( p_1(x) \) sigue siendo inversible y su término de mayor grado es \( a_k x^k \). Repitiendo el argumento anterior, se concluye que \( a_k \) es nilpotente.
	
	Por inducción, todos los coeficientes \( a_1, \dots, a_n \) son nilpotentes.
	
	
	
	
	
	
	
	
	
	
	
	
	
	
	
	
	
	
	
	
	
	
	
	
	\medskip
	
	\noindent \textbf{($\impliedby$) Suficiencia:} \\
	Supongamos que \( a_0 \) es inversible en \( R \) y que \( a_1, \dots, a_n \) son nilpotentes. Escribimos:
	\[
	p(x) = a_0 + N(x),
	\]
	donde \( N(x) = a_1x + a_2x^2 + \cdots + a_nx^n \).
	
	Observemos lo siguiente:
	\begin{itemize}
		\item Cada término \( a_i x^i \) es nilpotente en \( R[x] \), ya que \( a_i \) lo es y \( R[x] \) es conmutativo.
		\item La suma \( N(x) \) es nilpotente en \( R[x] \), pues la suma finita de nilpotentes es nilpotente.
		\item En un anillo conmutativo, si \( u \) es inversible y \( n \) es nilpotente, entonces \( u + n \) es inversible. En este caso:
		\[
		p(x) = a_0 (1 + a_0^{-1} N(x)).
		\]
		Como \( N(x) \) es nilpotente, \( a_0^{-1} N(x) \) también lo es, y por tanto \( 1 + a_0^{-1} N(x) \) es inversible.
		\item El producto de inversibles es inversible, por lo que \( p(x) \) es inversible en \( R[x] \).
	\end{itemize}
	
	\medskip
	
	
\end{proof}









\begin{proposicion}
	\upshape{
		Sea \( R \) un anillo conmutativo con identidad. Son equivalentes:
		\begin{enumerate}
			\item El polinomio \( p(x) = a_nx^n + \cdots + a_1x + a_0 \in R[x] \) es inversible.
			\item Se cumplen:
			\begin{itemize}
				\item \( a_0 \) es un elemento inversible de \( R \).
				\item Los coeficientes \( a_1, \dots, a_n \) son elementos nilpotentes de \( R \).
			\end{itemize}
		\end{enumerate}
	}
\end{proposicion}

\begin{proof}
	Demostramos ambas implicaciones por separado.
	
	\medskip
	
	\noindent \textbf{[1 $\Rightarrow$ 2]} Supongamos que \( p(x) \) es inversible en \( R[x] \), es decir, existe \( q(x) \in R[x] \) tal que \( p(x)q(x) = 1 \).
	
	Del término constante del producto se obtiene \( a_0 b_0 = 1 \), de modo que \( a_0 \) es inversible en \( R \).
	
	A continuación, verificaremos sucesivamente que los coeficientes \( a_1, \dots, a_n \) son nilpotentes, comenzando por el de mayor grado.
	
	\medskip
	
\medskip

\noindent \textit{Coeficiente principal \( a_n \):} Sea \( q(x) = b_m x^m + \cdots + b_0 \in R[x] \) tal que \( p(x) q(x) = 1 \). Como el lado derecho tiene grado 0, y el producto de dos polinomios no nulos tiene grado al menos \( n + m \), el coeficiente de grado \( n + m \) en \( p(x) q(x) \) debe ser cero. Dicho coeficiente es \( a_n b_m \), por lo tanto:


\[
a_n b_m = 0.
\]



Multiplicamos ambos lados de la igualdad \( p(x) q(x) = 1 \) por \( a_n \):


\[
a_n p(x) q(x) = a_n.
\]


Expandiendo el lado izquierdo:


\[
(a_n p(x)) q(x) = a_n a_n x^n q(x) + \text{términos de menor grado}.
\]


El coeficiente de grado \( n + m \) de \( a_n p(x) q(x) \) es entonces \( a_n^2 b_m \), que también debe anularse debido a que $a_nb_m=0$. Siguiendo este proceso:


\[
a_n^2 p(x) q(x) = a_n^2, \quad a_n^3 p(x) q(x) = a_n^3, \quad \ldots
\]


Para cada \( r \in \mathbb{N} \), el coeficiente de grado \( n + m \) del producto \( a_n^r p(x) q(x) \) es \( a_n^{r} b_m \). Como ya vimos que \( a_n b_m = 0 \), se deduce que:


\[
a_n^r b_m = 0 \quad \text{para todo } r \geq 1.
\]



Como los términos dominantes del producto desaparecen, y el lado izquierdo sigue siendo igual a \( a_n^r \), debe ocurrir que, en algún momento, \( a_n^r = 0 \). Por tanto, \( a_n \) es nilpotente.


	
	\medskip
	
	\noindent \textit{Coeficientes restantes:} Ahora restamos el término \( a_n x^n \) de \( p(x) \), y consideramos el polinomio \( p_1(x) = p(x) - a_n x^n \), que también es inversible (ya que restamos un nilpotente). Su coeficiente principal ahora es \( a_{n-1} \). Repetimos el mismo argumento con \( p_1(x) \) en lugar de \( p(x) \), concluyendo que \( a_{n-1} \) es nilpotente.
	
	Repitiendo este procedimiento con los coeficientes restantes, concluimos que \( a_1, \dots, a_n \) son nilpotentes.
	
	\bigskip
	
	\noindent \textbf{[2 $\Rightarrow$ 1]} Supongamos ahora que \( a_0 \in R^\times \) y que \( a_1, \dots, a_n \) son nilpotentes. Definimos:
	
	
	\[
	N(x) := a_1 x + a_2 x^2 + \cdots + a_n x^n,
	\]
	
	
	por lo que \( p(x) = a_0 + N(x) = a_0 \left( 1 + a_0^{-1} N(x) \right) \).
	
	Dado que cada \( a_i \) es nilpotente y \( R \) es conmutativo, \( N(x) \) es nilpotente en \( R[x] \). Por tanto, \( a_0^{-1} N(x) \) es nilpotente, y en consecuencia, \( 1 + a_0^{-1} N(x) \) es invertible, con inverso dado por una serie geométrica finita:
	
	
	\[
	\left( 1 + a_0^{-1} N(x) \right)^{-1} = \sum_{k=0}^{r-1} (-a_0^{-1} N(x))^k,
	\]
	
	
	donde \( r \in \mathbb{N} \) es tal que \( (a_0^{-1} N(x))^r = 0 \).
	
	Finalmente, como \( a_0 \) es invertible y el producto de elementos invertibles es invertible, se concluye que \( p(x) \) es invertible en \( R[x] \).
\end{proof}


































































\begin{ejemplo}
	\upshape{
		Sea \( R \) un anillo conmutativo con identidad. Entonces
		\[
		\left( R[x_1, x_2, \dots, x_n] \right)^{\times} = R^{\times}.
		\]
	}
\end{ejemplo}

\begin{proof}
	Sea \( f(x_1, \dots, x_n) \in R[x_1, \dots, x_n] \) tal que existe \( g(x_1, \dots, x_n) \in R[x_1, \dots, x_n] \) con
	\[
	f(x_1, \dots, x_n) \cdot g(x_1, \dots, x_n) = 1.
	\]
	Sea \( f = a + p \), donde \( a \in R \) es el término constante y \( p \) es un polinomio sin término constante (es decir, todos los términos de \( p \) tienen grado al menos 1).
	
	De igual forma, sea \( g = b + q \), con \( b \in R \) el término constante y \( q \) un polinomio sin término constante.
	
	Calculando:
	\[
	(a + p)(b + q) = ab + aq + bp + pq = 1.
	\]
	El término constante de este producto es \( ab \), y los demás términos tienen grado al menos 1. Por lo tanto:
	\[
	ab = 1.
	\]
	Esto implica que \( a \in R^\times \) y \( b = a^{-1} \).
	
	Además, como los demás términos tienen grado al menos 1 y su suma es 0, necesariamente:
	\[
	aq + bp + pq = 0.
	\]
	Por lo tanto, el polinomio \( f \) no tiene términos no constantes, es decir, \( p = 0 \).
	
	Por lo tanto:
	\[
	f = a, \quad g = b.
	\]
	Es decir, \( f \) es constante e invertible en \( R \).
	
	Recíprocamente, si \( a \in R^\times \), entonces \( f = a \) y su inverso es \( a^{-1} \), que también pertenece a \( R[x_1, \dots, x_n] \).
	
	Por lo tanto:
	\[
	\left( R[x_1, \dots, x_n] \right)^\times = R^\times.
	\]
\end{proof}



	%\chapter{Lista de ejercicios I}
	%\begin{ejercicio}
	\upshape{
	Diga si los siguientes enunciados son verdaderos (V) o falsos (F). Justifique su respuesta.
	\begin{enumerate}[label=\alph*)]
		\item El complemento de todo submonoide finitamente generado de $(\mathbb{N},+)$ es finito.
		\item Dados dos números naturales $m,n \geq 2$, la presentación del grupo $\mathbb{Z}/m\mathbb{Z} \times \mathbb{Z}/n\mathbb{Z}$ es
		\[
		\langle x, y : x^m = y^n = e \rangle.
		\]
	\end{enumerate}
	}
\end{ejercicio}
\begin{proof}[Sulución]
	{\color{white} Texto en azul.}\\
	
	\noindent $a)$ Es falso: En el monoide $\left(\mathbb{N},+\right)$ considere el sobmonoide $M=\{0;2;4;4;6;\cdots\}$. Veamos que $M$ es finitamente generado. En efecto, $$\langle 2\rangle=\left\{n2 : n\in\mathbb{N}\right\}=\left\{2n:n\in\mathbb{N}\right\}=M$$ Ahora notemos que $\mathbb{N}\setminus M=\{1;2;3; \cdot\}$ es un conjunto infinito.\\
	
	\noindent $b)$  Consideremos la función 
	\[
	f : \{x,y\} \to \mathbb{Z}/m\mathbb{Z} \times \mathbb{Z}/n\mathbb{Z}
	\]
	definida por
	\[
	x \mapsto (\overline{1}, \overline{0}), \quad y \mapsto (\overline{0}, \overline{1}).
	\] Por la propiedad universal del grupo libre, existe un único homomorfismo
	\[
	\varphi : F\{x,y\} \to \mathbb{Z}/m\mathbb{Z} \times \mathbb{Z}/n\mathbb{Z}
	\]
	tal que el siguiente diagrama conmuta:
	\[
	{\large \xymatrix{
			X \ar[ddrr]_{f} \ar[rr]^{\iota}& & F\{x,y\} \ar[dd]^{\varphi} \\
			&&\\
			&& \mathbb{Z}/m\mathbb{Z} \times \mathbb{Z}/n\mathbb{Z}
	}}
	\]
	donde \(\iota : \{x,y\} \to F\{x,y\}\) es la inclusión canónica.\\
	
	\noindent Observemos que
	\[
	\varphi(x^m) = f(x)^m = (\overline{1}, \overline{0})^m = (\overline{0}, \overline{0}),
	\]
	\[
	\varphi(y^n) = f(y)^n = (\overline{0}, \overline{1})^n = (\overline{0}, \overline{0}),
	\]
	y además, dado que \(\mathbb{Z}/m\mathbb{Z} \times \mathbb{Z}/n\mathbb{Z}\) es abeliano,
	\[
	\varphi(x y y^{-1} x^{-1}) = f(x) f(y) f(y)^{-1} f(x)^{-1} = (\overline{0}, \overline{0}).
	\] Luego, el núcleo de \(\varphi\) contiene el subgrupo normal generado por los elementos
	\[
	x^m, \quad y^n, \quad xyx^{-1}y^{-1}.
	\]
	
	Por el teorema fundamental de homomorfismos, se tiene el isomorfismo
	\[
	F\{x,y\} / \langle\langle x^m, y^n, x y y^{-1} x^{-1} \rangle\rangle \cong \mathbb{Z}/m\mathbb{Z} \times \mathbb{Z}/n\mathbb{Z}.
	\]
	
	En consecuencia, una presentación del grupo \(\mathbb{Z}/m\mathbb{Z} \times \mathbb{Z}/n\mathbb{Z}\) es
	\[
	\langle x, y : x^m = e, \quad y^n = e, \quad x y = y x \rangle.
	\]
	
\end{proof}

\begin{observacion}
	\upshape{
	Sea \( X \) un conjunto. Entonces existe un grupo \( F(X) \) y una función de conjuntos
	\[
	\iota : X \to F(X)
	\]
	tal que para todo grupo \( G \) y toda función de conjuntos
	\[
	f : X \to G,
	\]
	existe un único homomorfismo de grupos
	\[
	\varphi : F(X) \to G
	\]
	tal que el siguiente diagrama conmuta:
	\[
	{\large \xymatrix{
		X \ar[dr]_{f} \ar[r]^{\iota} & G \\
		 & F(X) \ar@{-->}[u]_{\varphi}
	}}
	\]
	Es decir, se cumple:
	\[
	\varphi \circ \iota = f.
	\]
	}
\end{observacion}






\begin{observacion}
	\upshape{
		Sea \( X \) un conjunto. Entonces existe un grupo \( F(X) \) y una función de conjuntos
		\[
		\iota : X \to F(X)
		\]
		tal que para todo grupo \( G \) y toda función de conjuntos
		\[
		f : X \to G,
		\]
		existe un único homomorfismo de grupos
		\[
		\varphi : F(X) \to G
		\]
		tal que el siguiente diagrama conmuta:
		\[
		{\large \xymatrix{
				X \ar[dr]_{f} \ar[r]^{\iota} & G \\
				& F(X) \ar@{-->}[u]_{\varphi}
		}}
		\]
		Es decir, se cumple:
		\[
		\varphi \circ \iota = f.
		\]
		
		\medskip
		
		\noindent
		\textbf{Recuerde:} El grupo \( F(X) \) es el \emph{grupo libre generado por el conjunto \( X \)}, formado por palabras reducidas en los símbolos de \( X \) y sus inversos, con la operación de concatenación seguida de reducción. La función \(\iota\) es la inclusión natural que asigna a cada \( x \in X \) el generador correspondiente en \( F(X) \). Esta inclusión es solo una función de conjuntos, no un homomorfismo, pues \( X \) no tiene estructura de grupo.
	}
\end{observacion}


\begin{observacion}[Teorema Fundamental del Homomorfismo de grupos]
	\upshape{
		Sea \( f : G \to H \) un homomorfismo de grupos. Entonces el kernel
		\[
		\ker(f) = \{ g \in G : f(g) = e_H \}
		\]
		es un subgrupo normal de \( G \), y \( f \) induce un isomorfismo de grupos
		\[
		\overline{f} : G / \ker(f) \xrightarrow{\ \cong \ } \mathrm{Im}(f),
		\]
		es decir,
		\[
		G / \ker(f) \cong \mathrm{Im}(f).
		\]
	}
\end{observacion}

\begin{ejercicio}
	Dados dos conjuntos $X$ e $Y$. Probar que los grupos libres $F_X$ y $F_Y$ son isomorfos si y solo si $X$ e $Y$ son biyectivos.
\end{ejercicio}
\begin{proof}
	{\color{white} Texto en azul.}\\
	
	\noindent $[\Rightarrow]$
	
	\noindent   $[\Rightarrow]$
\end{proof}
\begin{ejercicio}
	\upshape{
	Dados dos conjuntos \(X\) e \(Y\). Probar que los grupos libres \(F_X\) y \(F_Y\) son isomorfos si y solo si \(X\) e \(Y\) son biyectivos.
	}
\end{ejercicio}




\begin{proof}
	\noindent\emph{($\Rightarrow$)}  Supongamos que \(F_X \cong F_Y\). Como los grupos libres están determinados por la cardinalidad de sus conjuntos generadores libres, esto implica que \(X\) e \(Y\) tienen la misma cardinalidad. Por tanto, existe una biyección entre \(X\) e \(Y\).
	
	\vspace{1ex}
	\noindent\emph{($\Leftarrow$)} Supongamos ahora que existe una biyección de conjuntos
	\[
	f: X \to Y.
	\]
	Entonces podemos considerar la función
	\[
	\iota_X : X \longrightarrow F_X
	\quad \text{y} \quad
	\iota_Y : Y \longrightarrow F_Y
	\]
	que son las inclusiones canónicas. Componiendo con \(f\), obtenemos una función
	\[
	f \circ \iota_X : X \to F_Y.
	\]
	Por la propiedad universal del grupo libre \(F_X\), existe un único homomorfismo de grupos
	\[
	\varphi : F_X \to F_Y
	\]
	tal que \(\varphi \circ \iota_X = \iota_Y \circ f\). Análogamente, como \(f\) es biyectiva, su inversa \(f^{-1} : Y \to X\) induce, por la misma propiedad universal, un único homomorfismo de grupos
	\[
	\psi : F_Y \to F_X
	\]
	tal que \(\psi \circ \iota_Y = \iota_X \circ f^{-1}\). Entonces se verifica que
	\[
	\psi \circ \varphi \circ \iota_X = \psi \circ \iota_Y \circ f = \iota_X \circ f^{-1} \circ f = \iota_X,
	\]
	y como \(\iota_X\) genera \(F_X\), se deduce que \(\psi \circ \varphi = \operatorname{id}_{F_X}\). De forma similar, se verifica que \(\varphi \circ \psi = \operatorname{id}_{F_Y}\). Por tanto, \(\varphi\) es un isomorfismo de grupos.
	
	\vspace{1ex}
	\noindent En conclusión, \(F_X \cong F_Y\) si y solo si \(X\) e \(Y\) son biyectivos.
\end{proof}


\begin{ejercicio}
	\upshape{
		Probar que todo subgrupo de un grupo libre es también libre.
	}
\end{ejercicio}

\begin{proof}[Solución]
Esta afirmación se conoce como el \emph{teorema de Nielsen-Schreier}. 
\end{proof}

\begin{teorema}[Teorema de Nielsen-Schreier]
	\upshape{
	Todo subgrupo de un grupo libre es también libre.
	}
\end{teorema}
\begin{proof}
	Sea \( F = F_X \) el grupo libre generado por un conjunto \( X \). Por definición, los elementos de \( F \) son palabras reducidas formadas por elementos de \( X \) y sus inversos, y la operación es concatenación seguida de reducción.
	
	Sea \( H \leq F \) un subgrupo de \( F \).
	
	\begin{enumerate}
		\item \textbf{Descomposición en clases laterales:}
		
		El grupo \( F \) se descompone en clases laterales derechas respecto a \( H \):
		\[
		F = \bigsqcup_{t \in T} H t,
		\]
		donde \( T \subseteq F \) es un conjunto de representantes de las clases laterales, conteniendo exactamente un representante por clase lateral derecha.
		
		\item \textbf{Definición de generadores de \( H \):}
		
		Para cada \( t \in T \) y cada generador \( x \in X \), el producto \( t x \) pertenece a alguna clase lateral \( H s \) para cierto \( s \in T \). Esto implica que existe un elemento
		\[
		h_{t,x} := t x s^{-1} \in H.
		\]
		
		Definimos el conjunto
		\[
		S := \{ h_{t,x} = t x s^{-1} \mid t, s \in T, x \in X, \text{ y } t x \in H s \}.
		\]
		
		\item \textbf{Generación de \( H \) por \( S \):}
		
		Sea \( h \in H \). Entonces \( h \) puede escribirse como una palabra reducida en los generadores \( X \):
		\[
		h = x_1^{\epsilon_1} x_2^{\epsilon_2} \cdots x_n^{\epsilon_n},
		\]
		donde \( x_i \in X \) y \( \epsilon_i = \pm 1 \).
		
		Definimos una sucesión \( (t_0, t_1, \ldots, t_n) \) en \( T \) con \( t_0 = e \) (la identidad) tal que
		\[
		t_{i-1} x_i^{\epsilon_i} \in H t_i,
		\]
		para cada \( i = 1, \ldots, n \).
		
		De la definición de \( h_{t,x} \), para cada \( i \) existe \( h_i \in S \cup S^{-1} \) tal que
		\[
		t_{i-1} x_i^{\epsilon_i} = h_i t_i.
		\]
		
		Multiplicando en orden,
		\[
		h = h_1 h_2 \cdots h_n,
		\]
		con \( h_i \in \langle S \rangle \).
		
		Así, \( h \in \langle S \rangle \) y por tanto
		\[
		H = \langle S \rangle.
		\]
		
		\item \textbf{Conclusión:}
		
		Como \( F \) es libre y la construcción de \( S \) no introduce relaciones adicionales, se concluye que \( H \) es libre con conjunto de generadores \( S \).
		
	\end{enumerate}
	
	Por lo tanto, todo subgrupo de un grupo libre es libre.
\end{proof}
\begin{observacion}
	\upshape{
	Sea \( G \) un grupo y \( H \leq G \) un subgrupo de \( G \). Para cada elemento \( g \in G \), definimos la \emph{clase lateral derecha} de \( H \) en \( G \) determinada por \( g \) como el subconjunto
	\[
	H g = \{ h g : h \in H \}.
	\]
	
	Análogamente, la \emph{clase lateral izquierda} de \( H \) determinada por \( g \) es
	\[
	g H = \{ g h : h \in H \}.
	\]
	
	Estas clases laterales son subconjuntos de \( G \) que particionan \( G \), es decir,
	\[
	G = \bigsqcup_{t \in T} H t,
	\]
	donde \( T \subseteq G \) es un conjunto de representantes, es decir, cada elemento de \( G \) pertenece exactamente a una clase lateral derecha.
	
	Las clases laterales sirven para estudiar la estructura del grupo \( G \) en relación con el subgrupo \( H \), y son fundamentales en resultados como el Teorema de Lagrange.
	}
\end{observacion}



\begin{ejercicio}
	\upshape{
	Sea \( m \in \mathbb{N} \) con \( m \geq 1 \). Entonces
	\[
	(\mathbb{Z}/2^m\mathbb{Z})^\times = \{ [a] \in \mathbb{Z}/2^m\mathbb{Z} : a \text{ es impar} \}.
	\]
	Además, si \( m \geq 3 \), entonces
	\[
	(\mathbb{Z}/2^m\mathbb{Z})^\times = \langle [-1], [5] \rangle \cong \mathbb{Z}/2\mathbb{Z} \times \mathbb{Z}/2^{m-2}\mathbb{Z}.
	\]
	}
\end{ejercicio}
\begin{proof}
		\leavevmode
		\begin{enumerate}
			\item \textbf{Descripción del grupo multiplicativo:} 
			
			Sea \( a \in \mathbb{Z} \). El elemento \( [a] \in \mathbb{Z}/2^m\mathbb{Z} \) es invertible si y solo si \( \gcd(a, 2^m) = 1 \). Esto ocurre si y solo si \( a \) es impar, ya que \( 2^m \) es una potencia de \( 2 \). Por tanto:
			\[
			(\mathbb{Z}/2^m\mathbb{Z})^\times = \{ [a] : a \text{ es impar} \}.
			\]
			
			\item \textbf{Cardinalidad:}
			
			El número de enteros entre \( 1 \) y \( 2^m \) que son primos con \( 2^m \) es \( \varphi(2^m) = 2^m - 2^{m-1} = 2^{m-1} \), pues la mitad de los números entre \( 1 \) y \( 2^m \) son impares. Por lo tanto, el grupo tiene orden \( 2^{m-1} \).
			
			\item \textbf{Estructura del grupo multiplicativo para \( m \geq 3 \):}
			
			Se sabe que para \( m \geq 3 \), el grupo \( (\mathbb{Z}/2^m\mathbb{Z})^\times \) no es cíclico, pero sí es isomorfo a
			\[
			\mathbb{Z}/2\mathbb{Z} \times \mathbb{Z}/2^{m-2}\mathbb{Z}.
			\]
			Esto se justifica constructivamente:
			
			\begin{itemize}
				\item El elemento \( [-1] \) tiene orden \( 2 \), ya que:
				\[
				[-1]^2 = [1].
				\]
				
				\item Sea \( m \geq 3 \) y consideremos \( [5] \in \mathbb{Z}/2^m\mathbb{Z} \). Se puede demostrar que \( [5] \) tiene orden \( 2^{m-2} \). Por ejemplo:
				\[
				5^1 = 5,\quad 5^2 = 25,\quad 5^3 = 125,\quad \dots
				\]
				En general, se puede mostrar que:
				\[
				5^{2^{m-2}} \equiv 1 \pmod{2^m},\quad \text{pero no antes.}
				\]
				Por lo tanto, el orden de \( [5] \) es \( 2^{m-2} \).
				
				\item Los elementos \( [-1] \) y \( [5] \) generan un subgrupo de orden:
				\[
				\operatorname{mcm}(2, 2^{m-2}) = 2 \cdot 2^{m-2} = 2^{m-1},
				\]
				que es el orden de todo el grupo. Así, se concluye que
				\[
				(\mathbb{Z}/2^m\mathbb{Z})^\times = \langle [-1], [5] \rangle.
				\]
				Y por el Teorema de estructura de grupos abelianos finitos, se deduce que:
				\[
				(\mathbb{Z}/2^m\mathbb{Z})^\times \cong \mathbb{Z}/2\mathbb{Z} \times \mathbb{Z}/2^{m-2}\mathbb{Z}.
				\]
			\end{itemize}
		\end{enumerate}
\end{proof}


\begin{observacion}
	\upshape{
	El grupo multiplicativo \( (\mathbb{Z}/2^m\mathbb{Z})^\times \) es un grupo abeliano finito de orden \( \varphi(2^m) = 2^{m-1} \), pues contiene exactamente las clases de los enteros impares módulo \( 2^m \).
	
	Por el \textbf{teorema de estructura de grupos abelianos finitos}, todo grupo abeliano finito de orden \( n \) es isomorfo a un producto directo de grupos cíclicos de la forma
	\[
	\mathbb{Z}/n_1\mathbb{Z} \times \mathbb{Z}/n_2\mathbb{Z} \times \cdots \times \mathbb{Z}/n_r\mathbb{Z},
	\]
	con \( n_1 \mid n_2 \mid \cdots \mid n_r \). Esta descomposición es única salvo isomorfismo.
	
	Aplicando esto a \( (\mathbb{Z}/2^m\mathbb{Z})^\times \), se obtiene que para \( m \geq 3 \), este grupo es isomorfo a
	\[
	\mathbb{Z}/2\mathbb{Z} \times \mathbb{Z}/2^{m-2}\mathbb{Z},
	\]
	ya que:
	
	\begin{itemize}
		\item La clase \( [-1] \) tiene orden 2, y
		\item La clase \( [5] \) tiene orden \( 2^{m-2} \).
	\end{itemize}
	
	Por lo tanto,
	\[
	(\mathbb{Z}/2^m\mathbb{Z})^\times = \langle [-1], [5] \rangle \cong \mathbb{Z}/2\mathbb{Z} \times \mathbb{Z}/2^{m-2}\mathbb{Z}.
	\]
	}
\end{observacion}


\begin{ejemplo}
	\upshape{
	Sea $C^\infty(\mathbb{R})$ el anillo de funciones reales de clase $C^\infty$, es decir, funciones \( f : \mathbb{R} \to \mathbb{R} \) derivables infinitamente. Consideremos:
	
	\begin{enumerate}
		\item \textbf{¿Es \( C^\infty(\mathbb{R}) \) un dominio de integridad?}
		
		\begin{proof}
			Sea \( f, g \in C^\infty(\mathbb{R}) \) tales que
			\[
			fg = 0,
			\]
			es decir,
			\[
			(fg)(x) = f(x) g(x) = 0 \quad \text{para todo } x \in \mathbb{R}.
			\]
			
			Queremos probar que esto implica que \( f = 0 \) o \( g = 0 \) (la función nula).
			
			Supongamos, sin pérdida de generalidad, que \( f \neq 0 \). Entonces existe algún punto \( a \in \mathbb{R} \) tal que \( f(a) \neq 0 \).
			
			Como \( f \) es continua, existe un intervalo abierto \( I \) alrededor de \( a \) tal que para todo \( x \in I \), \( f(x) \neq 0 \).
			
			Pero en \( I \), debido a que \( f(x) g(x) = 0 \), necesariamente debe cumplirse que
			\[
			g(x) = 0 \quad \text{para todo } x \in I.
			\]
			
			Dado que \( g \in C^\infty(\mathbb{R}) \) y es nula en un intervalo abierto \( I \), por unicidad de funciones analíticas (o usando que las funciones suaves que se anulan en un intervalo abierto son idénticamente cero en ese intervalo y por extensión nula en todo \(\mathbb{R}\)) concluimos que
			\[
			g(x) = 0 \quad \text{para todo } x \in \mathbb{R}.
			\]
			
			Por tanto, \( g = 0 \).
			
			Esto demuestra que si \( fg = 0 \), entonces \( f=0 \) o \( g=0 \).
			
			Así, \( C^\infty(\mathbb{R}) \) no tiene divisores de cero y es un dominio de integridad.
		\end{proof}
		\begin{observacion}
			\upshape{
			Dado que \( g \in C^\infty(\mathbb{R}) \) y es nula en un intervalo abierto \( I \subseteq \mathbb{R} \), es decir,
			\[
			g(x) = 0 \quad \text{para todo } x \in I,
			\]
			podemos concluir que \( g \) es la función nula en todo \(\mathbb{R}\).
			
			Esto se debe a que las funciones de clase \( C^\infty \) tienen una propiedad de *unicidad de extensión analítica local*: si una función suave se anula en un intervalo abierto, entonces todas sus derivadas en cualquier punto de ese intervalo también se anulan.
			
			En efecto, para cada \( a \in I \), como \( g \equiv 0 \) en un entorno de \( a \), se cumple que
			\[
			g^{(n)}(a) = 0 \quad \text{para todo } n \in \mathbb{N}_0,
			\]
			donde \( g^{(n)} \) denota la \( n \)-ésima derivada de \( g \).
			
			Esto implica que el desarrollo en serie de Taylor de \( g \) alrededor de \( a \) es idénticamente cero:
			\[
			g(x) = \sum_{n=0}^\infty \frac{g^{(n)}(a)}{n!} (x - a)^n = 0.
			\]
			
			Como la función coincide con su serie de Taylor en un entorno de \( a \), \( g \) es nula en un entorno de cualquier punto \( a \in I \). Dado que \( I \) es un intervalo abierto, la función \( g \) es nula en una región abierta y conexa.
			
			Por continuidad e infinitas derivadas nulas en todo \( I \), \( g \) debe ser idénticamente cero en toda \(\mathbb{R}\), ya que una función suave con un conjunto abierto de ceros tiene que ser la función nula en todo su dominio (por la unicidad de la extensión analítica).
			
			Por tanto,
			\[
			g(x) = 0 \quad \text{para todo } x \in \mathbb{R}.
			\]
			
			Así, \( g \) es la función nula.
			}
		\end{observacion}
		
		
		\item \textbf{Sea \( R \subseteq C^\infty(\mathbb{R}) \) el subanillo generado por las funciones \( \sin \) y \( \cos \). Probar que todo elemento de \( R \) se puede escribir como}
		\[
		f(x) = a_0 + \sum_{m=1}^{n} (a_m \cos(mx) + b_m \sin(mx)), \quad a_0, a_m, b_m \in \mathbb{R}.
		\]
		
		\begin{proof}
			Sea \( R \subseteq C^\infty(\mathbb{R}) \) el subanillo generado por \( \sin(x) \) y \( \cos(x) \). Por definición, esto significa que \( R \) es el conjunto más pequeño que:
			
			\begin{itemize}
				\item contiene \( \sin(x) \) y \( \cos(x) \),
				\item es cerrado bajo suma y producto finitos,
				\item contiene constantes reales (ya que \( R \subseteq C^\infty(\mathbb{R}) \) es un subanillo).
			\end{itemize}
			
			Queremos probar que:
			\[
			\forall f \in R, \quad f(x) = a_0 + \sum_{m=1}^n (a_m \cos(mx) + b_m \sin(mx)),
			\]
			para ciertos \( a_0, a_1, \dots, a_n, b_1, \dots, b_n \in \mathbb{R} \) y algún \( n \in \mathbb{N} \).
			
			\medskip
			\textbf{Paso 1. Base: } probamos que \( \cos(x) \) y \( \sin(x) \) están en la forma requerida.
			
			Esto es inmediato:
			\[
			\cos(x) = 0 + 1 \cdot \cos(1x) + 0 \cdot \sin(1x), \quad
			\sin(x) = 0 + 0 \cdot \cos(1x) + 1 \cdot \sin(1x).
			\]
			Entonces están en la forma exigida.
			
			\medskip
			\textbf{Paso 2. Inducción:} probamos que si \( \cos(kx), \sin(kx) \in R \) para todo \( k \leq n \), entonces \( \cos((n+1)x), \sin((n+1)x) \in R \).
			
			Utilizamos las identidades trigonométricas para deducir las siguientes fórmulas de recurrencia:
			\[
			\cos((n+1)x) = 2 \cos(x)\cos(nx) - \cos((n-1)x),
			\]
			\[
			\sin((n+1)x) = 2 \cos(x)\sin(nx) - \sin((n-1)x).
			\]
			
			**Demostración de que esas fórmulas generan elementos en \( R \):**
			
			- Por hipótesis inductiva, \( \cos(nx) \), \( \cos((n-1)x) \in R \).
			- También sabemos que \( \cos(x) \in R \).
			- Como \( R \) es un subanillo, entonces:
			\[
			\cos(x)\cos(nx) \in R, \quad 2\cos(x)\cos(nx) \in R, \quad \cos((n+1)x) \in R.
			\]
			porque la suma y producto de elementos en \( R \) están en \( R \).
			
			- Lo mismo se aplica a:
			\[
			\sin((n+1)x) = 2\cos(x)\sin(nx) - \sin((n-1)x).
			\]
			porque:
			\[
			\sin(nx), \sin((n-1)x), \cos(x) \in R \quad \Rightarrow \quad \sin((n+1)x) \in R.
			\]
			
			Por lo tanto, por inducción matemática, tenemos que:
			\[
			\forall m \in \mathbb{N}, \quad \cos(mx), \sin(mx) \in R.
			\]
			
			\medskip
			\textbf{Paso 3. Combinaciones lineales:}
			
			Ahora, \( R \) es un anillo generado por \( \sin(x) \) y \( \cos(x) \), así que cualquier \( f \in R \) es suma finita de productos finitos de funciones de la forma \( \cos(kx), \sin(kx) \). Por identidades trigonométricas (producto a suma), sabemos que los productos como \( \cos(mx)\cos(nx) \), \( \sin(mx)\sin(nx) \), \( \sin(mx)\cos(nx) \) pueden escribirse como combinaciones lineales de \( \cos(kx) \) o \( \sin(kx) \) con \( k \in \mathbb{N} \).
			
			Por ejemplo:
			\[
			\cos(mx)\cos(nx) = \frac{1}{2}\left( \cos((m-n)x) + \cos((m+n)x) \right),
			\]
			\[
			\sin(mx)\cos(nx) = \frac{1}{2} \left( \sin((m+n)x) + \sin((m-n)x) \right),
			\]
			\[
			\sin(mx)\sin(nx) = \frac{1}{2}\left( \cos((m-n)x) - \cos((m+n)x) \right).
			\]
			
			Estas identidades muestran que todo producto de funciones trigonométricas de múltiplos de \( x \) puede ser reescrito como combinación lineal finita de \( \cos(kx), \sin(kx) \).
			
			\medskip
			\textbf{Conclusión:}
			
			Cualquier elemento \( f \in R \) se escribe como suma finita de productos finitos de \( \cos(kx), \sin(kx) \), y por lo anterior, estas se reducen a combinación lineal finita de funciones del tipo:
			\[
			\{ \cos(mx), \sin(mx) \mid m \in \mathbb{N} \} \cup \{1\}.
			\]
			Por tanto,
			\[
			f(x) = a_0 + \sum_{m=1}^{n} \bigl(a_m \cos(mx) + b_m \sin(mx)\bigr),
			\]
			para ciertos \( a_0, a_m, b_m \in \mathbb{R} \), y algún \( n \in \mathbb{N} \). Esto concluye la prueba.
		\end{proof}
		
		\item \textbf{Probar que \( R \) no tiene divisores de cero.}
		
			
		\begin{proof}
			Queremos probar que el subanillo \( R \subseteq C^\infty(\mathbb{R}) \), generado por \( \sin(x) \) y \( \cos(x) \), no tiene divisores de cero. Es decir, si \( f(x), g(x) \in R \) y 
			\[
			f(x) \cdot g(x) = 0 \quad \text{para todo } x \in \mathbb{R},
			\]
			entonces necesariamente \( f = 0 \) o \( g = 0 \).
			
			\begin{enumerate}
				\item \textbf{Forma de los elementos de \( R \):}  
				Por una prueba previa, sabemos que todo \( f \in R \) puede escribirse como
				\[
				f(x) = a_0 + \sum_{m=1}^n \bigl(a_m \cos(mx) + b_m \sin(mx)\bigr),
				\]
				donde \( a_0, a_m, b_m \in \mathbb{R} \), y la suma es finita. Es decir, los elementos de \( R \) son combinaciones lineales finitas de las funciones \( \sin(mx) \) y \( \cos(mx) \).
				
				\item \textbf{Funciones trigonométricas como funciones analíticas:}  
				Cada función \( \cos(mx) \), \( \sin(mx) \) es analítica en \( \mathbb{R} \), y por tanto también lo es cualquier combinación finita de ellas. Entonces, todo elemento de \( R \) es una función analítica real de \( \mathbb{R} \).
				
				\item \textbf{Propiedad de unicidad del producto nulo para funciones analíticas:}  
				Supongamos que \( f(x), g(x) \in R \) y que \( f(x)\cdot g(x) = 0 \) para todo \( x \in \mathbb{R} \). Como \( f \cdot g \) es analítica y se anula en todos los reales, resulta que su conjunto de ceros tiene un punto de acumulación (de hecho, todo \( \mathbb{R} \)). Por el principio de identidad de las funciones analíticas, se deduce que:
				\[
				f(x) = 0 \quad \text{o} \quad g(x) = 0 \quad \text{en todo } \mathbb{R}.
				\]
				
				\item \textbf{Conclusión:}  
				Por tanto, si \( f \cdot g = 0 \) en \( \mathbb{R} \), entonces \( f = 0 \) o \( g = 0 \) como funciones reales. Es decir, no hay divisores de cero en \( R \), por lo que:
				\[
				R \text{ es un dominio de integridad}.
				\]
			\end{enumerate}
		\end{proof}
		
		
		\begin{observacion}
			\upshape{
				No todo subanillo de un dominio de integridad es también un dominio de integridad. 
				
				Para ilustrar esto, consideremos el siguiente ejemplo rápido:
				
				Sea 
				\[
				M = \mathbb{Z}[X]
				\]
				el anillo de polinomios en una variable \(X\) con coeficientes enteros. Sabemos que \(M\) es un dominio de integridad porque \(\mathbb{Z}\) es un dominio de integridad y el anillo de polinomios sobre un dominio de integridad también es un dominio de integridad. Esto significa que en \(M\) no existen divisores de cero, es decir, si \(f, g \in M\) y \(f \neq 0, g \neq 0\), entonces \(f \cdot g \neq 0\).
				
				Ahora, definimos un subanillo de \(M\) de la siguiente manera:
				\[
				N = \{ a + bX : a, b \in \mathbb{Z} \text{ y } X^2 = 0 \}.
				\]
				Este anillo \(N\) es isomorfo al cociente
				\[
				\mathbb{Z}[X]/(X^2),
				\]
				donde se impone la relación \(X^2 = 0\).
				
				En el subanillo \(N\), el elemento \(X\) es distinto de cero pero satisface
				\[
				X \cdot X = X^2 = 0,
				\]
				lo cual muestra que \(X\) es un divisor de cero en \(N\).
				
				Por lo tanto, \(N\) no es un dominio de integridad, a pesar de ser un subanillo de \(M\), que sí lo es.
				
				Este ejemplo muestra que la propiedad de ser dominio de integridad no se hereda necesariamente a los subanillos.
			}
		\end{observacion}
		
		
		
		\item \textbf{Probar que \( \sin \) y \( 1 - \cos \) son irreducibles en \( R \). Deducir que \( R \) no es un dominio de factorización única.}
		
		\begin{proof}
			Recordemos que 
			\[
			R = \text{el subanillo de } C^\infty(\mathbb{R}) \text{ generado por } \sin \text{ y } \cos,
			\]
			y que todo elemento de \(R\) puede expresarse como combinación lineal finita de funciones \(\sin(mx)\) y \(\cos(mx)\).
			
			\begin{enumerate}
				\item \textbf{Irreducibilidad de \(\sin\):}
				
				Supongamos que \(\sin\) es reducible en \(R\), es decir, existen \(f, g \in R\), no invertibles en \(R\), tales que
				\[
				\sin = f \cdot g.
				\]
				
				Como \(R \subset C^\infty(\mathbb{R})\), esto implica que \(f\) y \(g\) son funciones suaves. 
				Además, dado que \(\sin\) tiene ceros infinitos (por ejemplo, en \(k\pi\), \(k \in \mathbb{Z}\)), si \(f\) y \(g\) son factores no invertibles, entonces cada uno debe tener al menos un cero.
				
				Sin embargo, por la expresión en series de Fourier, cualquier función en \(R\) tiene ceros aislados o finitos, y la estructura del anillo \(R\) no permite factores no triviales que reproduzcan la periodicidad y ceros infinitos de \(\sin\) sin ser invertibles.
				
				Más formalmente, si \(f\) o \(g\) tuvieran ceros distintos a los de \(\sin\), el producto tendría ceros distintos o multiplicidad mayor, lo que no es el caso.
				
				Por lo tanto, \(\sin\) no se puede factorizar no trivialmente en \(R\), es decir, es irreducible.
				
				\item \textbf{Irreducibilidad de \(1 - \cos\):}
				
				Análogamente, supongamos que
				\[
				1 - \cos = f \cdot g,
				\]
				con \(f, g \in R\) no invertibles.
				
				Notemos que \(1 - \cos(x)\) es una función que se anula en múltiplos enteros de \(2\pi\), y que es no negativa y periódica.
				
				Por la misma razón que en el caso de \(\sin\), si \(f\) y \(g\) tuvieran ceros adicionales o diferente comportamiento, el producto no coincidiría exactamente con \(1 - \cos\).
				
				Dado que \(1 - \cos\) es función elemental y no se puede escribir como producto no trivial en \(R\), concluimos que es irreducible.
				
				\item \textbf{No unicidad de factorización en \(R\):}
				
				Consideremos ahora la función
				\[
				f(x) = 2(1 - \cos x) = 4 \sin^2\left(\frac{x}{2}\right).
				\]
				
				Esta función puede factorizarse de dos maneras distintas en \(R\):
				\[
				2(1 - \cos x) = 2(1 - \cos x),
				\]
				o
				\[
				4 \sin^2\left(\frac{x}{2}\right) = (2 \sin(x/2))^2.
				\]
				
				Notemos que \(2\) es invertible en \(R\) porque es una constante no nula.
				
				Esto indica que
				\[
				1 - \cos x = \left(2 \sin\left(\frac{x}{2}\right)\right)^2 \cdot \frac{1}{2}.
				\]
				
				Sin embargo, la factorización en términos de \(\sin\) y \(1-\cos\) da diferentes factorizaciones irreducibles para la misma función, lo cual muestra que la factorización no es única.
				
				Por lo tanto, \(R\) no es un dominio de factorización única.
			\end{enumerate}
		\end{proof}
		\begin{observacion}\upshape{
				Un \emph{dominio de factorización única} (DFU) es un dominio de integridad \(R\) tal que todo elemento no nulo y no invertible de \(R\) puede escribirse como un producto finito de elementos irreducibles, y además esta factorización es única salvo el orden de los factores y las unidades asociadas.
				
				Formalmente, para todo \(a \in R\), \(a \neq 0\) y que no sea invertible, existen irreducibles \(p_1, p_2, \ldots, p_n \in R\) tales que
				\[
				a = p_1 p_2 \cdots p_n,
				\]
				y si existe otra factorización
				\[
				a = q_1 q_2 \cdots q_m,
				\]
				con irreducibles \(q_j \in R\), entonces \(m = n\) y existe una permutación \(\sigma\) de \(\{1, \ldots, n\}\) tal que cada \(p_i\) es asociado a \(q_{\sigma(i)}\) (es decir, difieren por un factor invertible).
				
				Este concepto generaliza la factorización única en los enteros y es fundamental en la teoría de anillos y factorización.
			}
		\end{observacion}
		
	\end{enumerate}
	}
\end{ejemplo}









\begin{ejercicio}
	
\upshape{

Sea $R$ un anillo conmutativo y $M$ un $R$-módulo. Un elemento $m \in M$ es de torsión si existe un elemento no nulo $a \in R$ tal que $am = 0$.
\begin{enumerate}
	\item Probar que el conjunto de elementos de torsión de $M$ forma un submódulo de $M$, llamado submódulo de torsión y denotado por $M_{\text{tor}}$.
	\item Supongamos que $R$ es un dominio de ideales principales (DIP) y que $M$ es finitamente generado. Probar que $M/M_{\text{tor}}$ es un módulo libre, y deducir que existe un submódulo libre $F$ de $M$ tal que
	\[
	M \cong M_{\text{tor}} \oplus F.
	\]
\end{enumerate}
}
\end{ejercicio}
	
	\begin{enumerate}
		\item \textbf{Probar que el conjunto de elementos de torsión de $M$ forma un submódulo de $M$, llamado submódulo de torsión y denotado por $M_{\text{tor}}$.}
		
		\begin{proof}
			Definimos el conjunto
			\[
			M_{\text{tor}} := \{ m \in M : \exists a \in R \setminus \{0\} \text{ tal que } am = 0 \}.
			\]
			Queremos probar que $M_{\text{tor}}$ es un submódulo de $M$. Para ello, verificamos las tres propiedades:
			
			\textbf{(i) $0 \in M_{\text{tor}}$:} Para $m = 0 \in M$, tenemos $a \cdot 0 = 0$ para todo $a \in R$, en particular para cualquier $a \neq 0$. Por tanto, $0 \in M_{\text{tor}}$.
			
			\textbf{(ii) Cerrado bajo suma:} Sean $m_1, m_2 \in M_{\text{tor}}$. Entonces existen $a_1, a_2 \in R \setminus \{0\}$ tales que
			\[
			a_1 m_1 = 0 \quad \text{y} \quad a_2 m_2 = 0.
			\]
			Como $R$ es conmutativo, el producto $a := a_1 a_2 \in R \setminus \{0\}$. Entonces,
			\begin{align*}
				a(m_1 + m_2) &= a_1 a_2 (m_1 + m_2) \\
				&= a_2 (a_1 m_1) + a_1 (a_2 m_2) \\
				&= a_2 \cdot 0 + a_1 \cdot 0 = 0.
			\end{align*}
			Por tanto, $m_1 + m_2 \in M_{\text{tor}}$.
			
			\textbf{(iii) Cerrado bajo multiplicación por escalares:} Sea $r \in R$ y $m \in M_{\text{tor}}$, con $a \in R \setminus \{0\}$ tal que $am = 0$. Entonces,
			\[
			a(rm) = r(am) = r \cdot 0 = 0,
			\]
			y por tanto $rm \in M_{\text{tor}}$.
			
			Como $M_{\text{tor}}$ contiene al cero, es cerrado bajo suma y multiplicación por escalares, concluimos que es un submódulo de $M$.
		\end{proof}
		
		\item \textbf{Supongamos que $R$ es un dominio de ideales principales (DIP) y que $M$ es finitamente generado. Probar que $M/M_{\text{tor}}$ es un módulo libre, y deducir que existe un submódulo libre $F$ de $M$ tal que}
		\[
		M \cong M_{\text{tor}} \oplus F.
		\]
		
	\begin{proof}
		Como $R$ es un dominio de ideales principales (DIP), y $M$ es un $R$-módulo finitamente generado, se aplica el \emph{teorema de estructura de módulos sobre un DIP}, el cual afirma que:
		\[
		M \cong \bigoplus_{i=1}^r R \oplus \bigoplus_{j=1}^s R/(a_j),
		\]
		donde $r, s \geq 0$, $a_j \neq 0$ y $a_j \mid a_{j+1}$ para todo $j$. El primer sumando es un módulo libre de rango $r$, mientras que el segundo es un módulo de torsión, pues todo elemento $x_j \in R/(a_j)$ cumple que $a_j x_j = 0$.
		
		Denotemos 
		\[
		F := \bigoplus_{i=1}^r R \quad \text{y} \quad T := \bigoplus_{j=1}^s R/(a_j),
		\]
		entonces 
		\[
		M \cong F \oplus T,
		\]
		donde $F$ es libre y $T$ es de torsión.
		
		\textbf{Afirmación:} $T = M_{\text{tor}}$.\\
		Todo elemento de $T$ es de torsión por construcción, y si $m \in M_{\text{tor}}$, entonces existe $a \neq 0$ tal que $am = 0$. Como $F$ es libre, y en un dominio toda relación de torsión es trivial en un módulo libre, se deduce que $m \in T$. Así, $M_{\text{tor}} = T$.
		
		Consideremos el cociente:
		\[
		M / M_{\text{tor}} = (F \oplus T)/T \cong F.
		\]
		Como $F \cong R^r$ es libre, se concluye que $M / M_{\text{tor}}$ es libre.
		
		Como todo módulo libre es proyectivo, la sucesión exacta corta
		\[
		0 \to M_{\text{tor}} \xrightarrow{\iota} M \xrightarrow{\pi} M/M_{\text{tor}} \to 0
		\]
		se escinde. Es decir, existe un morfismo $s: M/M_{\text{tor}} \to M$ tal que $\pi \circ s = \mathrm{id}$. Entonces la imagen de $s$ es un submódulo complementario de $M_{\text{tor}}$ en $M$.
		
		Sea $F := \operatorname{im}(s)$. Entonces $F \cong M / M_{\text{tor}} \cong R^r$ es libre, y se tiene
		\[
		M = M_{\text{tor}} \oplus F.
		\]
		
		Por tanto, hemos demostrado que si $M$ es un $R$-módulo finitamente generado sobre un DIP, entonces se descompone como suma directa de su submódulo de torsión y un submódulo libre complementario.
	\end{proof}
	\end{enumerate}
	
	\begin{observacion} \upshape
		Un $R$-módulo $M$ se dice \emph{libre} si admite una base, es decir, existe un conjunto $\{e_i\}_{i \in I} \subseteq M$ tal que todo elemento $m \in M$ se escribe de manera única como combinación finita
		\[
		m = \sum_{i \in I} r_i e_i, \quad r_i \in R,
		\]
		donde todos salvo una cantidad finita de los coeficientes $r_i$ son cero. La unicidad de esta representación implica que los $e_i$ son linealmente independientes y generan el módulo.\\
		
	\noindent 	En particular, si $M$ es libre de rango $n$, entonces $M \cong R^n$ como $R$-módulos, es decir, $M$ es isomorfo al producto cartesiano $R \times R \times \cdots \times R$ ($n$ veces), con la estructura natural de módulo componente a componente.
	\end{observacion}
	
	
	\begin{ejercicio}
		\upshape{
		Sea $X$ un conjunto no vacío. Sea $M$ el $\mathbb{Z}$-módulo libre generado por $X$, y sea $N \ne 0$ un submódulo de $M$.
		\begin{enumerate}
			\item ¿Será cierto que $M/N$ es libre?
			\item Si $M/N$ es finitamente generado, ¿es posible determinar su parte libre y de torsión en términos de las clases $\bar{x} = x + N$, con $x \in X$?
		\end{enumerate}
		}
	\end{ejercicio}

\begin{proof}
	Sea \( M \) el \(\mathbb{Z}\)-módulo libre generado por un conjunto no vacío \( X \). Esto significa que
	\[
	M \cong \bigoplus_{x \in X} \mathbb{Z} x,
	\]
	donde cada \( x \in X \) actúa como generador libre.
	
	Sea \( N \subseteq M \) un submódulo distinto de cero, y consideremos el cociente
	\[
	M / N.
	\]
	
	\textbf{(1) ¿Es \( M/N \) necesariamente libre?}
	
	No necesariamente. Para mostrarlo, consideremos el caso particular en que \( M = \mathbb{Z} \) (es decir, \( X = \{1\} \)) y
	\[
	N = 2 \mathbb{Z} = \{ 2k : k \in \mathbb{Z} \}.
	\]
	
	Entonces el cociente es
	\[
	M/N = \mathbb{Z} / 2\mathbb{Z}.
	\]
	
	Veamos si \( \mathbb{Z} / 2\mathbb{Z} \) es un módulo libre sobre \(\mathbb{Z}\):
	
	Supongamos que existe un generador \( \bar{1} = 1 + 2\mathbb{Z} \) que hace libre al módulo. Pero,
	\[
	2 \cdot \bar{1} = 2 + 2\mathbb{Z} = 0 + 2\mathbb{Z} = \bar{0}.
	\]
	
	Esto implica que \( \bar{1} \) es un elemento de torsión, pues \(2 \bar{1} = \bar{0}\) con \(2 \neq 0\). Por definición, un módulo libre no puede tener elementos de torsión distintos de cero. Por lo tanto,
	\[
	\mathbb{Z} / 2\mathbb{Z} \text{ no es un módulo libre}.
	\]
	
	Así, el cociente \( M/N \) no es libre en general.
	
	\bigskip
	
	\textbf{(2) Si \(M/N\) es finitamente generado, ¿cómo determinar su parte libre y de torsión?}
	
	Como \(\mathbb{Z}\) es un dominio de ideales principales (DIP), el teorema de estructura de módulos finitamente generados sobre \(\mathbb{Z}\) garantiza que
	\[
	M/N \cong \mathbb{Z}^r \oplus \left(\bigoplus_{i=1}^t \mathbb{Z}/n_i \mathbb{Z}\right),
	\]
	con \(r \geq 0\), \(t \geq 0\), y \(n_i \geq 2\) enteros tales que \(n_i \mid n_{i+1}\).
	
	Ahora, para ver esto en términos de las clases \(\bar{x} = x + N\), con \(x \in X\), procedemos así:
	
	\medskip
	
	Para cada generador \(x \in X\), consideramos la clase \(\bar{x} \in M/N\). 
	
	\textbf{Caso 1:} \(\bar{x}\) es un elemento de torsión.
	
	Esto significa que existe un entero \(n \geq 1\) tal que
	\[
	n \bar{x} = \overline{n x} = \bar{0} \quad \Longleftrightarrow \quad n x \in N.
	\]
	
	Por ejemplo, si \(n=3\) y \(3x \in N\), entonces el orden de \(\bar{x}\) divide a 3, y por lo tanto \(\bar{x}\) genera un factor \(\mathbb{Z}/3\mathbb{Z}\) en la descomposición.
	
	\textbf{Cálculo:} Si \(x \in X\) y se encuentra \(n \neq 0\) con \(n x \in N\), entonces el elemento \(\bar{x}\) satisface
	\[
	n \bar{x} = \bar{0},
	\]
	es decir, \(\bar{x}\) es torsión de orden (divisor de) \(n\).
	
	\medskip
	
	\textbf{Caso 2:} \(\bar{x}\) es libre.
	
	Esto ocurre si no existe ningún entero \(n \neq 0\) tal que
	\[
	n x \in N,
	\]
	es decir, la clase \(\bar{x}\) no es anulada por ningún entero distinto de cero. En este caso,
	\[
	\mathbb{Z} \cdot \bar{x} \cong \mathbb{Z},
	\]
	es un submódulo libre de \(M/N\).
	
	\medskip
	
	\textbf{Ejemplo explícito:}
	
	Supongamos
	\[
	M = \mathbb{Z} e_1 \oplus \mathbb{Z} e_2,
	\]
	y
	\[
	N = \langle 2 e_1, 3 e_2 \rangle,
	\]
	es decir, \(N\) es generado por \(2 e_1\) y \(3 e_2\).
	
	Entonces, consideramos
	\[
	M/N = \frac{\mathbb{Z} e_1 \oplus \mathbb{Z} e_2}{\langle 2 e_1, 3 e_2 \rangle}.
	\]
	
	Las clases \(\bar{e}_1, \bar{e}_2\) en \(M/N\) cumplen:
	\[
	2 \bar{e}_1 = \overline{2 e_1} = \bar{0} \implies \bar{e}_1 \text{ es torsión de orden 2},
	\]
	y
	\[
	3 \bar{e}_2 = \overline{3 e_2} = \bar{0} \implies \bar{e}_2 \text{ es torsión de orden 3}.
	\]
	
	Por lo tanto,
	\[
	M/N \cong \mathbb{Z}/2\mathbb{Z} \oplus \mathbb{Z}/3\mathbb{Z}.
	\]
	
	En este caso, no hay parte libre.
	
	\medskip
	
	Si en cambio tomamos
	\[
	N = \langle 2 e_1 \rangle,
	\]
	entonces
	\[
	M/N = \frac{\mathbb{Z} e_1 \oplus \mathbb{Z} e_2}{\langle 2 e_1 \rangle}.
	\]
	
	Las clases \(\bar{e}_1, \bar{e}_2\) cumplen:
	\[
	2 \bar{e}_1 = \bar{0} \implies \bar{e}_1 \text{ es torsión de orden 2},
	\]
	pero
	\[
	n \bar{e}_2 = \overline{n e_2} = \bar{0} \implies n e_2 \in N = \langle 2 e_1 \rangle,
	\]
	y como \(e_2\) y \(e_1\) son generadores libres, ningún múltiplo no nulo de \(e_2\) está en \(N\). Por tanto, no existe \(n \neq 0\) con \(n \bar{e}_2 = \bar{0}\), y \(\bar{e}_2\) es libre.
	
	Por lo tanto,
	\[
	M/N \cong \mathbb{Z}/2\mathbb{Z} \oplus \mathbb{Z},
	\]
	donde
	\[
	M_{\text{tor}} = \langle \bar{e}_1 \rangle \cong \mathbb{Z}/2\mathbb{Z}, \quad F = \langle \bar{e}_2 \rangle \cong \mathbb{Z}.
	\]
	
\end{proof}



	
	\appendix
	\chapter{}
	
\section{Relación de orden parcial y Lema de Zorn}



\begin{definicion}[Relación de orden parcial]\label{lkjbhgcdb}
	\upshape{
		Dada una relación  ''\(\prec\)`` en un conjunto \(\mathcal{F}\), decimos que \(\prec\) es una \textit{relación de orden parcial} si cumple:
		\begin{enumerate}
			\item \textit{Reflexividad}: Para todo \(a \in \mathcal{F}\), se cumple que \(a \prec a\). 
			\item \textit{Antisimetría}: Para todos \(a, b \in \mathcal{F}\), si \(a \prec b\) y \(b \prec a\), entonces \(a = b\).
			\item \textit{Transitividad}: Para todos \(a, b, c \in \mathcal{F}\), si \(a \prec b\) y \(b \prec c\), entonces \(a \prec c\).
		\end{enumerate}
	}
\end{definicion}
\begin{observacion}
	\upshape{
	Sean $a,b\in\mathcal{F}$. Decimos que $a$ es \textit{mayor} que $b$ si $b\prec a$.
	}
\end{observacion}


\begin{proposicion}\label{ibvdshkjcsd}
	\upshape{
		La relación de inclusión ``$\subset$'' es una relación de orden parcial en el conjunto de subconjuntos de un conjunto dado.
	}
\end{proposicion}
\begin{proof}
	Sea $X$ un conjunto y consideremos $\mathcal{F}$, un subconjunto del conjunto potencia de $X$ con la relación de inclusión ``$\subset$''. Entonces
	
	\begin{itemize}
		\item \textit{Reflexividad}: Para cualquier subconjunto $A \subset X$, se cumple que $A \subset A$.
		
		\item \textit{Antisimetría}: Para cualesquiera  subconjuntos $A, B \subset X$, si $A \subset B$ y $B \subset A$, entonces $A = B$. 
		
		\item \textit{Transitividad}: Para cualesquiera subconjunto $A, B, C \subset X$, si $A \subset B$ y $B \subset C$, entonces $A \subset C$. 
	\end{itemize}
	Concluimos que ``$\subset$'' es una relación de orden parcial en $\mathcal{F}$.
\end{proof}



\begin{definicion}[Orden total]
	\upshape{
	
		Sea $\mathcal{F}$ un conjunto con una relación de orden parcial $\prec$. Decimos que $\prec$ es un \textit{orden total} (o \textit{orden lineal}) en $\mathcal{F}$ si cumple la siguiente propiedad:
		\begin{enumerate}
			\item [$4.$] \textit{Comparabilidad}: para cualquier par de elementos $a, b \in \mathcal{F}$, se cumple que $a \prec b$ o $b \prec a$ (o ambos).
		\end{enumerate} 
	}
\end{definicion}
\begin{definicion}
	\upshape{
	
	Sea $\mathcal{F}$ un conjunto parcialmente ordenado con una relación de orden parcial ''$\prec$''. Un subconjunto $\mathcal{C} \subseteq \mathcal{F}$ se llama \textit{cadena} si $ \mathcal{C}$ es totalmente ordenado con la relación ''$\prec$''.
	
	}
\end{definicion}
\begin{observacion}
	\upshape{
	Sea $d\in\mathcal{F}$. Decimos que $d$ es una \textit{cota superior} de $\mathcal{C}$ si $a\prec d$ para todo $a\in\mathcal{C}$.
	}
\end{observacion}

\begin{lema}[Lema de Zorn]\label{lemazorn}
	\upshape{
	Supongamos que $\mathcal{F}$ es un conjunto parcialmente ordenado tal que cada cadena en $\mathcal{F}$ tiene una cota superior en $\mathcal{F}$. Entonces, $\mathcal{F}$ contiene al menos un elemento maximal(no existe otro elemento que sea estrictamente mayor que él con la relación de orden).
	}
\end{lema}








\backmatter
	\begin{thebibliography}{99}
		
		\bibitem{Dorronsoro1996}
		Dorronsoro, J. \& Hernández, E. (1996).
		\textit{Números, grupos y anillos}.
		Universidad Autónoma de Madrid.
		
		\bibitem{FernadoGuvea}
		Gouvêa, F. (2020). \textit{$p-$adict Numbers: An Introduction}. Springer. 
			
		
	\end{thebibliography}
	
	
\end{document}









%\section*{Demostración}


%\noindent En efecto. Fijado $p\in\mathbb{N}$, primo. Sea la susesión $\left(S_n\right)_{n\in\mathbb{N}}$ en $\mathbb{Q}$, donde $S_n:=\sum_{k=0}^{n} p^k$

%\begin{itemize}

%	\item $\lim_{n\to\infty}p^{n}=0$. Sea $\varepsilon>0$, luego por el Principio arquimediano existe $m_0\in \mathbb{N}\in $ tal que $p^{m_0}\varepsilon>1$, entonces  para todo $n\geq m_0$ se tiene $$|p^n - 0|_p=p^{-v_p(p^n)}=p^{-n}<\varepsilon $$
	
	

%	\item $\left(S_n\right)_{n\in\mathbb{N}}$ es Cauchy. Podemos suponer que $n<m$ entonces

%	\begin{align*}

%		\left| S_m-S_n\right|_p&=\left|\, \sum_{k=n+1}^{m}p^k\,\right|_p=\left|p^{n+1}\left(1+p+p^2+\dots+ p^{m-n-1}\right)\right|_p\\

%		&\\

%		&=\left|p^{n+1}\right|_p \left| 1+p+p^2+\dots + p^{m-n-1} \right|_p

%	\end{align*}

%	Como $p$ no divide a $1+p+p^2+\dots + p^{m-n-1}$, luego $v_p\left(1+p+p^2+\dots + p^{m-n-1} \right)=0.$ Así\\

%	$$\left| 1+p+p^2+\dots + p^{m-n-1} \right|_p=p^{-v_p\left(1+p+p^2+\dots + p^{m-n-1} \right)}=p^{-0}=1.$$ 
	

%	\noindent Así. Dado $\varepsilon>0$ entonces existe $n_0$ tal que si $n,m\geq n_o$, entonces $$\left| S_m-S_n\right|_p=|p^{n+1}|<\varepsilon.$$

%	\item Notemos que $$S_n=\sum_{k=0}^{n}p^k=\dfrac{1-p^{n+1}}{1-p}$$ de donde $$S_n \xrightarrow{\scalebox{1}{$|\cdot|_p$}} \dfrac{1}{1-p}$$
	

%\end{itemize}








